% Generated mechanically from frozen input p07/ko/pione.tex.
% Do not edit this generated file; edit the bound locale source instead.

% OK, start here.
%

\setcounter{chapter}{57}
\chapter{스킴의 기본군}



\phantomsection
\label{pione-section-phantom}


\section{서론}
\label{pione-section-introduction}

\noindent
이 장에서는 Grothendieck의 스킴 기본군과 그 응용을 논한다.
기초 문헌은 \cite{SGA1}이다.
좋은 입문서로 \cite{Lenstra}가 있다.
그 밖의 참고문헌으로 \cite{Murre-lectures}와
\cite{Grothendieck-Murre}가 있다.










\section{한 점 위에서 에탈인 스킴}
\label{pione-section-schemes-etale-point}

\noindent
이 절에서는 체의 스펙트럼 위에서 에탈인 스킴을 기술한다.
결과를 서술하기 전에 위상군 $G$에 대한 $G$-집합의 범주를 도입한다.

\begin{definition}
\label{pione-definition-G-set-continuous}
위상군 $G$를 잡자.
{\it $G$-집합} 또는 {\it 이산 $G$-집합}이란 집합 $X$와
왼쪽 작용 $a : G \times X \to X$의 쌍으로서, $X$에는 이산위상을 주고
$G \times X$에는 곱위상을 주었을 때 $a$가 연속인 것을 말한다.
{\it $G$-집합의 사상} $f : X \to Y$란 단순히 $G$-등변인 사상,
즉 $X$에서 $Y$로 가는 사상을 말한다.
$G$-집합의 범주는 {\it $G\textit{-Sets}$}로 나타낸다.
\end{definition}

\noindent
$a : G \times X \to X$가 연속이라는 조건은 단순히 임의의
$x \in X$의 안정자가 $G$의 열린 부분군이라는 뜻이다.
% P07-ERR-1144: frozen-source duplicate group symbol retained; see sidecar.
$G$가 추상군 $G$(즉, 위상군이 아닌 군)인 경우, $G$에
이산위상을 주면 앞의 정의와 일치한다(예를 들어
사이트, 예 \ref{sites-example-site-on-group}을 보라).

\medskip\noindent
$L/K$가 무한 Galois 확대이면 Galois 군
$G = \text{Gal}(L/K)$에는 표준적인 위상이 주어짐을 상기하자.
체, 절 \ref{fields-section-infinite-galois}을 보라.

\begin{lemma}
\label{pione-lemma-sheaves-point}
체 $K$를 잡자. $K^{sep}$를 $K$의 분리 폐포라 하자.
프로유한군 $G = \text{Gal}(K^{sep}/K)$를 생각하자. 함자
$$
\begin{matrix}
\text{스킴(}K\text{ 위에서 에탈)} &
\longrightarrow &
G\textit{-Sets} \\
X/K & \longmapsto &
\Mor_{\Spec(K)}(\Spec(K^{sep}), X)
\end{matrix}
$$
는 범주의 동치이다.
\end{lemma}

\begin{proof}
스킴 $X$가 $K$ 위에 있고 $K$ 위에서 에탈일 필요충분조건은
$X \cong \coprod_{i\in I} \Spec(K_i)$로 쓸 수 있고 각 $K_i$가 $K$의
유한 분리 확대인 것이다
(사상, 보조정리 \ref{morphisms-lemma-etale-over-field}).
보조정리의 함자는 $X$에 $G$-집합
$$
\coprod\nolimits_i \Hom_K(K_i, K^{sep})
$$
를 그 자연스러운 왼쪽 $G$-작용과 함께 대응시킨다. $G$ 위상의 정의에 따라
각 원소의 안정자는 열려 있다. 반대로 임의의 $G$-집합 $S$는 그 궤도들의
서로소 합이다. $S = \coprod S_i$라 쓰자. $s_i \in S_i$를 고르고 그
열린 안정자를 $G_i \subset G$로 나타내자. Galois 이론
(체, 정리 \ref{fields-theorem-inifinite-galois-theory})에 의해 체
$(K^{sep})^{G_i}$는 $K$의 유한 분리 확대이다. 따라서 스킴
$$
\coprod\nolimits_i \Spec((K^{sep})^{G_i})
$$
은 $K$ 위에서 에탈이다. 이로써 보조정리의 함자에 대한 역함자를 얻는다.
일부 세부사항은 생략한다.
\end{proof}

\begin{remark}
\label{pione-remark-covering-surjective}
보조정리 \ref{pione-lemma-sheaves-point}의 대응 아래에서 체 $K$의 작은 에탈 사이트
$\Spec(K)_\etale$의 덮개들은
$G\textit{-Sets}$에서 사상들의
전사족에 대응한다.
\end{remark}










\section{Galois 범주}
\label{pione-section-galois}

\noindent
이 절에서는 \cite[Expos\'e V, Sections 4, 5, and 6]{SGA1}에 실린
내용의 일부를 논한다.

\medskip\noindent
$F : \mathcal{C} \to \textit{Sets}$를 함자라 하자.
우리의 규약상 범주의 대상들은 집합을 이루며, 임의의 두 대상 사이의 사상들도
집합을 이룸을 상기하자. 표준적인 단사 사상
\begin{equation}
\label{pione-equation-embedding-product}
\text{Aut}(F)
\longrightarrow
\prod\nolimits_{X \in \Ob(\mathcal{C})} \text{Aut}(F(X))
\end{equation}
이 있다. 집합 $E$에 대해 $\text{Aut}(E)$에는 콤팩트-열린 위상을 준다
(위상수학, 예 \ref{topology-example-automorphisms-of-a-set}).
$E$가 유한이면 이는 물론 이산위상이며, 이 절에서 관심을 두는 경우가
바로 이것이다\footnote{pro-\'etale 기본군을 논할 때에는 일반적인 경우가 중요하다.}.
$\text{Aut}(F)$에는 (\ref{pione-equation-embedding-product})의 오른쪽에 있는 곱위상으로부터
유도되는 위상을 준다. 특히 작용 사상
$$
\text{Aut}(F) \times F(X) \longrightarrow F(X)
$$
은 $F(X)$에 이산위상을 주면 연속이다. 실제로 임의의 집합 $E$에 대해 작용 사상
$\text{Aut}(E) \times E \to E$가 연속이기 때문이다.
$\text{Aut}(F)$에 준 위상의 보편 성질은 다음과 같다. $G$를 위상군이라 하고,
$G \to \text{Aut}(F)$를 군 준동형으로서 유도된 작용
$G \times F(X) \to F(X)$가 모든 $X \in \Ob(\mathcal{C})$에 대해
$F(X)$에 이산위상을 주었을 때 연속이라고 하자. 그러면
$G \to \text{Aut}(F)$는 연속이다.

\medskip\noindent
다음 보조정리는 유한 집합의 범주로 가는 함자의 자기동형군이 자동으로
프로유한군이 됨을 말해 준다.

\begin{lemma}
\label{pione-lemma-aut-inverse-limit}
$\mathcal{C}$를 범주, $F : \mathcal{C} \to \textit{Sets}$를 함자라 하자.
사상 (\ref{pione-equation-embedding-product})은 $\text{Aut}(F)$를
$\prod_{X \in \Ob(\mathcal{C})} \text{Aut}(F(X))$의 닫힌 부분군과 동일시한다.
특히 모든 $X$에 대해 $F(X)$가 유한이면 $\text{Aut}(F)$는
프로유한군이다.
\end{lemma}

\begin{proof}
$\xi = (\gamma_X) \in \prod \text{Aut}(F(X))$가
$\text{Aut}(F)$에 속하지 않는 원소라 하자. 그러면 $\mathcal{C}$의 사상
$f : X \to X'$와 원소 $x \in F(X)$가 존재하여
$F(f)(\gamma_X(x)) \not = \gamma_{X'}(F(f)(x))$.
를 만족한다. $\gamma_X$의 열린 근방
$U = \{\gamma \in \text{Aut}(F(X)) \mid \gamma(x) = \gamma_X(x)\}$
와 열린 근방
$U' = \{\gamma' \in \text{Aut}(F(X')) \mid \gamma'(F(f)(x)) =
\gamma_{X'}(F(f)(x))\}$.
를 생각하자. 그러면
$U \times U' \times \prod_{X'' \not = X, X'} \text{Aut}(F(X''))$
은 $\xi$의 열린 근방이며 $\text{Aut}(F)$와 만나지 않는다.
마지막 주장은 각 $F(X)$가 유한이면 $\prod \text{Aut}(F(X))$가
프로유한 공간이라는 사실에서 따른다.
\end{proof}

\begin{example}
\label{pione-example-galois-category-G-sets}
$G$를 위상군이라 하자. 중요한 예는 망각 함자
\begin{equation}
\label{pione-equation-forgetful}
\textit{Finite-}G\textit{-Sets} \longrightarrow \textit{Sets}
\end{equation}
이다. 여기서 $\textit{Finite-}G\textit{-Sets}$는 유한 $G$-집합을 대상으로 하는
$G\textit{-Sets}$의 충만 부분범주이다. $G$-집합의 범주 $G\textit{-Sets}$는
정의 \ref{pione-definition-G-set-continuous}에서 정의하였다.
\end{example}

\noindent
$G$를 위상군이라 하자. $G$의 {\it 프로유한 완비화}란 프로유한 위상을 갖춘
프로유한군
$$
G^\wedge =
\lim_{U \subset G\text{ open, normal, finite index}} G/U
$$
을 말한다. 유한 지수의 열린 정규 부분군들을 유한 개 교차해도 같은 성질을 가지므로
이 극한은 공여과적이다. 프로유한 완비화의 보편 성질은 프로유한군 $H$로 가는 임의의 연속 사상
$G \to H$가 표준적으로
$G \to G^\wedge \to H$로 분해된다는 것이다.

\begin{lemma}
\label{pione-lemma-single-out-profinite}
$G$를 위상군이라 하자. 함자 (\ref{pione-equation-forgetful})의 자기동형군에 보조정리
\ref{pione-lemma-aut-inverse-limit}의 프로유한 위상을 준 것은 $G$의 프로유한 완비화이다.
\end{lemma}

\begin{proof}
(\ref{pione-equation-forgetful})의 함자를 $F_G$라 쓰자. $\textit{Finite-}G\textit{-Sets}$의
임의의 사상 $X \to Y$는 $G$의 작용과 가환한다.
따라서 임의의 $g \in G$는 $F_G$의 자기동형을 정하며, 표준적인 군 준동형
$G \to \text{Aut}(F_G)$를 얻는다. 임의의 유한 $G$-집합 $X$는 표준적인 $G$-작용을
갖는 꼴 $G/H_i$의 $G$-집합들의 유한 서로소 합이다. 여기서 $H_i \subset G$는
유한 지수의 열린 부분군이다. 이때 $U_i = \bigcap gH_ig^{-1}$는 열리고 정규이며
유한 지수를 갖는다. 또한 $U_i$는 $G/H_i$에 자명하게 작용하므로
$U = \bigcap U_i$는 $F_G(X)$에 자명하게 작용한다. 따라서 작용
$G \times F_G(X) \to F_G(X)$는 연속이다. $\text{Aut}(F_G)$ 위상의
보편 성질에 의해 사상 $G \to \text{Aut}(F_G)$는 연속이다.
보조정리 \ref{pione-lemma-aut-inverse-limit}과 프로유한 완비화의 보편 성질에 의해 유도된
연속 군 준동형
$$
G^\wedge \longrightarrow \text{Aut}(F_G)
$$
이 존재한다. 또한 $G/U$가 $G/U$에 충실하게 작용하므로 이 사상은 단사이다.
그 상이 조밀하면 이 사상은 전사이고, 따라서 위상수학, 보조정리
\ref{topology-lemma-bijective-map}에 의해 위상동형이다.

\medskip\noindent
$\gamma \in \text{Aut}(F_G)$와 $X \in \Ob(\mathcal{C})$를 잡자.
$\gamma$와 $g$가 $F_G(X)$에 같은 작용을 유도하는 $g \in G$가
존재함을 보이면 증명이 끝난다. 앞에서와 같이 $X$는 $G/H_i$들의 유한 서로소 합이다.
위와 같은 $U_i$와 $U$를 쓰면, 유한 $G$-집합 $Y = G/U$는 모든 $i$에 대해
$G/H_i$ 위로 전사한다. 따라서 $\gamma$와 $g$가
$F_G(G/U) = G/U$에 같은 작용을 유도하는 $g \in G$를 찾으면 충분하다.
$e \in G$를 항등원이라 하고, 어떤 $g_0 \in G$에 대해 $\gamma(eU) = g_0U$라 쓰자.
임의의 $g_1 \in G$에 대해 사상
$$
R_{g_1} : G/U \longrightarrow G/U,\quad gU \longmapsto gg_1U
$$
은 $\textit{Finite-}G\textit{-Sets}$의 사상이며 $\gamma$의 작용과 가환한다. 따라서
$$
\gamma(g_1U) = \gamma(R_{g_1}(eU)) = R_{g_1}(\gamma(eU)) =
R_{g_1}(g_0U) = g_0g_1U
$$
이다. 그러므로 $g = g_0$로 잡으면 된다.
\end{proof}

\noindent
완전 함자란 모든 유한 극한 및 유한 여극한과 가환하는 함자이다.
특히 그러한 함자는 등화자, 쌍대등화자, 섬유곱, 푸시아웃 등과 가환한다.

\begin{lemma}
\label{pione-lemma-second-fundamental-functor}
$G$를 위상군이라 하자.
$F : \textit{Finite-}G\textit{-Sets} \to \textit{Sets}$를
완전 함자라 하고, 모든 $X$에 대해 $F(X)$가 유한하다고 하자.
그러면 $F$는 함자 (\ref{pione-equation-forgetful})와 동형이다.
\end{lemma}

\begin{proof}
$X$를 $\textit{Finite-}G\textit{-Sets}$의 공집합이 아닌 대상이라 하자. 도식
$$
\xymatrix{
X \ar[r] \ar[d] & \{*\} \ar[d] \\
\{*\} \ar[r] & \{*\}
}
$$
은 쌍대카르테시안이다. 따라서 $F(X)$는 공집합이 아니다.
$U \subset G$를 유한 지수의 열린 정규 부분군이라 하자. 등식
$$
G/U \times G/U = \coprod\nolimits_{gU \in G/U} G/U
$$
에서 $gU$에 대응하는 성분은 왼쪽의 $(eU, gU)$의 궤도에 대응한다. 따라서
$$
F(G/U) \times F(G/U) = F(G/U \times G/U) = \coprod\nolimits_{gU \in G/U} F(G/U)
$$
$|F(G/U)| = |G/U|$이다. 실제로 $F(G/U)$는 공집합이 아니다. 그러므로
$$
% P07-ERR-1145: frozen-source spelling retained; see sidecar.
\lim_{U \subset G\text{ open, normal, finite idex}} F(G/U)
$$
은 공집합이 아니다(범주론, 보조정리 \ref{categories-lemma-nonempty-limit}).
이 극한의 원소 $\gamma = (\gamma_U)$를 하나 잡자.
함자 (\ref{pione-equation-forgetful})를 $F_G$라 쓰자. $F_G$를 함자
$$
X \longmapsto \colim_U \Mor(G/U, X)
$$
와 동일시할 수 있다. 여기서 $f : G/U \to X$는 $f(eU) \in X = F_G(X)$에 대응한다
(세부사항은 생략한다). 따라서 원소 $\gamma$는 잘 정의된 사상
$$
t : F_G \longrightarrow F
$$
를 정한다. 실제로 $x \in X$가 주어지면 $U$와 $eU$를 $x$로 보내는
$f : G/U \to X$를 고르고 $t_X(x) = F(f)(\gamma_U)$로 놓는다.
임의의 $U$에 대해 $t$가 전단사 사상
$t_{G/U} : F_G(G/U) \to F(G/U)$를 유도함을 보이자.
그러면 $t$가 동형이라는 것은 곧바로 따른다(세부사항은 생략한다).
$|F_G(G/U)| = |F(G/U)|$이므로 $t_{G/U}$가 전사임을 보이면 충분하다.
그 상은 적어도 원소
$t_{G/U}(eU) = F(\text{id}_{G/U})(\gamma_U) = \gamma_U$.
를 포함한다. $g \in G$에 대해 오른쪽 곱셈을 $R_g : G/U \to G/U$로 나타내자.
$F(R_g) : F(G/U) \to F(G/U)$의 고정점 집합은 $F(\emptyset) = \emptyset$과 같다.
이는 $g \not \in U$이면 $F$가 등화자와 가환하기 때문이다.
따라서 $g_1, \ldots, g_{|G/U|}$가 $G/U$의 대표원계이면 원소
$F(R_{g_i})(\gamma_U)$들은 서로 다르며, 따라서 $F(G/U)$ 전체를 이룬다. 이때
$$
t_{G/U}(g_iU) = F(R_{g_i})(\gamma_U)
$$
이므로 증명이 끝난다.
\end{proof}

\begin{example}
\label{pione-example-from-C-F-to-G-sets}
$\mathcal{C}$를 범주, $F : \mathcal{C} \to \textit{Sets}$를 함자라 하고,
모든 $X \in \Ob(\mathcal{C})$에 대해 $F(X)$가 유한하다고 하자.
보조정리 \ref{pione-lemma-aut-inverse-limit}에 의해 $G = \text{Aut}(F)$에는 표준적으로
프로유한 위상군의 구조가 주어진다. 함자
\begin{equation}
\label{pione-equation-remember}
\mathcal{C} \longrightarrow \textit{Finite-}G\textit{-Sets},\quad
X \longmapsto F(X)
\end{equation}
를 얻는다. 여기서 $F(X)$에는 $G$의 유도된 작용을 준다. 이 작용은
$\text{Aut}(F)$ 위상의 구성에 의해 연속이다.
\end{example}

\noindent
Galois 범주를 정의하는 목적은 함자 (\ref{pione-equation-remember})가 동치가 되는
쌍 $(\mathcal{C}, F)$를 특징짓는 데 있다. Galois 범주를 다음과 같이 정의한다.

\begin{definition}
\label{pione-definition-galois-category}
\begin{reference}
\cite[Expos\'e V, Definition 5.1]{SGA1}의 정의와는 다르다.
\cite[Definition 7.2.1]{BS}와 비교하라.
\end{reference}
$\mathcal{C}$를 범주, $F : \mathcal{C} \to \textit{Sets}$를 함자라 하자.
쌍 $(\mathcal{C}, F)$가 {\it Galois 범주}라는 것은 다음 조건들을 만족한다는 뜻이다.
\begin{enumerate}
\item $\mathcal{C}$는 유한 극한과 유한 여극한을 갖는다.
\item
\label{pione-item-connected-components}
$\mathcal{C}$의 모든 대상은 연결 대상들의 유한 쌍대곱(공집합일 수도 있다)이다.
\item 모든 $X \in \Ob(\mathcal{C})$에 대해 $F(X)$는 유한하다.
\item $F$는 동형을 반영하고\footnote{즉, $\mathcal{C}$의 사상 $f$에 대해
$F(f)$가 동형이면 $f$도 동형이다.}, 완전하다\footnote{이는 $F$가
유한 극한 및 유한 여극한과 가환한다는 뜻이다. 범주론, 절
\ref{categories-section-exact-functor}을 보라.}.
\end{enumerate}
여기서 $X \in \Ob(\mathcal{C})$가 연결이라는 것은 이 대상이 시작 대상이 아니며,
임의의 단사 사상 $Y \to X$에 대해 $Y$가 시작 대상이거나
$Y \to X$가 동형이라는 뜻이다.
\end{definition}

\noindent
{\bf 주의:} 이 정의는 대부분의 문헌에 나오는 정의와 같지 않다(뒤에서 두 정의가
동치임을 보일 것이다). 실제로 \cite[Expos\'e V, Definition 5.1]{SGA1}에서는
Galois 범주를 어떤 프로유한군 $G$에 대한 $\textit{Finite-}G\textit{-Sets}$와
동치인 범주로 정의한다. 이어서 Grothendieck은 위 공리들보다 약한 공리
(G1)--(G6)의 목록으로 Galois 범주를 특징짓는다. 여기의 정의를 택한 이유는
유한 극한과 유한 여극한의 존재 및 함자 $F$의 완전성을 강조하기 위해서이다.
그 대가로 뒤에서 이 절의 결과를 적용할 때 조금 더 많은 작업이 필요하다.

\begin{lemma}
\label{pione-lemma-epi-mono}
$(\mathcal{C}, F)$를 Galois 범주라 하고,
$X \to Y \in \text{Arrows}(\mathcal{C})$라 하자. 그러면 다음이 성립한다.
\begin{enumerate}
\item $F$는 충실하다.
\item $X \to Y$가 단사 사상이다.
$\Leftrightarrow F(X) \to F(Y)$가 단사 사상이다.
\item $X \to Y$가 전사 사상이다.
$\Leftrightarrow F(X) \to F(Y)$가 전사 사상이다.
\item 대상 $A$가 $\mathcal{C}$의 시작 대상일 필요충분조건은
$F(A) = \emptyset$이다.
\item 대상 $Z$가 $\mathcal{C}$의 끝 대상일 필요충분조건은
$F(Z)$가 한원소 집합인 것이다.
\item $X$와 $Y$가 연결이면 $X \to Y$는 전사 사상이다.
\item
\label{pione-item-one-element}
$X$가 연결이고 $a, b : X \to Y$가 두 사상이라고 하자.
$F(a)$와 $F(b)$가 $F(X)$의 한 원소에서 일치하면 $a = b$이다.
\item $X = \coprod_{i = 1, \ldots, n} X_i$이고
$Y = \coprod_{j = 1, \ldots, m} Y_j$이며 $X_i$, $Y_j$가 연결이라고 하자.
그러면 사상 $\alpha : \{1, \ldots, n\} \to \{1, \ldots, m\}$가 존재하여
$X \to Y$는 사상들의 모음 $X_i \to Y_{\alpha(i)}$로부터 얻어진다.
\end{enumerate}
\end{lemma}

\begin{proof}
(1)의 증명. $a, b : X \to Y$이고 $F(a) = F(b)$라고 하자.
$E$를 $a$와 $b$의 등화자라 하자. 그러면 $F(E) = F(X)$이고,
$F$가 동형을 반영하므로 $E = X$이다.

\medskip\noindent
(2)의 증명. $F$는 사상 $X \to X \times_Y X$를 사상
$F(X) \to F(X) \times_{F(Y)} F(X)$로 보내고, $F$는 동형을 반영하기 때문이다.

\medskip\noindent
(3)의 증명. $F$는 사상 $Y \amalg_X Y \to Y$를 사상
$F(Y) \amalg_{F(X)} F(Y) \to F(Y)$로 보내고, $F$는 동형을 반영하기 때문이다.

\medskip\noindent
(4)의 증명. 시작 대상 $A$가 존재하고 분명히 $F(A) = \emptyset$이다.
한편 대상 $X$가 $F(X) = \emptyset$을 만족하면 유일한 사상 $A \to X$는
전단사 $F(A) \to F(X)$를 유도하므로 $A \to X$는 동형이다.

\medskip\noindent
(5)의 증명. 끝 대상 $Z$가 존재하고 분명히 $F(Z)$는 한원소 집합이다.
한편 대상 $X$에 대해 $F(X)$가 한원소 집합이면 유일한 사상 $X \to Z$는
전단사 $F(X) \to F(Z)$를 유도하므로 $X \to Z$는 동형이다.

\medskip\noindent
(6)의 증명. 두 사상 $Y \to Y \amalg_X Y$의 등화자 $E$는
$\mathcal{C}$의 시작 대상이 아니다. 실제로 $X \to Y$가 $E$를 거쳐 분해되고
$F(X) \not = \emptyset$이기 때문이다. 따라서 $E = Y$이며 주장이 따른다.

\medskip\noindent
(\ref{pione-item-one-element})의 증명.
$a$와 $b$의 등화자 $E$에는 단사 사상 $E \to X$가 주어지며,
$F(E) \subset F(X)$는 $F(a)$와 $F(b)$가 일치하는 원소들의 집합이다.
$E$가 시작 대상이거나 $E = X$이므로 주장이 따른다.

\medskip\noindent
(8)의 증명. 각 $i, j$에 대해 $E_{ij} = X_i \times_Y Y_j$는 시작 대상이거나
$X_i$와 같다. $s \in F(X_i)$를 하나 고르면 $E_{ij} = X_i$일 필요충분조건은
$s$가 $F(Y_j) \subset F(Y)$의 원소로 보내지는 것이다.
따라서 이는 유일한 $j = \alpha(i)$에 대해 일어난다.
\end{proof}

\noindent
위 보조정리에 의해 Galois 범주 $(\mathcal{C}, F)$의 연결 대상 $X$가 주어지면
자기동형군 $\text{Aut}(X)$의 위수는 $|F(X)|$ 이하이다. 실제로
$s \in F(X)$와 $g \in \text{Aut}(X)$에 대해
$F(g)(s) = s$일 필요충분조건은 (\ref{pione-item-one-element})에 의해
$g = \text{id}_X$인 것이다. 등호가 성립할 때 $X$를
{\it Galois}라고 한다. 동치로, $X$가 Galois라는 것은
연결이고 $\text{Aut}(X)$가 $F(X)$에 추이적으로 작용한다는 뜻이다.

\begin{lemma}
\label{pione-lemma-galois}
$(\mathcal{C}, F)$를 Galois 범주라 하자. $\mathcal{C}$의 임의의 연결 대상 $X$에 대해
Galois 대상 $Y$와 사상 $Y \to X$가 존재한다.
\end{lemma}

\begin{proof}
이하에서는 보조정리 \ref{pione-lemma-epi-mono}의 결과를 별도의 언급 없이 사용한다.
$n = |F(X)|$라 하자. $X^n$에 $S_n$의 자연스러운 작용을 준 것을 생각하자.
$$
X^n = \coprod\nolimits_{t \in T} Z_t
$$
를 연결 대상들로의 분해라 하자. $F(Z_t)$가 서로 다른 $s_i$들로 이루어진 원소
$(s_1, \ldots, s_n)$를 포함하도록 $t$를 하나 고르자. 다른 원소
$(s'_1, \ldots, s'_n) \in F(Z_t)$에 대해서도 $s'_i$들이 서로 다르다고 주장한다.
% P07-ERR-1146: frozen-source article error recorded in sidecar.
그렇지 않고, 예를 들어 $s'_i = s'_j$라면 $Z_t$는 $X^{n - 1}$의 어떤 연결 성분을
대각 사상
$$
\Delta_{ij} : X^{n - 1} \longrightarrow X^n
$$
으로 보낸 상이다. 연결 대상 사이의 사상은 전사 사상이고, $F$를 적용하면 전사 사상을
유도하므로 $s_i = s_j$가 따라야 하지만 이는 성립하지 않는다.

\medskip\noindent
$G \subset S_n$을 $g(Z_t) = Z_t$를 만족하는 원소들로 이루어진 부분군이라 하자.
$S_n$의
$$
F(X)^n = F(X^n) = \coprod\nolimits_{t' \in T} F(Z_{t'})
$$
위의 작용을 살펴보면 $G = \{g \in S_n \mid g(s_1, \ldots, s_n) \in F(Z_t)\}$이다.
이제 두 번째 원소 $(s'_1, \ldots, s'_n) \in F(Z_t)$를 고르자. 위에서 $s'_i$들이
서로 다름을 보였다. 따라서
$g(s_1, \ldots, s_n) = (s'_1, \ldots, s'_n)$를 만족하는 $g \in S_n$이 존재한다.
즉 $G$의 $F(Z_t)$ 위의 작용은 추이적이며 증명이 끝난다.
\end{proof}

\noindent
다음은 핵심 보조정리이다.

\begin{lemma}
\label{pione-lemma-tame}
\begin{reference}
\cite[Definition 7.2.4]{BS}와 비교하라.
\end{reference}
$(\mathcal{C}, F)$를 Galois 범주라 하자. $G = \text{Aut}(F)$는 예
\ref{pione-example-from-C-F-to-G-sets}의 것이라 하자. $\mathcal{C}$의 임의의 연결 대상 $X$에 대해
$G$의 $F(X)$ 위의 작용은 추이적이다.
\end{lemma}

\begin{proof}
이하에서는 보조정리 \ref{pione-lemma-epi-mono}의 결과를 별도의 언급 없이 사용한다.
$I$를 $\mathcal{C}$에서 Galois 대상들의 동형류로 이루어진 집합이라 하자.
각 $i \in I$에 대해 그 동형류의 대표를 $X_i$라 하자.
각 $i \in I$에 대해 $\gamma_i \in F(X_i)$를 고른다. $I$ 위의 부분순서를
$i \geq i'$일 필요충분조건이 사상 $f_{ii'} : X_i \to X_{i'}$의
존재가 되도록 정의한다. 그러한 사상이 주어지면 자기동형
$X_{i'} \to X_{i'}$를 뒤에 합성하여
$F(f_{ii'})(\gamma_i) = \gamma_{i'}$가 되게 할 수 있다. 이 정규화 아래에서
사상 $f_{ii'}$는 유일하다. $I$가 유향 부분순서집합임에 유의하자
(범주론, 정의 \ref{categories-definition-directed-set}). 실제로
$i_1, i_2 \in I$이면 Galois 대상 $Y$와 사상
$Y \to X_{i_1} \times X_{i_2}$가 존재한다. 이는 보조정리 \ref{pione-lemma-galois}를
$X_{i_1} \times X_{i_2}$의 연결 성분에 적용하면 된다. 이때
% P07-ERR-1147: frozen-source index case retained; see sidecar.
어떤 $i \in I$에 대해 $Y \cong X_i$이고, $i \geq i_1$, $i \geq I_2$이다.

\medskip\noindent
함자 $F$는 $X$를
$$
F'(X) = \colim_I \Mor_\mathcal{C}(X_i, X)
$$
로 보내는 함자 $F'$와 동형이라고 주장한다. 이 동형은 다음과 같이 정의되는 함자 변환
$t : F' \to F$로 주어진다. $f : X_i \to X$가 주어지면
$t_X(f) = F(f)(\gamma_i)$로 놓는다. (\ref{pione-item-one-element})를 사용하면
$t_X$는 단사이다. 전사임을 보이기 위해 $\gamma \in F(X)$를 잡자.
Galois 범주의 정의에 의해 곧바로 $X$가 연결인 경우로 환원할 수 있다. 더 나아가 보조정리
\ref{pione-lemma-galois}에 의해 $X$가 Galois라고 가정해도 된다. 이 경우 어떤 $i$에 대해 $X$는
$X_i$와 동형이다. 동형 $X_i \to X$는 $\gamma_i$가 $\gamma$로
보내지도록 고를 수 있다(Galois 대상의 정의에 의한다). 따라서 $t$는 동형이다.

\medskip\noindent
$A_i = \text{Aut}(X_i)$로 놓자. $i \geq i'$일 때 모든 $a \in A_i$에 대해 도식
$$
\xymatrix{
X_i \ar[d]_a \ar[r]_{f_{ii'}} & X_{i'} \ar[d]^{h_{ii'}(a)} \\
X_i \ar[r]^{f_{ii'}} & X_{i'}
}
$$
이 가환하도록 하는 표준적인 사상 $h_{ii'} : A_i \to A_{i'}$가 존재한다고 주장한다.
실제로 $h_{ii'}(a) = a' : X_{i'} \to X_{i'}$를
$F(a')(\gamma_{i'}) = F(f_{ii'} \circ a)(\gamma_i)$를 만족하는 유일한 자기동형으로 잡으면 된다.
앞에서와 같이 이로써 도식은 가환하며, 그 선택도 유일하다. 따라서
$h_{i'i''} \circ h_{ii'} = h_{ii''}$
가 $i \geq i' \geq i''$일 때 성립한다.
$F(X_i) \to F(X_{i'})$가 전사이므로 $A_i \to A_{i'}$도 전사이다.
역극한을 취하여 군
$$
A = \lim_I A_i
$$
를 얻는다. 자기동형군들이 유한하므로 이는 프로유한군이다. 범주론, 보조정리
\ref{categories-lemma-nonempty-limit}에 의해 사상 $A \to A_i$는 모든 $i$에 대해
전사이다.

\medskip\noindent
$A$의 원소들은 역계 $X_i$에 작용하므로 앞합성을 통해 $F'$ 위의 $A$의 오른쪽 작용을 얻는다.
달리 말하면 준동형 $A^{opp} \to G$를 얻는다. $A \to A_i$가 전사이므로
$G$는 모든 $i$에 대해 $F(X_i)$에 추이적으로 작용한다. 모든 연결 대상은
어떤 $X_i$에 의해 지배되므로 보조정리가 따른다.
\end{proof}

\begin{proposition}
\label{pione-proposition-galois}
\begin{reference}
이는 \cite[Expos\'e V]{SGA1}의 약한 형태이다.
증명은 \cite[Theorem 7.2.5]{BS}에서 가져왔다.
\end{reference}
$(\mathcal{C}, F)$를 Galois 범주라 하자. $G = \text{Aut}(F)$는 예
\ref{pione-example-from-C-F-to-G-sets}의 것이라 하자. 함자
$F : \mathcal{C} \to \textit{Finite-}G\textit{-Sets}$
% P07-ERR-1148: frozen-source missing copula recorded in sidecar.
(\ref{pione-equation-remember})는 범주의 동치이다.
\end{proposition}

\begin{proof}
이하에서는 보조정리 \ref{pione-lemma-epi-mono}의 결과를 별도의 언급 없이 사용한다. 특히 함자가 충실함은
이미 알고 있다. 보조정리 \ref{pione-lemma-tame}에 의해 임의의 연결 대상 $X$에 대한
$G$의 $F(X)$ 위의 작용은 추이적이다. 따라서 일반적인 $X$에 대해 함자 $F$는
$X$의 연결 성분들(Galois 범주의 공리에 의해 존재한다)을 $F(X)$의 $G$-궤도들로 보낸다.
$X$와 $Y$를 대상이라 하고, $s : F(X) \to F(Y)$를 사상이라 하자. 그러면 $s$의 그래프
$\Gamma_s \subset F(X) \times F(Y)$는 연결 성분들의 합이다.
따라서 $X \times Y$의 연결 성분들의 합 $Z$가 존재하여 단사 사상 $Z \to X \times Y$를 갖고
$F(Z) = \Gamma_s$를 만족한다. $F(Z) \to F(X)$가 전단사이므로 $Z \to X$는
동형이다. 그러므로 $s = F(f)$이며, 여기서 $f : X \cong Z \to Y$는 합성이다.
따라서 $F$는 충만충실하다.

\medskip\noindent
증명을 끝내기 위해 $F$가 본질적으로 전사임을 보이자. 유한 지수의 임의의
열린 부분군 $H \subset G$에 대해 $G/H$가 본질적 상에 속함을 보이면 충분하다.
$G$ 위상의 정의에 의해 유한 개의 대상 $X_i$가 존재하여
$$
\Ker(G \longrightarrow \prod\nolimits_i \text{Aut}(F(X_i)))
$$
가 $H$에 포함된다. 모든 $i$에 대해 $X_i$가 연결이라고 가정해도 된다.
보조정리 \ref{pione-lemma-galois}를 사용하여 $\prod X_i$의 한 연결 성분으로 사상하는
Galois 대상 $Y$를 고를 수 있다. 어떤 열린 부분군 $U \subset G$에 대해
$G\textit{-sets}$에서 동형 $F(Y) = G/U$를 고르자. $Y$가 Galois이므로 군
$\text{Aut}(Y) = \text{Aut}_{G\textit{-Sets}}(G/U)$는 $F(Y) = G/U$에
추이적으로 작용한다. 이는 $U$가 정규임을 뜻한다. $F(Y)$가 각 $i$에 대해
$F(X_i)$ 위로 전사하므로 $U \subset H$이다.
$M \subset \text{Aut}(Y)$을 다음에 대응하는 유한 부분군이라 하자.
$$
(H/U)^{opp} \subset (G/U)^{opp} = \text{Aut}_{G\textit{-Sets}}(G/U)
= \text{Aut}(Y).
$$
$X = Y/M$으로 놓자. 즉 $X$는 사상 $m : Y \to Y$, $m \in M$들의
쌍대등화자이다. $F$가 완전하므로 $F(X) = G/H$이고 증명이 끝난다.
\end{proof}

\begin{lemma}
\label{pione-lemma-functoriality-galois}
$(\mathcal{C}, F)$와 $(\mathcal{C}', F')$를 Galois 범주라 하자.
$H : \mathcal{C} \to \mathcal{C}'$를 완전 함자라 하자. 동형
$t : F' \circ H \to F$가 존재한다. $t$의 선택은 연속 준동형
$h : G' = \text{Aut}(F') \to \text{Aut}(F) = G$와
$2$-가환 도식
$$
\xymatrix{
\mathcal{C} \ar[r]_H \ar[d] & \mathcal{C}' \ar[d] \\
\textit{Finite-}G\textit{-Sets} \ar[r]^h &
\textit{Finite-}G'\textit{-Sets}
}
$$
을 정한다. 사상 $h$는 $t$의 선택에 의존하지만 $G$의 내부 자기동형을 제외하면
유일하다. 반대로 연속 준동형 $h : G' \to G$가 주어지면 완전 함자
$H : \mathcal{C} \to \mathcal{C}'$와 동형 $t$가 존재하며, 위의 구성으로
$h$를 복원한다.
\end{lemma}

\begin{proof}
명제 \ref{pione-proposition-galois}와 보조정리 \ref{pione-lemma-single-out-profinite}에 의해
$\mathcal{C} = \textit{Finite-}G\textit{-Sets}$이고 $F$가 망각 함자라고 가정해도 되며,
$\mathcal{C}'$에 대해서도 마찬가지이다. 따라서 $t$의 존재는 보조정리
\ref{pione-lemma-second-fundamental-functor}에서 따른다. 사상 $h$는 $t$를 통한 구조의 운반으로 얻는다.
도식의 가환성은 명백하다. $G$의 원소에 의한 내부 공액을 제외하고 $h$가 유일하다는 것은
$G$의 원소를 제외하고 $t$의 선택이 유일하다는 사실에서 따른다. 마지막 주장은 자명하다.
\end{proof}





\section{함자와 준동형}
\label{pione-section-translation}

\noindent
$(\mathcal{C}, F)$, $(\mathcal{C}', F')$, $(\mathcal{C}'', F'')$를 Galois 범주라 하자.
$G = \text{Aut}(F)$, $G' = \text{Aut}(F')$, $G'' = \text{Aut}(F'')$로 놓자.
$H : \mathcal{C} \to \mathcal{C}'$와 $H' : \mathcal{C}' \to \mathcal{C}''$를
완전 함자라 하자. 보조정리 \ref{pione-lemma-functoriality-galois}와 같이 이들에 대응하는
연속 준동형을 각각 $h : G' \to G$와 $h' : G'' \to G'$라 하자. 이 절에서는 대응하는
$2$-가환 도식
\begin{equation}
\label{pione-equation-translation}
\vcenter{
\xymatrix{
\mathcal{C} \ar[r]_H \ar[d] &
\mathcal{C}' \ar[r]_{H'} \ar[d] &
\mathcal{C}'' \ar[d] \\
\textit{Finite-}G\textit{-Sets} \ar[r]^h &
\textit{Finite-}G'\textit{-Sets} \ar[r]^{h'} &
\textit{Finite-}G''\textit{-Sets}
}
}
\end{equation}
을 고찰하고, 열 $1 \to G'' \to G' \to G \to 1$의 완전성에 관한 성질을 함자 $H$와
$H'$의 성질에 연결한다.

\begin{lemma}
\label{pione-lemma-functoriality-galois-surjective}
도식 (\ref{pione-equation-translation})에서 다음 조건들은 동치이다.
\begin{enumerate}
\item $h : G' \to G$는 전사이다.
\item $H : \mathcal{C} \to \mathcal{C}'$는 충만충실하다.
\item $X \in \Ob(\mathcal{C})$가 연결이면 $H(X)$도 연결이다.
\item $X \in \Ob(\mathcal{C})$가 연결이고 $\mathcal{C}'$에 사상
$*' \to H(X)$가 존재하면 사상 $* \to X$가 존재한다.
\item $\mathcal{C}$의 임의의 대상 $X$에 대해 사상
$\Mor_\mathcal{C}(*, X) \to \Mor_{\mathcal{C}'}(*', H(X))$
은 전단사이다.
\end{enumerate}
여기서 $*$와 $*'$는 각각 $\mathcal{C}$와 $\mathcal{C}'$의 종대상이다.
\end{lemma}

\begin{proof}
함의 (5) $\Rightarrow$ (4)와 (2) $\Rightarrow$ (5)는 명백하다.

\medskip\noindent
(3)을 가정하자. $X$를 $\mathcal{C}$의 연결 대상으로 하고, $*' \to H(X)$를 사상이라 하자.
(3)에 의해 $H(X)$는 연결이므로 $*' \to H(X)$는 동형이다. 따라서
$H(X)$에 대응하는 $G'$-집합은 원소를 정확히 하나 가지며, 이는
$X$에 대응하는 $G$-집합이 원소를 하나 가짐을 뜻한다. 그러므로 $X$는 $\mathcal{C}$의 종대상과
동형이고, 특히 사상 $* \to X$가 존재한다. 이로써 (3) $\Rightarrow$ (4)이다.

\medskip\noindent
(1)이 성립하면 함자
$\textit{Finite-}G\textit{-Sets} \to \textit{Finite-}G'\textit{-Sets}$
는 충실충만이다. 실제로 이 경우 $G$-집합의 사상이 $G$의 작용과 가환하는 것과
$G'$의 작용과 가환하는 것은 동치이다. 따라서 (1) $\Rightarrow$ (2)이다.

\medskip\noindent
(1)이 성립하면 $G$-집합 $X$에서 $G$-궤도와 $G'$-궤도는 일치한다.
따라서 (1) $\Rightarrow$ (3)이다.

\medskip\noindent
증명을 끝내려면 (4)가 (1)을 함의함을 보이면 충분하다. (1)이 거짓이라고,
즉 $h$가 전사가 아니라고 하자. 그러면 $h(G')$를 포함하지만 $G$와 같지 않은
열린 부분군 $U \subset G$가 존재한다. 유한 $G$-집합 $M = G/U$에 대한
작용은 추이적이지만 $G'$는 고정점을 갖는다. $M$에 대응하는 $\mathcal{C}$의
대상 $X$는 (3)에 모순된다. 이로써 (3) $\Rightarrow$ (1)이고, 증명이 끝난다.
\end{proof}

\begin{lemma}
\label{pione-lemma-composition-trivial}
도식 (\ref{pione-equation-translation})에서 다음 조건들은 동치이다.
\begin{enumerate}
\item $h \circ h'$는 자명하다.
\item $H' \circ H$의 상은 종대상의 유한 쌍대곱과 동형인 대상들로 이루어진다.
\end{enumerate}
\end{lemma}

\begin{proof}
$H$와 $H'$를 $h$와 $h'$가 정하는 표준 함자
$\textit{Finite-}G\textit{-Sets} \to \textit{Finite-}G'\textit{-Sets}
\to \textit{Finite-}G''\textit{-Sets}$로 바꾸어도 된다. 그러면 주장은
모든 $G$-집합에 대한 $G''$의 작용이 자명한 것과 준동형 $G'' \to G$가
자명한 것이 동치라는 명제이며, 이는 명백하다.
\end{proof}

\begin{lemma}
\label{pione-lemma-functoriality-galois-ses}
도식 (\ref{pione-equation-translation})에서 다음 조건들은 동치이다.
\begin{enumerate}
\item 열 $G'' \xrightarrow{h'} G' \xrightarrow{h} G \to 1$은 다음 의미에서 완전하다.
$h$는 전사이고, $h \circ h'$는 자명하며, $\Ker(h)$는 $\Im(h')$를 포함하는 가장 작은 닫힌 정규 부분군이다.
\item $H$는 충실충만이고, $\mathcal{C}'$의 대상 $X'$가 $H$의 본질적 상에 속할 필요충분조건은
$H'(X')$가 종대상의 유한 쌍대곱과 동형인 것이다.
% P07-ERR-1149: frozen-source composite order retained; see sidecar.
\item $H$는 충실충만이고, $H \circ H'$는 모든 대상을 종대상의 유한 쌍대곱으로 보낸다.
더 나아가 $\mathcal{C}'$의 대상 $X'$가 $H'(X')$가 종대상의 유한 쌍대곱이라는 조건을
만족하면, $\mathcal{C}$의 대상 $X$와 전사 $H(X) \to X'$가 존재한다.
\end{enumerate}
\end{lemma}

\begin{proof}
보조정리 \ref{pione-lemma-functoriality-galois-surjective}와
\ref{pione-lemma-composition-trivial}에 의해 $H$는 충실충만이고, $h$는 전사이며,
$H' \circ H$는 대상을 종대상의 서로소 합으로 보내고, $h \circ h'$는 자명하다고 해도 된다.
$N \subset G'$를 $h'$의 상을 포함하는 가장 작은 닫힌 정규 부분군이라 하자. 명백히
$N \subset \Ker(h)$이다. 함자 $H$와 $H'$는 $h$와 $h'$가 정하는 표준 함자
$\textit{Finite-}G\textit{-Sets} \to \textit{Finite-}G'\textit{-Sets}
\to \textit{Finite-}G''\textit{-Sets}$라고 해도 된다.

\medskip\noindent
(2)를 가정하자. 이는 유한 $G'$-집합 $X'$에 $G''$가 자명하게 작용하면
$G'$의 작용이 $G$를 통해 분해됨을 뜻한다. 이를
$X' = G'/U'N$에 적용하자. 여기서 $U'$는 충분히 작은 $G'$의 열린 부분군이다.
그러면 모든 $U'$에 대해 $\Ker(h) \subset U'N$이다. $N$은 닫혀 있으므로
$\Ker(h) \subset N$이 따른다. 즉 (1)이 성립한다.

\medskip\noindent
(1)을 가정하자. 이는 $N = \Ker(h)$임을 뜻한다. $X'$를 유한 $G'$-집합으로 하고,
$G''$가 자명하게 작용한다고 하자. 이는 $\Ker(G' \to \text{Aut}(X'))$가 $\Im(h')$를 포함하는
닫힌 정규 부분군임을 뜻한다. 따라서 $N = \Ker(h)$는 이 부분군에 포함되고,
$X'$에 대한 $G'$의 작용은 $G$를 통해 분해된다. 즉 (2)가 성립한다.

\medskip\noindent
(3)을 가정하자. 이는 유한 $G'$-집합 $X'$에 $G''$가 자명하게 작용하면,
$X$가 $G$-집합인 $G'$-집합의 전사 $X \to X'$가 존재함을 뜻한다. 명백히 이는
$X'$에 대한 $G'$의 작용이 $G$를 통해 분해됨을 뜻한다. 즉 (2)가 성립한다.

\medskip\noindent
함의 (2) $\Rightarrow$ (3)은 곧바로 따른다. 이것으로 증명이 끝난다.
\end{proof}

\begin{lemma}
\label{pione-lemma-functoriality-galois-injective}
도식 (\ref{pione-equation-translation})에서 다음 조건들은 동치이다.
\begin{enumerate}
\item $h'$는 단사이다.
\item $\mathcal{C}''$의 임의의 연결 대상 $X''$에 대하여,
$\mathcal{C}'$의 대상 $X'$와 도식
% P07-ERR-0644: frozen-source functor symbol retained; see sidecar.
$$
X'' \leftarrow Y'' \rightarrow H(X')
$$
이 $\mathcal{C}''$에 존재하여, $Y'' \to X''$는 전사이고 $Y'' \to H(X')$는 단사이다.
\end{enumerate}
\end{lemma}

\begin{proof}
$H'$를 유한 $G'$-집합의 범주에서 유한 $G''$-집합의 범주로 가는 대응하는 함자로
바꾸어도 된다.

\medskip\noindent
$h' : G'' \to G'$가 단사라고 가정하자. $H'' \subset G''$를 열린 부분군이라 하자.
$G''$의 위상은 $G'$에서 유도된 위상이므로, $(h')^{-1}(H') \subset H''$를 만족하는
열린 부분군 $H' \subset G'$가 존재한다. 구하는 도식은
$$
G''/H'' \leftarrow G''/(h')^{-1}(H') \rightarrow G'/H'
$$
이다. 역으로 함자
$\textit{Finite-}G'\textit{-Sets} \to \textit{Finite-}G''\textit{-Sets}$.
에 대해 (2)가 성립한다고 가정하자. $g'' \in \Ker(h')$라 하자. 임의의 열린 부분군
$H'' \subset G''$를 택한다. 가정에 의해 유한 $G'$-집합 $X'$와 도식
$$
G''/H'' \leftarrow Y'' \rightarrow X'
$$
이 $G''$-집합의 범주에 존재하며, 왼쪽 화살표는 전사이고 오른쪽 화살표는 단사이다.
$g''$는 $h'$의 핵에 속하므로 $g''$는 $X'$에 자명하게 작용한다.
따라서 $g''$는 $Y''$에 자명하게 작용하고, 나아가 $G''/H''$에도 자명하게 작용한다.
그러므로 $g'' \in H''$이다. 이는 모든 열린 부분군에 대해 성립하므로
$g''$는 원하는 대로 항등원이다.
\end{proof}

\begin{lemma}
\label{pione-lemma-functoriality-galois-normal}
도식 (\ref{pione-equation-translation})에서 다음 조건들은 동치이다.
\begin{enumerate}
\item $h'$의 상은 정규이다.
\item $\mathcal{C}'$의 임의의 연결 대상 $X'$에 대하여, $\mathcal{C}''$의 종대상에서
$H'(X')$로 가는 사상이 존재하면 $H'(X')$는 종대상의 유한 쌍대곱과 동형이다.
\end{enumerate}
\end{lemma}

\begin{proof}
이는 연속 군 준동형 $h' : G'' \to G'$에 관한 다음 주장으로 옮겨진다.
$h'$의 상이 정규일 필요충분조건은 $h'(G'')$를 포함하는 임의의 열린 부분군 $U' \subset G'$가
$h'(G'')$의 모든 켤레도 포함하는 것이다. 이로부터 결과는 쉽게 따른다. 세부 일부는 생략한다.
\end{proof}






\section{유한 에탈 사상}
\label{pione-section-finite-etale}

\noindent
이 절에서는 에탈 기본군을 구성하는 데 충분한 유한 에탈 사상에 관한
기본 결과들을 증명한다.

\medskip\noindent
$X$를 스킴이라 하자. 기호 $\textit{F\'Et}_X$는 $X$ 위의 유한 에탈
스킴의 범주를 나타낸다. 따라서
\begin{enumerate}
\item $\textit{F\'Et}_X$의 대상은 공역이 $X$인 유한 에탈 사상 $Y \to X$이다.
\item $\textit{F\'Et}_X$에서 $Y \to X$로부터 $Y' \to X$로 가는 사상은
도식
$$
\xymatrix{
Y \ar[rr] \ar[rd] & &  Y' \ar[ld] \\
& X
}
$$
을 가환하게 하는 사상 $Y \to Y'$이다.
\end{enumerate}
$\textit{F\'Et}_X$의 대상을 흔히 $X$의 {\it 유한 에탈 덮개}라고 부른다
($Y$가 공집합인 경우도 포함한다). 스킴의 범주 위에는 스택
$p : \textit{F\'Et} \to \Sch$가 존재하며, 그 $X$ 위의 올은 방금 정의한 범주
$\textit{F\'Et}_X$이다. Examples of Stacks, 절
\ref{examples-stacks-section-finite-etale}를 보라.

\begin{example}
\label{pione-example-finite-etale-geometric-point}
$k$를 대수적으로 닫힌 체라 하고 $X = \Spec(k)$라 하자. 이 경우 $\textit{F\'Et}_X$는
유한 집합의 범주와 동치이다. 더 일반적으로 $k$가 분리적으로 대수적으로 닫힌 경우에도 성립한다.
이는 $k$ 위의 에탈 스킴이 $k$의 유한 분리 확대체들의 스펙트럼의
서로소 합이기 때문이다. Morphisms, 보조정리 \ref{morphisms-lemma-etale-over-field}를 보라.
\end{example}

\begin{lemma}
\label{pione-lemma-finite-etale-covers-limits-colimits}
$X$를 스킴이라 하자. 범주 $\textit{F\'Et}_X$는 유한 극한과 유한 여극한을 가지며,
임의의 사상 $X' \to X$에 대한 기저변환 함자
$\textit{F\'Et}_X \to \textit{F\'Et}_{X'}$는 완전하다.
\end{lemma}

\begin{proof}
유한 극한과 왼쪽 완전성. Categories, 보조정리 \ref{categories-lemma-finite-limits-exist}에 의해
$\textit{F\'Et}_X$가 종대상과 올곱을 가짐을 보이면 충분하다. 이는
$X$ 위의 모든 스킴의 범주가 종대상(즉 $X$)과 올곱을 갖는다는 사실에서
명백하다. 또한 $X$ 위의 유한 에탈 스킴들의 올곱도
$X$ 위에서 유한 에탈이다. 기저변환이 이 연산들과 가환함도 명백하므로
기저변환은 왼쪽 완전하다(Categories, 보조정리
\ref{categories-lemma-characterize-left-exact}).

\medskip\noindent
유한 여극한과 오른쪽 완전성. Categories, 보조정리 \ref{categories-lemma-colimits-exist}에 의해
$\textit{F\'Et}_X$가 유한 쌍대곱과 쌍대등화자를 가짐을 보이면 충분하다. 유한 쌍대곱은
서로소 합으로 주어진다(빈 쌍대곱은 빈 스킴이다). $a, b : Z \to Y$를
$\textit{F\'Et}_X$의 두 사상이라 하자. $Z \to X$와 $Y \to X$가 유한 에탈이므로
$Z = \underline{\Spec}(\mathcal{C})$ 및 $Y = \underline{\Spec}(\mathcal{B})$로 쓸 수 있다.
여기서 $\mathcal{C}$와 $\mathcal{B}$는 각각 $\mathcal{O}_X$ 위의 유한 국소 자유 대수이다.
사상 $a, b$는 두 사상 $a^\sharp, b^\sharp : \mathcal{B} \to \mathcal{C}$를 유도한다.
$\mathcal{A} = \text{Eq}(a^\sharp, b^\sharp)$를 그 등화자라 하자. 만일
$$
\underline{\Spec}(\mathcal{A}) \longrightarrow X
$$
가 유한 에탈이면, 이것이 쌍대등화자임은 명백하다(실제로
$\textit{F\'Et}_X$의 임의의 대상은 $\mathcal{O}_X$-대수의 층의 상대 스펙트럼으로 쓸 수 있다).
이 확인은 $X$를 에탈 덮개의 한 원소로 바꾼 뒤 수행해도 된다
(Descent, 보조정리 \ref{descent-lemma-descending-property-finite} 및
\ref{descent-lemma-descending-property-etale}). 따라서 \'Etale Morphisms, 보조정리
\ref{etale-lemma-finite-etale-etale-local}에 의해
$Y = \coprod_{i = 1, \ldots, n} X$ 및 $Z = \coprod_{j = 1, \ldots, m} X$
라고 해도 된다. 이때
$$
\mathcal{C} = \prod\nolimits_{1 \leq j \leq m} \mathcal{O}_X
\quad\text{이고}\quad
\mathcal{B} = \prod\nolimits_{1 \leq i \leq n} \mathcal{O}_X
$$
이다. 열린 덮개의 한 원소로 다시 바꾸면 $a, b$가 사상
$a_s, b_s : \{1, \ldots, m\} \to \{1, \ldots, n\}$에 대응한다고 해도 된다.
즉 $Z$에서 지표 $j$에 대응하는 성분 $X$는 사상 $a$, $b$에 의해
$Y$에서 각각 지표 $a_s(j)$, $b_s(j)$에 대응하는 성분 $X$로 보내진다.
$\{1, \ldots, n\} \to T$를 $a_s, b_s$의 쌍대등화자라 하자. 그러면
$$
\mathcal{A} = \prod\nolimits_{t \in T} \mathcal{O}_X
$$
이고, 그 스펙트럼은 명백히 $X$ 위에서 유한 에탈이다. 이것이 기저변환과
양립한다는 확인은 생략한다. 따라서 기저변환은 오른쪽 완전 함자이다.
\end{proof}

\begin{remark}
\label{pione-remark-colimits-commute-forgetful}
$X$를 스킴이라 하자. 자연스러운 함자
$F_1 : \textit{F\'Et}_X \to \Sch$와 $F_2 : \textit{F\'Et}_X \to \Sch/X$
를 생각하자. 그러면
\begin{enumerate}
\item 함자 $F_1$과 $F_2$는 유한 여극한과 가환한다.
\item 함자 $F_2$는 유한 극한과 가환한다.
\item 함자 $F_1$은 연결된 유한 극한, 즉 등화자와 올곱과 가환한다.
\end{enumerate}
극한에 관한 결과는 보조정리 \ref{pione-lemma-finite-etale-covers-limits-colimits}의 증명에서 한 논의와
Categories, 보조정리 \ref{categories-lemma-connected-limit-over-X}에서 곧바로 따른다.
$F_1$과 $F_2$가 유한 쌍대곱과 가환함은 명백하다. Categories, 보조정리
\ref{categories-lemma-characterize-left-exact}의 쌍대에 의해, $F_1$과 $F_2$가
쌍대등화자와 가환함을 보이면 충분하다. 보조정리
\ref{pione-lemma-finite-etale-covers-limits-colimits}의 증명에서 $\textit{F\'Et}_X$의
쌍대등화자는 에탈 국소적으로 다음 꼴임을 보았다.
$$
\xymatrix{
\coprod_{j \in J} U \ar@<1ex>[r]^a \ar@<-1ex>[r]_b &
\coprod_{i \in I} U \ar[r] &
\coprod_{t \in \text{Coeq}(a, b)} U
}
$$
이는 명백히 스킴의 범주에서 쌍대등화자이다. 따라서 주장은 쌍대등화자라는 성질이
fpqc 국소적이라는 사실에서 따른다. 정확한 정식화는
Descent, 보조정리 \ref{descent-lemma-coequalizer-fpqc-local}에 있다.
\end{remark}

\begin{lemma}
\label{pione-lemma-internal-hom-finite-etale}
$X$를 스킴이라 하자. $U, V$가 $X$ 위에서 유한 에탈이면,
$X$ 위에서 유한 에탈인 스킴 $W$가 존재하여
$$
\Mor_X(X, W) = \Mor_X(U, V)
$$
를 만족하고, 이 등식은 임의의 기저변환 뒤에도 성립한다.
\end{lemma}

\begin{proof}
More on Morphisms, 보조정리
\ref{more-morphisms-lemma-hom-from-finite-locally-free-separated-lqf}에 의해
$\mathit{Mor}_X(U, V)$를 표현하는 스킴 $W$가 존재한다(여기서 에탈 사상은 국소 준유한이라는
Morphisms, 보조정리 \ref{morphisms-lemma-etale-locally-quasi-finite}와
유한 사상은 분리적이라는 사실을 사용한다). 이 스킴이 임의의 기저변환 뒤에도 등식을
만족함은 명백하다. 증명을 끝내려면 $W \to X$가 유한 에탈임을 보여야 한다.
이는 $X$를 에탈 덮개의 한 원소로 바꾼 뒤 보여도 된다
(Descent, 보조정리 \ref{descent-lemma-descending-property-finite} 및
\ref{descent-lemma-descending-property-separated}). 따라서 \'Etale Morphisms, 보조정리
\ref{etale-lemma-finite-etale-etale-local}에 의해
$U = \coprod_{i = 1, \ldots, n} X$ 및 $V = \coprod_{j = 1, \ldots, m} X$
라고 해도 된다. 이 경우 직접 확인하면
$W = \coprod_{\alpha : \{1, \ldots, n\} \to \{1, \ldots, m\}} X$
이다(세부는 생략한다). 이것으로 증명이 끝난다.
\end{proof}

\noindent
$X$를 스킴이라 하자. $X$의 {\it 기하점}은 $k$가 대수적으로 닫힌 체인 사상
$\Spec(k) \to X$이다. 이러한 점은 보통 윗줄을 붙인 소문자
$\overline{x}$로 나타낸다. $\overline{x}$는 사상뿐 아니라 스킴 $\Spec(k)$도 나타내며,
$\kappa(\overline{x})$로 $k$를 나타낸다. $\overline{x}$가 $x$의 {\it 위에 놓인다}는 것은
$x \in X$가 $\overline{x}$의 상이라는 뜻이다. 이에 대해서는
\'Etale Cohomology, 절 \ref{etale-cohomology-section-stalks}에서 더 논의한다.
$\overline{x}$와 에탈 사상 $U \to X$가 주어지면 다음을 생각할 수 있다.
$$
|U_{\overline{x}}| : \text{스킴의 점들로 이루어진 바탕 집합 }U_{\overline{x}} = U \times_X \overline{x}
$$
$U_{\overline{x}}$는 $\overline{x}$ 위의 스킴으로서 $\overline{x}$의 복사본들의
서로소 합이므로(Morphisms, 보조정리 \ref{morphisms-lemma-etale-over-field}),
이 집합은 다음과 같이도 기술할 수 있다.
$$
|U_{\overline{x}}| =
\left\{
\begin{matrix}
\text{가환} \\
\text{도식}
\end{matrix}
\vcenter{
\xymatrix{
\overline{x} \ar[rd]_{\overline{x}} \ar[r]_{\overline{u}} & U \ar[d] \\
& X
}
}
\right\}
$$
대응 $U \mapsto |U_{\overline{x}}|$는 함자이며, 흔히
$F_{\overline{x}}$로 쓴다.

\begin{lemma}
\label{pione-lemma-finite-etale-connected-galois-category}
$X$를 연결 스킴이라 하고 $\overline{x}$를 기하점이라 하자. 함자
$$
F_{\overline{x}} : \textit{F\'Et}_X \longrightarrow \textit{Sets},\quad
Y \longmapsto |Y_{\overline{x}}|
$$
는 Galois 범주를 정한다(정의 \ref{pione-definition-galois-category}).
\end{lemma}

\begin{proof}
$\textit{F\'Et}_{\overline{x}}$를 유한 집합의 범주와 동일시하면
(예 \ref{pione-example-finite-etale-geometric-point}), 함자 $F_{\overline{x}}$는 사상
$\overline{x} \to X$에 대한 기저변환 함자에 지나지 않는다. 따라서
$\textit{F\'Et}_X$는 유한 극한과 유한 여극한을 가지며, 보조정리
\ref{pione-lemma-finite-etale-covers-limits-colimits}에 의해 $F_{\overline{x}}$는 완전하다.
또한 $\textit{F\'Et}_X$의 유한 극한은 $X$ 위의 스킴의 범주에서 대응하는
유한 극한과 일치한다는 사실도 사용한다. 주의
\ref{pione-remark-colimits-commute-forgetful}를 보라.

\medskip\noindent
$Y' \to Y$가 $\textit{F\'Et}_X$의 단사사상이면
$Y' \to Y' \times_Y Y'$는 동형이고, 따라서 $Y' \to Y$는 스킴의 단사사상이다.
그러므로 $Y' \to Y$는 열린 몰입이다(\'Etale Morphisms, 정리
\ref{etale-theorem-etale-radicial-open}). $Y'$는 $X$ 위에서 유한이고 $Y$는 $X$ 위에서 분리적이므로,
사상 $Y' \to Y$는 유한이며(Morphisms, 보조정리 \ref{morphisms-lemma-finite-permanence}),
따라서 닫혀 있다(Morphisms, 보조정리 \ref{morphisms-lemma-finite-proper}). 그러므로 이는
$Y$의 열린닫힌 부분스킴의 포함이다. 따라서 $Y$가 범주 $\textit{F\'Et}_X$의
연결 대상인 것(정의 \ref{pione-definition-galois-category}의 의미에서)과 $Y$가 스킴으로서
연결인 것은 동치이다. 이때 Topology, 보조정리
\ref{topology-lemma-finite-fibre-connected-components}에 의해 $Y$는 스킴으로서도,
정의 \ref{pione-definition-galois-category}의 의미에서도 그 연결 성분들의 유한 쌍대곱이다.

\medskip\noindent
$Y \to Z$를 $\textit{F\'Et}_X$의 사상으로서 전단사
$F_{\overline{x}}(Y) \to F_{\overline{x}}(Z)$를 유도한다고 하자.
$Y \to Z$가 동형임을 보여야 한다. 위의 논의에 의해 $Z$가 연결이라고 해도 된다.
$Y \to Z$는 유한 에탈, 따라서 유한 국소 자유이므로, $Y \to Z$가 차수 $1$인
유한 국소 자유 사상임을 보이면 충분하다. 이는 $Z$에서
$\overline{x}$ 위에 놓인 임의의 점의 근방에서 성립한다. $Z$는 연결이고 차수는 국소 상수이므로 결론이 따른다.
\end{proof}



\section{기본군}
\label{pione-section-fundamental-groups}

\noindent
이 절에서는 Grothendieck의 대수적 기본군을 정의한다. 다음 정의가 의미를 갖는 것은
보조정리 \ref{pione-lemma-finite-etale-connected-galois-category}에 의한다.

\begin{definition}
\label{pione-definition-fundamental-group}
$X$를 연결 스킴이라 하고 $\overline{x}$를 $X$의 기하점이라 하자.
기점 $\overline{x}$를 갖는 $X$의 {\it 기본군}은 군
$$
\pi_1(X, \overline{x}) = \text{Aut}(F_{\overline{x}})
$$
즉 올 함자
$F_{\overline{x}} : \textit{F\'Et}_X \to \textit{Sets}$
의 자기동형군에 보조정리 \ref{pione-lemma-aut-inverse-limit}의 표준 프로유한 위상을 준 것이다.
\end{definition}

\noindent
위의 결과와 절 \ref{pione-section-galois}의 내용을 합치면 다음 정리를 얻는다.

\begin{theorem}
\label{pione-theorem-fundamental-group}
$X$를 연결 스킴이라 하고 $\overline{x}$를 $X$의 기하점이라 하자.
\begin{enumerate}
\item 올 함자 $F_{\overline{x}}$는 범주의 동치
$$
\textit{F\'Et}_X \longrightarrow
\textit{Finite-}\pi_1(X, \overline{x})\textit{-Sets}
$$
를 정한다.
\item 두 번째 기하점 $\overline{x}'$가 $X$에 주어지면 동형
$t : F_{\overline{x}} \to F_{\overline{x}'}$가 존재한다. 이는 (1)의 범주 동치와 양립하는 동형
$\pi_1(X, \overline{x}) \to \pi_1(X, \overline{x}')$을 준다. 이 동형은 내부 켤레를 제외하면
$t$에 의존하지 않는다.
\item 연결 스킴의 사상 $f : X \to Y$가 주어졌을 때
$\overline{y} = f \circ \overline{x}$라 쓰자. 표준 연속 준동형
$$
f_* : \pi_1(X, \overline{x}) \to \pi_1(Y, \overline{y})
$$
이 존재하고, 도식
$$
\xymatrix{
\textit{F\'Et}_Y \ar[r]_{\text{기저변환}} \ar[d]_{F_{\overline{y}}} &
\textit{F\'Et}_X \ar[d]^{F_{\overline{x}}} \\
\textit{Finite-}\pi_1(Y, \overline{y})\textit{-Sets} \ar[r]^{f_*} &
\textit{Finite-}\pi_1(X, \overline{x})\textit{-Sets}
}
$$
은 가환한다.
\end{enumerate}
\end{theorem}

\begin{proof}
(1)은 보조정리 \ref{pione-lemma-finite-etale-connected-galois-category}와 명제
\ref{pione-proposition-galois}에서 따른다. (2)는 보조정리 \ref{pione-lemma-functoriality-galois}의 특별한 경우이다.
(3)에 대해 도식
$$
\xymatrix{
\textit{F\'Et}_Y \ar[r] \ar[d]_{F_{\overline{y}}} &
\textit{F\'Et}_X \ar[d]^{F_{\overline{x}}} \\
\textit{Sets} \ar@{=}[r] & \textit{Sets}
}
$$
은 가환한다(단지 $2$-가환인 것이 아니라 실제로 가환한다). 이는
$\overline{y} = f \circ \overline{x}$이기 때문이다. 따라서 암묵적인 함자 사이의 변환과 함께
보조정리 \ref{pione-lemma-functoriality-galois}를 적용하면 (3)을 얻는다.
\end{proof}

\begin{lemma}
\label{pione-lemma-fundamental-group-Galois-group}
$K$를 체라 하고 $X = \Spec(K)$라 하자. $\overline{K}$를 대수적 폐포라 하고, 대응하는 기하점을
$\overline{x} : \Spec(\overline{K}) \to X$라 쓰자. $K^{sep} \subset \overline{K}$를
분리 폐포라 하자.
\begin{enumerate}
\item 보조정리 \ref{pione-lemma-sheaves-point}의 함자는 범주의 동치
$$
\textit{F\'Et}_X \longrightarrow
\textit{Finite-}\text{Gal}(K^{sep}/K)\textit{-Sets}.
$$
를 유도하며, $F_{\overline{x}}$ 및 함자
$\textit{Finite-}\text{Gal}(K^{sep}/K)\textit{-Sets} \to \textit{Sets}$와 양립한다.
\item 이는 프로유한 위상군 사이의 표준 동형
$$
\text{Gal}(K^{sep}/K) \longrightarrow \pi_1(X, \overline{x})
$$
을 유도한다.
\end{enumerate}
\end{lemma}

\begin{proof}
보조정리 \ref{pione-lemma-sheaves-point}의 함자는 함자 $F_{\overline{x}}$와 같다. 실제로
$X$ 위에서 에탈인 임의의 $Y$에 대하여
$$
\Mor_X(\Spec(\overline{K}), Y) = \Mor_X(\Spec(K^{sep}), Y)
$$
이다. 즉 보조정리 \ref{pione-lemma-sheaves-point}의 증명에서 보았듯이
$Y = \coprod_{i \in I} \Spec(L_i)$이고, $L_i/K$는 $K$의 유한 분리 확대이다.
따라서 임의의 $K$-대수 준동형 $L_i \to \overline{K}$는 $K^{sep}$를 통해 분해된다.
또한 $F_{\overline{x}}(Y)$가 유한인 것, $I$가 유한인 것, 그리고
$Y \to X$가 유한 에탈인 것은 서로 동치이다. 이것으로 (1)이 증명되었다.

\medskip\noindent
(2)는 (1), 보조정리 \ref{pione-lemma-functoriality-galois}, 그리고 보조정리
\ref{pione-lemma-single-out-profinite}의 형식적 귀결이다(아래의 주의도 보라).
\end{proof}

\begin{remark}
\label{pione-remark-variance}
보조정리 \ref{pione-lemma-fundamental-group-Galois-group}의 상황에서 동형
$\text{Gal}(K^{sep}/K) \to
\pi_1(X, \overline{x}) = \text{Aut}(F_{\overline{x}})$.
을 더 명시적으로 구성하자. 등식
$\text{Gal}(K^{sep}/K) = \text{Aut}(\overline{K}/K)$
이 성립한다. 실제로 $\overline{K}$는 $K^{sep}$의 완전 폐포이다.
$F_{\overline{x}}(Y) = \Mor_X(\Spec(\overline{K}), Y)$이므로 사상
$$
\text{Aut}(\overline{K}/K) \times F_{\overline{x}}(Y) \to F_{\overline{x}}(Y),
\quad
(\sigma, \overline{y}) \mapsto
\sigma \cdot \overline{y} = \overline{y} \circ \Spec(\sigma)
$$
을 생각할 수 있다. 이는 작용이다. 실제로
$$
\sigma\tau \cdot \overline{y} =
\overline{y} \circ \Spec(\sigma\tau) =
\overline{y} \circ \Spec(\tau) \circ \Spec(\sigma) =
\sigma \cdot (\tau \cdot \overline{y})
$$
이다. 이 작용은 $Y \in \textit{F\'Et}_X$에 관해 함자적이며, 원하는 사상을 얻는다.
\end{remark}





\section{연결 스킴의 Galois 덮개}
\label{pione-section-finite-etale-under-galois}

\noindent
$X$를 기하점 $\overline{x}$를 갖는 연결 스킴이라 하자.
$F_{\overline{x}} : \textit{F\'Et}_X \to \textit{Sets}$는 Galois 범주이므로
(보조정리 \ref{pione-lemma-finite-etale-connected-galois-category}), 절 \ref{pione-section-galois}의
내용을 적용할 수 있다. 이 절에서는 그 용어와 결과의 일부를 스킴과 유한 에탈 사상의
상황으로 명시적으로 옮긴다.

\medskip\noindent
유한 에탈 사상 $Y \to X$가 {\it Galois 덮개}라는 것은 $Y$가
$\textit{F\'Et}_X$의 Galois 대상을 정한다는 뜻이다. 유한 에탈 사상 $Y \to X$에 대해
$G = \text{Aut}_X(Y)$라 하면 다음 조건들은 동치이다.
\begin{enumerate}
\item $Y$는 $X$의 Galois 덮개이다.
\item $Y$는 연결이고, $|G|$는 $Y \to X$의 차수와 같다.
\item $Y$는 연결이고, $G$는 $F_{\overline{x}}(Y)$에 추이적으로 작용한다.
\item $Y$는 연결이고, $G$는 $F_{\overline{x}}(Y)$에 단순 추이적으로 작용한다.
\end{enumerate}
이는 절 \ref{pione-section-galois}의 논의에서 곧바로 따른다.

\medskip\noindent
$Y$가 연결인 임의의 유한 에탈 사상 $f : Y \to X$에 대하여, $Y$를 지배하는
유한 에탈 Galois 덮개 $Y' \to X$가 존재한다(보조정리 \ref{pione-lemma-galois}).

\medskip\noindent
$\textit{F\'Et}_X$의 Galois 대상들은 동치
$$
F_{\overline{x}} : \textit{F\'Et}_X \to
\textit{Finite-}\pi_1(X, \overline{x})\textit{-Sets}
$$
(정리 \ref{pione-theorem-fundamental-group})를 통해 $G = \pi_1(X, \overline{x})/H$ 꼴의 유한
$\pi_1(X, \overline{x})\textit{-Sets}$에 대응한다. 여기서 $H$는 열린 정규 부분군이다.
동치인 표현으로, $G$가 유한군이고 $\pi_1(X, \overline{x}) \to G$가 연속 전사이면,
$G$를 $\pi_1(X, \overline{x})$-집합으로 본 것이 Galois 덮개에 대응한다.

\medskip\noindent
$Y_i \to X$, $i = 1, 2$가 Galois 군 $G_i$를 갖는 유한 에탈 Galois 덮개이면,
그 Galois 군이 $G_1 \times G_2$의 부분군인 유한 에탈 Galois 덮개
$Y \to X$가 존재한다. 실제로 대응하는 연속 준동형
$\pi_1(X, \overline{x}) \to G_i$들을 택하고, $G$를 유도된 연속 준동형
$\pi_1(X, \overline{x}) \to G_1 \times G_2$의 상으로 두면 된다.









\section{기본군의 위상적 불변성}
\label{pione-section-topological-invariance}

\noindent
이 절의 주결과는 연결 스킴의 보편 위상동형사상이 기본군의 동형을 유도한다는 것이다.
명제 \ref{pione-proposition-universal-homeomorphism}을 보라.

\medskip\noindent
두 스킴의 기본군이 같음을 직접 증명하는 대신, 그 유한 에탈 덮개의
범주가 같음을 증명하는 경우가 많다. 두 스킴이 연결이면 물론 이로부터 기본군이
같다는 결론이 따른다.

\begin{lemma}
\label{pione-lemma-what-equivalence-gives}
$f : X \to Y$를 준콤팩트 준분리 스킴의 사상이라 하고, 기저변환 함자
$\textit{F\'Et}_Y \to \textit{F\'Et}_X$가 범주의 동치라고 하자. 그러면
\begin{enumerate}
\item $f$는 위상동형사상 $\pi_0(X) \to \pi_0(Y)$를 유도한다.
\item $X$, 또는 동치인 조건으로 $Y$가 연결이면
$\pi_1(X, \overline{x}) = \pi_1(Y, \overline{y})$이다.
\end{enumerate}
\end{lemma}

\begin{proof}
$Y = Y_0 \amalg Y_1$을 공집합이 아닌 열린닫힌 부분스킴들로의 분해라 하자.
$f(X)$가 두 $Y_i$ 모두와 만난다고 주장한다. 그렇지 않고, 예를 들어
$f(X) \subset Y_1$이면 유한 에탈 사상 $V = Y_1 \to Y$를 생각할 수 있다.
이는 동형이 아니지만 $V \times_Y X \to X$는 동형이 되어 모순이다.

\medskip\noindent
$X = X_0 \amalg X_1$을 열린닫힌 부분스킴들로의 분해라 하자. 유한 에탈 사상
$U = X_1 \to X$를 생각하자. 그러면 $U = X \times_Y V$이며, 여기서 $V \to Y$는
어떤 유한 에탈 사상이다. 사상 $V \to Y$의 차수는 국소 상수이므로, 열린닫힌 부분스킴들로의 분해
$Y = \coprod_{d \geq 0} Y_d$로서 $V \to Y$가 $Y_d$ 위에서 차수 $d$를 갖는 것을 얻는다.
$d > 1$이면 $f^{-1}(Y_d) = \emptyset$이므로, 위의 논의에 의해
$d > 1$이면 $Y_d = \emptyset$이다. 따라서
$i = 0, 1$에 대하여 $f^{-1}(Y_i) = X_i$이다.

\medskip\noindent
그러므로 $f^{-1}$은 $Y$의 열린닫힌 부분집합 전체와 $X$의 열린닫힌 부분집합 전체 사이의
전단사를 유도한다. $X$와 $Y$가 스펙트럼 공간임에 주의하자
(Properties, 보조정리 \ref{properties-lemma-quasi-compact-quasi-separated-spectral}).
Topology, 보조정리 \ref{topology-lemma-connected-component-intersection}에 의해
스펙트럼 공간의 열린닫힌 부분집합들의 격자는 연결 성분들의 집합을 결정한다. 그러므로
$\pi_0(X) \to \pi_0(Y)$는 전단사이다. $\pi_0(X)$와 $\pi_0(Y)$는 프로유한 공간이므로
(Topology, 보조정리 \ref{topology-lemma-pi0-profinite}), Topology, 보조정리
\ref{topology-lemma-bijective-map}에 의해 $\pi_0(X) \to \pi_0(Y)$는 위상동형사상이다.
이것으로 (1)이 증명되었다. (2)는 곧바로 따른다.
\end{proof}

\noindent
다음 보조정리는 Hensel 쌍의 기본군이 그 닫힌 부분집합의 기본군임을 보인다.

\begin{lemma}
\label{pione-lemma-gabber}
$(A, I)$를 Hensel 쌍이라 하자. $X = \Spec(A)$ 및 $Z = \Spec(A/I)$라 두자. 함자
$$
\textit{F\'Et}_X \longrightarrow \textit{F\'Et}_Z,\quad
U \longmapsto U \times_X Z
$$
는 범주의 동치이다.
\end{lemma}

\begin{proof}
이는 More on Algebra, 보조정리 \ref{more-algebra-lemma-finite-etale-equivalence}를
옮겨 쓴 것이다.
\end{proof}

\noindent
다음 보조정리는 두꺼워짐의 기본군이 원래의 기본군과 같음을 보인다. 이는 아래의
보편 위상동형사상에 관한 더 강한 명제를 증명할 때 사용한다.

\begin{lemma}
\label{pione-lemma-thickening}
$X \subset X'$를 스킴의 두꺼워짐이라 하자. 함자
$$
\textit{F\'Et}_{X'} \longrightarrow \textit{F\'Et}_X,\quad
U' \longmapsto U' \times_{X'} X
$$
는 범주의 동치이다.
\end{lemma}

\begin{proof}
두꺼워짐에 관한 논의는 More on Morphisms, 절 \ref{more-morphisms-section-thickenings}를 보라.
$U' \to X'$를 에탈 사상이라 하고, $U = U' \times_{X'} X \to X$가 유한 에탈이라고 하자.
그러면 $U' \to X'$도 유한 에탈이다. 이는 예를 들어 More on Morphisms, 보조정리
\ref{more-morphisms-lemma-properties-that-extend-over-thickenings}에서 따른다.
이제 $X \subset X'$가 유한 차수의 두꺼워짐이면, 이 주의와 \'Etale Morphisms, 정리
\ref{etale-theorem-remarkable-equivalence}를 합쳐 보조정리를 얻는다. 아래에서는 일반적인 두꺼워짐에
대해 보조정리를 증명하지만, 독자에게는 이 증명을 건너뛸 것을 권한다.

\medskip\noindent
아핀 열린 덮개를 $X' = \bigcup X_i'$라 하자. 다음과 같이 놓는다.
$X_i = X \times_{X'} X_i'$, $X_{ij}' = X'_i \cap X'_j$,
$X_{ij} = X \times_{X'} X_{ij}'$, $X_{ijk}' = X'_i \cap X'_j \cap X'_k$,
$X_{ijk} = X \times_{X'} X_{ijk}'$.
두꺼워짐 $X_i \subset X'_i$, $X_{ij} \subset X_{ij}'$, 그리고
$X_{ijk} \subset X_{ijk}'$ 각각에 대해 정리를 증명할 수 있다고 하자. 그러면
스킴의 상대적 이어붙이기에 의해 $X \subset X'$에 대한 결과가 따른다.
Constructions, 절 \ref{constructions-section-relative-glueing}를 보라. 스킴
$X_i'$, $X_{ij}'$, $X_{ijk}'$는 각각 아핀 스킴의 열린 부분스킴으로서
분리적임에 주의하자. 이 논의를 한 번 더 반복하면 스킴
$X'_i$, $X_{ij}'$, $X_{ijk}'$가 아핀인 경우로 귀착된다.

\medskip\noindent
아핀인 경우 $X' = \Spec(A')$ 및 $X = \Spec(A'/I')$이고, $I'$는 국소 멱영 아이디얼이다.
이때 $(A', I')$는 Hensel 쌍이며(More on Algebra, 보조정리
\ref{more-algebra-lemma-locally-nilpotent-henselian}), 결과는 보조정리 \ref{pione-lemma-gabber}에서 따른다
(이 경우에는 훨씬 더 쉽다).
\end{proof}

\noindent
다음 명제를 증명하는 ``올바른'' 방법은 에탈 사이트의 불변성에서 이를 이끌어 내는 것일 것이다.
\'Etale Cohomology, 정리 \ref{etale-cohomology-theorem-topological-invariance}를 보라.

\begin{proposition}
\label{pione-proposition-universal-homeomorphism}
$f : X \to Y$를 스킴의 보편 위상동형사상이라 하자. 그러면
$$
\textit{F\'Et}_Y \longrightarrow \textit{F\'Et}_X,\quad
V \longmapsto V \times_Y X
$$
는 범주의 동치이다. 따라서 $X$와 $Y$가 연결이면 $f$는 기본군의 동형
$\pi_1(X, \overline{x}) \to \pi_1(Y, \overline{y})$을 유도한다.
\end{proposition}

\begin{proof}
보편 위상동형사상이라는 것은 정수적이고 보편 단사이며 전사인 사상과 같다.
Morphisms, 보조정리 \ref{morphisms-lemma-universal-homeomorphism}을 보라. 특히 대각사상
$\Delta : X \to X \times_Y X$는 Morphisms, 보조정리
\ref{morphisms-lemma-universally-injective}에 의해 두꺼워짐이다. 따라서 보조정리
\ref{pione-lemma-thickening}에 의해 유한 에탈 사상 $U \to X$가 주어지면 유일한 동형
$$
\varphi : U \times_Y X \to X \times_Y U
$$
이 존재하는데, 이는 $X \times_Y X$ 위의 유한 에탈 스킴들의 동형이고 $\Delta$에 의한 당김은
$X$ 위에서 $\text{id} : U \to U$를 준다. $X \to X \times_Y X \times_Y X$도
두꺼워짐이므로(이는 전단사인 닫힌 몰입이다), $(U, \varphi)$는 $X/Y$에 관한
내림 데이터이다. \'Etale Morphisms, 명제 \ref{etale-proposition-effective}에 의해
$U = X \times_Y V$이다. 여기서 $V \to Y$는 준콤팩트이고 분리적이며 에탈이다.
$V \to Y$가 유한임을 보이는 증명은 생략한다(힌트: 사상 $U \to V$는 전사이고
$U \to Y$는 정수적이다). 따라서 $\textit{F\'Et}_Y \to \textit{F\'Et}_X$는
본질적으로 전사이다.

\medskip\noindent
위와 마찬가지로 논하면, $V_1 \to Y$ 및 $V_2 \to Y$가 $\textit{F\'Et}_Y$의 대상일 때
임의의 사상 $a : X \times_Y V_1 \to X \times_Y V_2$는 $X$ 위의
표준 내림 데이터와 양립한다. 따라서 \'Etale Morphisms, 보조정리
\ref{etale-lemma-fully-faithful-cases}에 의해 $a$는 $Y$ 위의 사상 $V_1 \to V_2$로 내려간다.
\end{proof}










\section{고유 스킴의 유한 에탈 덮개}
\label{pione-section-finite-etale-over-proper}

\noindent
이 절에서는 Hensel 국소환 위의 연결 고유 스킴의 기본군이 그 특수
올의 기본군과 같음을 보인다. 또한 이 결과의 Hensel 쌍에 대한 변형도
증명한다.

\medskip\noindent
더 나아가 대수적으로 닫힌 체 $k$ 위의 연결 고유 스킴의 기본군은 $k$를 대수적으로 닫힌 확대체로
바꾸어도 변하지 않음을 보인다.

\medskip\noindent
연결인 경우에 결과를 진술하고 증명하는 대신 일반적인 경우에 결과를 증명하고,
기본군에 관한 결론은 보조정리
\ref{pione-lemma-what-equivalence-gives}를 사용하여 독자가 이끌어 내도록 한다.

\begin{lemma}
\label{pione-lemma-finite-etale-on-proper-over-henselian}
$A$를 Hensel 국소환이라 하고 $X$를 $A$ 위의 고유 스킴, 닫힌 올을
$X_0$이라 하자. 그러면 함자
$$
\textit{F\'Et}_X \to \textit{F\'Et}_{X_0},\quad
U \longmapsto U_0 = U \times_X X_0
$$
는 범주의 동치이다.
\end{lemma}

\begin{proof}
여기서의 증명은 대수화와 근사를 적용하는 한 예이다. 몇 단계로 나누어 진행한다.

\medskip\noindent
\medskip\noindent
$A$가 완비 Noether 국소환인 경우의 본질적 전사성.
$X_n = X \times_{\Spec(A)} \Spec(A/\mathfrak m^{n + 1})$이라 놓자.
\'Etale Morphisms, 정리 \ref{etale-theorem-remarkable-equivalence}에 의해 포함
$$
X_0 \to X_1 \to X_2 \to \ldots
$$
은 $X_0$ 위의 에탈 스킴의 범주와 $X_n$ 위의 에탈 스킴의 범주 사이의
범주 동치를 유도한다.
또한 $U_n \to X_n$이 유한 에탈 사상 $U_0 \to X_0$에 대응하면
$U_n \to X_n$도 유한이다. 예를 들어 More on Morphisms, 보조정리
\ref{more-morphisms-lemma-thicken-property-morphisms-cartesian}.
이 경우 사상 $U_0 \to \Spec(A/\mathfrak m)$는 $X_0$가 $A/\mathfrak m$ 위에서 고유이므로 고유이다. 따라서 Grothendieck의 대수화 정리를
(다음 형태로) 적용할 수 있다.
Cohomology of Schemes, 보조정리
\ref{coherent-lemma-algebraize-formal-scheme-finite-over-proper})
이에 의해 유한 사상 $U \to X$가 존재하며, 이를 $X_0$로 제한하면 $U_0$를 되찾는다.
More on Morphisms, 보조정리
\ref{more-morphisms-lemma-check-smoothness-on-infinitesimal-nbhds}
에 의해 $U \to X$는 $U_0$의 모든 점에서 에탈이다.
그러나 $U$의 각 점은 $U_0$의 한 점으로 특수화되므로($U$는 $A$ 위에서 고유이다),
$U \to X$가 에탈이라는 결론을 얻는다. 이로써 함자가 본질적으로 전사임이 따른다.

\medskip\noindent
\medskip\noindent
$A$가 완비 Noether 국소환인 경우의 충실충만성.
$U \to X$와 $V \to X$를 유한 에탈 사상이라 하고, $\varphi_0 : U_0 \to V_0$를
$X_0$ 위의 사상이라 하자. 다음 사상을 생각한다.
$$
\Gamma_{\varphi_0} : U_0 \longrightarrow U_0 \times_{X_0} V_0
$$
이 사상은 유한 에탈 사상인 동시에 닫힌 몰입이다.
$X = U \times_X V$에 본질적 전사성을 적용하면 유한 에탈 사상
$W \to U \times_X V$를 얻으며, 그 특수 올은 $\Gamma_{\varphi_0}$와 동형이다.
사영 $W \to U$를 생각하자. 이는 구성에 의해 유한 에탈이고 $U_0$ 위에서는 동형이다.
\'Etale Morphisms, 보조정리
\ref{etale-lemma-finite-etale-one-point}
$W \to U$는 $U_0$의 열린 근방에서 동형이다. 따라서 이는 동형이고, 합성
$\varphi : U \cong W \to V$는 $\varphi_0$의 원하는 올림이다.

\medskip\noindent
\medskip\noindent
$A$가 Hensel Noether 국소 G-환인 경우의 본질적 전사성.
$U_0 \to X_0$를 유한 에탈 사상이라 하자.
$A^\wedge$를 극대 아이디얼에 관한 $A$의 완비화라 하고, $X^\wedge$를
$X$의 $A^\wedge$로의 기저변환이라 하자. 위의 결과에 의해 특수 올이 $U_0$인
유한 에탈 사상 $V \to X^\wedge$가 존재한다. $A^\wedge = \colim A_i$로 쓰되,
$A \to A_i$가 유한형이라고 하자.
Limits, 보조정리 \ref{limits-lemma-descend-finite-presentation}에 의해 어떤 $i$와 유한 표시 사상
$U_i \to X_{A_i}$가 존재하여, 이를 $X^\wedge$로 기저변환하면 $V$가 된다. $i$를 충분히 크게 잡으면
$U_i \to X_{A_i}$가 유한 에탈이라고 가정할 수 있다
(Limits, 보조정리 \ref{limits-lemma-descend-finite-finite-presentation} 및
\ref{limits-lemma-descend-etale}). 다음과 같이 쓰자.
$$
A_i = A[x_1, \ldots, x_n]/(f_1, \ldots, f_m)
$$
환사상 $A_i \to A^\wedge$는 방정식계 $f_j = 0$의 해
$(a_1, \ldots, a_n)$를 $A^\wedge$ 안에서 택한 것으로 해석할 수 있다.
Smoothing Ring Maps, 정리 \ref{smoothing-theorem-approximation-property}에 의해
이 해를 $A$ 안의 해 $(b_1, \ldots, b_n)$로 (예를 들어 차수 $11$까지) 근사할 수 있다.
이를 다시 해석하면 $A$-대수 사상 $A_i \to A$를 얻으며, 이는 원래 사상
$A_i \to A^\wedge$와 같은 닫힌점을 준다($11 > 1$이기 때문이다). 따라서 이 환사상에 의한
$U \to X$는 $V \to X_{A_i}$의 기저변환이고, 특수 올이 $U_0$와 동형인 유한 에탈 사상이다.

\medskip\noindent
\medskip\noindent
$A$가 Hensel Noether 국소 G-환인 경우의 충실충만성.
이는 $A$가 완비 Noether 환인 경우와 완전히 같은 방식으로 본질적 전사성에서 도출된다.

\medskip\noindent
\medskip\noindent
일반적인 경우. $(A, \mathfrak m)$을 Hensel 국소환이라 하자. $S = \Spec(A)$라 두고
닫힌점을 $s \in S$로 쓰자. Limits, 보조정리
\ref{limits-lemma-proper-limit-of-proper-finite-presentation-noetherian}
에 의해 $X \to \Spec(A)$를 고유 사상 $X_i \to S_i$들의 공여과 극한으로 쓸 수 있다.
$S_i$는 $\mathbf{Z}$ 위에서 유한형이다. 각 $i$에 대해 $s_i \in S_i$를 $s$의 상이라 하자.
$S = \lim S_i$이고 $A = \mathcal{O}_{S, s}$이므로
$A = \colim \mathcal{O}_{S_i, s_i}$이다. 환 $A_i = \mathcal{O}_{S_i, s_i}$는
Noether 국소 G-환이다(More on Algebra, 명제
\ref{more-algebra-proposition-ubiquity-G-ring}).
More on Algebra, 보조정리 \ref{more-algebra-lemma-henselization-colimit}에 의해
$A = \colim A_i^h$이다. More on Algebra, 보조정리
\ref{more-algebra-lemma-henselization-G-ring}에 의해 환 $A_i^h$는 G-환이다.
따라서 $A = \colim A_i^h$이고
$$
X = \lim (X_i \times_{S_i} \Spec(A_i^h))
$$
가 스킴으로서 성립한다. $X$ 위의 유한 에탈 스킴의 범주는
$X_i \times_{S_i} \Spec(A_i^h)$ 위의 유한 에탈 스킴의 범주들의 극한이다(Limits, 보조정리
\ref{limits-lemma-descend-finite-presentation},
\ref{limits-lemma-descend-finite-finite-presentation} 및
\ref{limits-lemma-descend-etale})
에 대해서도 마찬가지이다. $X_0 = \lim (X_i \times_{S_i} s_i)$ 위의 유한 에탈 스킴에도
같은 결과가 성립한다. 따라서 위에서 다룬
$X / \Spec(A)$에 대한 결과를, 위에서 다룬
$(X_i \times_{S_i} \Spec(A_i^h)) / \Spec(A_i^h)$에 대한 결과에서 형식적으로 이끌어 낼 수 있다.
\end{proof}

\begin{lemma}
\label{pione-lemma-finite-etale-on-proper-over-henselian-pair}
$(A, I)$를 Hensel 쌍이라 하고 $X$를 $A$ 위의 고유 스킴이라 하자.
$X_0 = X \times_{\Spec(A)} \Spec(A/I)$라 두자. 그러면 함자
$$
\textit{F\'Et}_X \to \textit{F\'Et}_{X_0},\quad
U \longmapsto U_0 = U \times_X X_0
$$
는 범주의 동치이다.
\end{lemma}

\begin{proof}
이 보조정리의 증명은 보조정리
\ref{pione-lemma-finite-etale-on-proper-over-henselian}의 증명과 완전히 같다.

\medskip\noindent
\medskip\noindent
$A$가 Noether 환이고 $I$-진 완비인 경우의 본질적 전사성.
$X_n = X \times_{\Spec(A)} \Spec(A/I^{n + 1})$이라 두자.
\'Etale Morphisms, 정리 \ref{etale-theorem-remarkable-equivalence}에 의해 포함
$$
X_0 \to X_1 \to X_2 \to \ldots
$$
은 $X_0$ 위의 에탈 스킴의 범주와 $X_n$ 위의 에탈 스킴의 범주 사이의
범주 동치를 유도한다.
또한 $U_n \to X_n$이 유한 에탈 사상 $U_0 \to X_0$에 대응하면
$U_n \to X_n$도 유한이다. 예를 들어 More on Morphisms, 보조정리
\ref{more-morphisms-lemma-thicken-property-morphisms-cartesian}.
이 경우 사상 $U_0 \to \Spec(A/I)$는 $X_0$가 $A/I$ 위에서 고유이므로 고유이다.
따라서 Grothendieck의 대수화 정리를 (다음 형태로) 적용할 수 있다.
Cohomology of Schemes, 보조정리
\ref{coherent-lemma-algebraize-formal-scheme-finite-over-proper})
이에 의해 유한 사상 $U \to X$가 존재하며, 이를 $X_0$로 제한하면 $U_0$를 되찾는다.
More on Morphisms, 보조정리
\ref{more-morphisms-lemma-check-smoothness-on-infinitesimal-nbhds}
에 의해 $U \to X$는 $U_0$의 모든 점에서 에탈이다. 그러나 $U$의 각 점은
$U_0$의 한 점으로 특수화되므로($U$는 $A$ 위에서 고유이다), $U \to X$는 에탈이다.
이로써 함자가 본질적으로 전사임이 따른다.

\medskip\noindent
\medskip\noindent
$A$가 Noether 환이고 $I$-진 완비인 경우의 충실충만성.
$U \to X$와 $V \to X$를 유한 에탈 사상이라 하고, $\varphi_0 : U_0 \to V_0$를
$X_0$ 위의 사상이라 하자. 다음 사상을 생각한다.
$$
\Gamma_{\varphi_0} : U_0 \longrightarrow U_0 \times_{X_0} V_0
$$
이 사상은 유한 에탈 사상인 동시에 닫힌 몰입이다.
$X = U \times_X V$에 본질적 전사성을 적용하면 유한 에탈 사상
$W \to U \times_X V$를 얻으며, 그 특수 올은 $\Gamma_{\varphi_0}$와 동형이다.
사영 $W \to U$를 생각하자. 이는 구성에 의해 유한 에탈이고 $U_0$ 위에서는 동형이다.
\'Etale Morphisms, 보조정리
\ref{etale-lemma-finite-etale-one-point}
$W \to U$는 $U_0$의 열린 근방에서 동형이다. 따라서 이는 동형이고, 합성
$\varphi : U \cong W \to V$는 $\varphi_0$의 원하는 올림이다.

\medskip\noindent
\medskip\noindent
$(A, I)$가 Hensel 쌍이고 $A$가 Noether G-환인 경우의 본질적 전사성.
$U_0 \to X_0$를 유한 에탈 사상이라 하자. $A^\wedge$를 $I$에 관한 $A$의 완비화라 하자.
$A^\wedge$는 Noether 환이고 $IA^\wedge$-진 완비임에 주의하자. Algebra, 보조정리
\ref{algebra-lemma-completion-complete} 및
\ref{algebra-lemma-completion-Noetherian-Noetherian}을 보라. $X^\wedge$를 $X$의 $A^\wedge$로의
기저변환이라 하자. 위의 결과에 의해 특수 올이 $U_0$인 유한 에탈 사상
$V \to X^\wedge$가 존재한다. $A^\wedge = \colim A_i$로 쓰되
$A \to A_i$가 유한형이라고 하자. Limits, 보조정리 \ref{limits-lemma-descend-finite-presentation}에 의해
어떤 $i$와 유한 표시 사상 $U_i \to X_{A_i}$가 존재하여, 이를 $X^\wedge$로 기저변환하면 $V$가 된다.
$i$를 충분히 크게 잡으면 $U_i \to X_{A_i}$가 유한 에탈이라고 가정할 수 있다
(Limits, 보조정리 \ref{limits-lemma-descend-finite-finite-presentation} 및
\ref{limits-lemma-descend-etale}). 다음과 같이 쓰자.
$$
A_i = A[x_1, \ldots, x_n]/(f_1, \ldots, f_m)
$$
환사상 $A_i \to A^\wedge$는 방정식계 $f_j = 0$의 해
$(a_1, \ldots, a_n)$를 $A^\wedge$ 안에서 택한 것으로 해석할 수 있다. Smoothing Ring Maps, 보조정리
\ref{smoothing-lemma-henselian-pair}에 의해 이 해를 $A$ 안의 해
$(b_1, \ldots, b_n)$로 (예를 들어 차수 $11$까지) 근사할 수 있다. 이를 다시 해석하면 $A$-대수 사상 $A_i \to A$를 얻으며, 이는 원래 사상
$A_i \to A^\wedge$와 같은 닫힌점을 준다($11 > 1$이기 때문이다). 따라서 이 환사상에 의한
$U \to X$는 $V \to X_{A_i}$의 기저변환이고, 특수 올이 $U_0$와 동형인 유한 에탈 사상이다.

\medskip\noindent
\medskip\noindent
$(A, I$가 Hensel 쌍이고 $A$가 Noether G-환인 경우의 충실충만성.
이는 $A$가 완비 Noether 환인 경우와 완전히 같은 방식으로 본질적 전사성에서 도출된다.

\medskip\noindent
\medskip\noindent
일반적인 경우. $(A, I)$를 Hensel 쌍이라 하자. $S = \Spec(A)$, $S_0 = \Spec(A/I)$라 두자.
Limits, 보조정리
\ref{limits-lemma-proper-limit-of-proper-finite-presentation-noetherian}
에 의해 $X \to \Spec(A)$를 고유 사상 $X_i \to S_i$들의 공여과 극한으로 쓸 수 있다.
여기서 $S_i$는 $\mathbf{Z}$ 위의 아핀 유한형이다. $S_i = \Spec(A_i)$라 쓰고,
$I_i \subset A_i$를 사상 $A_i \to A$에 의한 $I$의 역상이라 하자.
$S_{i, 0} = \Spec(A_i/I_i)$라 두자. $S = \lim S_i$이므로 $A = \colim A_i$이고,
따라서 $I = \colim I_i$ 및 $A/I = \colim A_i/I_i$도 성립한다.
환 $A_i$는 Noether G-환이다(More on Algebra, 명제
\ref{more-algebra-proposition-ubiquity-G-ring}).
$(A_i^h, I_i^h)$를 $(A_i, I_i)$의 Hensel화라 쓰자. More on Algebra, 보조정리
\ref{more-algebra-lemma-henselization-colimit}에 의해 $A = \colim A_i^h$이다.
More on Algebra, 보조정리 \ref{more-algebra-lemma-henselization-pair-G-ring}에 의해 환 $A_i^h$
는 G-환이다. 따라서 $A = \colim A_i^h$이고
$$
X = \lim (X_i \times_{S_i} \Spec(A_i^h))
$$
가 스킴으로서 성립한다. $X$ 위의 유한 에탈 스킴의 범주는
$X_i \times_{S_i} \Spec(A_i^h)$ 위의 유한 에탈 스킴의 범주들의 극한이다(Limits, 보조정리
\ref{limits-lemma-descend-finite-presentation},
\ref{limits-lemma-descend-finite-finite-presentation} 및
\ref{limits-lemma-descend-etale})
에 대해서도 마찬가지이다. $X_0 = \lim (X_i \times_{S_i} S_{i, 0})$ 위의 유한 에탈 스킴에도
같은 결과가 성립한다. 따라서 위에서 다룬
따라서 $X / \Spec(A)$에 대한 결과를, 위에서 다룬
$(X_i \times_{S_i} \Spec(A_i^h)) / \Spec(A_i^h)$에 대한 결과에서 형식적으로 이끌어 낼 수 있다.
\end{proof}

\begin{lemma}
\label{pione-lemma-finite-etale-invariant-over-proper}
$k'/k$를 대수적으로 닫힌 체들의 확대라 하자. $X$를 $k$ 위의 고유 스킴이라 하자. 그러면 함자
$$
U \longmapsto U_{k'}
$$
는 $X$ 위의 유한 에탈 스킴의 범주와 $X_{k'}$ 위의 유한 에탈 스킴의 범주 사이의
범주 동치이다.
\end{lemma}

\begin{proof}
함자가 본질적으로 전사임을 보이자. $U' \to X_{k'}$를 유한 에탈 사상이라 하자.
$k' = \colim A_i$를 유한형 $k$-대수들의 여과 여극한으로 쓰자.
Limits, 보조정리 \ref{limits-lemma-descend-finite-presentation}에 의해 어떤 $i$와
유한 표시 사상 $U_i \to X_{A_i}$가 존재하여, 이를 $X_{k'}$로 밑변환하면
$U'$가 된다. $i$를 충분히 크게 잡으면 $U_i \to X_{A_i}$가 유한 에탈이라고
가정할 수 있다(Limits, 보조정리
\ref{limits-lemma-descend-finite-finite-presentation} 및
\ref{limits-lemma-descend-etale}).
$k$는 대수적으로 닫혀 있으므로 $\Spec(A_i)$ 안에서 $k$-값점 $t$를 찾을 수 있다.
$U = (U_i)_t$를 $t$ 위에서의 $U_i$의 올이라 하자. 점 $t$에 대응하는 극대 아이디얼을
$\mathfrak m$이라 하고, $(A_i)_{\mathfrak m}$의 헨젤화를 $A_i^h$라 하자.
보조정리 \ref{pione-lemma-finite-etale-on-proper-over-henselian}에 의해
$X_{A_i^h}$ 위의 스킴으로서
$(U_i)_{A_i^h} = U \times \Spec(A_i^h)$이다.
이제 $A_i^h$는 $A_i$ 위에서 대수적이고(예를 들어 Smoothing Ring Maps, 예
\ref{smoothing-example-describe-henselian}의 논의를 보라), $k'$는 대수적으로 닫혀
있으므로, 주어진 포함 $A_i \subset k'$를 연장하는 환 사상
$A_i^h \to k'$를 찾을 수 있다. 따라서 $U'$는 $U$의 밑변환과 동형이다.
충실충만성의 증명도 완전히 같다.
\end{proof}








\section{국소 연결성}
\label{pione-section-unibranch}

\noindent
이 절에서는 스킴 $X$의 조밀 열린집합 $U$에 대하여 언제
$\pi_1(U) \to \pi_1(X)$가 전사인지 묻는다. 대략 말해, 이는
$x \in X \setminus U$ 주위의 임의의 작은 “공” $B$에 대하여
$U \cap B$가 연결일 때 성립함을 보일 것이다.

\begin{lemma}
\label{pione-lemma-dense-faithful}
$f : X \to Y$를 스킴의 사상이라 하자. $f(X)$가 $Y$에서 조밀하면 밑변환 함자
$\textit{F\'Et}_Y \to \textit{F\'Et}_X$는 충실하다.
\end{lemma}

\begin{proof}
유한 에탈 덮개의 범주는 내부 Hom을 가지므로(보조정리
\ref{pione-lemma-internal-hom-finite-etale}), 다음을 보이면 충분하다.
$W$를 $Y$ 위의 유한 에탈 스킴이라 하고, $s : X \to W$를 $X$ 위의 사상이라 하자.
$s = t \circ f$를 만족하는 단면 $t : Y \to W$는 많아야 하나이다.
$s = t_1 \circ f = t_2 \circ f$를 만족하는 두 단면
$t_1, t_2 : Y \to W$를 생각하자. $t_1$과 $t_2$의 등화자는 $Y$에서 닫혀 있고
(Schemes, 보조정리 \ref{schemes-lemma-where-are-they-equal}),
$f(X)$는 $Y$에서 조밀하므로 $t_1$과 $t_2$는 $Y_{red}$ 위에서 일치한다.
따라서 $t_1$과 $t_2$의 상은 같고, 이 상은 $W$의 열린닫힌 부분스킴이며 $Y$로 동형으로 사상된다
(\'Etale Morphisms, 명제 \ref{etale-proposition-properties-sections}).
그러므로 두 단면은 같다.
\end{proof}

\noindent
다음 보조정리에서 엄밀 헨젤화의 천공 스펙트럼이 연결이라는 조건은, 예를 들어
국소환이 기하학적으로 일분기라는 가정에서 따른다. More on Algebra, 보조정리
\ref{more-algebra-lemma-geometrically-unibranch}를 보라. 그 부분적 역은
Properties, 보조정리 \ref{properties-lemma-geometrically-unibranch}에 있다.

\begin{lemma}
\label{pione-lemma-same-etale-extensions}
$(A, \mathfrak m)$을 국소환이라 하자. $X = \Spec(A)$라 하고
$U = X \setminus \{\mathfrak m\}$라 하자. $A$의 엄밀 헨젤화의 천공 스펙트럼이
연결이면
$$
\textit{F\'Et}_X \longrightarrow \textit{F\'Et}_U,\quad
Y \longmapsto Y \times_X U
$$
는 충실충만 함자이다.
\end{lemma}

\begin{proof}
먼저 $A$가 엄밀 헨젤 환이라고 가정하자. 이 경우 $X$의 임의의 유한 에탈 덮개 $Y$는
$X$의 유한 개 서로소 복사본의 합과 동형이다. 따라서 $U$ 위의 임의의 사상
$U \to U \amalg \ldots \amalg U$가 $X$ 위의 사상
$X \to X \amalg \ldots \amalg X$로 유일하게 연장됨을 보이면 충분하다.
$U$가 연결이면(특히 공집합이 아니면) 이는 참이다.

\medskip\noindent
이제 일반적인 경우를 보자. 유한 에탈 덮개의 범주는 내부 Hom을 가지므로(보조정리
\ref{pione-lemma-internal-hom-finite-etale}), 다음을 보이면 충분하다. $Y$를 $X$ 위의
유한 에탈 스킴이라 할 때, $X$ 위의 임의의 사상 $s : U \to Y$는 $X$ 위의 사상
$t : X \to Y$로 연장된다. $A$의 엄밀 헨젤화를 $A^{sh}$라 하고
$X^{sh} = \Spec(A^{sh})$, $U^{sh} = U \times_X X^{sh}$,
$Y^{sh} = Y \times_X X^{sh}$라 쓰자. 첫 문단과 $A$에 대한 가정에 의해
$s$의 밑변환 $s^{sh} : U^{sh} \to Y^{sh}$를
$t^{sh} : X^{sh} \to Y^{sh}$로 연장할 수 있다.
$A' = A^{sh} \otimes_A A^{sh}$라 하자. $t^{sh}$를
$X' = \Spec(A')$로 당긴 두 사상 $t'_1, t'_2$는
$s$를 $U' = U \times_X X'$로 당긴 사상 $s'$의 연장이다.
$A \to A'$가 평탄하므로, $A \to A'$에 대한 내려가기
(Algebra, 보조정리 \ref{algebra-lemma-flat-going-down})에 의해
$U' \subset X'$는 위상적으로 조밀하다. 따라서 보조정리
\ref{pione-lemma-dense-faithful}에 의해 $t'_1 = t'_2$이다.
그러므로 예를 들어 Descent, 보조정리
\ref{descent-lemma-fpqc-universal-effective-epimorphisms}에 의해
$t^{sh}$는 사상 $t : X \to Y$로 내려간다.
\end{proof}

\noindent
보조정리 \ref{pione-lemma-same-etale-extensions}에 비추어 보면, 환과 그 엄밀 헨젤화의
천공 스펙트럼이 언제 연결인지 아는 것이 흥미롭다. 충분조건을 주는 Hartshorne의
유명한 보조정리에 대해서는 Local Cohomology, 보조정리
\ref{local-cohomology-lemma-depth-2-connected-punctured-spectrum}를 보라.

\begin{lemma}
\label{pione-lemma-quasi-compact-dense-open-connected-at-infinity-Noetherian}
$X$를 스킴이라 하고 $U \subset X$를 조밀 열린집합이라 하자. 다음을 가정하자.
\begin{enumerate}
\item $X$의 바탕 위상공간은 뇌터 공간이고,
\item 모든 $x \in X \setminus U$에 대하여 $\mathcal{O}_{X, x}$의
엄밀 헨젤화의 천공 스펙트럼이 연결이다.
\end{enumerate}
그러면 $\textit{F\'Et}_X \to \textit{F\'et}_U$는 충실충만하다.
\end{lemma}

\begin{proof}
$Y_1, Y_2$를 $X$ 위의 유한 에탈 스킴이라 하고
$\varphi : (Y_1)_U \to (Y_2)_U$를 $U$ 위의 사상이라 하자.
$\varphi$가 $X$ 위의 사상 $Y_1 \to Y_2$로 유일하게 올라감을 보여야 한다.
유일성은 보조정리 \ref{pione-lemma-dense-faithful}에서 따른다.

\medskip\noindent
$x \in X \setminus U$를 $X \setminus U$의 한 기약 성분의 일반점이라 하자.
$V = U \times_X \Spec(\mathcal{O}_{X, x})$라 하자. $x$의 선택에 의해 이는
$\Spec(\mathcal{O}_{X, x})$의 천공 스펙트럼이다. 보조정리
\ref{pione-lemma-same-etale-extensions}에 의해 사상
$\varphi_V : (Y_1)_V \to (Y_2)_V$를 유일하게 사상
$(Y_1)_{\Spec(\mathcal{O}_{X, x})} \to (Y_2)_{\Spec(\mathcal{O}_{X, x})}$로
연장할 수 있다. Limits, 보조정리 \ref{limits-lemma-glueing-near-point}에 의해
$x$를 포함하는 열린집합 $U \subset U'$와 $\varphi$의 연장
$\varphi' : (Y_1)_{U'} \to (Y_2)_{U'}$를 얻는다.
$X$의 바탕 위상공간은 뇌터 공간이므로, $\varphi$가 정의된 열린집합의 여집합에 대한
뇌터 귀납법으로 증명이 끝난다.
\end{proof}

\begin{lemma}
\label{pione-lemma-retrocompact-dense-open-connected-at-infinity-closed}
$X$를 스킴이라 하고 $U \subset X$를 조밀 열린집합이라 하자. 다음을 가정하자.
\begin{enumerate}
\item $U \to X$는 준콤팩트이고,
\item $X \setminus U$의 모든 점은 닫힌점이며,
\item 모든 $x \in X \setminus U$에 대하여 $\mathcal{O}_{X, x}$의
엄밀 헨젤화의 천공 스펙트럼이 연결이다.
\end{enumerate}
그러면 $\textit{F\'Et}_X \to \textit{F\'et}_U$는 충실충만하다.
\end{lemma}

\begin{proof}
$Y_1, Y_2$를 $X$ 위의 유한 에탈 스킴이라 하고
$\varphi : (Y_1)_U \to (Y_2)_U$를 $U$ 위의 사상이라 하자.
$\varphi$가 $X$ 위의 사상 $Y_1 \to Y_2$로 유일하게 올라감을 보여야 한다.
유일성은 보조정리 \ref{pione-lemma-dense-faithful}에서 따른다.

\medskip\noindent
$x \in X \setminus U$라 하고 $V = U \times_X \Spec(\mathcal{O}_{X, x})$라 하자.
$X \setminus U$의 모든 점이 닫힌점이므로 $V$는
$\Spec(\mathcal{O}_{X, x})$의 천공 스펙트럼이다. 보조정리
\ref{pione-lemma-same-etale-extensions}에 의해 사상
$\varphi_V : (Y_1)_V \to (Y_2)_V$를 유일하게 사상
$(Y_1)_{\Spec(\mathcal{O}_{X, x})} \to (Y_2)_{\Spec(\mathcal{O}_{X, x})}$로
연장할 수 있다. Limits, 보조정리 \ref{limits-lemma-glueing-near-point}
(여기서는 $U$가 $X$에서 레트로콤팩트임을 사용한다)에 의해
$x$를 포함하는 열린집합 $U \subset U'_x$와 $\varphi$의 연장
$\varphi'_x : (Y_1)_{U'_x} \to (Y_2)_{U'_x}$를 얻는다.
두 점 $x, x' \in X \setminus U$가 주어지면, $U$가
$U'_x \cap U'_{x'}$에서 조밀하므로 사상 $\varphi'_x$와 $\varphi'_{x'}$는
그 교집합 위에서 일치한다(보조정리 \ref{pione-lemma-dense-faithful}).
따라서 원하는 대로 $\varphi$를 $\bigcup U'_x = X$까지 연장할 수 있다.
\end{proof}

\begin{lemma}
\label{pione-lemma-quasi-compact-dense-open-connected-at-infinity}
$X$를 스킴이라 하고 $U \subset X$를 조밀 열린집합이라 하자. 다음을 가정하자.
\begin{enumerate}
\item $X$의 모든 준콤팩트 열린집합은 유한 개의 기약 성분을 가지며,
\item 모든 $x \in X \setminus U$에 대하여 $\mathcal{O}_{X, x}$의
엄밀 헨젤화의 천공 스펙트럼이 연결이다.
\end{enumerate}
그러면 $\textit{F\'Et}_X \to \textit{F\'et}_U$는 충실충만하다.
\end{lemma}

\begin{proof}
$Y_1, Y_2$를 $X$ 위의 유한 에탈 스킴이라 하고
$\varphi : (Y_1)_U \to (Y_2)_U$를 $U$ 위의 사상이라 하자.
$\varphi$가 $X$ 위의 사상 $Y_1 \to Y_2$로 유일하게 올라감을 보여야 한다.
유일성은 보조정리 \ref{pione-lemma-dense-faithful}에서 따른다.
존재성을 보이기 위해 $U \not = X$일 때 $U$를 더 크게 만들 수 있음을 보인 뒤
초른의 보조정리를 적용한다.

\medskip\noindent
$x \in X \setminus U$를 $X \setminus U$의 한 기약 성분의 일반점이라 하고
$V = U \times_X \Spec(\mathcal{O}_{X, x})$라 하자. $x$의 선택에 의해 이는
$\Spec(\mathcal{O}_{X, x})$의 천공 스펙트럼이다. 보조정리
\ref{pione-lemma-same-etale-extensions}에 의해 사상
$\varphi_V : (Y_1)_V \to (Y_2)_V$를 유일하게 사상
$(Y_1)_{\Spec(\mathcal{O}_{X, x})} \to (Y_2)_{\Spec(\mathcal{O}_{X, x})}$로
연장할 수 있다. $W \subset X$를 $x$의 아핀 근방으로 택하자.
$U \cap W$는 $W$에서 조밀하므로 $W$의 일반점
$\eta_1, \ldots, \eta_n$을 모두 포함한다. $\eta_1, \ldots, \eta_n$을 포함하는 아핀 열린집합
$W' \subset W \cap U$를 택하고
$V' = W' \times_X \Spec(\mathcal{O}_{X, x})$라 하자.
Limits, 보조정리 \ref{limits-lemma-glueing-near-point}를
$x \in W \supset W'$에 적용하면, $x \in W''$인 열린집합
$W' \subset W'' \subset W$와 사상
$\varphi'' : (Y_1)_{W''} \to (Y_2)_{W''}$를 얻으며, 이 사상은
$W'$ 위에서 $\varphi$와 일치한다. $W'$는 $W'' \cap U$에서 조밀하므로
보조정리 \ref{pione-lemma-dense-faithful}에 의해 $\varphi$와 $\varphi''$는
$U \cap W'$ 위에서 일치한다. 따라서 $\varphi$와 $\varphi''$를 붙여서
$U' = U \cup W''$ 위의 사상 $\varphi'$를 얻으며, 이는 $U$ 위에서
$\varphi$와 일치한다. $x \in U'$이므로 $\varphi$를 진정으로 더 큰
열린집합으로 연장한 것이다.

\medskip\noindent
$U \subset U'$이고 $\varphi'$가 $\varphi$의 연장인 순서쌍
$(U', \varphi')$들의 집합을 $\mathcal{S}$라 하자. $\mathcal{S}$에 자명한 방식으로
부분순서를 주자. $(U'_i, \varphi'_i)$가 전순서 부분집합이면, 이는 상계
$(U', \varphi')$를 갖는다. 실제로 $U' = \bigcup U'_i$라 하고
$\varphi' : (Y_1)_{U'} \to (Y_2)_{U'}$를 각 $U'_i$ 위에서
$\varphi'_i$와 일치하는 사상으로 잡으면 된다. 따라서 초른의 보조정리를 적용할 수 있고,
$\mathcal{S}$는 극대 원소를 갖는다. 위의 논의에 의해 이 극대 원소는
$X$ 전체 위에서 정의된 $\varphi$의 연장이다.
\end{proof}

\begin{lemma}
\label{pione-lemma-local-exact-sequence}
$(A, \mathfrak m)$을 국소환이라 하자. $X = \Spec(A)$라 하고
$U = X \setminus \{\mathfrak m\}$라 하자. $A$의 엄밀 헨젤화 $A^{sh}$의
천공 스펙트럼을 $U^{sh}$라 하자. $U$가 준콤팩트이고 $U^{sh}$가 연결이라고
가정하자. 그러면 열
$$
\pi_1(U^{sh}, \overline{u}) \to \pi_1(U, \overline{u}) \to
\pi_1(X, \overline{u}) \to 1
$$
은 보조정리 \ref{pione-lemma-functoriality-galois-ses} (1)의 의미에서 완전하다.
\end{lemma}

\begin{proof}
사상 $\pi_1(U) \to \pi_1(X)$는 보조정리
\ref{pione-lemma-same-etale-extensions} 및
\ref{pione-lemma-functoriality-galois-surjective}에 의해 전사이다.

\medskip\noindent
$X^{sh} = \Spec(A^{sh})$라 쓰자. $Y \to X$를 유한 에탈 사상이라 하자.
그러면 $Y^{sh} = Y \times_X X^{sh} \to X^{sh}$도 유한 에탈 사상이다.
$A^{sh}$가 엄밀 헨젤 환이므로 $Y^{sh}$는 $X^{sh}$의 서로소 복사본들의 합과
동형이다. 따라서 $Y \times_X U^{sh}$에 대해서도 마찬가지이다. 그러므로 합성
$\pi_1(U^{sh}) \to \pi_1(U) \to \pi_1(X)$는 자명하다.
보조정리 \ref{pione-lemma-composition-trivial}를 보라.

\medskip\noindent
증명을 끝내려면 보조정리 \ref{pione-lemma-functoriality-galois-ses}에 따라 다음을 보이면
충분하다. $V \to U$가 유한 에탈 사상이고
$V \times_U U^{sh}$가 $U^{sh}$의 서로소 복사본들의 합이라고 하자.
그러면 유한 에탈 사상 $Y \to X$를 찾아
$U$ 위에서 $V \cong Y \times_X U$가 되게 할 수 있다.
가정에 의해 유한 에탈 사상 $Y^{sh} \to X^{sh}$와 동형
$V \times_U U^{sh} \cong Y^{sh} \times_{X^{sh}} U^{sh}$가 존재한다.
다음 도식을 생각하자.
$$
\xymatrix{
U \ar[d] & U^{sh} \ar[d] \ar[l] &
U^{sh} \times_U U^{sh} \ar[d] \ar@<1ex>[l] \ar@<-1ex>[l] &
U^{sh} \times_U U^{sh} \times_U U^{sh}
\ar[d] \ar@<1ex>[l] \ar[l] \ar@<-1ex>[l] \\
X & X^{sh} \ar[l] &
X^{sh} \times_X X^{sh} \ar@<1ex>[l] \ar@<-1ex>[l] &
X^{sh} \times_X X^{sh} \times_X X^{sh} \ar@<1ex>[l] \ar[l] \ar@<-1ex>[l]
}
$$
가정에 의해 $U \subset X$가 준콤팩트이므로 아래쪽을 향한 모든 화살표는
준콤팩트 열린 몰입이다. $U^{sh} \times_U U^{sh}$에 속하지 않는 점
$\xi \in X^{sh} \times_X X^{sh}$를 택하자. 그러면 $\xi$는
$X^{sh}$의 닫힌점 $x^{sh}$ 위에 놓인다. 첫째 사영
$X^{sh} \times_X X^{sh}$가 정하는 국소환 준동형
$$
A^{sh} = \mathcal{O}_{X^{sh}, x^{sh}} \to
\mathcal{O}_{X^{sh} \times_X X^{sh}, \xi}
$$
을 생각하자. 이는 국소화인 에탈 환 사상들로 이루어진 국소 준동형들의
여과 여극한이다. $A^{sh}$가 엄밀 헨젤 환이므로 이 사상은 동형이다.
여집합의 모든 $\xi$에 대해 이것이 성립하므로 이 점들 사이에는 특수화가 없고,
따라서 각 $\xi$는 닫힌점이다(이를 직접 보일 수도 있다).
$\xi$에서의 국소환은 $A^{sh}$와 동형이므로 엄밀 헨젤 환이며
연결인 천공 스펙트럼을 갖는다. 마찬가지로,
$U^{sh} \times_U U^{sh} \times_U U^{sh}$에 속하지 않는
$X^{sh} \times_X X^{sh} \times_X X^{sh}$의 점 $\xi$에 대해서도 같은 말이 성립한다.
보조정리 \ref{pione-lemma-retrocompact-dense-open-connected-at-infinity-closed}에 의해
수직 화살표를 따른 당김은 유한 에탈 스킴의 범주 사이의 충실충만 함자를 유도한다.
따라서 fpqc 덮개 $\{U^{sh} \to U\}$에 대한
$V \times_U U^{sh}$ 위의 표준 내림 데이터는 fpqc 덮개
$\{X^{sh} \to X\}$에 대한 $Y^{sh}$의 내림 데이터로 옮겨진다.
$Y^{sh} \to X^{sh}$는 유한, 따라서 아핀 사상이므로 이 내림 데이터는 유효하다
(Descent, 보조정리 \ref{descent-lemma-affine}).
그러므로 아핀 사상 $Y \to X$와 내림 데이터와 양립하는 동형
$Y \times_X X^{sh} \to Y^{sh}$를 얻는다. 내림 데이터의 충실충만성
(Descent, 보조정리 \ref{descent-lemma-refine-coverings-fully-faithful}의 경우와 같이)에 의해
동형 $V \to U \times_X Y$를 얻는다. 마지막으로
$Y^{sh} \to X^{sh}$가 유한 에탈이므로 $Y \to X$도 유한 에탈이다.
Descent, 보조정리 \ref{descent-lemma-descending-property-etale} 및
\ref{descent-lemma-descending-property-finite}를 보라.
\end{proof}

\noindent
$X$를 기약 스킴이라 하고 $\eta \in X$를 그 일반점이라 하자. 표준 사상
$\eta \to X$는 표준 사상
\begin{equation}
\label{pione-equation-inclusion-generic-point}
\text{Gal}(\kappa(\eta)^{sep}/\kappa(\eta)) = \pi_1(\eta, \overline{\eta})
\longrightarrow \pi_1(X, \overline{\eta})
\end{equation}
을 유도한다. 왼쪽의 동일시는 보조정리
\ref{pione-lemma-fundamental-group-Galois-group}에 따른다.

\begin{lemma}
\label{pione-lemma-irreducible-geometrically-unibranch}
$X$를 기약이고 기하학적으로 일분기인 스킴이라 하자. 공집합이 아닌 임의의 열린집합
$U \subset X$에 대하여 표준 사상
$$
\pi_1(U, \overline{u}) \longrightarrow \pi_1(X, \overline{u})
$$
은 전사이다. 사상 (\ref{pione-equation-inclusion-generic-point})
$\pi_1(\eta, \overline{\eta}) \to \pi_1(X, \overline{\eta})$도 전사이다.
\end{lemma}

\begin{proof}
보조정리 \ref{pione-lemma-thickening}에 의해 $X$를 그 축약으로 바꾸어도 된다.
따라서 $X$가 정역 스킴이라고 가정할 수 있다. 보조정리
\ref{pione-lemma-functoriality-galois-surjective}에 의해 이 보조정리의 주장은
함자 $\textit{F\'Et}_X \to \textit{F\'Et}_U$와
$\textit{F\'Et}_X \to \textit{F\'Et}_\eta$가 충실충만이라는 명제로 바뀐다.

\medskip\noindent
$\textit{F\'Et}_X \to \textit{F\'Et}_U$에 대한 결과는 보조정리
\ref{pione-lemma-quasi-compact-dense-open-connected-at-infinity}와, 기하학적으로 일분기인
국소환 $A$의 엄밀 헨젤화가 기약 스펙트럼을 갖는다는 사실에서 따른다.
More on Algebra, 보조정리
\ref{more-algebra-lemma-geometrically-unibranch}를 보라.

\medskip\noindent
잉여체 $\kappa(\eta) = \mathcal{O}_{X, \eta}$는 공집합이 아닌 아핀 열린집합
$U \subset X$에 대한 $\mathcal{O}_X(U)$의 여과 여극한이다.
따라서 $\textit{F\'Et}_\eta$는 그러한 $U$에 대한 범주
$\textit{F\'Et}_U$들의 여극한이다. Limits, 보조정리
\ref{limits-lemma-descend-finite-presentation},
\ref{limits-lemma-descend-finite-finite-presentation} 및
\ref{limits-lemma-descend-etale}를 보라. 이제 형식적인 논증으로
$\textit{F\'Et}_X \to \textit{F\'Et}_\eta$의 충실충만성은 함자
$\textit{F\'Et}_X \to \textit{F\'Et}_U$의 충실충만성에서 따른다.
\end{proof}

\begin{lemma}
\label{pione-lemma-exact-sequence-finite-nr-closed-pts}
$X$를 스킴이라 하자. $x_1, \ldots, x_n \in X$를 다음 조건을 만족하는
유한 개의 닫힌점이라 하자.
\begin{enumerate}
\item $U = X \setminus \{x_1, \ldots, x_n\}$는 연결이고 $X$의 레트로콤팩트 열린집합이며,
\item 각 $i$에 대하여 $\mathcal{O}_{X, x_i}$의 엄밀 헨젤화의 천공 스펙트럼
$U_i^{sh}$는 연결이다.
\end{enumerate}
그러면 사상 $\pi_1(U) \to \pi_1(X)$는 전사이고, 그 핵은
$i = 1, \ldots, n$에 대한 $\pi_1(U_i^{sh}) \to \pi_1(U)$의 상들을 모두 포함하는
$\pi_1(U)$의 가장 작은 닫힌 정규 부분군이다.
\end{lemma}

\begin{proof}
전사성은 보조정리
\ref{pione-lemma-retrocompact-dense-open-connected-at-infinity-closed} 및
\ref{pione-lemma-functoriality-galois-surjective}에서 따른다. 다음 사상들의 열을 생각하자.
$$
\pi_1(U)  \to \ldots \to
\pi_1(X \setminus \{x_1, x_2\}) \to \pi_1(X \setminus \{x_1\}) \to \pi_1(X)
$$
군론적 논증에 의해 핵에 관한 명제는 $n = 1$인 경우만 보이면 충분하다
(세부사항은 생략한다). $x = x_1$, $U^{sh} = U_1^{sh}$라 쓰고,
$A = \mathcal{O}_{X, x}$라 하며, 그 엄밀 헨젤화를 $A^{sh}$라 하자.
다음 도식을 생각하자.
$$
\xymatrix{
U \ar[d] &
\Spec(A) \setminus \{\mathfrak m\} \ar[l] \ar[d] &
U^{sh} \ar[d] \ar[l] \\
X & \Spec(A) \ar[l] & \Spec(A^{sh}) \ar[l]
}
$$
보조정리 \ref{pione-lemma-functoriality-galois-ses}에 의해,
$U^{sh}$의 자명한 덮개로 당겨지는 유한 에탈 사상 $V \to U$가
$X$ 위의 유한 에탈 스킴으로 연장됨을 보여야 한다.
보조정리 \ref{pione-lemma-local-exact-sequence}에 의해 $A$의 천공 스펙트럼 위의
유한 에탈 스킴에 대한 대응 명제를 알고 있다. 한편 Limits, 보조정리
\ref{limits-lemma-glueing-near-closed-point}에 의해 $X$ 위의 유한 표시 스킴은
$A$의 천공 스펙트럼 위에서 붙인, $U$ 위의 유한 표시 스킴과
$A$ 위의 유한 표시 스킴의 데이터와 같다. 이것으로 증명이 끝난다.
\end{proof}








\section{정규 스킴의 기본군}
\label{pione-section-normal}

\noindent
$X$를 정역이고 기하학적으로 일분기인 스킴이라 하자. 앞 절에서 $X$의 기본군은
$X$의 함수체 $K$의 갈루아 군의 몫임을 보았다. 사상이 연속이므로 그 핵은
갈루아 군의 닫힌 정규 부분군이다. 따라서 갈루아 이론
(Fields, 정리 \ref{fields-theorem-inifinite-galois-theory})에 의해 이 핵은
갈루아 확대 $M/K$에 대응한다. 이 절에서는 $X$가 정규 정역 스킴일 때
$M$을 결정한다.

\medskip\noindent
$X$를 함수체 $K$를 갖는 정규 정역 스킴이라 하자. $L/K$를 유한 확대라 하자.
Morphisms, 절 \ref{morphisms-section-normalization-X-in-Y}에서 정의한,
사상 $\Spec(L) \to X$에서의 $X$의 정규화 $Y \to X$를 생각하자.
이 상황에서 $Y \to X$가 스킴의 비분기 사상이면
{\it $X$는 $L$에서 비분기이다}라고 한다. 이 조건은 보조정리
\ref{pione-lemma-unramified}에서 더 자세히 설명한다. $Y \to X$의
$\Spec(K)$ 위 스킴론적 올은 $\Spec(L)$임에 주의하자.
따라서 $X$가 $L$에서 비분기이면 체 확대 $L/K$는 분리 확대이다.
Morphisms, 보조정리 \ref{morphisms-lemma-unramified-over-field}를 보라.

\begin{lemma}
\label{pione-lemma-unramified-in-L}
위의 상황에서 다음 조건들은 동치이다.
\begin{enumerate}
\item $X$는 $L$에서 비분기이다.
\item $Y \to X$는 에탈이다.
\item $Y \to X$는 유한 에탈이다.
\end{enumerate}
\end{lemma}

\begin{proof}
$Y \to X$는 정수적 사상임에 주의하자. 각 경우에서 정의에 의해
사상 $Y \to X$는 국소 유한형이다. 따라서 각 경우에
Morphisms, 보조정리 \ref{morphisms-lemma-finite-integral}을 적용하면
$Y \to X$는 유한이다. 특히 (2)와 (3)은 동치이다.
에탈 사상은 비분기이므로 (2)는 (1)을 함의한다.

\medskip\noindent
거꾸로 $Y \to X$가 비분기라고 가정하자. 정규 스킴은 기하학적으로 일분기이므로
(Properties, 보조정리 \ref{properties-lemma-normal-geometrically-unibranch}),
More on Morphisms, 보조정리
\ref{more-morphisms-lemma-global-unramified-dominant-unibranch-is-etale}에 의해
사상 $Y \to X$는 에탈이다. 다음 문단에서는 직접적인 증명도 제시한다.

\medskip\noindent
$x \in X$라 하자. 다음 사상이 서로소 닫힌 몰입들의 합이 되게 하는
에탈 근방 $(U, u) \to (X, x)$를 택할 수 있다.
$$
Y \times_X U = \coprod V_j \longrightarrow U
$$
\'Etale Morphisms, 보조정리
\ref{etale-lemma-finite-unramified-etale-local}를 보라. $U$를 축소하여
준콤팩트라고 가정해도 된다. 그러면 $U$는 유한 개의 기약 성분을 갖는다
(Descent, 보조정리 \ref{descent-lemma-locally-finite-nr-irred-local-fppf}).
$U$는 정규이고(Descent, 보조정리
\ref{descent-lemma-normal-local-smooth}), $U$의 기약 성분들은 열린닫힌집합이므로
(Properties, 보조정리 \ref{properties-lemma-normal-locally-finite-nr-irreducibles}),
$U$가 기약이라고 가정해도 된다. 그러면 $U$는 정역 스킴이고 그 일반점
$\xi$는 $X$의 일반점으로 사상된다. 한편 More on Morphisms, 보조정리
\ref{more-morphisms-lemma-normalization-smooth-localization}에 의해
$Y \times_X U$는 $\Spec(L) \times_X U$에서의 $U$의 정규화이다.
$\Spec(L) \times_X U$의 모든 점은 $\xi$로 사상된다. 따라서
Morphisms, 보조정리 \ref{morphisms-lemma-normalization-generic}에 의해
각 $V_j$는 $\xi$로 사상되는 점을 포함한다. $U = \overline{\{\xi\}}$이므로
$V_j \to U$는 동형이다. 따라서 $Y \times_X U \to U$는 에탈이다.
Descent, 보조정리 \ref{descent-lemma-descending-property-etale}에 의해
$Y \to X$는 $U \to X$의 상, 즉 $x$의 열린 근방 위에서 에탈이다.
\end{proof}

\begin{lemma}
\label{pione-lemma-finite-etale-covering-normal-unramified}
$X$를 함수체 $K$를 갖는 정규 정역 스킴이라 하자.
$Y \to X$를 유한 에탈 사상이라 하자. $Y$가 연결이면 $Y$는 정규 정역 스킴이고,
$Y$는 $Y$의 함수체에서의 $X$의 정규화이다.
\end{lemma}

\begin{proof}
Descent, 보조정리 \ref{descent-lemma-normal-local-smooth}에 의해 스킴 $Y$는 정규이다.
$Y \to X$가 평탄하므로 Morphisms, 보조정리
\ref{morphisms-lemma-generalizations-lift-flat}에 의해 $Y$의 각 일반점은
$X$의 일반점으로 사상된다. $Y \to X$가 유한이므로 $Y$는 유한 개의
기약 성분을 갖는다. 따라서 Properties, 보조정리
\ref{properties-lemma-normal-locally-finite-nr-irreducibles}에 의해 $Y$는
유한 개의 정규 정역 스킴의 서로소 합이다. 그러므로 $Y$가 연결이면
$Y$는 정규 정역 스킴이다.

\medskip\noindent
$L$을 $Y$의 함수체라 하고 $Y' \to X$를 $L$에서의 $X$의 정규화라 하자.
Morphisms, 보조정리 \ref{morphisms-lemma-characterize-normalization}에 의해
분해 $Y' \to Y \to X$를 얻고, $Y' \to Y$는 $L$에서의 $Y$의 정규화이다.
$Y$가 정규이므로 $Y' = Y$임은 분명하다(이는 Morphisms, 보조정리
\ref{morphisms-lemma-finite-birational-over-normal}에서도 따른다).
\end{proof}

\begin{proposition}
\label{pione-proposition-normal}
$X$를 함수체 $K$를 갖는 정규 정역 스킴이라 하자. 그러면 표준 사상
(\ref{pione-equation-inclusion-generic-point})
$$
\text{Gal}(K^{sep}/K) = \pi_1(\eta, \overline{\eta})
\longrightarrow \pi_1(X, \overline{\eta})
$$
은 몫사상
$\text{Gal}(K^{sep}/K) \to \text{Gal}(M/K)$와 동일시된다. 여기서
$M \subset K^{sep}$는 $X$가 $L$에서 비분기인 유한 부분확대 $L$들의 합집합이다.
\end{proposition}

\begin{proof}
$X$는 정규 스킴이므로 기하학적으로 일분기이다
(Properties, 보조정리 \ref{properties-lemma-normal-geometrically-unibranch}).
따라서 보조정리 \ref{pione-lemma-irreducible-geometrically-unibranch}를 $X$에 적용할 수 있다.
그러므로 $\pi_1(\eta, \overline{\eta}) \to \pi_1(X, \overline{\eta})$는 전사이고,
다음 가환 도식의 위쪽 수평 화살표는 충실충만하다.
$$
\xymatrix{
\textit{F\'Et}_X \ar[r] \ar[d] \ar[rd]_c & \textit{F\'Et}_\eta \ar[d] \\
\textit{Finite-}\pi_1(X, \overline{\eta})\textit{-sets} \ar[r] &
\textit{Finite-}\text{Gal}(K^{sep}/K)\textit{-sets}
}
$$
왼쪽 수직 화살표는 정리 \ref{pione-theorem-fundamental-group}의 동치이고,
오른쪽 수직 화살표는 보조정리
\ref{pione-lemma-fundamental-group-Galois-group}의 동치이다.
아래쪽 수평 화살표는 명제의 사상에서 유도된다. 보조정리
\ref{pione-lemma-unramified-in-L} 및
\ref{pione-lemma-finite-etale-covering-normal-unramified}에 의해
$c$의 본질적 상은 다음 꼴의 집합과 동형인
$\text{Gal}(K^{sep}/K)\textit{-Sets}$들로 이루어진다.
$$
S = \Hom_K(\prod\nolimits_{i = 1, \ldots, n} L_i, K^{sep}) =
\coprod\nolimits_{i = 1, \ldots, n} \Hom_K(L_i, K^{sep})
$$
여기서 $L_i/K$는 유한 분리 확대이고 $X$는 $L_i$에서 비분기이다.
따라서 $M \subset K^{sep}$가 보조정리의 명제에서와 같다면
$\text{Gal}(K^{sep}/M)$은 $\text{Gal}(K^{sep}/K)$의 부분군 가운데
$c$의 본질적 상에 속하는 모든 대상에 자명하게 작용하는 부분군과 정확히 같다.
한편 $c$의 본질적 상은 $\text{Gal}(K^{sep}/K)$의 작용이 전사
$\text{Gal}(K^{sep}/K) \to \pi_1(X, \overline{\eta})$를 통해 분해되는
$S$들의 범주와 정확히 같다. 따라서
$\text{Gal}(K^{sep}/M)$은 그 핵이다. 그러므로
$\text{Gal}(K^{sep}/M)$은 정규 부분군이고 $M/K$는 갈루아 확대이며,
다음 짧은 완전열을 얻는다.
$$
1 \to \text{Gal}(K^{sep}/M) \to
\text{Gal}(K^{sep}/K) \to
\text{Gal}(M/K) \to 1
$$
이는 갈루아 이론(Fields, 정리
\ref{fields-theorem-inifinite-galois-theory} 및
보조정리 \ref{fields-lemma-ses-infinite-galois})에 따른다. 증명이 끝났다.
\end{proof}

\begin{lemma}
\label{pione-lemma-local-exact-sequence-normal}
$(A, \mathfrak m)$을 정규 국소환이라 하자. $X = \Spec(A)$라 하자.
$A^{sh}$를 $A$의 엄밀 헨젤화라 하자. $K$와 $K^{sh}$를 각각
$A$와 $A^{sh}$의 분수체라 하자. 그러면 열
$$
\pi_1(\Spec(K^{sh})) \to \pi_1(\Spec(K)) \to \pi_1(X) \to 1
$$
은 보조정리 \ref{pione-lemma-functoriality-galois-ses} (1)의 의미에서 완전하다.
\end{lemma}

\begin{proof}
$A^{sh}$가 정규 정역임에 주의하자. More on Algebra, 보조정리
\ref{more-algebra-lemma-henselization-normal}를 보라.
$\pi_1(\Spec(K)) \to \pi_1(X)$는 명제
\ref{pione-proposition-normal}에 의해 전사이다.

\medskip\noindent
$X^{sh} = \Spec(A^{sh})$라 쓰자. $Y \to X$를 유한 에탈 사상이라 하자.
그러면 $Y^{sh} = Y \times_X X^{sh} \to X^{sh}$도 유한 에탈 사상이다.
$A^{sh}$가 엄밀 헨젤 환이므로 $Y^{sh}$는 $X^{sh}$의 서로소 복사본들의 합과
동형이다. 따라서 $Y \times_X \Spec(K^{sh})$도 마찬가지이다. 그러므로 합성
$\pi_1(\Spec(K^{sh})) \to \pi_1(X)$는 자명하다.
보조정리 \ref{pione-lemma-composition-trivial}를 보라.

\medskip\noindent
증명을 끝내기 위해 보조정리 \ref{pione-lemma-functoriality-galois-ses}에 따라 다음을
보이면 충분하다. $V \to \Spec(K)$를 유한 에탈 사상이라 하고,
$V \times_{\Spec(K)} \Spec(K^{sh})$가 $\Spec(K^{sh})$의 서로소 복사본들의
합이라고 하자. 그러면 유한 에탈 사상 $Y \to X$를 찾아
$\Spec(K)$ 위에서 $V \cong Y \times_X \Spec(K)$가 되게 할 수 있다.
$V = \Spec(L)$이라 쓰면 $L$은 $K$의 유한 분리 확대들의 유한곱이다.
$B \subset L$을 $L$에서의 $A$의 정수적 폐포라 하자.
$A \to B$가 에탈이면 $Y = \Spec(B)$로 잡아 증명을 끝낼 수 있다.
Algebra, 보조정리
\ref{algebra-lemma-integral-closure-commutes-smooth}와 생략한 극한 논증에 의해
$B \otimes_A A^{sh}$는
$L^{sh} = L \otimes_K K^{sh}$에서의 $A^{sh}$의 정수적 폐포이다.
가정에 의해 $L^{sh}$는 $K^{sh}$의 복사본들의 곱이고, 따라서
$B^{sh}$는 $A^{sh}$의 복사본들의 곱이다. 그러므로
$A^{sh} \to B^{sh}$는 에탈이다. $A \to A^{sh}$가 충실평탄하므로
Descent, 보조정리 \ref{descent-lemma-descending-property-etale}에 의해
$A \to B$는 에탈이며, 원하는 결론을 얻는다.
\end{proof}






\section{군 작용과 정수적 폐포}
\label{pione-section-group-actions-integral}

\noindent
이 절에서는 More on Algebra, 절
\ref{more-algebra-section-group-actions-integral}의 논의를 계속한다.
정의에 의해 정규 국소환은 정역임을 상기하자.

\begin{lemma}
\label{pione-lemma-get-algebraic-closure}
분수체 $K$가 분리대수적으로 닫힌 정규 정역 $A$를 택하자.
$\mathfrak p \subset A$를 영이 아닌 소 아이디얼이라 하자. 그러면 잉여체
$\kappa(\mathfrak p)$는 대수적으로 닫혀 있다.
\end{lemma}

\begin{proof}
보조정리가 거짓이라고 가정하여 모순을 얻자. 그러면 차수가 $d > 1$인
기약 일계수 다항식 $P(T) \in \kappa(\mathfrak p)[T]$가 존재한다.
분모를 없애기 위해 적절한 $a \in A$에 대하여 $P$를
$a^d P(a^{-1}T)$로 바꾼 뒤, $P$가 $A[T]$의 일계수 다항식
$Q$의 상이라고 가정할 수 있다. $Q$는 $K[T]$에서 기약임에 주의하자.
실제로 $K$ 위에서의 인수분해가 존재하면 Algebra, 보조정리
\ref{algebra-lemma-polynomials-divide}에 의해 $A$ 위에서의 인수분해가 생기고,
이를 $\mathfrak p$로 나누어 나머지를 취하면 $P$의 인수분해가 생긴다.
$K$가 분리적으로 닫혀 있으므로 $Q$는 분리 다항식이 아니다
(Fields, 정의 \ref{fields-definition-separable}). 따라서 $K$의 표수는
$p > 0$이고 $Q$의 일차항은 영이다
(Fields, 정의 \ref{fields-definition-separable}). 그러나 이때
$Q$를 $Q + a T$로 바꾸면 모순을 얻는다. 여기서
$a \in \mathfrak p$는 영이 아니다.
\end{proof}

\begin{lemma}
\label{pione-lemma-normal-local-domain-separablly-closed-fraction-field}
분리적으로 닫힌 분수체를 갖는 정규 국소환은 엄밀 헨젤 환이다.
\end{lemma}

\begin{proof}
$(A, \mathfrak m, \kappa)$를 분리적으로 닫힌 분수체 $K$를 갖는
정규 국소환이라 하자. $A = K$이면 끝났다. 그렇지 않으면 보조정리
\ref{pione-lemma-get-algebraic-closure}에 의해 잉여체 $\kappa$는 대수적으로 닫혀 있으므로,
$A$가 헨젤 환임을 확인하면 충분하다.
$f \in A[T]$를 일계수 다항식이라 하고, $a_0 \in \kappa$를
환원 $\overline{f} \in \kappa[T]$의 중복도 $1$인 근이라 하자.
$f = \prod f_i$를 $K[T]$에서의 인수분해라 하자.
Algebra, 보조정리 \ref{algebra-lemma-polynomials-divide}에 의해
$f_i \in A[T]$이다. 따라서 어떤 $i$에 대하여 $a_0$는 $f_i$의 근이다.
$f$를 $f_i$로 바꾸어 $f$가 기약이라고 가정해도 된다.
$a_0$가 단근이므로 도함수 $f'$는 $A[T]$에서 영일 수 없다.
$K$가 분리대수적으로 닫혀 있으므로 이로부터 $f$는 일차임을 얻는다.
따라서 $A$는 헨젤 환이다. Algebra, 정의
\ref{algebra-definition-henselian}을 보라.
\end{proof}

\begin{lemma}
\label{pione-lemma-inertia-base-change}
유한군 $G$가 환 $R$에 작용한다고 하자. $R^G \to A$를 환 준동형이라 하자.
$\mathfrak q' \subset A \otimes_{R^G} R$를 소 아이디얼
$\mathfrak q \subset R$ 위에 놓인 소 아이디얼이라 하자. 그러면
$$
I_\mathfrak q = \{\sigma \in G \mid
\sigma(\mathfrak q) = \mathfrak q\text{ 및 }
\sigma \bmod \mathfrak q = \text{id}_{\kappa(\mathfrak q)}\}
$$
는
$$
I_{\mathfrak q'} = \{\sigma \in G \mid
\sigma(\mathfrak q') = \mathfrak q'\text{ 및 }
\sigma \bmod \mathfrak q' = \text{id}_{\kappa(\mathfrak q')}\}
$$
와 같다.
\end{lemma}

\begin{proof}
$\mathfrak q$는 $\mathfrak q'$의 역상이고
$\kappa(\mathfrak q) \subset \kappa(\mathfrak q')$이므로
$I_{\mathfrak q'} \subset I_\mathfrak q$이다.
거꾸로 $\sigma \in I_\mathfrak q$이면 $\sigma$는 올환
$A \otimes_{R^G} \kappa(\mathfrak q)$에 자명하게 작용한다.
따라서 $\sigma$는 $\mathfrak q$ 위에 놓인 모든 소 아이디얼을 고정하고,
각각의 잉여체 위에서 항등사상을 유도한다.
\end{proof}

\begin{lemma}
\label{pione-lemma-inertia-invariants-etale}
유한군 $G$가 환 $R$에 작용한다고 하자. $\mathfrak q \subset R$를 소 아이디얼이라 하고
다음과 같이 놓자.
$$
I = \{\sigma \in G \mid \sigma(\mathfrak q) = \mathfrak q
\text{ 및 } \sigma \bmod \mathfrak q = \text{id}_\mathfrak q\}
$$
그러면 $R^G \to R^I$는 $R^I \cap \mathfrak q$에서 에탈이다.
\end{lemma}

\begin{proof}
증명의 전략은 에탈 국소화를 사용하여 $R \to R^I$가
$R^I \cap \mathfrak p$에서 국소 동형인 경우로 환원하는 것이다.
$R^G \to A$를 에탈 환 준동형이라 하자. $G$가
$A \otimes_{R^G} R$에 작용하는 경우와
$\mathfrak q$ 위에 놓인 $A \otimes_{R^G} R$의 어떤 소 아이디얼
$\mathfrak q'$에 대하여 결과가 성립하면 원래 결과도 성립한다고 주장한다.

\medskip\noindent
이를 확인하자. $R^G \to A$가 평탄하므로
$A = (A \otimes_{R^G} R)^G$이다. More on Algebra, 보조정리
\ref{more-algebra-lemma-base-change-invariants}를 보라.
보조정리 \ref{pione-lemma-inertia-base-change}에 의해 군 $I$는 바뀌지 않는다.
다시 More on Algebra, 보조정리
\ref{more-algebra-lemma-base-change-invariants}을 적용하면
$A \otimes_{R^G} R^I = (A \otimes_{R^G} R)^I$를 얻는다
($R^I \to A \otimes_{R^G} R^I$가 평탄하기 때문이다).
따라서
$$
\xymatrix{
\Spec((A \otimes_{R^G} R)^I) \ar[d] \ar[r] & \Spec(R^I) \ar[d] \\
\Spec(A) \ar[r] & \Spec(R^G)
}
$$
은 데카르트 도식이고 수평 화살표들은 에탈이다. 그러므로 왼쪽 수직 화살표가
$(A \otimes_{R^G} R)^I \cap \mathfrak q'$를 포함하는 어떤 열린 근방 $W$에서
에탈이면, 오른쪽 수직 화살표는 $\Spec(R^I)$에서 $W$의 열린 상에 속하는 점들에서
에탈이다. Descent, 보조정리
\ref{descent-lemma-smooth-permanence}를 보라. 특히 사상
$\Spec(R^I) \to \Spec(R^G)$는 $R^I \cap \mathfrak q$에서 에탈이다.

\medskip\noindent
$\mathfrak p = R^G \cap \mathfrak q$라 하자.
More on Algebra, 보조정리 \ref{more-algebra-lemma-one-orbit}에 의해
$\mathfrak p$ 위에서의 $\Spec(R) \to \Spec(R^G)$의 올은 유한하다.
더욱이 이 점들에서의 잉여체 확대들은 대수적이고 정규이며,
그 자기동형군들은 유한하다. 이는 More on Algebra, 보조정리
\ref{more-algebra-lemma-one-orbit-geometric}에 따른다. 따라서
More on Morphisms, 보조정리
\ref{more-morphisms-lemma-etale-makes-integral-split}을 정수적 환 준동형
$R^G \to R$와 소 아이디얼 $\mathfrak p$에 적용할 수 있다.
위의 주장과 합치면, $R = A_1 \times \ldots \times A_n$인 경우로 환원된다.
각 $A_i$는 유일한 소아이디얼 $\mathfrak q_i$를 가지며, 이는 $\mathfrak p$ 위에 놓이고,
잉여체 확대 $\kappa(\mathfrak q_i)/\kappa(\mathfrak p)$는 순수 비분리이다.
물론 $\mathfrak q$는 이 소아이디얼들 가운데 하나이므로,
$\mathfrak q = \mathfrak q_1$이라 해도 좋다.

\medskip\noindent
$G$가 인자 $A_i$들을 치환하지 않을 수도 있다(예를 들어 $A_i$의 스펙트럼이
연결이고 $R^G$가 국소환이면 치환한다). 이는 다음과 같이 바로잡을 수 있다.
독자에게 이 과정을 직접 생각해 보기를 권하며, 위상 대신 멱등원을 사용해도 좋다.
곱 분해는 $\Spec(R)$를 열린닫힌 부분집합 $U_i$들로 나타내는 서로소 합 분해를
준다는 점을 상기하자. $G$가 유한이므로 이 덮개를 유한 서로소 합 분해
$\Spec(R) = \coprod_{j \in J} W_j$로 세분할 수 있다. 여기서 열린닫힌
부분집합 $W_j$들에 대해, 각 $j \in J$마다
어떤 $j' \in J$가 존재하여 $\sigma(W_j) = W_{j'}$가 되게 할 수 있다.
이 $W_j$들 가운데 $\{\mathfrak q_1, \ldots, \mathfrak q_n\}$와 만나지 않는 것들의 합집합은
$\mathfrak p$ 위의 올과 만나지 않는 닫힌집합이다. 따라서 이 합집합은
$\Spec(R^G)$에서 $\mathfrak p$와 만나지 않는 닫힌집합으로 사상된다.
이는 $\Spec(R) \to \Spec(R^G)$가 닫힌 사상이기 때문이다.
그러므로 $R^G$를 주국소화로 바꾸고 나면(위의 주장에 의해 허용된다),
각 $W_j$가 어떤 $\mathfrak q_i$와 만난다고 가정할 수 있다.
이제 $U_i = W_j$로 둔다. 단, 이는 $\mathfrak q_i \in W_j$일 때이다.
이에 대응하는 곱 분해 $R = A_1 \times \ldots \times A_n$에서는
$G$가 인자 $A_i$들을 치환한다.

\medskip\noindent
따라서 곱 분해
$R = A_1 \times \ldots \times A_n$가 $G$-작용과 양립하고, 각 $A_i$는
유일한 소아이디얼 $\mathfrak q_i$를 가지며, 이는 $\mathfrak p$ 위에 놓이고,
체 확대 $\kappa(\mathfrak q_i)/\kappa(\mathfrak p)$가 순수 비분리라고 가정할 수 있다.
$A' = A_2 \times \ldots \times A_n$이라 쓰면
$$
R = A_1 \times A'
$$
이다. $\mathfrak q = \mathfrak q_1$이므로 모든
$\sigma \in I$는 위의 곱 분해를 보존한다. 따라서
$$
R^I = (A_1)^I \times (A')^I
$$
이다. $I = D = \{\sigma \in G \mid \sigma(\mathfrak q) = \mathfrak q\}$임에 유의하자.
이는 $\kappa(\mathfrak q)/\kappa(\mathfrak p)$가 순수 비분리이기 때문이다.
$G$의 $\mathfrak p$ 위 소아이디얼들에 대한 작용은 추이적이다
(More on Algebra, 보조정리 \ref{more-algebra-lemma-one-orbit}).
따라서 $I$의 $G$에서의 지수는 $n$이고,
$G = eI \amalg \sigma_2I \amalg \ldots \amalg \sigma_nI$로 쓸 수 있으며,
$A_i = \sigma_i(A_1)$이다($i = 2, \ldots, n$). 그러므로
$$
R^G = (A_1)^I.
$$
사상 $R^G \to R^I$는 $R^I \cap \mathfrak q$에서 에탈이고, 이로써 증명이 끝난다.
\end{proof}

\noindent
다음 보조정리는 More on Algebra, 보조정리
\ref{more-algebra-lemma-inertial-invariants-unramified}를 일반화한다.

\begin{lemma}
\label{pione-lemma-inertial-invariants-unramified}
$A$를 분수체 $K$를 갖는 정규 정역이라 하자. $L/K$를 (무한일 수도 있는)
갈루아 확대라 하자. $G = \text{Gal}(L/K)$로 두고, $B$를
$A$의 $L$에서의 정수적 폐포라 하자. $\mathfrak q \subset B$라 하자. 다음과 같이 둔다.
$$
I = \{\sigma \in G \mid
\sigma(\mathfrak q) = \mathfrak q \text{ 이고 }
\sigma \bmod \mathfrak q = \text{id}_{\kappa(\mathfrak q)}\}
$$
그러면 $(B^I)_{B^I \cap \mathfrak q}$는 에탈 $A$-대수들의
여과 여극한이다.
\end{lemma}

\begin{proof}
$L$을 $K$의 유한 갈루아 확대들의 여과 여극한으로 쓸 수 있다.
따라서 $L/K$가 유한 갈루아 확대인 경우에 이 보조정리를 증명하면 충분하다.
Algebra, 보조정리 \ref{algebra-lemma-colimit-colimit-etale}를 보라.
$A = B^G$이다. 실제로 $A$는 $K = L^G$에서 정수적으로 닫혀 있다.
따라서 결과는 보조정리 \ref{pione-lemma-inertia-invariants-etale}에서 따른다.
\end{proof}








\section{분기 이론}
\label{pione-section-ramification}

\noindent
이 절에서는 More on Algebra, 절 \ref{more-algebra-section-ramification}의 논의를
계속하여 스킴의 기본군에 관한 논의와 연결한다.

\medskip\noindent
$(A, \mathfrak m, \kappa)$를 분수체 $K$를 갖는 정규 국소환이라 하자.
분리 대수적 폐포 $K^{sep}$를 택한다. $A^{sep}$를 $A$의 $K^{sep}$에서의 정수적 폐포라 하자. 극대 아이디얼
$\mathfrak m^{sep} \subset A^{sep}$를 하나 택한다.
$A \subset A^h \subset A^{sh}$를 각각 헨젤화와 엄밀 헨젤화라 하자.
$A^h$와 $A^{sh}$도 정규환임에 유의하자
(More on Algebra, 보조정리 \ref{more-algebra-lemma-henselization-normal}).
각각의 분수체를 $K^h$와 $K^{sh}$로 쓴다.
보조정리 \ref{pione-lemma-normal-local-domain-separablly-closed-fraction-field}에 의해
$(A^{sep})_{\mathfrak m^{sep}}$는 엄밀 헨젤 환이므로,
$A$-대수 준동형 $A^{sh} \to (A^{sep})_{\mathfrak m^{sep}}$를 택할 수 있다.
구체적으로 먼저 $\kappa$-몰입\footnote{이는
$\kappa(\mathfrak m^{sh})$가 $\kappa$의 분리 대수적 폐포이고,
$\kappa(\mathfrak m^{sep})$가 보조정리 \ref{pione-lemma-get-algebraic-closure}에 의해
$\kappa$의 대수적 폐포이기 때문에 가능하다.}
$\kappa(\mathfrak m^{sh}) \to \kappa(\mathfrak m^{sep})$를 택한 다음,
Algebra, 보조정리 \ref{algebra-lemma-strictly-henselian-functorial}에 의해
이를 $A$-대수 준동형으로 (유일하게) 확장한다.
다음 도식을 얻는다.
$$
\xymatrix{
K^{sep} & K^{sh} \ar[l] & K^h \ar[l] & K \ar[l] \\
(A^{sep})_{\mathfrak m^{sep}} \ar[u] &
A^{sh} \ar[u] \ar[l] &
A^h \ar[u] \ar[l] &
A \ar[u] \ar[l]
}
$$
이 환들의 스펙트럼의 기본군을 생각할 수 있다.
물론 $K^{sep}$, $(A^{sep})_{\mathfrak m^{sep}}$, $A^{sh}$는 엄밀 헨젤이므로
이들에서는 자명군을 얻는다. 따라서 흥미로운 부분은 다음과 같다.
\begin{equation}
\label{pione-equation-inertia-diagram-pione}
\vcenter{
\xymatrix{
\pi_1(U^{sh}) \ar[r] \ar[rd]_1 & \pi_1(U^h) \ar[d] \ar[r] & \pi_1(U) \ar[d] \\
& \pi_1(X^h) \ar[r] & \pi_1(X)
}
}
\end{equation}
여기서 $X^h$와 $X$는 각각 $A^h$와 $A$의 스펙트럼이고,
$U^{sh}$, $U^h$, $U$는 각각 $K^{sh}$, $K^h$, $K$의 스펙트럼이다.
표지 $1$은 사상이 자명하다는 뜻이다. 이는 사상이 자명군
$\pi_1(X^{sh})$를 경유하기 때문이다.
한편 프로유한군 $G = \text{Gal}(K^{sep}/K)$는 $A^{sep}$에 작용하므로
다음을 정의한다.
$$
D = \{\sigma \in G \mid \sigma(\mathfrak m^{sep}) = \mathfrak m^{sep}\}
\supset
I = \{\sigma \in D \mid \sigma \bmod \mathfrak m^{sep} =
\text{id}_{\kappa(\mathfrak m^{sep})}\}
$$
이 군들은 특히 $A$가 이산 부치환환일 때 각각
{\it 분해군}과 {\it 관성군}이라 불리기도 한다.

\begin{lemma}
\label{pione-lemma-identify-inertia}
위의 상황에서 동형
$\pi_1(U) = \text{Gal}(K^{sep}/K)$를 통하여 도식
(\ref{pione-equation-inertia-diagram-pione})은 다음 도식으로 옮겨진다.
$$
\xymatrix{
I \ar[r] \ar[rd]_1 & D \ar[d] \ar[r] & \text{Gal}(K^{sep}/K) \ar[d] \\
& \text{Gal}(\kappa(\mathfrak m^{sh})/\kappa) \ar[r] & \text{Gal}(M/K)
}
$$
여기서 $K^{sep}/M/K$는 $A$에 관해 비분기인 최대 부분확대이다.
또한 수직 화살표들은 전사이고, 왼쪽 수직 화살표의 핵은 $I$이며,
오른쪽 수직 화살표의 핵은
$\text{Gal}(K^{sep}/K)$에서 $I$를 포함하는 최소 닫힌 정규부분군이다.
\end{lemma}

\begin{proof}
구성에 의해 군 $D$는 $(A^{sep})_{\mathfrak m^{sep}}$에 $A$ 위에서 작용한다.
잉여체 위의 사상을 지정했을 때의
$A^{sh} \to (A^{sep})_{\mathfrak m^{sep}}$의 유일성
(Algebra, 보조정리 \ref{algebra-lemma-strictly-henselian-functorial})으로부터,
$A^{sh} \to (A^{sep})_{\mathfrak m^{sep}}$의 상은
$((A^{sep})_{\mathfrak m^{sep}})^I$에 포함된다.
한편 보조정리 \ref{pione-lemma-inertial-invariants-unramified}에 의해
$((A^{sep})_{\mathfrak m^{sep}})^I$는 $A$의 에탈 확대들의 여과 여극한이다.
$A^{sh}$는 그러한 최대 확대이므로
$A^{sh} = ((A^{sep})_{\mathfrak m^{sep}})^I$라고 결론내릴 수 있다.
따라서 $K^{sh} = (K^{sep})^I$이다.

\medskip\noindent
$I$는 전사
$D \to \text{Aut}(\kappa(\mathfrak m^{sep})/\kappa)$의 핵임을 상기하자.
More on Algebra, 보조정리 \ref{more-algebra-lemma-one-orbit-geometric-galois}를 보라.
위에서 보았듯이
$\text{Aut}(\kappa(\mathfrak m^{sep})/\kappa) =
\text{Gal}(\kappa(\mathfrak m^{sh})/\kappa)$이고, 이 체들은 각각
$\kappa$의 대수적 폐포와 분리 대수적 폐포이다.
한편 Algebra, 보조정리 \ref{algebra-lemma-henselian-functorial}의 유일성에 의해
$A^{sh}$의 $A$ 위의 모든 자기동형은 $A^{sh}$의 $A^h$ 위의 자기동형이다.
더 나아가 $A^{sh}$는 유한 에탈 확대 $A^h \subset A'$들의 여극한이며,
이 확대들은 $1$ 대 $1$로 유한 분리 확대 $\kappa'/\kappa$와 대응한다.
Algebra, 주의 \ref{algebra-remark-construct-sh-from-h}를 보라. 따라서
$$
\text{Aut}(A^{sh}/A) = \text{Aut}(A^{sh}/A^h) =
\text{Gal}(\kappa(\mathfrak m^{sh})/\kappa)
$$
$\kappa'/\kappa$를 갈루아군 $G$를 갖는 유한 갈루아 확대라 하자.
Algebra, 보조정리 \ref{algebra-lemma-henselian-cat-finite-etale}에 의해
유한 에탈 확대 $A^h \subset A'$를 택한다. 이는
$\kappa \subset \kappa'$에 대응한다.
분수체와 차수를 살펴보면 $(A')^G = A^h$가 따른다
(작은 세부사항은 생략한다). 여극한을 취하면
$(A^{sh})^{\text{Gal}(\kappa(\mathfrak m^{sh})/\kappa)} = A^h$라고 결론내릴 수 있다.
이상을 모두 합치면 $A^h = ((A^{sep})_{\mathfrak m^{sep}})^D$를 얻는다.
따라서 $K^h = (K^{sep})^D$이다.

\medskip\noindent
$U$, $U^h$, $U^{sh}$는 각각 체 $K$, $K^h$, $K^{sh}$의 스펙트럼이므로,
보조정리 \ref{pione-lemma-fundamental-group-Galois-group}에 의해 두 도식의 윗줄들이 대응한다.
보조정리 \ref{pione-lemma-gabber}에 의해
$\pi_1(X^h) = \text{Gal}(\kappa(\mathfrak m^{sh})/\kappa)$이다.
열
$1 \to I \to D \to \text{Gal}(\kappa(\mathfrak m^{sh})/\kappa) \to 1$
이 완전하다는 것은 위에서 지적했다.
명제 \ref{pione-proposition-normal}에 의해
$\pi_1(X) = \text{Gal}(M/K)$이다.
마지막으로 사상
$\text{Gal}(K^{sep}/K) \to \text{Gal}(M/K) = \pi_1(X)$의 핵에 관한 주장은
보조정리 \ref{pione-lemma-local-exact-sequence-normal}에서 따른다.
이로써 증명이 끝난다.
\end{proof}

\noindent
$X$를 함수체 $K$를 갖는 정규 정역 스킴이라 하자.
$K^{sep}$를 $K$의 분리 대수적 폐포라 하자.
$X^{sep} \to X$를 $X$의 $K^{sep}$에서의 정규화라 하자.
$G = \text{Gal}(K^{sep}/K)$가 $K^{sep}$에 작용하므로,
$G$의 오른쪽 작용을 $X^{sep}$ 위에 얻는다. $y \in X^{sep}$에 대해 다음과 같이 정의한다.
$$
D_y = \{\sigma \in G \mid \sigma(y) = y\} \supset
I_y = \{\sigma \in D \mid \sigma \bmod \mathfrak m_y =
\text{id}_{\kappa(y)} \}
$$
이는 위와 같다. 한편 $x \in X$에 대해
$\mathcal{O}_{X, x}^{sh}$를 엄밀 헨젤화라 하고,
$K_x^{sh}$를 $\mathcal{O}_{X, x}^{sh}$의 분수체라 하며,
$K$-몰입 $K_x^{sh} \to K^{sep}$를 하나 택한다.

\begin{lemma}
\label{pione-lemma-normal-pione-quotient-inertia}
$X$를 함수체 $K$를 갖는 정규 정역 스킴이라 하자.
위의 표기 아래 다음 세
$\text{Gal}(K^{sep}/K) = \pi_1(\Spec(K))$의 부분군은 서로 같다.
\begin{enumerate}
\item 전사 $\text{Gal}(K^{sep}/K) \longrightarrow \pi_1(X)$의 핵.
\item $I_y$를 포함하는 최소 닫힌 정규부분군(모든 $y \in X^{sep}$에 대해).
\item $\text{Gal}(K^{sep}/K_x^{sh})$를 포함하는 최소 닫힌 정규부분군
(모든 $x \in X$에 대해).
\end{enumerate}
\end{lemma}

\begin{proof}
보조정리 \ref{pione-lemma-identify-inertia}에 의해 (2)와 (3)은 동치이다.
그 보조정리는 $I_y$가 $\text{Gal}(K^{sep}/K_x^{sh})$와 켤레임을 말하며,
이는 $y$가 $x$ 위에 있을 때이다.
보조정리 \ref{pione-lemma-local-exact-sequence-normal}에 의해
$\text{Gal}(K^{sep}/K_x^{sh})$는
$\pi_1(\Spec(\mathcal{O}_{X, x}))$로 자명하게 사상된다.
따라서 (2)와 (3)의 부분군
$N \subset G = \text{Gal}(K^{sep}/K)$는
$G \longrightarrow \pi_1(X)$의 핵에 포함된다.

\medskip\noindent
반대 포함관계를 보이자. $N$은 정규이므로, $N \subset U \subset G$이고
$U$가 열린 정규부분군일 때 몫사상 $G \to G/U$가
$\pi_1(X)$를 경유함을 보이면 충분하다.
달리 말해, $L/K$가 $U$에 대응하는 갈루아 확대라면
$X$가 $L$에서 비분기임을 보여야 한다
(절 \ref{pione-section-normal}, 특히 명제 \ref{pione-proposition-normal}).
$X$가 아핀인 경우만 보이면 충분하다
(증명의 나머지에서 대수학 결과들을 인용할 수 있도록 이렇게 한다).
$Y \to X$를 $X$의 $L$에서의 정규화라 하자.
포함 $L \subset K^{sep}$는 사상 $\pi : X^{sep} \to Y$를 유도한다.
$y \in X^{sep}$에 대해 $\pi(y)$의
$\text{Gal}(L/K)$에서의 관성군은 $I_y$의
$\text{Gal}(L/K)$에서의 상이다. 이는 More on Algebra, 보조정리
\ref{more-algebra-lemma-one-orbit-geometric-galois-compare}에서 따른다.
$N \subset U$이므로 이 모든 관성군은 자명하다.
보조정리 \ref{pione-lemma-inertia-invariants-etale}를 적용하면
$Y \to X$가 에탈임을 얻는다.
(다른 방법으로 보조정리 \ref{pione-lemma-local-exact-sequence-normal}을 사용하여
$Y$를 $\Spec(\mathcal{O}_{X, x})$로 당긴 것이
모든 $x \in X$에 대해 에탈임을 보인 뒤, 약간의 작업을 더하여 결론을 내릴 수도 있다.)
\end{proof}

\begin{example}
\label{pione-example-bigger-codim}
$X$를 함수체 $K$를 갖는 정규 정역 뇌터 스킴이라 하자.
분기 자취의 순수성(아래를 보라)에 의해 $X$가 정칙이면
보조정리 \ref{pione-lemma-normal-pione-quotient-inertia}에서
관성군 $I = \pi_1(\Spec(K_x^{sh}))$만 점 $x$, 즉 여차원이 $1$인
$X$의 점에 대해 생각하면 충분하다.
그러나 일반적으로는 이것만으로 충분하지 않다. 실제로
$Y = \mathbf{A}_k^n = \Spec(k[t_1, \ldots, t_n])$이라 하고,
$k$를 표수가 $2$가 아닌 체라 하자.
$G = \{\pm 1\}$을 위수가 $2$인 군이라 하고, 좌표들에 대한 곱셈으로
$Y$에 작용하게 하며,
$$
X = \Spec(k[t_it_j, i, j \in \{1, \ldots, n\}])
$$
로 둔다. 몰입 $k[t_it_j] \subset k[t_1, \ldots, t_n]$은
차수 $2$인 사상 $Y \to X$를 정한다. 이 사상은 극대 아이디얼
$\mathfrak m = (t_it_j)$ 위를 제외하면 모든 곳에서 비분기이다.
이 아이디얼은 여차원이 $n$인 $X$의 점이다.
\end{example}

\begin{lemma}
\label{pione-lemma-unramified}
$X$를 함수체 $K$를 갖는 정규 정역 스킴이라 하자.
$L/K$를 유한 확대라 하고, $Y \to X$를 $X$의 $L$에서의 정규화라 하자.
다음 조건들은 동치이다.
\begin{enumerate}
\item $X$가 절 \ref{pione-section-normal}의 의미에서 $L$에서 비분기이다.
\item $Y \to X$가 스킴의 비분기 사상이다.
\item $Y \to X$가 스킴의 에탈 사상이다.
\item $Y \to X$가 유한 에탈 사상이다.
\item $x \in X$에 대한 사상
$Y \times_X \Spec(\mathcal{O}_{X, x}) \to \Spec(\mathcal{O}_{X, x})$가 비분기이다.
\item (5)와 같되 $\mathcal{O}_{X, x}^h$를 사용한다.
\item (5)와 같되 $\mathcal{O}_{X, x}^{sh}$를 사용한다.
\item $x \in X$에 대한 스킴론적 올 $Y_x$가
$x$ 위에서 에탈이고 차수가 $\geq [L : K]$이다.
\end{enumerate}
더 나아가 $L/K$가 갈루아군 $G$를 갖는 갈루아 확대이면,
이 조건들은 다음 조건과도 동치이다.
\begin{enumerate}
\item[(9)] $y \in Y$에 대해 군
$I_y = \{g \in G \mid g(y) = y\text{ 이고 }
g \bmod \mathfrak m_y = \text{id}_{\kappa(y)}\}$가 자명하다.
\end{enumerate}
\end{lemma}

\begin{proof}
 (1)과 (2)의 동치성은 (1)의 정의이다.
(2), (3), (4)의 동치성은 보조정리 \ref{pione-lemma-unramified-in-L}에 따른다.
(4) $\Rightarrow$ (5), (5) $\Rightarrow$ (6), (6) $\Rightarrow$ (7)은 쉽게 증명된다.

\medskip\noindent
 (7)을 가정하자. $\mathcal{O}_{X, x}^{sh}$는 정규 국소 정역임에 유의하자
(More on Algebra, 보조정리 \ref{more-algebra-lemma-henselization-normal}).
$L^{sh} = L \otimes_K K_x^{sh}$로 둔다. 여기서 $K_x^{sh}$는
$\mathcal{O}_{X, x}^{sh}$의 분수체이다. 그러면
$L^{sh} = \prod_{i = 1, \ldots, n} L_i$이고, $L_i/K_x^{sh}$는 유한 분리 확대이다.
Algebra, 보조정리 \ref{algebra-lemma-integral-closure-commutes-smooth}
(그리고 여기서는 생략하는 극한 논증)에 의해
$Y \times_X \Spec(\mathcal{O}_{X, x}^{sh})$는
$\Spec(\mathcal{O}_{X, x}^{sh})$의 $L^{sh}$에서의 정수적 폐포이다.
따라서 보조정리 \ref{pione-lemma-unramified-in-L}을 인자
$L_i$, 즉 $L^{sh}$의 인자들에 적용하면
$Y \times_X \Spec(\mathcal{O}_{X, x}^{sh}) \to \Spec(\mathcal{O}_{X, x}^{sh})$는
유한 에탈이다. 일반점을 살펴보면 그 차수는 $[L : K]$와 같으므로 (8)이 성립한다.

\medskip\noindent
 (8)을 가정하자. $x \in X$라 하고, 스킴론적 올 $Y_x$가
$x$ 위에서 에탈이며 차수가 $\geq [L : K]$라고 하자. 이는 $Y$가
$\geq [L : K]$개의 기하점을 $x$ 위에 갖는다는 뜻이다.
$Y \to X$가 $x$의 어떤 근방 위에서 유한 에탈임을 보이겠다. 그러면 (1)이 따른다.
증명을 위해 $X = \Spec(R)$이고 점 $x$가 소아이디얼
$\mathfrak p \subset R$에 대응하며 $Y = \Spec(S)$라고 가정해도 좋다.
More on Morphisms, 보조정리
\ref{more-morphisms-lemma-etale-makes-integral-split}을 적용하면,
에탈 근방 $(U, u) \to (X, x)$로서
$Y \times_X U = V_1 \amalg \ldots \amalg V_m$이고 각 $V_i$가
유일한 점 $v_i$를 가지며, 이 점은 $u$ 위에 놓이고
$\kappa(v_i)/\kappa(u)$가 순수 비분리인 것을 얻는다.
필요하면 $U$를 축소하여 $U$가 일반점 $\xi$를 갖는 정규 정역 스킴이라고
가정할 수 있다(Descent, 보조정리
\ref{descent-lemma-locally-finite-nr-irred-local-fppf},
\ref{descent-lemma-normal-local-smooth}, 그리고 Properties, 보조정리
\ref{properties-lemma-normal-locally-finite-nr-irreducibles}를 사용한다).
기하점에 관해 위에서 한 관찰로부터 $m \geq [L : K]$이다.
한편 More on Morphisms, 보조정리
\ref{more-morphisms-lemma-normalization-smooth-localization}에 의해
$\coprod V_i \to U$는 $U$의 $\Spec(L) \times_X U$에서의 정규화이다.
$K \subset \kappa(\xi)$가 유한 분리 확대이므로
$\Spec(L) \times_X U = \Spec(\prod_{i = 1, \ldots, n} L_i)$로 쓸 수 있으며,
$L_i/\kappa(\xi)$는 유한이고 $[L : K] = \sum [L_i : \kappa(\xi)]$이다.
$V_j$가 각 $j$에 대해 공집합이 아니고 $m \geq [L : K]$이므로
$m = n$이고 $[L_i : \kappa(\xi)] = 1$이라고 결론내릴 수 있다
(모든 $i$에 대해).
따라서 $V_j \to U$는 동형, 특히 에탈이고, $Y \times_X U \to U$도 에탈이다.
Descent, 보조정리 \ref{descent-lemma-descending-property-etale}에 의해
$Y \to X$는 $U \to X$의 상, 즉 $x$의 열린 근방 위에서 에탈이다.

\medskip\noindent
 $L/K$가 갈루아 확대이고 (9)가 성립한다고 가정하자. 그러면 보조정리
\ref{pione-lemma-inertial-invariants-unramified}에 의해 $Y \to X$는 에탈이다.
(1)이 (9)를 함의한다는 증명은 생략한다.
\end{proof}

\noindent
이산 부치환환의 무한 갈루아 확대의 경우에는 조금 더 상세히 말할 수 있다.
이를 위해 다음 표기를 도입한다. 정수의 부분집합 $S \subset \mathbf{N}$가
{\it 곱셈적으로 유향}이라는 것은 $1 \in S$이고, $n, m \in S$에 대해
어떤 $k \in S$가 존재하여 $n | k$와 $m | k$를 만족한다는 뜻이다.
$S$ 위의 부분순서를 $n \geq_S m$이 $m | n$과 동치가 되도록 정의한다.
체 $\kappa$가 주어지면 유한군들의 역계
$\{\mu_n(\kappa)\}_{n \in S}$와 전이 사상
$$
\mu_n(\kappa) \longrightarrow \mu_m(\kappa),\quad
\zeta \longmapsto \zeta^{n/m}
$$
을 얻는다. 여기서 $n \geq_S m$이다. 따라서 프로유한군
$$
\lim_{n \in S} \mu_n(\kappa)
$$
을 만들 수 있다. 이 극한은 공여과형이다($S$가 유향이기 때문이다).
이 구성은 $\kappa$에 관해 함자적이며, 특히
$\text{Aut}(\kappa)$는 이 프로유한군에 작용한다.
예를 들어 $S = \{1, n\}$이면 $\mu_n(\kappa)$를 얻는다.
$S = \{1, \ell, \ell^2, \ell^3, \ldots\}$이고 $\ell$이
$\kappa$의 표수와 다른 소수이면 $\lim_n \mu_{\ell^n}(\kappa)$를 얻는다.
이는 흔히 $\ell$진 Tate 가군, 즉 $\kappa$의 곱셈군의 Tate 가군이라 불린다
(More on Algebra, 예 \ref{more-algebra-example-spectral-sequence-principal}와 비교하라).

\begin{lemma}
\label{pione-lemma-structure-decomposition}
$A$를 분수체 $K$를 갖는 이산 부치환환이라 하자.
$L/K$를 (무한일 수도 있는) 갈루아 확대라 하자.
$B$를 $A$의 $L$에서의 정수적 폐포라 하자. 극대 아이디얼
$\mathfrak m$을 하나 택한다. 이는 $B$의 극대 아이디얼이다.
$G = \text{Gal}(L/K)$,
$D = \{\sigma \in G \mid \sigma(\mathfrak m) = \mathfrak m\}$,
$I = \{\sigma \in D \mid \sigma \bmod \mathfrak m =
\text{id}_{\kappa(\mathfrak m)}\}$로 둔다.
분해군 $D$는 다음 표준 완전열에 들어간다.
$$
1 \to I \to D \to \text{Aut}(\kappa(\mathfrak m)/\kappa_A) \to 1
$$
관성군 $I$는 다음 표준 완전열에 들어간다.
$$
1 \to P \to I \to I_t \to 1
$$
여기서 다음이 성립한다.
\begin{enumerate}
\item $P$는 $D$의 정규부분군이다.
\item $P$는 프로-$p$-군이며, 이는 $\kappa_A$의 표수가 $p > 1$일 때이다.
$P = \{1\}$인 것은 $\kappa_A$의 표수가 0일 때이다.
\item 곱셈적으로 유향인 $S \subset \mathbf{N}$가 존재하여,
$\kappa(\mathfrak m)$는 원시 $n$제곱근을 각 $n \in S$에 대해 포함한다
(단, $S$의 원소들은 $p$와 서로소이다).
\item 다음 표준 전사가 존재한다.
$$
\theta_{can} : I \to \lim_{n \in S} \mu_n(\kappa(\mathfrak m))
$$
그 핵은 $P$이고,
$\theta_{can}(\tau \sigma \tau^{-1}) = \tau(\theta_{can}(\sigma))$
를 $\tau \in D$, $\sigma \in I$에 대해 만족하며,
동형 $I_t \to \lim_{n \in S} \mu_n(\kappa(\mathfrak m))$를 유도한다.
\end{enumerate}
\end{lemma}

\begin{proof}
이는 주로 More on Algebra, 절 \ref{more-algebra-section-ramification}에서 증명한
유한 갈루아 확대에 관한 결과를 다시 표현한 것이다.
사상 $D \to \text{Aut}(\kappa(\mathfrak m)/\kappa)$의 전사성은
More on Algebra, 보조정리 \ref{more-algebra-lemma-one-orbit-geometric-galois}에 따른다.
이로부터 첫 번째 완전열을 얻는다.

\medskip\noindent
두 번째 짧은 완전열을 구성하기 위해 유한 갈루아 부분확대들의 집합을
$\Lambda$라 하자. 즉 $\lambda \in \Lambda$는 $L/L_\lambda/K$에 대응한다.
$G_\lambda = \text{Gal}(L_\lambda/K)$로 둔다.
$G_\lambda$는 전사 전이 사상을 갖는 유한군들의 역계이고
$G = \lim_{\lambda \in \Lambda} G_\lambda$임을 상기하자.
Fields, 보조정리 \ref{fields-lemma-infinite-galois-limit}를 보라.
$B_\lambda$를 $A$의 $L_\lambda$에서의 정수적 폐포라 하자.
이제 $\mathfrak m_\lambda = \mathfrak m \cap B_\lambda$로 두고,
$P_\lambda, I_\lambda, D_\lambda$를 각각 $\mathfrak m_\lambda$의
야성 관성군, 관성군, 분해군이라 쓰자. More on Algebra, 보조정리
\ref{more-algebra-lemma-galois-inertia}를 보라.
$\lambda \geq \lambda'$일 때 제한 사상으로 다음 가환 도식을 얻는다.
$$
\xymatrix{
P_\lambda \ar[d] \ar[r] &
I_\lambda \ar[d] \ar[r] &
D_\lambda \ar[d] \ar[r] &
G_\lambda \ar[d] \\
P_{\lambda'} \ar[r] &
I_{\lambda'} \ar[r] &
D_{\lambda'} \ar[r] &
G_{\lambda'}
}
$$
수직 사상들은 전사이다. More on Algebra, 보조정리
\ref{more-algebra-lemma-compare-inertia}를 보라.

\medskip\noindent
정의로부터 위의 동형 아래
$I = \lim I_\lambda$와 $D = \lim D_\lambda$, 그리고 $G = \lim G_\lambda$임이
즉시 따른다.
$L = \colim L_\lambda$이므로 $B = \colim B_\lambda$이고
$\kappa(\mathfrak m) = \colim \kappa(\mathfrak m_\lambda)$이다.
$D_\lambda$ 계의 전이 사상들은
$D_\lambda \to \text{Aut}(\kappa(\mathfrak m_\lambda)/\kappa)$와 양립한다
(More on Algebra, 보조정리 \ref{more-algebra-lemma-compare-inertia}).
따라서 사상 $D \to \text{Aut}(\kappa(\mathfrak m)/\kappa)$는
사상 $D_\lambda \to \text{Aut}(\kappa(\mathfrak m_\lambda)/\kappa)$들의 극한이다.

\medskip\noindent
표준 사상
$$
\theta_{\lambda, can} :
I_\lambda
\longrightarrow
\mu_{n_\lambda}(\kappa(\mathfrak m_\lambda))
$$
이 존재한다. 여기서 $n_\lambda = |I_\lambda|/|P_\lambda|$이고,
$\mu_{n_\lambda}(\kappa(\mathfrak m_\lambda))$의 위수는 $n_\lambda$이며,
$\theta_{\lambda, can}(\tau \sigma \tau^{-1}) =
\tau(\theta_{\lambda, can}(\sigma))$
를 $\tau \in D_\lambda$와 $\sigma \in I_\lambda$에 대해 만족한다.
또한 다음 가환 도식을 얻는다.
$$
\xymatrix{
I_\lambda \ar[r]_-{\theta_{\lambda, can}} \ar[d] &
\mu_{n_\lambda}(\kappa(\mathfrak m_\lambda))
\ar[d]^{(-)^{n_\lambda/n_{\lambda'}}} \\
I_{\lambda'} \ar[r]^-{\theta_{\lambda', can}} &
\mu_{n_{\lambda'}}(\kappa(\mathfrak m_{\lambda'}))
}
$$
More on Algebra, 주의 \ref{more-algebra-remark-canonical-inertia-character}를 보라.

\medskip\noindent
$S \subset \mathbf{N}$를 정수 $n_\lambda$들의 집합이라 하자.
$\Lambda$가 유향이므로 $S$는 곱셈적으로 유향이다.
위에 나타낸 가환 도식들에 따라 사상 $\theta_{\lambda, can}$들의 극한을 취하여
$$
\theta_{can} : I \to \lim_{n \in S} \mu_n(\kappa(\mathfrak m)).
$$
를 얻는다. 이 사상은 연속이다(작은 세부사항은 생략한다).
$I_\lambda$ 계의 전이 사상들이 전사이고 $\Lambda$가 유향이므로,
사영 $I \to I_\lambda$는 전사이다. 각 $\lambda$에 대해 다음 도식
$$
\xymatrix{
I \ar[d] \ar[r]_-{\theta_{can}} &
\lim_{n \in S} \mu_n(\kappa(\mathfrak m)) \ar[d] \\
I_{\lambda} \ar[r]^-{\theta_{\lambda, can}} &
\mu_{n_\lambda}(\kappa(\mathfrak m_\lambda))
}
$$
은 가환이다. 따라서 $\theta_{can}$의 상은 유한군
$\mu_{n_\lambda}(\kappa(\mathfrak m)) =
\mu_{n_\lambda}(\kappa(\mathfrak m_\lambda))$ (위수 $n_\lambda$)으로 전사한다
(위를 보라). 그러므로 $\theta_{can}$의 상은 조밀하다.
한편 $\theta_{can}$은 연속이고 그 정의역은 프로유한군이다.
따라서 위상적 논증에 의해 $\theta_{can}$은 전사이다.

\medskip\noindent
$\tau \in D$, $\sigma \in I$에 대하여
$\theta_{can}(\tau \sigma \tau^{-1}) = \tau(\theta_{can}(\sigma))$
가 성립한다는 것은 사상 $\theta_{\lambda, can}$의 대응하는 성질과
사상 $D \to \text{Aut}(\kappa(\mathfrak m))$ 및
사상 $D_\lambda \to \text{Aut}(\kappa(\mathfrak m_\lambda))$ 사이의
양립성에서 따른다. $P = \Ker(\theta_{can})$로 놓으면, 이는
$P$가 $D$의 정규 부분군임을 뜻한다. $I_t = I/P$로 놓으면
$\theta_{can}$의 전사성으로부터 동형
$I_t \to \lim_{n \in S} \mu_n(\kappa(\mathfrak m))$을 얻는다.

\medskip\noindent
증명을 끝내기 위해 $P = \lim P_\lambda$임을 보이자. 이로써
$P$가 프로-$p$-군임이 증명된다. 순한 관성군
$I_{\lambda, t} = I_\lambda/P_\lambda$의 위수는 $n_\lambda$이다.
전이 사상 $P_\lambda \to P_{\lambda'}$가 전사이고 $\Lambda$가 유향이므로,
다음 짧은 완전열을 얻는다.
$$
1 \to \lim P_\lambda \to I \to \lim I_{\lambda, t} \to 1
$$
(세부 사항은 생략한다). 각 $\lambda$에 대하여 사상
$\theta_{\lambda, can}$이 동형
$I_{\lambda, t} \cong \mu_{n_\lambda}(\kappa(\mathfrak m))$
을 유도하므로 원하는 결과가 따른다.
\end{proof}

\begin{lemma}
\label{pione-lemma-structure-decomposition-separable-closure}
$A$를 분수체가 $K$인 이산 부치환환이라 하자.
$K^{sep}$를 $K$의 분리 폐포라 하자.
$A^{sep}$를 $K^{sep}$에서 $A$의 정수적 폐포라 하자.
$\mathfrak m^{sep}$를 $A^{sep}$의 극대 아이디얼이라 하자.
$\mathfrak m = \mathfrak m^{sep} \cap A$로 놓고
$\kappa = A/\mathfrak m$, 그리고
$\overline{\kappa} = A^{sep}/\mathfrak m^{sep}$로 놓자.
그러면 $\overline{\kappa}$는 $\kappa$의 대수적 폐포이다.
$G = \text{Gal}(K^{sep}/K)$,
$D = \{\sigma \in G \mid \sigma(\mathfrak m^{sep}) = \mathfrak m^{sep}\}$,
그리고
$I = \{\sigma \in D \mid \sigma \bmod \mathfrak m^{sep} =
\text{id}_{\kappa(\mathfrak m^{sep})}\}$로 놓자.
분해군 $D$는 다음 표준 완전열에 들어간다.
$$
1 \to I \to D \to \text{Gal}(\kappa^{sep}/\kappa) \to 1
$$
여기서 $\kappa^{sep} \subset \overline{\kappa}$는 $\kappa$의 분리 폐포이다.
관성군 $I$는 다음 표준 완전열에 들어간다.
$$
1 \to P \to I \to I_t \to 1
$$
그리고 다음이 성립한다.
\begin{enumerate}
\item $P$는 $D$의 정규 부분군이다.
\item $\kappa_A$의 표수가 $p > 1$이면 $P$는 프로-$p$-군이고,
$\kappa_A$의 표수가 0이면 $P = \{1\}$이다.
\item 다음 표준 전사 사상이 존재한다.
$$
\theta_{can} : I \to \lim_{n\text{ prime to }p} \mu_n(\kappa^{sep})
$$
그 핵은 $P$이고, $\tau \in D$, $\sigma \in I$에 대하여
$\theta_{can}(\tau \sigma \tau^{-1}) = \tau(\theta_{can}(\sigma))$
를 만족하며, 또한 동형
$I_t \to \lim_{n\text{ prime to }p} \mu_n(\kappa^{sep})$을 유도한다.
\end{enumerate}
\end{lemma}

\begin{proof}
보조정리 \ref{pione-lemma-get-algebraic-closure}에 의해
$\overline{\kappa}$는 $\kappa$의 대수적 폐포이다.
나머지 명제 대부분은 보조정리 \ref{pione-lemma-structure-decomposition}의
대응하는 부분에서 곧바로 따른다. 예를 들어
$\text{Aut}(\overline{\kappa}/\kappa) = \text{Gal}(\kappa^{sep}/\kappa)$
이므로 첫 번째 완전열을 얻는다.
이제 증명이 필요한 유일한 나머지 명제는 보조정리
\ref{pione-lemma-structure-decomposition}의 $S$에 대하여
극한 $\lim_{n \in S} \mu_n(\overline{\kappa})$가
$\lim_{n\text{ prime to }p} \mu_n(\kappa^{sep})$와 같다는 것이다.
이를 보이려면 $p$와 서로소인 모든 정수 $n$이 $S$의 어떤 원소를
나눈다는 것을 보이면 충분하다.
균등화원 $\pi \in A$를 택하고 다항식 $X^n - \pi$의 분해체를 $L$이라 하자.
이 다항식은 분리적이므로 $L$은 $K$의 유한 갈루아 확대이다.
매장 $L \to K^{sep}$ 하나를 택한다.
$B$를 $L$에서 $A$의 정수적 폐포라 하면
$A \to B_{\mathfrak m^{sep} \cap B}$의 분기 지수는
$n$으로 나누어떨어진다($\pi$의 $n$제곱근이 $B$에 있기 때문이다. 실제로
분기 지수는 $n$과 같지만 여기서는 필요하지 않다).
보조정리 \ref{pione-lemma-structure-decomposition}의 증명에서 사용한
$S$의 구성으로부터 $n$이 $S$의 어떤 원소를 나눈다는 결론이 따른다.
\end{proof}









\section{기하적 기본군과 산술적 기본군}
\label{pione-section-galois-action}

\noindent
이 절에서는 체 $k$ 위의 스킴 $X$의 기본군과
$X_{\overline{k}}$의 기본군을 비교할 때 어떤 일이 일어나는지 살펴본다.
여기서 $\overline{k}$는 $k$의 대수적 폐포이다.

\begin{lemma}
\label{pione-lemma-limit}
$I$를 유향 집합이라 하자. $X_i$를 준콤팩트 준분리 스킴들로 이루어진
$I$ 위의 역계라 하고, 그 전이 사상들은 아핀이라고 하자.
Limits, 절 \ref{limits-section-limits}에서와 같이 $X = \lim X_i$로 놓자.
그러면 범주의 동치
$$
\colim \textit{F\'Et}_{X_i} = \textit{F\'Et}_X
$$
가 존재한다. 충분히 큰 모든 $i$에 대하여 $X_i$가 연결이고
$\overline{x}$가 $X$의 기하점이면,
$$
\pi_1(X, \overline{x}) = \lim \pi_1(X_i, \overline{x})
$$
이다.
\end{lemma}

\begin{proof}
범주의 동치는 Limits, 보조정리
\ref{limits-lemma-descend-finite-presentation},
\ref{limits-lemma-descend-finite-finite-presentation}, 그리고
\ref{limits-lemma-descend-etale}에서 따른다.
두 번째 명제는 범주에 관한 명제로부터 형식적으로 따른다.
\end{proof}

\begin{lemma}
\label{pione-lemma-perfection}
$k$를 체라 하고 그 완전화를 $k^{perf}$라 하자.
$X$를 $k$ 위의 연결 스킴이라 하자.
그러면 $X_{k^{perf}}$는 연결이고
$\pi_1(X_{k^{perf}}) \to \pi_1(X)$는 동형이다.
\end{lemma}

\begin{proof}
이는 기본군의 위상적 불변성의 특수한 경우이다.
명제 \ref{pione-proposition-universal-homeomorphism}을 보라.
$\Spec(k^{perf}) \to \Spec(k)$가 보편 위상동형임을 보이려면
Algebra, 보조정리 \ref{algebra-lemma-radicial-integral-bijective}를 사용하면 된다.
\end{proof}

\begin{lemma}
\label{pione-lemma-ses-field}
$k$를 체라 하고 그 대수적 폐포를 $\overline{k}$라 하자.
$X$를 체 $k$ 위의 준콤팩트 준분리 스킴이라 하자.
기저변환 $X_{\overline{k}}$가 연결이면, 다음은 프로유한 위상군들의 짧은 완전열이다.
$$
1 \to \pi_1(X_{\overline{k}}) \to \pi_1(X) \to \pi_1(\Spec(k)) \to 1
$$
\end{lemma}

\begin{proof}
$\textit{F\'Et}_{\Spec(k)}$의 연결 대상은
$\Spec(k') \to \Spec(k)$ 꼴이며, 여기서 $k'/k$는 유한 분리 확대이다.
그러면 $X_{\Spec{k'}}$는 연결이다. 실제로 사상
$X_{\overline{k}} \to X_{\Spec(k')}$가 전사이고
가정에 의해 $X_{\overline{k}}$가 연결이기 때문이다. 따라서
보조정리 \ref{pione-lemma-functoriality-galois-surjective}에 의해
$\pi_1(X) \to \pi_1(\Spec(k))$는 전사이다.

\medskip\noindent
계속하기 전에 $k$가 완전체라고 가정해도 됨에 주의하자.
실제로 보조정리 \ref{pione-lemma-perfection}에 의해
$\pi_1(X_{k^{perf}}) = \pi_1(X)$ 및
$\pi_1(\Spec(k^{perf})) = \pi_1(\Spec(k))$이다.

\medskip\noindent
함자의 합성
$\textit{F\'Et}_{\Spec(k)} \to \textit{F\'Et}_X \to
\textit{F\'Et}_{X_{\overline{k}}}$이 대상을
$X_{\Spec(\overline{k})}$의 복사본들의 서로소 합으로 보냄은 명백하다.
따라서 합성
$\pi_1(X_{\overline{k}}) \to \pi_1(X) \to \pi_1(\Spec(k))$
은 보조정리 \ref{pione-lemma-composition-trivial}에 의해 자명한 준동형이다.

\medskip\noindent
$U \to X$를 유한 에탈 사상이라 하고 $U$가 연결이라고 하자.
$U \times_X X_{\overline{k}} = U_{\overline{k}}$임에 주의하자.
$U_{\overline{k}} \to X_{\overline{k}}$가 단면
$s : X_{\overline{k}} \to U_{\overline{k}}$를 가진다고 가정하자.
그러면 $s(X_{\overline{k}})$는 $U_{\overline{k}}$의 열린 연결 성분이다.
$\sigma \in \text{Gal}(\overline{k}/k)$에 대하여
$\Spec(\sigma)$에 의한 $s$의 기저변환을 $s^\sigma$라 쓰자.
$U_{\overline{k}} \to X_{\overline{k}}$는 유한 에탈이므로
단면은 유한 개뿐이다. 따라서
$$
\overline{T} = \bigcup s^\sigma(X_{\overline{k}})
$$
는 유한 합이고, $\overline{T}$는
$\text{Gal}(\overline{k}/k)$-안정인 열린닫힌 부분집합이다.
Varieties, 보조정리 \ref{varieties-lemma-closed-fixed-by-Galois}에 의해
$\overline{T}$는 어떤 닫힌 부분집합 $T \subset U$의 역상이다.
$U_{\overline{k}} \to U$가 열린 사상이므로
(Morphisms, 보조정리 \ref{morphisms-lemma-scheme-over-field-universally-open}),
$T$도 열려 있다. $U$가 연결이므로 $T = U$이다.
따라서 $U_{\overline{k}}$는 $X_{\overline{k}}$의 복사본들의
(유한) 서로소 합이다. 보조정리
\ref{pione-lemma-functoriality-galois-normal}에 의해
$\pi_1(X_{\overline{k}}) \to \pi_1(X)$의 상은 정규 부분군이다.

\medskip\noindent
유한 에탈 덮개 $V \to X_{\overline{k}}$를 택하자.
$\overline{k}$가 $k$의 유한 분리 확대들의 합집합임을 상기하자.
보조정리 \ref{pione-lemma-limit}에 의해, 유한 분리 확대 $k'/k$와
유한 에탈 사상 $U \to X_{k'}$로서
$V = X_{\overline{k}} \times_{X_{k'}} U =
U \times_{\Spec(k')} \Spec(\overline{k})$
가 되는 것을 얻는다.
그러면 합성 $U \to X_{k'} \to X$는 유한 에탈이고,
$U \times_{\Spec(k)} \Spec(\overline{k})$는
$V = U \times_{\Spec(k')} \Spec(\overline{k})$를 열린닫힌 부분스킴으로 포함한다.
(실제로 $\Spec(\overline{k})$는 곱셈 사상
$k' \otimes_k \overline{k} \to \overline{k}$에 의해
$\Spec(k') \times_{\Spec(k)} \Spec(\overline{k})$의 열린닫힌 부분스킴이다.)
보조정리 \ref{pione-lemma-functoriality-galois-injective}에 의해
$\pi_1(X_{\overline{k}}) \to \pi_1(X)$는 단사이다.

\medskip\noindent
마지막으로, $U_{\overline{k}}$가 $X_{\overline{k}}$의 복사본들의 서로소 합이 되는
임의의 유한 에탈 사상 $U \to X$에 대하여, 유한 에탈 사상
$V \to \Spec(k)$와 전사
$V \times_{\Spec(k)} X \to U$가 존재함을 보여야 한다.
보조정리 \ref{pione-lemma-functoriality-galois-ses}를 보라.
위와 같이 보조정리 \ref{pione-lemma-limit}을 사용하면, 유한 분리 확대 $k'/k$로서
동형
$U_{k'} \cong \coprod_{i = 1, \ldots, n} X_{k'}$
이 존재하는 것을 얻는다.
따라서 $V = \coprod_{i = 1, \ldots, n} \Spec(k')$로 놓으면 된다.
\end{proof}







\section{호모토피 완전열}
\label{pione-section-homotopy-exact-sequence}

\noindent
이 절에서는 다음 결과를 논의한다.
$f : X \to S$를 유한 표시인 평탄 고유 사상이라 하고, 그 기하적 올들은
연결이고 기약이라고 하자.
$S$는 연결이고 $\overline{s}$는 $S$의 기하점이라고 가정하자.
그러면 기본군의 완전열
$$
\pi_1(X_{\overline{s}}) \to \pi_1(X) \to \pi_1(S) \to 1
$$
을 얻는다. 명제 \ref{pione-proposition-first-homotopy-sequence}을 보라.

\begin{lemma}
\label{pione-lemma-stein-factorization-etale}
\begin{reference}
\cite[Exposé X, 명제 1.2, 262쪽]{SGA1}.
\end{reference}
스킴의 고유 사상 $f : X \to S$를 택하자.
$X \to S' \to S$를 $f$의 Stein 인자분해라 하자. More on Morphisms, 정리
\ref{more-morphisms-theorem-stein-factorization-general}을 보라.
$f$가 유한 표시이고 평탄하며 기하적으로 기약인 올을 가지면,
$S' \to S$는 유한 에탈이다.
\end{lemma}

\begin{proof}
이는 Derived Categories of Schemes, 보조정리
\ref{perfect-lemma-proper-flat-geom-red}와 More on Morphisms, 정리
\ref{more-morphisms-theorem-stein-factorization-general}에 담긴 정보에서 따른다.
\end{proof}

\begin{proposition}
\label{pione-proposition-first-homotopy-sequence}
$f : X \to S$를 유한 표시인 평탄 고유 사상이라 하고, 그 기하적 올들은
연결이고 기약이라고 하자. $S$는 연결이고 $\overline{s}$는
$S$의 기하점이라고 가정하자. 그러면 기본군의 완전열
$$
\pi_1(X_{\overline{s}}) \to \pi_1(X) \to \pi_1(S) \to 1
$$
을 얻는다.
\end{proposition}

\begin{proof}
유한 에탈 사상 $Y \to X$를 택하자. 다음 Stein 인자분해를 생각하자.
$$
\xymatrix{
Y \ar[d] \ar[r] & X \ar[d] \\
T \ar[r] & S
}
$$
이는 $Y \to S$의 Stein 인자분해이다.
보조정리 \ref{pione-lemma-stein-factorization-etale}에 의해 사상 $T \to S$는 유한 에탈이다.
이와 같이 함자 $\textit{F\'Et}_X \to \textit{F\'Et}_S$를 얻는다.
임의의 유한 에탈 사상 $U \to S$에 대하여, $X$ 위의 사상
 $Y \to U \times_S X$는 $S$ 위의 사상
 $Y \to U$와 같고, 이러한 사상은 Stein 인자분해를 유일하게 경유한다.
즉, 유일한 사상 $T \to U$에 대응한다
(Stein 인자분해를 More on Morphisms, 보조정리
\ref{more-morphisms-lemma-stein-universally-closed}의 상대 정규화로
구성한 다음 Morphisms, 보조정리
\ref{morphisms-lemma-characterize-normalization}에 의해 인자화하면 된다).
따라서 함자
$\textit{F\'Et}_X \to \textit{F\'Et}_S$와
$\textit{F\'Et}_S \to \textit{F\'Et}_X$는 서로 수반이다.
또한 $U \times_S X \to S$의 Stein 인자분해가 $U$임에 주의하자.
이는 $U \times_S X \to U$의 올들이 기하적으로 연결이기 때문이다.

\medskip\noindent
위 논의와 Categories, 보조정리
\ref{categories-lemma-adjoint-fully-faithful}에 의해
$\textit{F\'Et}_S \to \textit{F\'Et}_X$
는 충실충만이다. 즉,
$\pi_1(X) \to \pi_1(S)$는 전사이다
(보조정리 \ref{pione-lemma-functoriality-galois-surjective}).

\medskip\noindent
합성
$\textit{F\'Et}_S \to \textit{F\'Et}_X \to \textit{F\'Et}_{X_{\overline{s}}}$
이 임의의 $U$를 $X_{\overline{s}}$의 복사본들의 서로소 합으로 보냄은 명백하다.
따라서 $\pi_1(X_{\overline{s}}) \to \pi_1(X) \to \pi_1(S)$는
보조정리 \ref{pione-lemma-composition-trivial}에 의해 자명하다.

\medskip\noindent
유한 에탈 사상 $Y \to X$에서 $Y$가 연결이고,
 $Y \times_X X_{\overline{s}}$가 $X_{\overline{s}}$와 동형인
 연결 성분 $Z$를 포함한다고 하자.
위의 Stein 인자분해 $T$를 생각하자.
$\overline{t} \in T_{\overline{s}}$를 올 $Z$에 대응하는 점이라 하자.
$T$는 연결 스킴의 상이므로 연결이고, 위의 전사성에 의해
$T \times_S X$도 연결이다.
다음 인자화를 생각하자.
$$
\pi : Y \longrightarrow T \times_S X
$$
임의의 닫힌점 $\overline{x} \in X_{\overline{s}}$를 택하자.
$\kappa(\overline{t}) = \kappa(\overline{s}) = \kappa(\overline{x})$
가 대수적으로 닫힌 체임에 주의하자.
그러면 $(\overline{t}, \overline{x})$ 위의 $\pi$의 올은 하나의 점,
즉 $\overline{z} \in Z$만으로 이루어진다. 이 점은 동형
$Z \to X_{\overline{s}}$ 아래에서
$\overline{x} \in X_{\overline{s}}$에 대응한다.
따라서 유한 에탈 사상 $\pi$는
$(\overline{t}, \overline{x})$의 근방에서 차수가 $1$이다.
$T \times_S X$가 연결이므로 모든 곳에서 차수가 $1$이고,
$Y \cong T \times_S X$를 얻는다. 따라서
$Y \times_X X_{\overline{s}}$는 완전히 분해된다.
이들을 종합하면 보조정리
\ref{pione-lemma-functoriality-galois-ses}와
\ref{pione-lemma-functoriality-galois-normal}을 모두 적용할 수 있고,
증명이 끝난다.
\end{proof}





\section{특수화 사상}
\label{pione-section-specialization-map}

\noindent
이 절에서는 특수화 사상을 구성한다.
$f : X \to S$를 기하적으로 연결인 올들을 갖는 스킴의 고유 사상이라 하자.
$s' \leadsto s$를 $S$의 점들의 특수화라 하자.
$\overline{s}$와 $\overline{s}'$를 각각 $s$와 $s'$ 위의
기하점이라 하자. 그러면 특수화 사상
$$
sp : \pi_1(X_{\overline{s}'}) \longrightarrow \pi_1(X_{\overline{s}})
$$
을 얻는다. 이 사상의 구성은 다음과 같다. $A$를
$\kappa(s) \subset \kappa(s)^{sep} \subset \kappa(\overline{s})$에 관한
$\mathcal{O}_{S, s}$의 엄밀 헨젤화라 하자.
Algebra, 정의 \ref{algebra-definition-henselization}을 보라.
$s' \leadsto s$이므로 점 $s'$에 대응하는
$\Spec(\mathcal{O}_{S, s})$의 점을 생각할 수 있다. 따라서 $\Spec(A)$에는
그 위에 놓이고 잉여체가 $\kappa(s')$의 분리 대수적 확대인 점 $s'$가
적어도 하나 존재한다(하나보다 많을 수도 있다).
$\kappa(\overline{s}')$가 대수적으로 닫혀 있으므로, 다음 가환 도식을 주는 사상
$\varphi : \overline{s}' \to \Spec(A)$를 택할 수 있다.
$$
\xymatrix{
\overline{s}' \ar[r]_-\varphi \ar[rd] &
\Spec(A) \ar[d] &
\overline{s} \ar[l] \ar[ld] \\
& S
}
$$
특수화 사상은 다음 합성이다.
$$
\pi_1(X_{\overline{s}'}) \longrightarrow
\pi_1(X_A) =
\pi_1(X_{\kappa(s)^{sep}}) =
\pi_1(X_{\overline{s}})
$$
여기서 첫 번째 등호는 보조정리
\ref{pione-lemma-finite-etale-on-proper-over-henselian}에 의한 것이고,
두 번째 등호는 보조정리 \ref{pione-lemma-perfection}과
\ref{pione-lemma-finite-etale-invariant-over-proper}에서 따른다.
구성상 특수화 사상은 다음 가환 도식에 들어간다.
$$
\xymatrix{
\pi_1(X_{\overline{s}'}) \ar[rr]_{sp} \ar[rd] & &
\pi_1(X_{\overline{s}}) \ar[ld] \\
& \pi_1(X)
}
$$
여기서는 $X$가 연결이라고 가정한다. 특수화 사상은 위에서 택한
$\varphi : \overline{s}' \to \Spec(A)$에 의존한다.
이 의존성을 나타낼 때는 $sp_\varphi$라고 쓴다.

\begin{lemma}
\label{pione-lemma-specialization-map-base-change}
다음 가환 도식을 생각하자.
$$
\xymatrix{
Y \ar[d]_g \ar[r] & X \ar[d]^f \\
T \ar[r] & S
}
$$
여기서 $f$와 $g$는 기하적으로 연결인 올들을 갖는 고유 사상이다.
$t' \leadsto t$를 $T$의 점들의 특수화라 하고, 위에서 구성한 특수화 사상
$sp : \pi_1(Y_{\overline{t}'}) \to \pi_1(Y_{\overline{t}})$를 생각하자.
그러면 다음 가환 도식이 존재한다.
$$
\xymatrix{
\pi_1(Y_{\overline{t}'}) \ar[r]_{sp} \ar[d] & \pi_1(Y_{\overline{t}}) \ar[d] \\
\pi_1(X_{\overline{s}'}) \ar[r]^{sp} & \pi_1(X_{\overline{s}})
}
$$
이는 특수화 사상들의 도식이며, $\overline{s}$와 $\overline{s}'$는
각각 $\overline{t}$와 $\overline{t}'$의 상이다.
\end{lemma}

\begin{proof}
 $B$를 $\kappa(t) \subset \kappa(t)^{sep} \subset \kappa(\overline{t})$에 관한
$\mathcal{O}_{T, t}$의 엄밀 헨젤화라 하자.
특수화 사상의 구성에서와 같이 $\psi : \overline{t}' \to \Spec(B)$를 택하자.
이는 $\overline{t}' \to T$의 올림이다.
$s$와 $s'$를 각각 $t$와 $t'$의 $S$에서의 상이라 하자.
$A$를 $\kappa(s) \subset \kappa(s)^{sep} \subset \kappa(\overline{s})$에 관한
$\mathcal{O}_{S, s}$의 엄밀 헨젤화라 하자.
$\kappa(\overline{s}) = \kappa(\overline{t})$이므로,
엄밀 헨젤화의 함자성(Algebra, 보조정리
\ref{algebra-lemma-strictly-henselian-functorial})에 의해 다음 가환 도식에 들어가는
환 준동형 $A \to B$를 얻는다.
$$
\xymatrix{
\overline{t}' \ar[r]_-\psi \ar[d] & \Spec(B) \ar[d] \ar[r] & T \ar[d] \\
\overline{s}' \ar[r]^-\varphi & \Spec(A) \ar[r] & S
}
$$
여기서 사상 $\varphi : \overline{s}' \to \Spec(A)$는 단순히 합성
$\overline{t}' \to \Spec(B) \to \Spec(A)$로 잡는다.
기저변환을 취하면 다음 가환 도식을 얻는다.
$$
\xymatrix{
Y_{\overline{t}'} \ar[r] \ar[d] & Y_B \ar[d] \\
X_{\overline{s}'} \ar[r] & X_A
}
$$
특수화 사상의 구성에 따라 이 도식의 가환성은 보조정리에 나온 도식의 가환성을 뜻한다.
\end{proof}

\begin{lemma}
\label{pione-lemma-specialization-map-composition}
$f : X \to S$를 기하적으로 연결인 올들을 갖는 고유 사상이라 하자.
$s'' \leadsto s' \leadsto s$를 $S$의 점들의 특수화라 하자.
특수화 사상들의 합성
$\pi_1(X_{\overline{s}''}) \to \pi_1(X_{\overline{s}'}) \to
\pi_1(X_{\overline{s}})$은 특수화 사상
$\pi_1(X_{\overline{s}''}) \to \pi_1(X_{\overline{s}})$이다.
\end{lemma}

\begin{proof}
 $\mathcal{O}_{S, s} \to A$를 $\kappa(s) \to \kappa(\overline{s})$를 사용해
구성한 엄밀 헨젤화라 하자.
첫 번째 특수화 사상을 구성할 때 사용한 사상을
$A \to \kappa(\overline{s}')$라 하자.
$\mathcal{O}_{S, s'} \to A'$를 $\kappa(s') \subset \kappa(\overline{s}')$를 사용해
구성한 엄밀 헨젤화라 하자.
엄밀 헨젤화의 함자성에 의해 사상
$A \to A'$가 존재하여, $A' \to \kappa(\overline{s}')$와의 합성이 주어진 사상이 된다.
(Algebra, 보조정리 \ref{algebra-lemma-map-into-henselian-colimit}).
다음으로 두 번째 특수화 사상을 구성할 때 사용한 사상을
$A' \to \kappa(\overline{s}'')$라 하자. 그러면 첫 번째와 두 번째 특수화 사상의
합성이 사상
$A \to A' \to \kappa(\overline{s}'')$를 사용해 구성한 특수화 사상
$\pi_1(X_{\overline{s}''}) \to \pi_1(X_{\overline{s}})$
임은 명백하다.
\end{proof}

\noindent
$X \to S$를 기하적으로 연결인 올들을 갖는 고유 사상이라 하자.
$R$을 분수체가 대수적으로 닫힌 엄밀 헨젤 부치환환이라 하고
사상 $\Spec(R) \to S$를 택하자.
$\eta, s \in \Spec(R)$를 각각 일반점과 닫힌점이라 하자.
그러면 기저변환 $X_R/\Spec(R)$에 대하여 특수화 사상
$$
sp_R : \pi_1(X_\eta) \to \pi_1(X_s)
$$
을 생각할 수 있다. $\eta$와 $s$의 잉여체가 모두 대수적으로 닫혀 있으므로
이 사상은 잘 정의된다.

\begin{lemma}
\label{pione-lemma-specialization-map-valuation-ring}
$f : X \to S$를 기하적으로 연결인 올들을 갖는 고유 사상이라 하자.
$s' \leadsto s$를 $S$의 점들의 특수화라 하고
$sp : \pi_1(X_{\overline{s}'}) \to \pi_1(X_{\overline{s}})$를
특수화 사상이라 하자. 그러면 분수체가 대수적으로 닫힌 $S$ 위의
엄밀 헨젤 부치환환 $R$로서, $sp$가 위에서 정의한 $sp_R$과 동형인 것이 존재한다.
\end{lemma}

\begin{proof}
$\mathcal{O}_{S, s} \to A$를 $\kappa(s) \to \kappa(\overline{s})$를 사용해
구성한 엄밀 헨젤화라 하자. $sp$는 구성에 사용된 사상
$A \to \kappa(\overline{s}')$으로부터 얻어진다.
$R \subset \kappa(\overline{s}')$을 분수체가 $\kappa(\overline{s}')$이고
$A$의 상을 지배하는 부치환환이라 하자.
Algebra, 보조정리 \ref{algebra-lemma-dominate}를 보라.
예를 들어 보조정리 \ref{pione-lemma-normal-local-domain-separablly-closed-fraction-field}와
Algebra, 보조정리 \ref{algebra-lemma-valuation-ring-normal}에 의해
$R$이 엄밀 헨젤임에 주의하자. 따라서 보조정리가 따른다.
\end{proof}

\noindent
$X \to S$를 기하적으로 연결인 올들을 갖는 고유 사상이라 하자.
$R$을 엄밀 헨젤 이산 부치환환이라 하고 사상 $\Spec(R) \to S$를 택하자.
$\eta, s \in \Spec(R)$를 각각 일반점과 닫힌점이라 하자. 그러면
기저변환 $X_R/\Spec(R)$에 대하여 특수화 사상
$$
sp_R : \pi_1(X_{\overline{\eta}}) \to \pi_1(X_s)
$$
을 생각할 수 있다. $s$의 잉여체가 대수적으로 닫혀 있으므로
이 사상은 잘 정의된다.

\begin{lemma}
\label{pione-lemma-specialization-map-discrete-valuation-ring}
$f : X \to S$를 기하적으로 연결인 올들을 갖는 고유 사상이라 하자.
$s' \leadsto s$를 $S$의 점들의 특수화라 하고,
$sp : \pi_1(X_{\overline{s}'}) \to \pi_1(X_{\overline{s}})$를
특수화 사상이라 하자. $S$가 뇌터이면, $S$ 위의 엄밀 헨젤 이산 부치환환
$R$로서 $sp$가 위에서 정의한 $sp_R$과 동형인 것이 존재한다.
\end{lemma}

\begin{proof}
$\mathcal{O}_{S, s} \to A$를 $\kappa(s) \to \kappa(\overline{s})$를 사용해
구성한 엄밀 헨젤화라 하자. $sp$는 구성에 사용된 사상
$A \to \kappa(\overline{s}')$으로부터 얻어진다.
$R \subset \kappa(\overline{s}')$을 $A$의 상을 지배하는 이산 부치환환이라 하자.
Algebra, 보조정리
\ref{algebra-lemma-exists-dvr}를 보라. 다음 체의 도식에서
$$
\xymatrix{
\kappa(\overline{s}) \ar[r] & k \\
A/\mathfrak m_A \ar[r] \ar[u] & R/\mathfrak m_R \ar[u]
}
$$
$k$가 대수적으로 닫히도록 택하자. $R^{sh}$를 $R \to k$를 사용해
$R$로부터 구성한 엄밀 헨젤화라 하자. 사상 $A \to R^{sh}$를 얻는다.
Algebra, 보조정리
\ref{algebra-lemma-strictly-henselian-functorial}.
More on Algebra, 보조정리 \ref{more-algebra-lemma-henselization-dvr}에 의해
환 $R^{sh}$는 이산 부치환환이다. 그 일반점과 닫힌점을 $\eta, o$라 쓰자.
다음 $\Spec(R^{sh})$의 스킴 도식은 가환이다.
$$
\xymatrix{
\overline{\eta} \ar[d] \ar[r] & \Spec(R^{sh}) \ar[d] &
\Spec(k) \ar[d] \ar[l] \\
\overline{s}' \ar[r] & \Spec(A) & \overline{s} \ar[l]
}
$$
따라서 다음 가환 도식을 얻는다.
$$
\xymatrix{
\pi_1(X_{\overline{\eta}}) \ar[d] \ar[r]_{sp_{R^{sh}}} & \pi_1(X_o) \ar[d] \\
\pi_1(X_{\overline{s}'}) \ar[r]^{sp} & X_{\overline{s}}
}
$$
사상들의 구성에 의해 이는 특수화 사상들의 도식이다.
세로 화살표들이 동형이므로(보조정리
\ref{pione-lemma-finite-etale-invariant-over-proper}), 보조정리가 증명된다.
\end{proof}









\section{닫힌 부분스킴으로의 제한}
\label{pione-section-lefschetz}

\noindent
이 절에서는 제한 함자에 관한 몇 가지 결과를 증명한다.
$$
\textit{F\'Et}_X \longrightarrow \textit{F\'Et}_Y,\quad
U \longmapsto V = U \times_X Y
$$
여기서 $X$는 스킴이고 $Y$는 닫힌 부분스킴이다.
기본군의 위상적 불변성을 사용하면 이 함자의 연구를
유한 국소 자유 가군에 대한 완비화 함자와 연관시킬 수 있다.

\medskip\noindent
이하의 보조정리에서는 Cohomology of Schemes, 절
\ref{coherent-section-coherent-formal}에서 정의한 연접 형식 가군의 개념을 사용한다.
뇌터 스킴과 준연접 아이디얼층 $\mathcal{I} \subset \mathcal{O}_X$가 주어졌을 때,
대상 $(\mathcal{F}_n)$을 $\textit{Coh}(X, \mathcal{I})$의 대상이라 하자.
이 대상이 {\it 유한 국소 자유}라는 것은 각 $\mathcal{F}_n$이
유한 국소 자유 $\mathcal{O}_X/\mathcal{I}^n$-가군이라는 뜻이다.

\begin{lemma}
\label{pione-lemma-restriction-fully-faithful}
 $X$를 뇌터 스킴이라 하고, $Y \subset X$를 아이디얼층
$\mathcal{I} \subset \mathcal{O}_X$로 정의되는 닫힌 부분스킴이라 하자.
완비화 함자
$$
\textit{Coh}(\mathcal{O}_X)
\longrightarrow 
\textit{Coh}(X, \mathcal{I}),\quad
\mathcal{F} \longmapsto \mathcal{F}^\wedge
$$
가 유한 국소 자유 대상들의 충만 부분범주 위에서 충실충만이라고 가정하자(위 설명을 보라).
그러면 제한 함자 $\textit{F\'Et}_X \to \textit{F\'Et}_Y$는 충실충만이다.
\end{lemma}

\begin{proof}
유한 에탈 덮개들의 범주에는 내부 Hom이 있으므로(보조정리
\ref{pione-lemma-internal-hom-finite-etale}), 다음을 보이면 충분하다.
$U$를 $X$ 위의 유한 에탈 스킴이라 하고 $X$ 위의 사상
$t : Y \to U$를 택하면, $t = s|_Y$를 만족하는 유일한 단면
$s : X \to U$가 존재한다. 도식은 다음과 같다.
$$
\xymatrix{
& U \ar[d]^f \\
Y \ar[r] \ar[ru] & X \ar@{..>}@/^1em/[u]
}
$$
점선 화살표 $s$를 찾는 것은 다음 $\mathcal{O}_X$-대수 사상을 찾는 것과 동치이다.
$$
s^\sharp : f_*\mathcal{O}_U \longrightarrow \mathcal{O}_X
$$
이 사상은 $Y$의 아이디얼층으로 나눈 뒤 주어진 대수 사상
$t^\sharp : f_*\mathcal{O}_U \to \mathcal{O}_Y$가 되어야 한다.
보조정리 \ref{pione-lemma-thickening}에 의해 $t$는 유일하게 양립하는 사상들의 계
$t_n : Y_n \to U$로 올라가며, 따라서 다음 사상을 얻는다.
$$
\lim t_n^\sharp : f_*\mathcal{O}_U \longrightarrow \lim \mathcal{O}_{Y_n}
$$
이는 $X$ 위의 대수층 사상이다.
$f_*\mathcal{O}_U$가 유한 국소 자유 $\mathcal{O}_X$-가군이므로,
유일한 $\mathcal{O}_X$-가군 사상
$\sigma : f_*\mathcal{O}_U \to \mathcal{O}_X$를 얻으며, 그 완비화는
$\lim t_n^\sharp$이다.
$\sigma$가 대수 준동형임을 보이려면 다음 도식의 가환성을 확인하면 된다.
$$
\xymatrix{
f_*\mathcal{O}_U \otimes_{\mathcal{O}_X} f_*\mathcal{O}_U
\ar[r] \ar[d]_{\sigma \otimes \sigma} &
f_*\mathcal{O}_U \ar[d]^\sigma \\
\mathcal{O}_X \otimes_{\mathcal{O}_X} \mathcal{O}_X \ar[r] &
\mathcal{O}_X
}
$$
임의의 $n$에 대하여 $Y_n$으로 제한한 뒤, 즉 완비화 함자를 적용한 뒤에는
이 도식이 가환함을 알 수 있다. 따라서 완비화 함자의 충실성으로부터
해당 사상이 대수 준동형이라는 결론이 따른다.
\end{proof}

\begin{lemma}
\label{pione-lemma-restriction-equivalence}
 $X$를 뇌터 스킴이라 하고, $Y \subset X$를 아이디얼층
$\mathcal{I} \subset \mathcal{O}_X$로 정의되는 닫힌 부분스킴이라 하자.
완비화 함자
$$
\textit{Coh}(\mathcal{O}_X)
\longrightarrow 
\textit{Coh}(X, \mathcal{I}),\quad
\mathcal{F} \longmapsto \mathcal{F}^\wedge
$$
가 유한 국소 자유 대상들의 충만 부분범주 사이의 동치라고 가정하자(위 설명을 보라).
그러면 제한 함자 $\textit{F\'Et}_X \to \textit{F\'Et}_Y$는 동치이다.
\end{lemma}

\begin{proof}
제한 함자가 충실충만이라는 것은 보조정리
\ref{pione-lemma-restriction-fully-faithful}에서 따른다.

\medskip\noindent
유한 에탈 사상 $U_1 \to Y$를 택하자. 증명을 끝내기 위해
$U_1$이 제한 함자의 본질적 상에 속함을 보이자.

\medskip\noindent
 $n \geq 1$에 대하여 $Y_n$을 $Y$의 $n$차 무한소 근방이라 하자.
보조정리 \ref{pione-lemma-thickening}에 의해 유일한 유한 에탈 사상
$\pi_n : U_n \to Y_n$이 존재하며, $Y = Y_1$로의 기저변환은
$U_1 \to Y_1$을 복원한다. 층
$\mathcal{F}_n = \pi_{n, *}\mathcal{O}_{U_n}$을 생각하자.
$\mathcal{F}_n$을 $X$ 위의 $\mathcal{O}_X$-가군으로 보면, 국소적으로
$(\mathcal{O}_X/f^{n + 1}\mathcal{O}_X)^{\oplus r}$와 동형이다.
이 $(\mathcal{F}_n)$은 $\textit{Coh}(X, \mathcal{I})$의 유한 국소 자유 대상이다.
가정에 의해 유한 국소 자유 $\mathcal{O}_X$-가군 $\mathcal{F}$와
양립하는 동형들의 계
$$
\mathcal{F}/\mathcal{I}^n\mathcal{F} \to \mathcal{F}_n
$$
을 얻는다. 이들은 $\mathcal{O}_X$-가군의 동형이다.

\medskip\noindent
 $\mathcal{F}$ 위의 대수 구조를 구성하기 위해 곱셈 사상
$\mathcal{F}_n \otimes_{\mathcal{O}_X} \mathcal{F}_n \to \mathcal{F}_n$을 생각하자.
이는 $\mathcal{F}_n = \pi_{n, *}\mathcal{O}_{U_n}$이 대수층이므로 주어진다.
이 사상들은 다음 사상을 정한다.
$$
(\mathcal{F}\otimes_{\mathcal{O}_X} \mathcal{F})^\wedge
\longrightarrow
\mathcal{F}^\wedge
$$
이는 범주 $\textit{Coh}(X, \mathcal{I})$의 사상이다. 따라서 가정에 의해
$\mu : \mathcal{F}\otimes_{\mathcal{O}_X} \mathcal{F} \to \mathcal{F}$를 얻으며,
그 $Y_n$으로의 제한들이 위의 곱셈 사상들을 주도록 할 수 있다.
보조정리의 명제에 나오는 함자의 충실성에 의해 $\mu$는
$\mathcal{F}$ 위의 $\mathcal{O}_X$-대수 구조를 정하며,
$\mathcal{F}_n$ 위의 주어진 대수 구조와 양립한다.
다음과 같이 놓자.
$$
U = \underline{\Spec}_X((\mathcal{F}, \mu))
$$
그러면 유한 국소 자유 스킴 $\pi : U \to X$를 얻고, 그 $Y$로의 제한은
$U_1$과 동형이다.
$\pi$의 판별식은 다음 단면의 영점 스킴이다.
$$
\det(Q_\pi) :
\mathcal{O}_X
\longrightarrow
\wedge^{top}(\pi_*\mathcal{O}_U)^{\otimes -2}
$$
이는 Discriminants, 절 \ref{discriminant-section-discriminant}에서 구성한 것이다.
모든 $n$에 대하여 그 $Y_n$으로의 제한이 동형이므로, Discriminants, 보조정리
\ref{discriminant-lemma-discriminant}에 의해 그 자체도 동형이다. 따라서 $\pi$는
Discriminants, 보조정리 \ref{discriminant-lemma-discriminant}에 의해 에탈이다.
\end{proof}

\begin{lemma}
\label{pione-lemma-restriction-fully-faithful-general}
 $X$를 뇌터 스킴이라 하고, $Y \subset X$를 아이디얼층
$\mathcal{I} \subset \mathcal{O}_X$로 정의되는 닫힌 부분스킴이라 하자.
$\mathcal{V}$를 $Y$를 포함하는 열린 부분스킴 $V \subset X$들의 집합이라 하고
역포함으로 순서화하자. 완비화 함자
$$
\colim_\mathcal{V} \textit{Coh}(\mathcal{O}_V)
\longrightarrow
\textit{Coh}(X, \mathcal{I}),
\quad
\mathcal{F} \longmapsto \mathcal{F}^\wedge
$$
가 유한 국소 자유 대상들의 충만 부분범주 위에서 충실충만이라고 가정하자(위 설명을 보라).
그러면 제한 함자
$\colim_\mathcal{V} \textit{F\'Et}_V \to \textit{F\'Et}_Y$
는 충실충만이다.
\end{lemma}

\begin{proof}
 $\mathcal{V}$는 유향 집합이므로 여극한은 Categories, 절
\ref{categories-section-directed-colimits}에서 설명한 것이다.
나머지 논증은 보조정리 \ref{pione-lemma-restriction-fully-faithful}의 증명과
거의 같으므로 독자는 건너뛰어도 된다.

\medskip\noindent
유한 에탈 덮개들의 범주에는 내부 Hom이 있으므로(보조정리
\ref{pione-lemma-internal-hom-finite-etale}), 다음을 보이면 충분하다.
$U$를 어떤 $V \in \mathcal{V}$ 위의 유한 에탈 스킴이라 하고,
$V$ 위의 사상 $t : Y \to U$를 택하면, 어떤 $V' \geq V$와
사상 $s : V' \to U$가 존재하여 $V$ 위에서 $t = s|_Y$가 된다.
도식은 다음과 같다.
$$
\xymatrix{
& & U \ar[d]^f \\
Y \ar[r] \ar[rru] & V' \ar@{..>}[ru] \ar[r] & V
}
$$
점선 화살표 $s$를 찾는 것은 다음
$\mathcal{O}_{V'}$-대수 사상을 찾는 것과 동치이다.
$$
s^\sharp : f_*\mathcal{O}_U|_{V'} \longrightarrow \mathcal{O}_{V'}
$$
이 사상은 $Y$의 아이디얼층으로 나눈 뒤 주어진 대수 사상
$t^\sharp : f_*\mathcal{O}_U \to \mathcal{O}_Y$가 되어야 한다.
보조정리 \ref{pione-lemma-thickening}에 의해 $t$는 유일하게 양립하는 사상들의 계
$t_n : Y_n \to U$로 올라가며, 따라서 다음 사상을 얻는다.
$$
\lim t_n^\sharp : f_*\mathcal{O}_U \longrightarrow \lim \mathcal{O}_{Y_n}
$$
이는 $V$ 위의 대수층 사상이다.
$f_*\mathcal{O}_U$는 유한 국소 자유 $\mathcal{O}_V$-가군이다.
따라서 어떤 $V' \geq V$와 사상
$\sigma : f_*\mathcal{O}_U|_{V'} \to \mathcal{O}_{V'}$
를 얻으며, 그 완비화는 $\lim t_n^\sharp$이다.
$\sigma$가 대수 준동형임을 보이려면 다음 도식
$$
\xymatrix{
(f_*\mathcal{O}_U \otimes_{\mathcal{O}_V} f_*\mathcal{O}_U)|_{V'}
\ar[r] \ar[d]_{\sigma \otimes \sigma} &
f_*\mathcal{O}_U|_{V'} \ar[d]^\sigma \\
\mathcal{O}_{V'} \otimes_{\mathcal{O}_{V'}} \mathcal{O}_{V'} \ar[r] &
\mathcal{O}_{V'}
}
$$
이 가환임을 확인하면 된다. 임의의 $n$에 대하여 $Y_n$으로 제한한 뒤,
즉 완비화 함자를 적용한 뒤에는 이 도식이 가환함을 알 수 있다.
따라서 완비화 함자의 충실성으로부터, 어떤 $V'' \geq V'$가 존재하여
$\sigma|_{V''}$가 대수 준동형이 된다는 결론을 얻는다.
\end{proof}

\begin{lemma}
\label{pione-lemma-restriction-equivalence-general}
 $X$를 뇌터 스킴이라 하고, $Y \subset X$를 아이디얼층
$\mathcal{I} \subset \mathcal{O}_X$로 정의되는 닫힌 부분스킴이라 하자.
$\mathcal{V}$를 $Y$를 포함하는 열린 부분스킴 $V \subset X$들의 집합이라 하고
역포함으로 순서화하자. 완비화 함자
$$
\colim_\mathcal{V} \textit{Coh}(\mathcal{O}_V)
\longrightarrow
\textit{Coh}(X, \mathcal{I}),
\quad
\mathcal{F} \longmapsto \mathcal{F}^\wedge
$$
가 유한 국소 자유 대상들의 충만 부분범주 사이의 동치를 정한다고 가정하자
(위 설명을 보라). 그러면 제한 함자
$$
\colim_\mathcal{V} \textit{F\'Et}_V \to \textit{F\'Et}_Y
$$
는 동치이다.
\end{lemma}

\begin{proof}
 $\mathcal{V}$는 유향 집합이므로 여극한은 Categories, 절
\ref{categories-section-directed-colimits}에서 설명한 것이다.
나머지 논증은 보조정리 \ref{pione-lemma-restriction-equivalence}의 증명과
거의 같으므로 독자는 건너뛰어도 된다.

\medskip\noindent
제한 함자가 충실충만이라는 것은 보조정리
\ref{pione-lemma-restriction-fully-faithful-general}에서 따른다.

\medskip\noindent
유한 에탈 사상 $U_1 \to Y$를 택하자. 증명을 끝내기 위해
$U_1$이 제한 함자의 본질적 상에 속함을 보이자.

\medskip\noindent
 $n \geq 1$에 대하여 $Y_n$을 $Y$의 $n$차 무한소 근방이라 하자.
보조정리 \ref{pione-lemma-thickening}에 의해 유일한 유한 에탈 사상
$\pi_n : U_n \to Y_n$이 존재하며, 그 $Y = Y_1$로의 기저변환은
$U_1 \to Y_1$을 복원한다. 층
$\mathcal{F}_n = \pi_{n, *}\mathcal{O}_{U_n}$을 생각하자.
$\mathcal{F}_n$을 $X$ 위의 $\mathcal{O}_X$-가군으로 보면, 국소적으로
$(\mathcal{O}_X/f^{n + 1}\mathcal{O}_X)^{\oplus r}$.
이 $(\mathcal{F}_n)$은 $\textit{Coh}(X, \mathcal{I})$의 유한 국소 자유 대상이다.
가정에 의해 어떤 $V \in \mathcal{V}$, 유한 국소 자유
$\mathcal{O}_V$-가군 $\mathcal{F}$, 그리고 양립하는 동형들의 계
$$
\mathcal{F}/\mathcal{I}^n\mathcal{F} \to \mathcal{F}_n
$$
이 존재한다. 이들은 $\mathcal{O}_V$-가군의 동형이다.

\medskip\noindent
 $\mathcal{F}$ 위의 대수 구조를 구성하기 위해 곱셈 사상
$\mathcal{F}_n \otimes_{\mathcal{O}_V} \mathcal{F}_n \to \mathcal{F}_n$을 생각하자.
이는 $\mathcal{F}_n = \pi_{n, *}\mathcal{O}_{U_n}$이 대수층이므로 주어진다.
이들은 다음 사상을 정한다.
$$
(\mathcal{F}\otimes_{\mathcal{O}_V} \mathcal{F})^\wedge
\longrightarrow
\mathcal{F}^\wedge
$$
이는 범주 $\textit{Coh}(X, \mathcal{I})$의 사상이다. 따라서 가정에 의해
$V$를 축소한 뒤 다음 사상이 존재한다고 가정할 수 있다.
$\mu : \mathcal{F}\otimes_{\mathcal{O}_V} \mathcal{F} \to \mathcal{F}$
그 $Y_n$으로의 제한들은 위의 곱셈 사상들을 준다.
필요하다면 열린 부분을 더 축소하여 $\mu$가
$\mathcal{F}$ 위에 가환 $\mathcal{O}_V$-대수 구조를 정하고
$\mathcal{F}_n$ 위의 주어진 대수 구조와 양립한다고 가정할 수 있다.
다음과 같이 놓자.
$$
U = \underline{\Spec}_V((\mathcal{F}, \mu))
$$
그러면 $V$ 위의 유한 국소 자유 스킴을 얻고, 그 $Y$로의 제한은
$U_1$과 동형이다. 따라서 $U \to V$는 $Y$ 위에 놓이는 모든 점에서
에탈이다. More on Morphisms, 보조정리
\ref{more-morphisms-lemma-check-smoothness-on-infinitesimal-nbhds}.
에 의해 $V$를 한 번 더 축소하여 $U \to V$가 유한 에탈이라고
가정할 수 있다. 이로써 증명이 끝난다.
\end{proof}

\begin{lemma}
\label{pione-lemma-restriction-faithful}
 $X$를 스킴이라 하고 $Y \subset X$를 닫힌 부분스킴이라 하자.
$X$의 모든 연결 성분이 $Y$와 만나면
제한 함자 $\textit{F\'Et}_X \to \textit{F\'Et}_Y$는 충실하다.
\end{lemma}

\begin{proof}
 $a, b : U \to U'$를 $X$ 위의 유한 에탈 스킴들 사이의 사상이라 하고,
그 $Y$로의 제한들이 같다고 하자. $U$의 연결 성분의 상은
$X$의 연결 성분이다. 이는 $U \to X$를 $X$의 연결 성분 위로 제한한 뒤
Topology, 보조정리
\ref{topology-lemma-finite-fibre-connected-components}를 적용하면 알 수 있다.
따라서 가정에 의해 $U$의 모든 연결 성분의 상은 $Y$와 만난다.
그러므로 \'Etale Morphisms, 명제
\ref{etale-proposition-equality}에 의해 $a = b$이다. 이 결론은
$U$의 각 연결 성분에 제한하여 얻는다. $a$와 $b$의 등화자가
$U$의 열린 부분스킴이므로($U'$의 대각 사상이 $X$ 위에서 열린 사상이기 때문이다)
결론이 따른다.
\end{proof}

\begin{lemma}
\label{pione-lemma-restriction-fully-faithful-special}
 $X$를 뇌터 스킴이라 하고 $Y \subset X$를 닫힌 부분스킴이라 하자.
$Y_n \subset X$를 $X$에서 $Y$의 $n$차 무한소 근방이라 하자.
다음 중 하나가 성립한다고 가정하자.
\begin{enumerate}
\item $X$가 준아핀이고
$\Gamma(X, \mathcal{O}_X) \to \lim \Gamma(Y_n, \mathcal{O}_{Y_n})$
가 동형이거나,
\item $X$가 풍부한 가역 가군 $\mathcal{L}$을 가지며,
$\Gamma(X, \mathcal{L}^{\otimes m}) \to
\lim \Gamma(Y_n, \mathcal{L}^{\otimes m}|_{Y_n})$
가 모든 $m \gg 0$에 대하여 동형이거나,
\item 모든 유한 국소 자유 $\mathcal{O}_X$-가군 $\mathcal{E}$에 대하여 사상
$\Gamma(X, \mathcal{E}) \to \lim \Gamma(Y_n, \mathcal{E}|_{Y_n})$
가 동형이다.
\end{enumerate}
그러면 제한 함자 $\textit{F\'Et}_X \to \textit{F\'Et}_Y$는 충실충만이다.
\end{lemma}

\begin{proof}
이 보조정리는 보조정리 \ref{pione-lemma-restriction-fully-faithful}와
Algebraic and Formal Geometry, 보조정리
\ref{algebraization-lemma-completion-fully-faithful}에서 형식적으로 따른다.
\end{proof}

\begin{lemma}
\label{pione-lemma-restriction-fully-faithful-general-special}
 $X$를 뇌터 스킴이라 하고 $Y \subset X$를 닫힌 부분스킴이라 하자.
$Y_n \subset X$를 $X$에서 $Y$의 $n$차 무한소 근방이라 하자.
$\mathcal{V}$를 $Y$를 포함하는 열린 부분스킴 $V \subset X$들의 집합이라 하고
역포함으로 순서화하자. 다음 중 하나가 성립한다고 가정하자.
\begin{enumerate}
\item $X$가 준아핀이고
$$
\colim_\mathcal{V} \Gamma(V, \mathcal{O}_V)
\longrightarrow
\lim \Gamma(Y_n, \mathcal{O}_{Y_n})
$$
가 동형이거나,
\item $X$가 풍부한 가역 가군 $\mathcal{L}$을 가지며,
$$
\colim_\mathcal{V} \Gamma(V, \mathcal{L}^{\otimes m})
\longrightarrow
\lim \Gamma(Y_n, \mathcal{L}^{\otimes m}|_{Y_n})
$$
가 모든 $m \gg 0$에 대하여 동형이거나,
\item 모든 $V \in \mathcal{V}$와 모든 유한 국소 자유
$\mathcal{O}_V$-가군 $\mathcal{E}$에 대하여 사상
$$
\colim_{V' \geq V} \Gamma(V', \mathcal{E}|_{V'})
\longrightarrow
\lim \Gamma(Y_n, \mathcal{E}|_{Y_n})
$$
가 동형이다.
\end{enumerate}
그러면 함자
$$
\colim_\mathcal{V} \textit{F\'Et}_V \to \textit{F\'Et}_Y
$$
는 충실충만이다.
\end{lemma}

\begin{proof}
이 보조정리는 보조정리 \ref{pione-lemma-restriction-fully-faithful-general}과
Algebraic and Formal Geometry, 보조정리
\ref{algebraization-lemma-completion-fully-faithful-general}에서 형식적으로 따른다.
\end{proof}







\section{밂과 기본군}
\label{pione-section-pushouts}

\noindent
주요 결과는 다음과 같다.

\begin{lemma}
\label{pione-lemma-pushout-along-closed-immersion-and-integral}
More on Morphisms, 상황
\ref{more-morphisms-situation-pushout-along-closed-immersion-and-integral}
에서 예를 들어 $Z \to Y$와 $Z \to X$가 스킴들의 닫힌 몰입이면,
범주의 동치
$$
\textit{F\'Et}_{Y \amalg_Z X}
\longrightarrow
\textit{F\'Et}_Y
\times_{\textit{F\'Et}_Z}
\textit{F\'Et}_X
$$
가 존재한다.
\end{lemma}

\begin{proof}
More on Morphisms, 명제
\ref{more-morphisms-proposition-pushout-along-closed-immersion-and-integral}.
에 의해 밂은 존재한다. 이 함자는 밂 위의 유한 에탈 스킴 $U$를
기저변환 $Y' = U \times_{Y \amalg_Z X} Y$와
$X' = U \times_{Y \amalg_Z X} X$, 그리고 $Z$ 위의 자연스러운 동형
$Y' \times_Y Z \to X' \times_X Z$로 보냄으로써 주어진다.
이 함자가 동치임을 보이기 위해 More on Morphisms, 보조정리
\ref{more-morphisms-lemma-pushout-functor-equivalence-flat}
를 사용하여 준역함자를 구성한다. 이제 다음만 보이면 된다.
분리이고 준유한이며 에탈인 사상 $U \to Y \amalg_Z X$가 주어지고
$X' \to X$와 $Y' \to Y$가 유한이면, $U \to Y \amalg_Z X$는 유한이다.
이는 대응하는 대수적 사실(More on Algebra, 보조정리
\ref{more-algebra-lemma-finite-module-over-fibre-product})
에서 따른다. 또는
$$
X' \amalg Y' \to U
$$
가 전사이고 $X'$와 $Y'$가 $Y \amalg_Z X$ 위에서 고유임을 이용할 수도 있다
(여기서는 More on Morphisms, 명제
\ref{more-morphisms-proposition-pushout-along-closed-immersion-and-integral})
의 밂에 대한 기술을 사용한다). 이제 Morphisms, 보조정리
\ref{morphisms-lemma-scheme-theoretic-image-is-proper}를 적용하면
$U$가 $Y \amalg_Z X$ 위에서 고유임을 알 수 있다.
준유한 고유 사상은 유한이므로(More on Morphisms, 보조정리
\ref{more-morphisms-lemma-characterize-finite}), 결론이 따른다.
\end{proof}






\section{천공 스펙트럼의 유한 에탈 덮개 I}
\label{pione-section-pi1-punctured-spec}

\noindent
먼저 Lefschetz형 결과 몇 가지를 증명한다.

\begin{situation}
\label{pione-situation-local-lefschetz}
 $(A, \mathfrak m)$을 뇌터 국소환이라 하고 $f \in \mathfrak m$이라 하자.
$X = \Spec(A)$ 및 $X_0 = \Spec(A/fA)$로 놓고,
$U = X \setminus \{\mathfrak m\}$ 및
$U_0 = X_0 \setminus \{\mathfrak m\}$를 각각
$A$ 및 $A/fA$의 천공 스펙트럼이라 하자.
\end{situation}

\noindent
스킴 $X$에 대하여 $X$ 위의 유한 에탈 스킴들의 범주를
$\textit{F\'Et}_X$로 나타냈음을 상기하자. 절
\ref{pione-section-finite-etale}을 보라.
상황 \ref{pione-situation-local-lefschetz}에서는 다음 기저변환 함자들을 조사한다.
$$
\xymatrix{
\textit{F\'Et}_X \ar[d] \ar[r] & \textit{F\'Et}_U \ar[d] \\
\textit{F\'Et}_{X_0} \ar[r] & \textit{F\'Et}_{U_0}
}
$$
많은 경우 오른쪽 세로 화살표는 충실하다.

\begin{lemma}
\label{pione-lemma-faithful}
상황 \ref{pione-situation-local-lefschetz}에서 다음 중 하나가
성립한다고 가정하자.
\begin{enumerate}
\item $f \not \in \mathfrak p$를 만족하는 모든 극소 소아이디얼
$\mathfrak p \subset A$에 대하여 $\dim(A/\mathfrak p) \geq 2$이거나,
\item $U$의 모든 연결 성분이 $U_0$와 만난다.
\end{enumerate}
그러면
$$
\textit{F\'Et}_U \longrightarrow \textit{F\'Et}_{U_0},\quad
V \longmapsto V_0 = V \times_U U_0
$$
는 충실 함자이다.
\end{lemma}

\begin{proof}
경우 (2)는 보조정리 \ref{pione-lemma-restriction-faithful}에서 곧바로 따른다.
가정 (1) 아래에서는 $U$의 모든 기약 성분이 $U_0$와 만난다.
Algebra, 보조정리 \ref{algebra-lemma-one-equation}을 보라.
따라서 (1)은 (2)로부터 따른다.
\end{proof}

\noindent
더 흥미로운 결과를 증명하기 전에 몇 가지 보조정리가 필요하다.

\begin{lemma}
\label{pione-lemma-fill-in-missing}
상황 \ref{pione-situation-local-lefschetz}에서 $V \to U$를 유한 사상이라 하자.
$A^\wedge$를 $A$의 $\mathfrak m$-진 완비화라 하고
$X' = \Spec(A^\wedge)$로 놓자. 또한 $U'$와 $V'$를 각각
$U$와 $V$의 $X'$로의 기저변환이라 하자. 유한 사상 $Y' \to X'$가
$V' = Y' \times_{X'} U'$를 만족하면, 유한 사상 $Y \to X$로서
$V = Y \times_X U$ 및 $Y' = Y \times_X X'$를 만족하는 것이 존재한다.
\end{lemma}

\begin{proof}
이는 More on Algebra, 명제
\ref{more-algebra-proposition-equivalence}의 직접적인 응용이다.
실제로 $f_1, \ldots, f_t$를 $\mathfrak m$의 생성원들로 택하자.
각 $i$에 대하여 $V \times_U D(f_i) = \Spec(B_i)$로 쓰자.
$1 \leq i, j \leq n$에 대하여 두 스펙트럼 모두
$V \times_U D(f_if_j)$를 나타내므로, $A_{f_if_j}$-대수들 사이의 동형
$\alpha_{ij} : (B_i)_{f_j} \to (B_j)_{f_i}$를 얻는다.
$Y' = \Spec(B')$로 쓰자. $V \times_U U' = Y \times_{X'} U'$이므로
동형 $\alpha_i : B'_{f_i} \to B_i \otimes_A A^\wedge$를 얻는다.
직접적인 논증으로 $(B', B_i, \alpha_i, \alpha_{ij})$가
$\text{Glue}(A \to A^\wedge, f_1, \ldots, f_t)$의 대상임을 알 수 있다.
More on Algebra, 주의 \ref{more-algebra-remark-glueing-data}를 보라.
위에서 인용한 명제를 적용하여(대수 구조를 얻으려면 More on Algebra, 주의
\ref{more-algebra-remark-formal-glueing-algebras}를 사용한다), $A$-대수 $B$로서
$\text{Can}(B)$가 $(B', B_i, \alpha_i, \alpha_{ij})$와 동형인 것을 얻는다.
$Y = \Spec(B)$로 놓으면 $Y \to X$가 사상이고, 원하는 양립하는 동형
$V \cong Y \times_X U$ 및 $Y' = Y \times_X X'$를 갖는다는 것을 알 수 있다.
\end{proof}

\begin{lemma}
\label{pione-lemma-fully-faithful-henselian-completion}
상황 \ref{pione-situation-local-lefschetz}에서 $A$가 헨젤이거나,
더 일반적으로 $(A, (f))$가 헨젤 쌍이라고 가정하자.
$A^\wedge$를 $A$의 $\mathfrak m$-진 완비화라 하고
$X' = \Spec(A^\wedge)$로 놓자. 또한 $U'$와 $U'_0$를 각각
$U$와 $U_0$의 $X'$로의 기저변환이라 하자.
$\textit{F\'Et}_{U'} \to \textit{F\'Et}_{U'_0}$가 충실충만이면
$\textit{F\'Et}_U \to \textit{F\'Et}_{U_0}$도 충실충만이다.
\end{lemma}

\begin{proof}
 $\textit{F\'Et}_{U'} \longrightarrow \textit{F\'Et}_{U'_0}$가
충실충만이라고 가정하자. $X' \to X$가 충실 평탄이므로, 함자
$V \to V_0 = V \times_U U_0$가 충실임은 곧바로 알 수 있다.
유한 에탈 덮개들의 범주에는 내부 Hom이 있으므로(보조정리
\ref{pione-lemma-internal-hom-finite-etale}), $V$가 $U$ 위에서 유한 에탈일 때
다음을 보이면 충분하다.
$$
\Mor_U(U, V) = \Mor_{U_0}(U_0, V_0)
$$
따라서 $U_0$ 위의 사상 $s_0 : U_0 \to V_0$가 주어졌다고 하고,
$U$ 위의 사상 $s : U \to V$를 구성하자.

\medskip\noindent
가정에 의해 사상 $s' : U' \to V'$로서 그 $V'_0$로의 제한이
$s'_0$, 즉 $s_0$의 기저변환인 것이 존재한다.
$V' \to U'$가 유한 에탈이므로 $V' = s'(U') \amalg W'$이고
$W' \to U'$는 유한 에탈이다. 유한 사상 $Z' \to X'$로서
$W' = Z' \times_{X'} U'$를 만족하는 것을 택하자. 이는 More on Morphisms, 보조정리
\ref{more-morphisms-lemma-quasi-finite-separated-pass-through-finite}
에 서술된 형태의 Zariski 주정리에 의해 가능하다
(한 가지 세부 사항은 생략한다). 그러면
$$
V' = s'(U') \amalg W' \longrightarrow X' \amalg Z' = Y'
$$
는 열린 몰입이고 $V' = Y' \times_{X'} U'$를 만족한다.
보조정리 \ref{pione-lemma-fill-in-missing}에 의해 유한 사상 $Y \to X$로서
$V = Y \times_X U$ 및 $Y' = Y \times_X X'$를 만족하는 것을 얻는다.
$Y = \Spec(B)$로 쓰면 $Y' = \Spec(B \otimes_A A^\wedge)$이다.
이때 $B \otimes_A A^\wedge$는 멱등원 $e'$을 갖는다. 이는
$Y' = X' \amalg Z'$에서 열린닫힌 부분스킴 $X'$에 대응한다.

\medskip\noindent
먼저 $A$가 헨젤인 경우를 생각하자(이쪽이 조금 더 쉽다).
상 $\overline{e}$, 즉 $e'$의
$B \otimes_A \kappa(\mathfrak m) = B/\mathfrak mB$에서의 상은
$B$의 멱등원 $e$로 올라간다. 이는 $A$가 헨젤이기 때문이다
(Algebra, 보조정리 \ref{algebra-lemma-characterize-henselian}에 의해
$B$는 국소환들의 곱이다). 멱등원 올림의 유일성에 의해 $e$는
$e'$로 사상된다(여기서는 $B \otimes_A A^\wedge$가 국소환들의 곱임을 사용한다).
$e$에 대응하는 열린닫힌 부분스킴을 $Y_1 \subset Y$라 하자.
그러면 $Y_1 \times_X X' = s'(X')$이고, 충실 평탄 내림에 의해
$Y_1 \to X$가 동형임이 따른다. 이것이 원하는 단면을 준다.

\medskip\noindent
다음으로 $(A, (f))$가 헨젤 쌍인 경우를 생각하자. 여기서는
$s'$가 $s'_0$의 올림임을 사용한다. 구체적으로,
$Y_{0, 1} \subset Y_0 = Y \times_X X_0$를
$s_0(U_0) \subset V_0 = Y_0 \times_{X_0} U_0$의 폐포라 하자.
$X' \to X$가 평탄이므로 기저변환 $Y'_{0, 1} \subset Y'_0$는
$s'_0(U'_0)$의 폐포이며, 이는 $X'_0 \subset Y'_0$와 같다
(Morphisms, 보조정리
\ref{morphisms-lemma-flat-base-change-scheme-theoretic-image}를 보라).
$Y'_0 \to Y_0$는 상사상이므로(Morphisms, 보조정리
\ref{morphisms-lemma-fpqc-quotient-topology}), $Y_{0, 1}$은
$Y_0$의 열린닫힌 부분스킴이다. 대응하는 멱등원을
$e_0 \in B/fB$라 하자. More on Algebra, 보조정리
\ref{more-algebra-lemma-characterize-henselian-pair}
에 의해 $e_0$를 멱등원 $e \in B$로 올릴 수 있다.
나머지는 앞의 경우와 같이 결론을 얻는다.
\end{proof}

\noindent
상황 \ref{pione-situation-local-lefschetz}에서 제한 함자
$\textit{F\'Et}_U \longrightarrow \textit{F\'Et}_{U_0}$
의 충실충만성은 상당히 약한 가정 아래에서도 성립한다.
특히 그 가정들은 $U$가 연결 스킴임을 보장하지 않지만, 결론은
$U$와 $U_0$의 연결 성분 수가 같음을 보장한다.

\begin{lemma}
\label{pione-lemma-fully-faithful-simple}
상황 \ref{pione-situation-local-lefschetz}에서 다음을 가정하자.
\begin{enumerate}
\item[(a)] $A$는 쌍대화 복합체를 갖는다.
\item[(b)] 쌍 $(A, (f))$는 헨젤이다.
\item[(c)] 다음 중 하나가 성립한다.
\begin{enumerate}
\item[(i)] $A_f$는 $(S_2)$를 만족하고, $X_0$에 포함되지 않는
$X$의 모든 기약 성분은 차원이 $\geq 3$이다.
\item[(ii)] $f \not \in \mathfrak p$를 만족하는 임의의 소아이디얼
$\mathfrak p \subset A$에 대하여
$\text{depth}(A_\mathfrak p) + \dim(A/\mathfrak p) > 2$.
\end{enumerate}
\end{enumerate}
그러면 제한 함자
$\textit{F\'Et}_U \longrightarrow \textit{F\'Et}_{U_0}$
는 충실충만이다.
\end{lemma}

\begin{proof}
$A'$을 $\mathfrak m$-진 $A$-완비화라 하자. 가정들이 $A'$에 대해서도
성립함을 보이자. 조건 (a)와 (b)는 명백하다. 조건 (c)(ii)는
Local Cohomology, 보조정리 \ref{local-cohomology-lemma-change-completion}
에 의해 보존된다. 이제 (c)(i)가 성립한다고 가정하자. $A$가 보편적 카테너리이므로
(Dualizing Complexes, 보조정리 \ref{dualizing-lemma-universally-catenary}),
$\Spec(A')$의 기약 성분 가운데 $V(f)$에 포함되지 않는 것은 모두 차원이
$\geq 3$이다. More on Algebra, 명제
\ref{more-algebra-proposition-ratliff}를 보라.
$A \to A'$가 Gorenstein 올을 갖는 평탄 사상이므로,
$A_f$가 $(S_2)$를 만족하면 $A'_f$도 $(S_2)$를 만족한다.
여기서 사용한 참조는 Dualizing Complexes, 절
\ref{dualizing-section-formal-fibres}, More on Algebra, 절
\ref{more-algebra-section-properties-formal-fibres}, 그리고 Algebra, 보조정리
\ref{algebra-lemma-Sk-goes-up}이다. 따라서 보조정리
\ref{pione-lemma-fully-faithful-henselian-completion}에 의해 $A$가
완비 뇌터 국소환이라고 가정해도 된다.

\medskip\noindent
다른 가정들에 더하여 $A$가 완비 국소환이라고 가정하자.
보조정리 \ref{pione-lemma-restriction-fully-faithful}에 의해 결과는
Algebraic and Formal Geometry, 보조정리
\ref{algebraization-lemma-fully-faithful-simple-two}에서 따른다.
\end{proof}

\begin{lemma}
\label{pione-lemma-fully-faithful-minimal}
\begin{reference}
\cite[Corollary 1.11]{Bhatt-local}
\end{reference}
상황 \ref{pione-situation-local-lefschetz}에서 다음을 가정하자.
\begin{enumerate}
\item $H^1_\mathfrak m(A)$와 $H^2_\mathfrak m(A)$는
$f$의 어떤 거듭제곱에 의해 소멸하고,
\item $A$는 헨젤이거나, 더 일반적으로 $(A, (f))$는 헨젤 쌍이다.
\end{enumerate}
그러면 제한 함자
$\textit{F\'Et}_U \longrightarrow \textit{F\'Et}_{U_0}$
는 충실충만이다.
\end{lemma}

\begin{proof}
보조정리 \ref{pione-lemma-fully-faithful-henselian-completion}에 의해
$A$가 완비 뇌터 국소환이라고 가정해도 된다.
(가정들은 그대로 보존된다. Dualizing Complexes, 보조정리
\ref{dualizing-lemma-torsion-change-rings}를 사용한다.)
보조정리 \ref{pione-lemma-restriction-fully-faithful}에 의해 결과는
Algebraic and Formal Geometry, 보조정리
\ref{algebraization-lemma-fully-faithful-alternative}에서 따른다.
\end{proof}

\begin{lemma}
\label{pione-lemma-fully-faithful}
상황 \ref{pione-situation-local-lefschetz}에서 $A$의 깊이가 $\geq 3$이고,
$A$는 헨젤이거나, 더 일반적으로 $(A, (f))$가 헨젤 쌍이라고 가정하자.
그러면 제한 함자
$\textit{F\'Et}_U \to \textit{F\'Et}_{U_0}$
는 충실충만이다.
\end{lemma}

\begin{proof}
깊이에 대한 가정으로부터
$H^1_\mathfrak m(A) = H^2_\mathfrak m(A) = 0$이다.
Dualizing Complexes, 보조정리 \ref{dualizing-lemma-depth}를 보라.
따라서 보조정리 \ref{pione-lemma-fully-faithful-minimal}를 적용할 수 있다.
\end{proof}



























\section{국소 경우의 순수성 I}
\label{pione-section-local-purity}

\noindent
$(A, \mathfrak m)$을 뇌터 국소환이라 하자. $X = \Spec(A)$로 놓고,
$U = X \setminus \{\mathfrak m\}$를 천공 스펙트럼이라 하자.
{\it $(A, \mathfrak m)$에 대해 순수성이 성립한다}는 것은 제한 함자
$$
\textit{F\'Et}_X \longrightarrow \textit{F\'Et}_U
$$
가 본질적으로 전사라는 뜻이다. 이 절에서는 $X$에서 초곡면 $X_0$, 즉 $X$ 안의
초곡면으로 옮길 때 이 문제가 어떻게 달라지는지 이해하고자 한다.
달리 말하면 천공 스펙트럼의 기본군에 대한 일종의 국소 Lefschetz 성질을 연구한다.
이 결과들은 국소 순수성에 관한 주요 결과, 즉 보조정리
\ref{pione-lemma-local-purity}, 명제
\ref{pione-proposition-purity-complete-intersection}, 그리고 명제
\ref{pione-proposition-purity-smooth-over-depth2}의 증명을 차원에 대한 귀납법으로
진행하는 데 유용하다.

\begin{lemma}
\label{pione-lemma-sections-over-punctured-spec}
$(A, \mathfrak m)$을 뇌터 국소환이라 하자. $X = \Spec(A)$로 놓고,
$U = X \setminus \{\mathfrak m\}$로 놓자. $\pi : Y \to X$를 유한 사상이라 하고,
$\text{depth}(\mathcal{O}_{Y, y}) \geq 2$가 모든 닫힌점
$y \in Y$에 대하여 성립한다고 가정하자.
그러면 $Y$는 $B = \mathcal{O}_Y(\pi^{-1}(U))$의 스펙트럼이다.
\end{lemma}

\begin{proof}
$V = \pi^{-1}(U)$로 놓고 $\pi' : V \to U$를 $\pi$의 제한이라 하자.
다음 $\mathcal{O}_X$-가군 사상을 생각하자.
$$
\pi_*\mathcal{O}_Y \longrightarrow j_*\pi'_*\mathcal{O}_V
$$
여기서 $j : U \to X$는 포함 사상이다. 이 사상에 Divisors, 보조정리
\ref{divisors-lemma-depth-2-hartog}를 적용할 수 있다고 주장한다.
그렇다면 $B = \Gamma(Y, \mathcal{O}_Y)$이고 보조정리가 성립한다.
$x \in X$를 닫힌점이라 하자. 다음을 보이면 충분하다.
$\text{depth}((\pi_*\mathcal{O}_Y)_x) \geq 2$.
$y_1, \ldots, y_n \in Y$를 $x$로 사상되는 점들이라 하자.
Algebra, 보조정리 \ref{algebra-lemma-depth-goes-down-finite}에 의해
$\text{depth}(\mathcal{O}_{Y, y_i}) \geq 2$가 $i = 1, \ldots, n$에 대하여
성립함을 보이면 충분하다.
이는 바로 보조정리의 가정이므로 증명이 끝난다.
\end{proof}

\begin{lemma}
\label{pione-lemma-reformulate-purity}
$(A, \mathfrak m)$을 뇌터 국소환이라 하자. $X = \Spec(A)$로 놓고,
$U = X \setminus \{\mathfrak m\}$로 놓자. $V$를 $U$ 위의 유한 에탈 스킴이라 하자.
$A$의 깊이가 $\geq 2$라고 가정하자. 다음은 동치이다.
\begin{enumerate}
\item $V = Y \times_X U$인 유한 에탈 사상 $Y \to X$가 존재한다.
\item $B = \Gamma(V, \mathcal{O}_V)$는 $A$ 위 유한 에탈이다.
\end{enumerate}
\end{lemma}

\begin{proof}
주어진 유한 에탈 사상을 $\pi : V \to U$라 하자.
(1)의 $Y$가 존재한다고 가정하자. $x \in X$를 $\mathfrak m$에 대응하는 점이라 하자.
$y \in Y$를 $x$로 사상되는 점이라 하자. 다음을 주장한다.
$\text{depth}(\mathcal{O}_{Y, y}) \geq 2$.
이는 $Y \to X$가 에탈이고, 따라서
$A = \mathcal{O}_{X, x}$와 $\mathcal{O}_{Y, y}$의 깊이가 같기 때문이다
(Algebra, 보조정리 \ref{algebra-lemma-apply-grothendieck}).
그러므로 보조정리 \ref{pione-lemma-sections-over-punctured-spec}를 적용할 수 있고
$Y = \Spec(B)$이다.

\medskip\noindent
함의 (2) $\Rightarrow$ (1)은 더 쉬우므로 세부 사항을 생략한다.
\end{proof}

\begin{lemma}
\label{pione-lemma-reformulate-purity-normal}
$(A, \mathfrak m)$을 뇌터 국소환이라 하자. $X = \Spec(A)$로 놓고,
$U = X \setminus \{\mathfrak m\}$로 놓자. $A$가 차원 $\geq 2$인 정규환이라고 가정하자.
함자
$$
\textit{F\'Et}_U \longrightarrow
\left\{
\begin{matrix}
\text{유한 정규 }A\text{-대수 }B\text{ 중} \\
\Spec(B) \to X\text{가 }U\text{ 위에서 에탈인 것}
\end{matrix}
\right\},
\quad
V \longmapsto \Gamma(V, \mathcal{O}_V)
$$
는 동치이다. 또한 $V = Y \times_X U$인 유한 에탈 사상 $Y \to X$가
존재할 필요충분조건은 $B = \Gamma(V, \mathcal{O}_V)$가 $A$ 위 유한 에탈인 것이다.
\end{lemma}

\begin{proof}
$\text{depth}(A) \geq 2$임에 주의하자. 이는 $A$가 정규이기 때문이다
(정규성에 관한 Serre의 판정법, Algebra, 보조정리
\ref{algebra-lemma-criterion-normal}).
따라서 마지막 명제는 보조정리 \ref{pione-lemma-reformulate-purity}에서 따른다.
유한 에탈 사상 $\pi : V \to U$가 주어졌을 때
$B = \Gamma(V, \mathcal{O}_V)$로 놓자. $B$가 정규이고 $A$ 위 유한임을 보이면
표시된 함자를 얻는다. 자명한 준역함자가 있으므로 보일 것은 이것뿐이다.

\medskip\noindent
$A$가 정규이므로 스킴 $V$도 정규이다
(Descent, 보조정리 \ref{descent-lemma-normal-local-smooth}).
따라서 $V$는 정역 스킴들의 유한 서로소 합이다
(Properties, 보조정리 \ref{properties-lemma-normal-Noetherian}).
그러므로 $V$가 정역 스킴이라고 가정해도 된다. 이 경우 함수체 $L$, 즉 $V$의 함수체는
(Morphisms, 절 \ref{morphisms-section-rational-maps})
$A$의 분수체의 유한 분리 확대이다
(이는 $V \to U$의 일반올을 살펴보고 Morphisms, 보조정리
\ref{morphisms-lemma-etale-over-field}를 사용하면 알 수 있다).
Algebra, 보조정리
\ref{algebra-lemma-Noetherian-normal-domain-finite-separable-extension}에 의해
$B' \subset L$을 $A$의 $L$ 안에서의 정수적 폐포라 하면, 이는 $A$ 위 유한이다.
More on Algebra, 보조정리 \ref{more-algebra-lemma-integral-closure-reflexive}에 의해
$B'$은 반사 $A$-가군이고, 따라서 More on Algebra, 보조정리
\ref{more-algebra-lemma-reflexive-over-normal}에 의해
$\text{depth}_A(B') \geq 2$이다.

\medskip\noindent
$f \in \mathfrak m$이라 하자. 그러면
$B_f = \Gamma(V \times_U D(f), \mathcal{O}_V)$이다
(Properties, 보조정리 \ref{properties-lemma-invert-f-sections}).
따라서 $B'_f = B_f$이다. 실제로 위에서 본 대로 $B_f$는 정규이고,
$A_f$ 위 유한이며 분수체가 $L$이다. 따라서 $V = \Spec(B') \times_X U$이다.
더 나아가 $B = B'$이다. 이는 보조정리
\ref{pione-lemma-sections-over-punctured-spec}를 $\Spec(B') \to X$에 적용하면 따른다.
이 보조정리를 적용할 수 있는 까닭은 국소화 $B'_{\mathfrak m'}$, 즉
$B'$의 극대 아이디얼 $\mathfrak m' \subset B'$ 중 $\mathfrak m$ 위에 놓이는 것에서
취한 국소화가 Algebra, 보조정리
\ref{algebra-lemma-depth-goes-down-finite}와 앞 문단의 깊이에 관한 관찰에 의해,
깊이 $\geq 2$를 갖기 때문이다.
\end{proof}

\begin{lemma}
\label{pione-lemma-purity-and-completion}
$(A, \mathfrak m)$을 뇌터 국소환이라 하자. $X = \Spec(A)$로 놓고,
$U = X \setminus \{\mathfrak m\}$로 놓자. $V$를 $U$ 위의 유한 에탈 스킴이라 하자.
$A^\wedge$를 $\mathfrak m$-진 $A$-완비화라 하고 $X' = \Spec(A^\wedge)$로 놓자.
$U'$와 $V'$를 각각 $U$와 $V$의 $X'$로의 기저변환이라 하자. 다음은 동치이다.
\begin{enumerate}
\item $V = Y \times_X U$인 유한 에탈 사상 $Y \to X$가 존재한다.
\item $V' = Y' \times_{X'} U'$인 유한 에탈 사상 $Y' \to X'$가 존재한다.
\end{enumerate}
\end{lemma}

\begin{proof}
함의 (1) $\Rightarrow$ (2)는 해 $Y \to X$의 기저변환을 취하면 따른다.
$Y' \to X'$를 (2)의 것으로 하자. 보조정리 \ref{pione-lemma-fill-in-missing}에 의해, 유한 사상
$Y \to X$로서 $V = Y \times_X U$이고 $Y' = Y \times_X X'$인 것을 취할 수 있다.
내림에 의해 $Y \to X$는 유한 에탈이다
(대수학의 보조정리 \ref{algebra-lemma-descend-properties-modules}와
\ref{algebra-lemma-etale}). 이것으로 증명이 끝난다.
\end{proof}

\noindent
다음 두 보조정리의 요점은 그 가정들만으로는 $A$의 깊이가 $\geq 3$임이
따르지 않는다는 데 있다. 예를 들어 $A$가 차원 $\geq 3$인 완비 정규
국소 정역이고 $f \in \mathfrak m$이 영이 아니면 이 가정들이 성립한다.

\begin{lemma}
\label{pione-lemma-lift-simple}
상황 \ref{pione-situation-local-lefschetz}에서 $V$를 $U$ 위의 유한 에탈 스킴이라 하자.
다음을 가정한다.
\begin{enumerate}
\item[(a)] $A$는 쌍대화 복합체를 가진다.
\item[(b)] 쌍 $(A, (f))$는 헨젤 쌍이다.
\item[(c)] 다음 중 하나가 성립한다.
\begin{enumerate}
\item[(i)] $A_f$는 $(S_2)$이고, $X_0$에 포함되지 않는 $X$의 모든 기약 성분의
차원은 $\geq 3$이다.
\item[(ii)] 모든 소 아이디얼 $\mathfrak p \subset A$에 대하여,
$f \not \in \mathfrak p$이면
$\text{depth}(A_\mathfrak p) + \dim(A/\mathfrak p) > 2$이다.
\end{enumerate}
\item[(d)] 어떤 유한 에탈 사상 $Y_0 \to X_0$에 대하여
$V_0 = V \times_U U_0$가 $Y_0 \times_{X_0} U_0$와 같다.
\end{enumerate}
그러면 어떤 유한 에탈 사상 $Y \to X$에 대하여 $V = Y \times_X U$이다.
\end{lemma}

\begin{proof}
보조정리 \ref{pione-lemma-purity-and-completion}을 사용하여 완비인 경우로 환원한다.
(가정들도 그대로 이어진다. 보조정리 \ref{pione-lemma-fully-faithful-simple}의 증명을 보라.)

\medskip\noindent
완비인 경우, 대수학 심화의 보조정리
\ref{more-algebra-lemma-finite-etale-equivalence}에 의해 $Y_0 \to X_0$를
유한 에탈 사상 $Y \to X$로 올릴 수 있다. 또한 대수학 심화의 보조정리
\ref{more-algebra-lemma-complete-henselian}에 의해
$(A, fA)$가 헨젤 쌍임에 유의하라. 이제 보조정리
\ref{pione-lemma-fully-faithful-simple}을 사용하면 $V$가 $Y \times_X U$와 동형임을 알 수 있다.
이것으로 증명이 끝난다.
\end{proof}

\begin{lemma}
\label{pione-lemma-lift-purity-general}
상황 \ref{pione-situation-local-lefschetz}에서 $V$를 $U$ 위의 유한 에탈 스킴이라 하자.
다음을 가정한다.
\begin{enumerate}
\item $H^1_\mathfrak m(A)$와 $H^2_\mathfrak m(A)$는
$f$의 어떤 거듭제곱에 의해 소멸한다.
\item 어떤 유한 에탈 사상 $Y_0 \to X_0$에 대하여
$V_0 = V \times_U U_0$가 $Y_0 \times_{X_0} U_0$와 같다.
\end{enumerate}
그러면 어떤 유한 에탈 사상 $Y \to X$에 대하여 $V = Y \times_X U$이다.
\end{lemma}

\begin{proof}
보조정리 \ref{pione-lemma-purity-and-completion}을 사용하여 완비인 경우로 환원한다.
(가정들도 그대로 이어진다. 쌍대화 복합체의 보조정리
\ref{dualizing-lemma-torsion-change-rings}를 사용하라.)

\medskip\noindent
완비인 경우, 대수학 심화의 보조정리
\ref{more-algebra-lemma-finite-etale-equivalence}에 의해 $Y_0 \to X_0$를
유한 에탈 사상 $Y \to X$로 올릴 수 있다. 또한 대수학 심화의 보조정리
\ref{more-algebra-lemma-complete-henselian}에 의해
$(A, fA)$가 헨젤 쌍임에 유의하라. 이제 보조정리
\ref{pione-lemma-fully-faithful-minimal}을 사용하면 $V$가 $Y \times_X U$와 동형임을 알 수 있다.
이것으로 증명이 끝난다.
\end{proof}

\begin{lemma}
\label{pione-lemma-lift-purity}
상황 \ref{pione-situation-local-lefschetz}에서 $V$를 $U$ 위의 유한 에탈 스킴이라 하자.
다음을 가정한다.
\begin{enumerate}
\item $A$의 깊이는 $\geq 3$이다.
\item 어떤 유한 에탈 사상 $Y_0 \to X_0$에 대하여
$V_0 = V \times_U U_0$가 $Y_0 \times_{X_0} U_0$와 같다.
\end{enumerate}
그러면 어떤 유한 에탈 사상 $Y \to X$에 대하여 $V = Y \times_X U$이다.
\end{lemma}

\begin{proof}
깊이에 관한 가정으로부터 $H^1_\mathfrak m(A) = H^2_\mathfrak m(A) = 0$이 따른다.
쌍대화 복합체의 보조정리 \ref{dualizing-lemma-depth}를 보라. 따라서 보조정리
\ref{pione-lemma-lift-purity-general}을 적용할 수 있다.
\end{proof}











\section{분기 자취의 순수성}
\label{pione-section-purity}

\noindent
유한 국소 자유 사상의 판별식을 사용할 것이다. 판별식의 절
\ref{discriminant-section-discriminant}을 보라.

\begin{lemma}
\label{pione-lemma-find-point-codim-1}
$(A, \mathfrak m)$을 $\dim(A) \geq 1$인 뇌터 국소환이라 하자.
$f \in \mathfrak m$이라 하자. 그러면 $\mathfrak p \in V(f)$이고
$\dim(A_\mathfrak p) = 1$인 것이 존재한다.
\end{lemma}

\begin{proof}
$\dim(A)$에 관한 귀납법을 쓴다. $\dim(A) = 1$이면
$\mathfrak p = \mathfrak m$으로 두면 된다. $\dim(A) > 1$이면,
$Z \subset \Spec(A)$를 차원 $> 1$인 기약 성분이라 하자. 그러면
$V(f) \cap Z$의 차원은 $> 0$이다
(대수학의 보조정리 \ref{algebra-lemma-one-equation}). 소 아이디얼
$\mathfrak q \in V(f) \cap Z$로서 $\mathfrak q \not = \mathfrak m$이고
$A$의 천공 스펙트럼의 닫힌점에 대응하는 것을 택한다.
이는 성질의 보조정리 \ref{properties-lemma-complement-closed-point-Jacobson}에 의해 가능하다.
이때 $\mathfrak q$는 $Z$의 일반점이 아니다. 따라서
$0 < \dim(A_\mathfrak q) < \dim(A)$이고 $f \in \mathfrak q A_\mathfrak q$이다.
차원에 관한 귀납법으로 $f \in \mathfrak p \subset A_\mathfrak q$이고
$\dim((A_\mathfrak q)_\mathfrak p) = 1$인 것을 찾을 수 있다. 그러면
$\mathfrak p \cap A$가 원하는 것이다.
\end{proof}

\begin{lemma}
\label{pione-lemma-ramification-quasi-finite-flat}
$f : X \to Y$를 국소 뇌터 스킴의 사상이라 하고 $x \in X$라 하자. 다음을 가정한다.
\begin{enumerate}
\item $f$는 평탄하다.
\item $f$는 $x$에서 준유한이다.
\item $x$는 $X$의 어느 기약 성분의 일반점도 아니다.
\item 특수화 $x' \leadsto x$로서 $\dim(\mathcal{O}_{X, x'}) = 1$인 것마다
$f$는 $x'$에서 비분기이다.
\end{enumerate}
그러면 $f$는 $x$에서 에탈이다.
\end{lemma}

\begin{proof}
$f$가 비분기인 점들의 집합은 $f$가 에탈인 점들의 집합과 같으며, 이 집합은 열린집합이다.
스킴의 사상의 정의 \ref{morphisms-definition-unramified},
\ref{morphisms-definition-etale} 및 보조정리
\ref{morphisms-lemma-flat-unramified-etale}를 보라. $f$가 $x$에서 에탈인지 확인할 때에는
밑과 공역 양쪽에서 에탈 국소적으로 작업해도 된다
(내림의 보조정리 \ref{descent-lemma-descending-property-etale} 및
\ref{descent-lemma-etale-etale-local-source}). 따라서 스킴의 사상 심화의 보조정리
\ref{more-morphisms-lemma-etale-makes-quasi-finite-finite-at-point}을 적용하여
$f : X \to Y$가 유한이고, $x$가 $X$에서 $y = f(x)$ 위에 놓이는 유일한 점이라고
가정할 수 있다. 그러면 $f$는 유한 국소 자유이다
(스킴의 사상의 보조정리 \ref{morphisms-lemma-finite-flat}).

\medskip\noindent
$f$가 유한 국소 자유이고 $x$가 $X$에서 $y = f(x)$ 위에 놓이는 유일한 점이라고 가정하자.
판별식의 보조정리 \ref{discriminant-lemma-discriminant}에 의해 국소 주인 닫힌 부분스킴
$D_\pi \subset Y$로서 다음 성질을 갖는 것을 얻는다. 즉, $y' \in D_\pi$인 것과
어떤 $x' \in X$에 대하여 $f(x') = y'$이고 $f$가 $x'$에서 분기하는 것이 동치이다.
따라서 $y \not \in D_\pi$임을 보이면 충분하다. 모순을 얻기 위해 $y \in D_\pi$라고 가정하자.

\medskip\noindent
조건 (3)에 의해 $\dim(\mathcal{O}_{X, x}) \geq 1$이다. 대수학의 보조정리
\ref{algebra-lemma-dimension-base-fibre-equals-total}에 의해
$\dim(\mathcal{O}_{X, x}) = \dim(\mathcal{O}_{Y, y})$이다. 보조정리
\ref{pione-lemma-find-point-codim-1}에 의해 $y' \in D_\pi$로서 $y$로 특수화되고
$\dim(\mathcal{O}_{Y, y'}) = 1$인 것을 택할 수 있다.
$f(x') = y'$이고 $f$가 분기하는 $x' \in X$를 택하자. $f$는 유한이므로 닫힌 사상이고,
따라서 $x' \leadsto x$이다. 앞에서와 마찬가지로
$\dim(\mathcal{O}_{X, x'}) = \dim(\mathcal{O}_{Y, y'}) = 1$이다.
이는 조건 (4)에 모순이다.
\end{proof}

\begin{lemma}
\label{pione-lemma-local-purity}
$(A, \mathfrak m)$을 차원 $d \geq 2$인 정칙 국소환이라 하자.
$X = \Spec(A)$ 및 $U = X \setminus \{\mathfrak m\}$이라 두자. 그러면
\begin{enumerate}
\item 함자 $\textit{F\'Et}_X \to \textit{F\'Et}_U$는 본질적으로 전사이다.
즉, $A$에 대하여 순수성이 성립한다.
\item $A \to B$가 유한이고 $B$가 정규이며 천공 스펙트럼 위에서 유한 에탈 사상을
유도하면, 이 사상은 에탈이다.
\end{enumerate}
\end{lemma}

\begin{proof}
대수학의 보조정리 \ref{algebra-lemma-regular-normal}에 의해 정칙 국소환은 정규임을
상기하자. 따라서 보조정리 \ref{pione-lemma-reformulate-purity-normal}에 의해 (1)과 (2)는
동치이다. $d$에 관한 귀납법으로 보조정리를 증명한다.

\medskip\noindent
$d = 2$인 경우. 이 경우 $A \to B$는 평탄하다. 실제로 대수학의 명제
\ref{algebra-proposition-going-down-normal-integral}에 의해 $A \to B$에 대하여 하강 정리가
성립한다. 그러면 대수학의 보조정리 \ref{algebra-lemma-dimension-base-fibre-equals-total}에 의해
모든 극대 아이디얼 $\mathfrak m' \subset B$에 대하여
$\dim(B_{\mathfrak m'}) = 2$이다. 또한 대수학의 보조정리
\ref{algebra-lemma-criterion-normal}에 의해 $B_{\mathfrak m'}$는 Cohen--Macaulay이다.
따라서, 이것이 핵심인데, 대수학의 보조정리
\ref{algebra-lemma-CM-over-regular-flat}을 적용하면 $A \to B_{\mathfrak m'}$가 평탄임을
알 수 있다. 이어서 대수학의 보조정리 \ref{algebra-lemma-flat-localization}에 의해
$A \to B$가 평탄이다. 그러므로 보조정리 \ref{pione-lemma-ramification-quasi-finite-flat}
(또는 더 쉬운 판별식의 보조정리
\ref{discriminant-lemma-discriminant}을 써서 직접 논증할 수도 있다)을 적용하면
$A \to B$가 에탈임을 알 수 있다.

\medskip\noindent
$d \geq 3$인 경우. $V \to U$를 유한 에탈이라 하자.
$f \in \mathfrak m_A$이고 $f \not \in \mathfrak m_A^2$인 원소를 택한다. 그러면
$A/fA$는 차원 $d - 1 \geq 2$인 정칙 국소환이다. 대수학의 보조정리
\ref{algebra-lemma-regular-ring-CM}을 보라. $U_0$를 $A/fA$의 천공 스펙트럼이라 하고,
$V_0 = V \times_U U_0$라 두자. 보조정리 \ref{pione-lemma-lift-purity}에 의해 $V_0$가 함자
$\textit{F\'Et}_{\Spec(A/fA)} \to \textit{F\'Et}_{U_0}$의 본질적 상에 속함을
보이면 충분하다. 이는 귀납 가정에서 따른다.
\end{proof}

\begin{lemma}[분기 자취의 순수성]
\label{pione-lemma-purity}
\begin{reference}
\cite{Nagata-Purity} 및 \cite[Exp. X, Thm. 3.1]{SGA1}
\end{reference}
\begin{history}
이 결과는 먼저 Zariski가 기하학적 형태로 진술하고 증명하였다
\cite{Zariski-Purity}. Nagata가 비완전체인 기초체로 일반화한 결과는 같은 권의
미국 국립과학원 회보에서 바로 다음 논문으로 \cite{Nagata-Remarks-Purity}에 발표되었다.
이듬해 Nagata는 뇌터 국소환에 대한 결과를 \cite{Nagata-Purity}에서 증명하였다.
그 증명은 완비 국소 정역에 관한 Bertini 정리인 Chow의 결과를 사용한다.
\cite{Chow-Bertini}를 보라. Bertini 정리들의 역사는 Kleiman의 역사 논문
\cite{Kleiman-Bertini}에서 논의된다. 몇 년 뒤 Auslander가 전혀 다른 증명을 얻었다.
\cite{Auslander-Purity}를 보라.
\end{history}
$f : X \to Y$를 국소 뇌터 스킴의 사상이라 하자. $x \in X$라 하고
$y = f(x)$라 두자. 다음을 가정한다.
\begin{enumerate}
\item $\mathcal{O}_{X, x}$는 정규이다.
\item $\mathcal{O}_{Y, y}$는 정칙이다.
\item $f$는 $x$에서 준유한이다.
\item $\dim(\mathcal{O}_{X, x}) = \dim(\mathcal{O}_{Y, y}) \geq 1$
\item 특수화 $x' \leadsto x$로서 $\dim(\mathcal{O}_{X, x'}) = 1$인 것마다
$f$는 $x'$에서 비분기이다.
\end{enumerate}
그러면 $f$는 $x$에서 에탈이다.
\end{lemma}

\begin{proof}
$d = \dim(\mathcal{O}_{X, x}) = \dim(\mathcal{O}_{Y, y})$에 관한 귀납법으로 증명한다.

\medskip\noindent
$d = 1$인 경우는 특별할 것이 없다. 이 경우 $f$가 $x$에서 비분기이고,
$\mathcal{O}_{Y, y}$가 이산 부치환환이라고 가정하였다
(대수학의 보조정리 \ref{algebra-lemma-characterize-dvr}). 그러면
$\mathcal{O}_{X, x}$는 $\mathcal{O}_{Y, y}$ 위에서 평탄하다
(그렇지 않으면 사상은 $x$에서 준유한일 수 없다). 따라서 $f$는 $x$에서 평탄이다.
평탄 $+$ 비분기인 사상은 에탈이므로 결론을 얻는다(몇몇 세부는 생략한다).

\medskip\noindent
$d \geq 2$인 경우. $d$에 관한 귀납법을 사용하여 보조정리
\ref{pione-lemma-local-purity}에서 논의한 경우로 환원한다. $f$가 $x$에서 에탈인지 확인할 때에는
밑과 공역 양쪽에서 에탈 국소적으로 작업해도 된다
(내림의 보조정리 \ref{descent-lemma-descending-property-etale} 및
\ref{descent-lemma-etale-etale-local-source}). 따라서 스킴의 사상 심화의 보조정리
\ref{more-morphisms-lemma-etale-makes-quasi-finite-finite-at-point}을 적용하여
$f : X \to Y$가 유한이고, $x$가 $X$에서 $y$ 위에 놓이는 유일한 점이라고 가정할 수 있다.
여기서는 국소환의 에탈 확대가 차원, 정규성, 정칙성을 바꾸지 않는다는 사실을 사용하였다.
대수학 심화의 절 \ref{more-algebra-section-permanence-etale} 및
에탈 사상의 절 \ref{etale-section-properties-permanence}을 보라.

\medskip\noindent
다음으로 $\Spec(\mathcal{O}_{Y, y})$에 의한 밑변환을 취하여
$Y$가 정칙 국소환의 스펙트럼이라고 가정할 수 있다. 실제로 $X$의 모든 점은
반드시 $x$로 특수화하므로 $X = \Spec(\mathcal{O}_{X, x})$이다.

\medskip\noindent
환 준동형 $\mathcal{O}_{Y, y} \to \mathcal{O}_{X, x}$는 유한이고, 차원이 같으므로
반드시 단사이다. 대수학의 명제 \ref{algebra-proposition-going-down-normal-integral}에 의해
$\mathcal{O}_{Y, y} \to \mathcal{O}_{X, x}$에 대하여 하강 정리가 성립한다
(대수학의 보조정리 \ref{algebra-lemma-regular-normal}에 의해 정칙환은 정규환이라는 사실도 사용한다).
$x' \in X$이고 $x' \not = x$이며 그 상이 $y' = f(x')$인 점을 택하자.
그러면 $\mathcal{O}_{X, x'}$는 정규 정역의 국소화이므로 정규이다. 마찬가지로
$\mathcal{O}_{Y, y'}$는 정칙이다(대수학의 보조정리
\ref{algebra-lemma-localization-of-regular-local-is-regular}를 보라). 대수학의 보조정리
\ref{algebra-lemma-dimension-base-fibre-equals-total}에 의해
$\dim(\mathcal{O}_{X, x'}) = \dim(\mathcal{O}_{Y, y'})$이다
(위에서 하강 정리를 확인하였다). 이 차원들은 물론 $d$보다 진정 작다. 실제로
$x' \not = x$이다.
$d$에 관한 귀납법으로 $f$가 $x'$에서 에탈임을 얻는다.

\medskip\noindent
이로써 다음 상황에 이르렀다. 차원 $d \geq 2$인 뇌터 국소환의 유한 국소 준동형
$A \to B$가 있고, $A$는 정칙이며 $B$는 정규이고, 이 준동형은 천공 스펙트럼 위에
유한 에탈 사상 $V \to U$를 유도한다. 목표는 $A \to B$가 에탈임을 보이는 것이다.
이는 보조정리 \ref{pione-lemma-local-purity}에서 따르므로 증명이 끝난다.
\end{proof}

\noindent
다음 보조정리는 유한 에탈 사상을 연장할 수 있는 최대 열린 부분집합을 찾을 때 유용하다.

\begin{lemma}
\label{pione-lemma-extend-S2}
$j : U \to X$를 국소 뇌터 스킴의 열린 몰입이라 하고,
모든 $x \not \in U$에 대하여 $\text{depth}(\mathcal{O}_{X, x}) \geq 2$라고 가정하자.
$\pi : V \to U$를 유한 에탈이라 하자. 그러면
\begin{enumerate}
\item $\mathcal{B} = j_*\pi_*\mathcal{O}_V$는 반사적인 연접
$\mathcal{O}_X$-대수이다. $Y = \underline{\Spec}_X(\mathcal{B})$라 두자.
\item $Y \to X$는 $V = Y \times_X U$이고
모든 $y \in Y \setminus V$에 대하여 $\text{depth}(\mathcal{O}_{Y, y}) \geq 2$인
유일한 유한 사상이다.
\item $Y \to X$가 $y$에서 에탈인 것과 $Y \to X$가
$y$에서 평탄인 것은 동치이다.
\item $Y \to X$가 에탈인 것과 $\mathcal{B}$가
$\mathcal{O}_X$-가군으로서 유한 국소 자유인 것은 동치이다.
\end{enumerate}
또한 (a) $\mathcal{B}$와 $Y \to X$의 구성은 국소 뇌터 스킴의 평탄 사상
$X' \to X$에 의한 밑변환과 가환한다. (b) $V' \to U'$가 유한 에탈 사상이고
$U \subset U' \subset X$가 열린집합이며 $U$ 위에서 $V \to U$와 일치하면,
$U'$ 위의 동형사상 $Y' \times_X U' = V'$가 유일하게 존재한다.
\end{lemma}

\begin{proof}
$\pi_*\mathcal{O}_V$는 유한 국소 자유 $\mathcal{O}_U$-가군이며, 특히 반사적이다.
인자의 보조정리 \ref{divisors-lemma-reflexive-S2-extend}에 의해
$j_*\pi_*\mathcal{O}_V$는 $X$ 위의 반사적 연접 가군 가운데
$U$ 위에서 $\pi_*\mathcal{O}_V$로 제한되는 유일한 것이다. 이로써 (1)이 증명되었다.

\medskip\noindent
구성에 의해 $Y \times_X U = V$이다. $\mathcal{B}$가 연접이므로 $Y \to X$는 유한이다.
인자의 보조정리 \ref{divisors-lemma-reflexive-S2}에 의해
모든 $x \in X \setminus U$에 대하여 $\text{depth}(\mathcal{B}_x) \geq 2$이다.
따라서 대수학의 보조정리 \ref{algebra-lemma-depth-goes-down-finite}에 의해
모든 $y \in Y \setminus V$에 대하여 $\text{depth}(\mathcal{O}_{Y, y}) \geq 2$이다.
거꾸로, $\pi' : Y' \to X$가 유한 사상이고 $V = Y' \times_X U$이며
모든 $y' \in Y' \setminus V$에 대하여
$\text{depth}(\mathcal{O}_{Y', y'}) \geq 2$라고 하자. 그러면
$\pi'_*\mathcal{O}_{Y'}$는 $U$ 위에서 $\pi_*\mathcal{O}_V$로 제한되고,
대수학의 보조정리 \ref{algebra-lemma-depth-goes-down-finite}에 의해
모든 $x \in X \setminus U$에 대하여
$\text{depth}((\pi'_*\mathcal{O}_{Y'})_x) \geq 2$이다. 그러면 예를 들어 인자의 보조정리
\ref{divisors-lemma-depth-2-hartog}에 의해 $\pi'_*\mathcal{O}_{Y'}$는
$j_*\pi_*\mathcal{O}_V$와 표준적으로 동형이다. 이로써 (2)가 증명되었다.

\medskip\noindent
$Y \to X$가 $y$에서 에탈이면 $Y \to X$는 $y$에서 평탄이다.
거꾸로 $Y \to X$가 $y$에서 평탄이라고 가정하자. $y \in V$이면
$Y \to X$는 $y$에서 에탈이다. $y \not \in V$이면 보조정리
\ref{pione-lemma-ramification-quasi-finite-flat}의 (1), (2), (3), (4)를 확인함으로써
$Y \to X$가 $y$에서 에탈임을 알 수 있다. (1)과 (2)는 명백하고,
$\text{depth}(\mathcal{O}_{Y, y}) \geq 2$이므로 (3)도 명백하다.
$y' \leadsto y$가 특수화이고 $\dim(\mathcal{O}_{Y, y'}) = 1$이면 $y' \in V$이다.
그렇지 않다면 대수학의 보조정리 \ref{algebra-lemma-bound-depth}에 의해 이 국소환의 깊이가
$2$가 되어 모순이기 때문이다. 따라서 $Y \to X$는 $y'$에서 에탈이고,
보조정리 \ref{pione-lemma-ramification-quasi-finite-flat}의 (4)도 성립한다. 이로써 (3)의 증명이 끝난다.

\medskip\noindent
(4)는 (3)과 다음 사실에서 따른다. $((Y \to X)_*\mathcal{O}_Y)_x$가 평탄한
$\mathcal{O}_{X, x}$-가군인 것과, $x$로 가는 모든 $y \in Y$에 대하여
$\mathcal{O}_{Y, y}$가 평탄한 $\mathcal{O}_{X, x}$-가군인 것은 동치이다.
대수학의 보조정리 \ref{algebra-lemma-flat-localization}을 보라. 또한 뇌터환 위의 유한 평탄 가군은
유한 국소 자유라는 사실도 사용한다. 대수학의 보조정리 \ref{algebra-lemma-finite-projective} 및
대수학의 보조정리 \ref{algebra-lemma-Noetherian-finite-type-is-finite-presentation}을 보라.

\medskip\noindent
보조정리의 마지막 주장들 가운데 (a)는 평탄 밑변환에서 따른다. 스킴의 코호몰로지의 보조정리
\ref{coherent-lemma-flat-base-change-cohomology}를 보라. (b)는 (2)의 유일성을 제한
$Y \times_X U'$에 적용하면 따른다.
\end{proof}

\begin{lemma}
\label{pione-lemma-extend-pure}
$j : U \to X$를 뇌터 스킴의 열린 몰입이라 하고, 모든 $x \not \in U$에 대하여
$\mathcal{O}_{X, x}$에서 순수성이 성립한다고 가정하자. 그러면
$$
\textit{F\'Et}_X \longrightarrow \textit{F\'Et}_U
$$
는 본질적으로 전사이다.
\end{lemma}

\begin{proof}
$V \to U$를 유한 에탈 사상이라 하자. 뇌터 귀납법에 의해 $V \to U$를
$X$의 진정 더 큰 열린 부분집합 위의 유한 에탈 사상으로 연장하면 충분하다.
$x \in X \setminus U$를 $X \setminus U$의 어떤 기약 성분의 일반점이라 하자.
이때 역상 $U_x$, 즉 $U$의 $\Spec(\mathcal{O}_{X, x})$ 안에서의 역상은
$\mathcal{O}_{X, x}$의 천공 스펙트럼이다. 가정에 의해 $V_x = V \times_U U_x$는
어떤 유한 에탈 사상 $Y_x \to \Spec(\mathcal{O}_{X, x})$를 $U_x$로 제한한 것이다.
극한의 보조정리 \ref{limits-lemma-glueing-near-point}에 의해, $x$를 포함하는 열린 부분스킴
$U \subset U' \subset X$와 유한 표시 사상 $V' \to U'$로서, $U$로 제한하면
$V \to U$를 복원하고 $\Spec(\mathcal{O}_{X, x})$로 제한하면 $Y_x$를 복원하는 것을 얻는다.
마지막으로 필요하면 극한의 보조정리
\ref{limits-lemma-glueing-near-point-properties}에 의해 $U'$를 더 작은 열린 부분집합으로 축소하여,
사상 $V' \to U'$가 유한 에탈이 되게 할 수 있다.
\end{proof}
















\section{천공 스펙트럼의 유한 에탈 덮개 II}
\label{pione-section-pi1-punctured-spec-II}

\noindent
이 절에서는 절 \ref{pione-section-pi1-punctured-spec}에서 논의한 내용의 몇 가지 변형을 증명한다.
뇌터 국소환 $(A, \mathfrak m)$과 $f \in \mathfrak m$이 주어졌다고 하자.
$X = \Spec(A)$ 및 $X_0 = \Spec(A/fA)$라 두고, $U = X \setminus \{\mathfrak m\}$와
$U_0 = X_0 \setminus \{\mathfrak m\}$를 각각 $A$와 $A/fA$의 천공 스펙트럼이라 하자.
이는 정확히 상황 \ref{pione-situation-local-lefschetz}와 같다. 다른 점은 제한 함자
$$
\colim_{U_0 \subset U' \subset U\text{ open}} \textit{F\'Et}_{U'}
\longrightarrow
\textit{F\'Et}_{U_0}
$$
를 고려한다는 것이다. 다시 말해, $U_0$의 유한 에탈 덮개를 $U$ 전체로 올리려는 것이 아니라
어떤 열린 근방 $U'$, 즉 $U$ 안에서 $U_0$의 열린 근방으로만 올리려 한다.

\begin{lemma}
\label{pione-lemma-faithful-general}
상황 \ref{pione-situation-local-lefschetz}에서 $U' \subset U$를 $U_0$를 포함하는 열린 부분집합이라 하자.
$\mathfrak p \subset A$가 $\mathfrak p \in U'$ 및 $\mathfrak p \not \in U_0$를 만족하는
극소 소 아이디얼이면 $\dim(A/\mathfrak p) \geq 2$라고 가정하자. 그러면
$$
\textit{F\'Et}_{U'} \longrightarrow \textit{F\'Et}_{U_0},\quad
V' \longmapsto V_0 = V' \times_{U'} U_0
$$
는 충실한 함자이다. 또한 이 가정을 만족하는 $U'$가 존재하며, $U_0$를 포함하는
더 작은 임의의 열린 부분집합 $U'' \subset U'$도 이 가정을 만족한다. 특히 제한 함자
$$
\colim_{U_0 \subset U' \subset U\text{ open}} \textit{F\'Et}_{U'}
\longrightarrow
\textit{F\'Et}_{U_0}
$$
는 충실하다.
\end{lemma}

\begin{proof}
대수학의 보조정리 \ref{algebra-lemma-one-equation}에 의해
$\dim(A/\mathfrak p) \geq 2$인 $A$의 모든 소 아이디얼 $\mathfrak p$에 대하여
$V(\mathfrak p)$가 $U_0$와 만난다. 따라서 보조정리 \ref{pione-lemma-restriction-faithful}에 의해
표시한 함자는 명제에 나온 것과 같은 $U$에 대하여 충실하다. 그런 $U'$가 존재함을 보이기 위해,
$\mathfrak p \subset A$가 $\mathfrak p \in U$, $\mathfrak p \not \in U_0$,
$\dim(A/\mathfrak p) = 1$을 만족하면 $\mathfrak p$는 $U$의 닫힌점에 대응하고, 따라서
$V(\mathfrak p) \cap U_0 = \emptyset$임에 유의하자. 그러므로 $U'$로서
$U_0$와 만나지 않으며 차원이 $1$인 $X$의 기약 성분들의 여집합을 취할 수 있다.
\end{proof}

\begin{lemma}
\label{pione-lemma-fully-faithful-general-better}
상황 \ref{pione-situation-local-lefschetz}에서 다음을 가정한다.
\begin{enumerate}
\item $A$는 쌍대화 복합체를 가지며 $f$-진 완비이다.
\item $X_0$에 포함되지 않는 $X$의 모든 기약 성분의 차원은 $\geq 3$이다.
\end{enumerate}
그러면 제한 함자
$$
\colim_{U_0 \subset U' \subset U\text{ open}} \textit{F\'Et}_{U'}
\longrightarrow
\textit{F\'Et}_{U_0}
$$
는 충만충실하다.
\end{lemma}

\begin{proof}
기본군의 위상적 불변성에 의해 증명할 때 $A$를 그 축약으로 바꾸어도 된다.
보조정리 \ref{pione-lemma-thickening}을 보라. 그러면 결과는 보조정리
\ref{pione-lemma-restriction-fully-faithful-general}과 대수적 및 형식적 기하의 보조정리
\ref{algebraization-lemma-fully-faithful-general}에서 따른다.
\end{proof}

\begin{lemma}
\label{pione-lemma-fully-faithful-general}
상황 \ref{pione-situation-local-lefschetz}에서 다음을 가정한다.
\begin{enumerate}
\item $A$는 $f$-진 완비이다.
\item $f$는 영인자가 아니다.
\item $H^1_\mathfrak m(A/fA)$는 유한 $A$-가군이다.
\end{enumerate}
그러면 제한 함자
$$
\colim_{U_0 \subset U' \subset U\text{ open}} \textit{F\'Et}_{U'}
\longrightarrow
\textit{F\'Et}_{U_0}
$$
는 충만충실하다.
\end{lemma}

\begin{proof}
보조정리 \ref{pione-lemma-restriction-fully-faithful-general}과
대수적 및 형식적 기하의 보조정리
\ref{algebraization-lemma-fully-faithful-general-alternative}에서 따른다.
\end{proof}








\section{천공 스펙트럼의 유한 에탈 덮개 III}
\label{pione-section-pi1-punctured-spec-III}

\noindent
이 절에서는 상황 \ref{pione-situation-local-lefschetz}에서 제한 함자
$$
\colim_{U_0 \subset U' \subset U\text{ open}} \textit{F\'Et}_{U'}
\longrightarrow
\textit{F\'Et}_{U_0}
$$
가 범주의 동치가 되는 조건을 조사한다.

\begin{lemma}
\label{pione-lemma-essentially-surjective-general-better}
상황 \ref{pione-situation-local-lefschetz}에서 다음을 가정한다.
\begin{enumerate}
\item $A$는 쌍대화 복합체를 가지며 $f$-진 완비이다.
\item 다음 중 하나가 성립한다.
\begin{enumerate}
\item $A_f$는 $(S_2)$이고, $X_0$에 포함되지 않는 $X$의 모든 기약 성분의
차원은 $\geq 4$이다.
\item $\mathfrak p \not \in V(f)$이고
$V(\mathfrak p) \cap V(f) \not = \{\mathfrak m\}$이면
$\text{depth}(A_\mathfrak p) + \dim(A/\mathfrak p) > 3$이다.
\end{enumerate}
\end{enumerate}
그러면 제한 함자
$$
\colim_{U_0 \subset U' \subset U\text{ open}} \textit{F\'Et}_{U'}
\longrightarrow
\textit{F\'Et}_{U_0}
$$
는 동치이다.
\end{lemma}

\begin{proof}
보조정리 \ref{pione-lemma-restriction-equivalence-general}과
대수적 및 형식적 기하의 보조정리
\ref{algebraization-lemma-equivalence-better}에서 따른다.
\end{proof}

\begin{lemma}
\label{pione-lemma-essentially-surjective-general}
상황 \ref{pione-situation-local-lefschetz}에서 다음을 가정한다.
\begin{enumerate}
\item $A$는 $f$-진 완비이다.
\item $f$는 영인자가 아니다.
\item $H^1_\mathfrak m(A/fA)$ 및 $H^2_\mathfrak m(A/fA)$는
유한 $A$-가군이다.
\end{enumerate}
그러면 제한 함자
$$
\colim_{U_0 \subset U' \subset U\text{ open}} \textit{F\'Et}_{U'}
\longrightarrow
\textit{F\'Et}_{U_0}
$$
는 동치이다.
\end{lemma}

\begin{proof}
보조정리 \ref{pione-lemma-restriction-equivalence-general}과
대수적 및 형식적 기하의 보조정리
\ref{algebraization-lemma-equivalence}에서 따른다.
\end{proof}

\begin{remark}
\label{pione-remark-combine}
$(A, \mathfrak m)$을 완비 뇌터 국소환이라 하고, $f \in \mathfrak m$을 영이 아닌 원소라 하자.
$A_f$가 $(S_2)$이고 $\Spec(A)$의 모든 기약 성분의 차원이 $\geq 4$라고 가정하자.
그러면 보조정리 \ref{pione-lemma-essentially-surjective-general-better}에 의해 범주
$$
\colim\nolimits_{U' \subset U\text{ open, }U_0 \subset U}
\text{유한 에탈 }U'\text{-스킴의 범주}
$$
는 유한 에탈 $U_0$-스킴의 범주와 동치이다. 예를 들어 $A$가 차원 $\geq 4$인 정규 정역이면
이 조건이 성립한다.
\end{remark}




\section{천공 스펙트럼의 유한 에탈 덮개 IV}
\label{pione-section-pi1-punctured-spec-IV}

\noindent
상황 \ref{pione-situation-local-lefschetz}에서와 같이 $X, X_0, U, U_0$를 잡는다.
이 절에서는 제한 함자
$$
\textit{F\'Et}_U
\longrightarrow
\textit{F\'Et}_{U_0}
$$
가 본질적으로 전사인 조건을 묻는다. 절 \ref{pione-section-pi1-punctured-spec-III}의 결과를 사용하고
나머지 틈을 순수성으로 메워 이를 보일 것이다. 뇌터 국소환 $(A, \mathfrak m)$에 대하여
제한 함자 $\textit{F\'Et}_X \to \textit{F\'Et}_U$가 본질적으로 전사인 경우를 생각하자.
여기서 $X = \Spec(A)$ 및 $U = X \setminus \{\mathfrak m\}$이다. 이때 이 환에서는
{\it 순수성이 성립한다}고 하였다.

\begin{lemma}
\label{pione-lemma-equivalence-better}
상황 \ref{pione-situation-local-lefschetz}에서 다음을 가정한다.
\begin{enumerate}
\item $A$는 쌍대화 복합체를 가지며 $f$-진 완비이다.
\item 다음 중 하나가 성립한다.
\begin{enumerate}
\item $A_f$는 $(S_2)$이고, $X_0$에 포함되지 않는 $X$의 모든 기약 성분의
차원은 $\geq 4$이다.
\item $\mathfrak p \not \in V(f)$이고
$V(\mathfrak p) \cap V(f) \not = \{\mathfrak m\}$이면
$\text{depth}(A_\mathfrak p) + \dim(A/\mathfrak p) > 3$이다.
\end{enumerate}
\item 모든 극대 아이디얼 $\mathfrak p \subset A_f$에 대하여 $(A_f)_\mathfrak p$에서 순수성이 성립한다.
\end{enumerate}
그러면 제한 함자 $\textit{F\'Et}_U \to \textit{F\'Et}_{U_0}$는 본질적으로 전사이다.
\end{lemma}

\begin{proof}
$V_0 \to U_0$를 유한 에탈 사상이라 하자. 보조정리
\ref{pione-lemma-essentially-surjective-general-better}에 의해 $U_0$를 포함하는 열린 부분집합
$U' \subset U$와 유한 에탈 사상 $V' \to U$가 존재하며, 그 $U_0$로의 밑변환은
$V_0 \to U_0$와 동형이다. $U' \supset U_0$이므로 $U \setminus U'$는 (3)에 나온
소 아이디얼 $\mathfrak p_1, \ldots, \mathfrak p_n$에 대응하는 점들로 이루어진다. 가정에 의해
유한 에탈 사상 $V'_i \to \Spec(A_{\mathfrak p_i})$로서 $V' \to U'$와
$U' \times_U \Spec(A_{\mathfrak p_i})$ 위에서 일치하는 것을 택할 수 있다. 극한의 보조정리
\ref{limits-lemma-glueing-near-closed-point}을 $n$번 적용하면 $V' \to U'$가
유한 에탈 사상 $V \to U$로 연장됨을 알 수 있다.
\end{proof}

\begin{lemma}
\label{pione-lemma-equivalence}
$(A, \mathfrak m)$을 뇌터 국소환이라 하고 $f \in \mathfrak m$이라 하자. 다음을 가정한다.
\begin{enumerate}
\item $A$는 $f$-진 완비이다.
\item $f$는 영인자가 아니다.
\item $H^1_\mathfrak m(A/fA)$ 및 $H^2_\mathfrak m(A/fA)$는 유한 $A$-가군이다.
\item 모든 극대 아이디얼 $\mathfrak p \subset A_f$에 대하여 $(A_f)_\mathfrak p$에서 순수성이 성립한다.
\end{enumerate}
그러면 제한 함자 $\textit{F\'Et}_U \to \textit{F\'Et}_{U_0}$는 본질적으로 전사이다.
\end{lemma}

\begin{proof}
증명은 보조정리 \ref{pione-lemma-equivalence-better}의 증명과 같고, 보조정리
\ref{pione-lemma-essentially-surjective-general-better} 대신 보조정리
\ref{pione-lemma-essentially-surjective-general}을 쓰면 된다.
\end{proof}



\section{국소적인 경우의 순수성 II}
\label{pione-section-local-purity-II}

\noindent
이 절은 절 \ref{pione-section-local-purity}의 계속이다. 뇌터 국소환 $(A, \mathfrak m)$에 대하여
제한 함자 $\textit{F\'Et}_X \to \textit{F\'Et}_U$가 본질적으로 전사인 경우를 생각하자.
여기서 $X = \Spec(A)$이고 $U = X \setminus \{\mathfrak m\}$이다. 이때 이 환에서는
{\it 순수성이 성립한다}고 하였다.

\begin{lemma}
\label{pione-lemma-purity-inherited-by-hypersurface-better}
$(A, \mathfrak m)$을 뇌터 국소환이라 하고, $f \in \mathfrak m$이라 하자. 다음을 가정한다.
\begin{enumerate}
\item $A$는 쌍대화 복합체를 가지며 $f$-진 완비이다.
\item 다음 중 하나가 성립한다.
\begin{enumerate}
\item $A_f$는 $(S_2)$이고, $X$의 모든 기약 성분 중 $X_0$에 포함되지 않는 것의
차원은 $\geq 4$이다.
\item $\mathfrak p \not \in V(f)$이고
$V(\mathfrak p) \cap V(f) \not = \{\mathfrak m\}$이면,
$\text{depth}(A_\mathfrak p) + \dim(A/\mathfrak p) > 3$이다.
\end{enumerate}
\item 모든 극대 아이디얼 $\mathfrak p \subset A_f$에 대하여 $(A_f)_\mathfrak p$에서 순수성이 성립한다.
\item $A$에서 순수성이 성립한다.
\end{enumerate}
그러면 $A/fA$에서 순수성이 성립한다.
\end{lemma}

\begin{proof}
$X = \Spec(A)$라 두고, $U = X \setminus \{\mathfrak m\}$를 천공 스펙트럼이라 하자.
마찬가지로 $X_0 = \Spec(A/fA)$ 및 $U_0 = X_0 \setminus \{\mathfrak m\}$를 얻는다.
$V_0 \to U_0$를 유한 에탈 사상이라 하자. 보조정리 \ref{pione-lemma-equivalence-better}에 의해,
유한 에탈 사상 $V \to U$로서 그 $U_0$로의 밑변환이 $V_0 \to U_0$와 동형인 것을 얻는다.
가정 (5)에 의해 $V \to U$는 유한 에탈 사상 $Y \to X$로 연장된다. 그러면 $Y$의
$X_0$로의 제한은 원하는 $V_0 \to U_0$의 연장이다.
\end{proof}

\begin{lemma}
\label{pione-lemma-purity-inherited-by-hypersurface}
$(A, \mathfrak m)$을 뇌터 국소환이라 하고, $f \in \mathfrak m$이라 하자. 다음을 가정한다.
\begin{enumerate}
\item $A$는 $f$-진 완비이다.
\item $f$는 영인자가 아니다.
\item $H^1_\mathfrak m(A/fA)$ 및 $H^2_\mathfrak m(A/fA)$는 유한 $A$-가군이다.
\item 모든 극대 아이디얼 $\mathfrak p \subset A_f$에 대하여 $(A_f)_\mathfrak p$에서 순수성이 성립한다.
\item $A$에서 순수성이 성립한다.
\end{enumerate}
그러면 $A/fA$에서 순수성이 성립한다.
\end{lemma}

\begin{proof}
증명은 보조정리 \ref{pione-lemma-purity-inherited-by-hypersurface-better}의 증명과 같고, 보조정리
\ref{pione-lemma-equivalence}을 보조정리 \ref{pione-lemma-equivalence-better} 대신 쓰면 된다.
\end{proof}

\noindent
지금까지의 결과를 거듭 적용하면 차원 $\geq 3$인 완전교차에서 순수성이 성립함을 증명할 수 있다.
뇌터 국소환의 완비화가 정칙 국소환을 정칙열이 생성하는 아이디얼로 나눈 몫일 때, 그 환을
완전교차라고 하였다. 나눗셈 거듭제곱 대수의 절 \ref{dpa-section-lci}에서 한 논의를 보라.

\begin{proposition}
\label{pione-proposition-purity-complete-intersection}
$(A, \mathfrak m)$을 뇌터 국소환이라 하자. $A$가 차원 $\geq 3$인 완전교차이면,
$A$에서 순수성이 성립한다. 즉, 천공 스펙트럼의 모든 유한 에탈 덮개는 연장된다.
\end{proposition}

\begin{proof}
보조정리 \ref{pione-lemma-purity-and-completion}에 의해 $A$가 완비 국소환이라고 가정해도 된다. 가정에 의해
$A = B/(f_1, \ldots, f_r)$로 쓸 수 있다. 여기서 $B$는 완비 정칙 국소환이고,
$f_1, \ldots, f_r$은 정칙열이다. $r$에 관한 귀납법으로 증명을 마무리한다. 기초 단계 $r = 0$은
차원 $\geq 2$인 정칙환에 적용되는 보조정리 \ref{pione-lemma-local-purity}에서 따른다.

\medskip\noindent
$A = B/(f_1, \ldots, f_r)$라 하고, 이 명제가 $r - 1$에 대하여 성립한다고 가정하자.
$A' = B/(f_1, \ldots, f_{r - 1})$라 두고, $f_r \in A'$에 보조정리
\ref{pione-lemma-purity-inherited-by-hypersurface}을 적용한다. 다음과 같은 이유로 이를 적용할 수 있다.
조건 (1)은 $f_1, \ldots, f_r$이 정칙열이므로 성립한다. 조건 (2)는 $B$, 따라서
$A'$가 완비이므로 성립한다. 조건 (3)은 $A = A'/f_r A'$가 차원
$\dim(A) \geq 3$인 Cohen--Macaulay 환이므로 성립한다. 쌍대화 복합체의 보조정리
\ref{dualizing-lemma-depth}를 보라. 조건 (4)는 귀납 가정으로부터 성립한다. 실제로
$\dim((A'_{f_r})_\mathfrak p) \geq 3$이다. 여기서 $\mathfrak p$는 $A'_{f_r}$의 극대 소 아이디얼이고, 또한
$(A'_{f_r})_\mathfrak p = B_\mathfrak q/(f_1, \ldots, f_{r - 1})$가 어떤 $\mathfrak q \subset B$에 대하여
성립한다. 조건 (5)도 귀납
가정으로부터 성립한다.
\end{proof}




\section{국소적인 경우의 순수성 III}
\label{pione-section-local-purity-III}

\noindent
이 절에서는 절 \ref{pione-section-local-purity} 및 \ref{pione-section-local-purity-II}의 논의를 계속한다.

\begin{lemma}
\label{pione-lemma-fully-faithful-power-series-over-depth2}
$(A, \mathfrak m)$을 깊이 $\geq 2$인 뇌터 국소환이라 하자.
$B = A[[x_1, \ldots, x_d]]$라 하고, $d \geq 1$이라 하자.
$Y = \Spec(B)$ 및 $Y_0 = V(x_1, \ldots, x_d)$라 두자.
$V \subset Y$를 열린 부분스킴이라 하고, $V_0 = V \cap Y_0$가
$Y_0 \setminus \{\mathfrak m_B\}$와 같다고 하자. 그러면 제한 함자
$$
\textit{F\'Et}_V \longrightarrow \textit{F\'Et}_{V_0}
$$
는 충만충실하다.
\end{lemma}

\begin{proof}
$I = (x_1, \ldots, x_d)$ 및 $X = \Spec(A)$라 두자.
사상 $Y \to X$를 사용하여 $Y_0$를 $X$와 동일시하면, $V_0$는 천공 스펙트럼 $U$, 즉
$A$의 천공 스펙트럼과 동일시된다. 이 아핀 사상으로 가군을 전진시키면
\begin{align*}
\lim_n \Gamma(V_0, \mathcal{O}_V/I^n\mathcal{O}_V)
& =
\lim_n \Gamma(V_0, \mathcal{O}_Y/I^n\mathcal{O}_Y) \\
& =
\lim_n \Gamma(U, \mathcal{O}_U[x_1, \ldots, x_d]/(x_1, \ldots, x_d)^n) \\
& =
\lim_n A[x_1, \ldots, x_d]/(x_1, \ldots, x_d)^n \\
& =
B
\end{align*}
실제로 $A$의 깊이가 $\geq 2$이므로 $\Gamma(U, \mathcal{O}_U) = A$이다. 국소 코호몰로지의 보조정리
\ref{local-cohomology-lemma-finiteness-pushforwards-and-H1-local}을 보라.
따라서 보조정리에 나오는 임의의 열린 부분집합 $V \subset Y$에 대하여
$$
B = \Gamma(Y, \mathcal{O}_Y) \to \Gamma(V, \mathcal{O}_V) \to
\lim_n \Gamma(V_0, \mathcal{O}_Y/I^n\mathcal{O}_Y) = B
$$
를 얻는다. 이로부터 두 화살표가 모두 동형임이 따른다(작은 세부사항은 생략한다).
대수적 및 형식적 기하의 보조정리 \ref{algebraization-lemma-completion-fully-faithful}에 의해,
$\textit{Coh}(\mathcal{O}_V) \to \textit{Coh}(V, I\mathcal{O}_V)$는
유한 국소 자유 대상들로 이루어진 충만 부분범주 위에서 충만충실하다. 따라서 보조정리
\ref{pione-lemma-restriction-fully-faithful}에서 결론이 따른다.
\end{proof}

\begin{lemma}
\label{pione-lemma-purity-power-series-over-depth2}
\begin{slogan}
유한 에탈 덮개에 대한 Ramanujam--Samuel
\end{slogan}
$(A, \mathfrak m)$을 깊이 $\geq 2$인 뇌터 국소환이라 하자.
$B = A[[x_1, \ldots, x_d]]$라 하고, $d \geq 1$이라 하자. 열린 부분집합
$V \subset Y = \Spec(B)$가 다음을 포함한다고 가정한다.
\begin{enumerate}
\item 모든 소 아이디얼 $\mathfrak q \subset B$로서 $\mathfrak q \cap A \not = \mathfrak m$인 것.
\item 소 아이디얼 $\mathfrak m B$.
\end{enumerate}
그러면 함자
$
\textit{F\'Et}_Y
\to
\textit{F\'Et}_V
$
는 동치이다. 특히 $B$에서 순수성이 성립한다.
\end{lemma}

\begin{proof}
소 아이디얼 $\mathfrak q \subset B$가 $V$에 포함되지 않으면 $\mathfrak m$ 위에 놓인다. 이 경우
$A \to B_\mathfrak q$는 평탄 국소 준동형이므로
$\text{depth}(B_\mathfrak q) \geq 2$이다(대수학의 보조정리
\ref{algebra-lemma-apply-grothendieck}). 따라서 보조정리
\ref{pione-lemma-quasi-compact-dense-open-connected-at-infinity-Noetherian}와 국소 코호몰로지의 보조정리
\ref{local-cohomology-lemma-depth-2-connected-punctured-spectrum}에 의해 이 함자는 충만충실하다.

\medskip\noindent
$W \to V$를 유한 에탈 사상이라 하자. $B \to C$를 다음 성질을 갖는 유일한 유한 환 준동형이라 하자.
즉, $\Spec(C) \to Y$는 보조정리 \ref{pione-lemma-extend-S2}에서 구성한, $W \to V$를 연장하는 유한 사상이다.
$C = \Gamma(W, \mathcal{O}_W)$임에 유의하라.

\medskip\noindent
$Y_0 = V(x_1, \ldots, x_d)$ 및 $V_0 = V \cap Y_0$라 두고, $X = \Spec(A)$라 두자.
사상 $Y \to X$를 사용하여 $Y_0$를 $X$와 동일시하면, $V_0$는 천공 스펙트럼 $U$, 즉
$A$의 천공 스펙트럼과 동일시된다. 따라서 $W_0 = W \times_Y Y_0$를 $U$ 위의 유한 에탈 스킴으로 볼 수 있다. 이때
$$
W_0 \times_U (U \times_X Y)
\quad\text{및}\quad
W \times_V (U \times_X Y)
$$
는 $U \times_X Y$ 위의 유한 에탈 스킴들이고, $V_0$ 위로 제한하면 서로 동형이다.
열린 부분집합 $U \times_X Y$에 보조정리 \ref{pione-lemma-fully-faithful-power-series-over-depth2}을 적용하여 동형
$$
W_0 \times_U (U \times_X Y) \longrightarrow W \times_V (U \times_X Y)
$$
을 $U \times_X Y$ 위에서 얻는다.

\medskip\noindent
$C_0 = \Gamma(W_0, \mathcal{O}_{W_0})$는 유한 $A$-대수임에 유의하라. 이는
$W_0 \to U \subset X$에 보조정리 \ref{pione-lemma-extend-S2}를 적용하면 따른다(위에서 $B \to C$에 대하여
한 것과 완전히 같다). 보조정리 \ref{pione-lemma-extend-S2}의 구성은 평탄 밑변환 및
열린집합의 변경과 가환하므로, 위의 동형은 동형
$$
\Psi : C \longrightarrow C_0 \otimes_A B
$$
을 유도하며, 이는 유한 $B$-대수의 동형이다. 한편 $\Spec(C) \to Y$는 $Y$의 적어도 한 점으로서
$\mathfrak m \in X$ 위에 놓이는 점의 위에 있는 모든 점에서 에탈임을 알고 있다.
$\Psi$가 동형이므로 $\Spec(C_0) \to X$는
$\mathfrak m$ 위에서 에탈이다(작은 세부사항은 생략한다). 이는 물론 $A \to C_0$가 유한 에탈이고,
따라서 $B \to C$도 유한 에탈임을 뜻한다.
\end{proof}

\begin{lemma}
\label{pione-lemma-purity-smooth-over-depth2}
$f : X \to S$를 스킴의 사상이라 하고, $U \subset X$를 열린 부분스킴이라 하자. 다음을 가정한다.
\begin{enumerate}
\item $f$는 매끄럽다.
\item $S$는 뇌터이다.
\item $s \in S$가 $\text{depth}(\mathcal{O}_{S, s}) \leq 1$을 만족하면 $X_s = U_s$이다.
\item $U_s \subset X_s$는 모든 $s \in S$에 대하여 조밀하다.
\end{enumerate}
그러면 $\textit{F\'Et}_X \to \textit{F\'Et}_U$는 동치이다.
\end{lemma}

\begin{proof}
보조정리 \ref{pione-lemma-quasi-compact-dense-open-connected-at-infinity-Noetherian}와 국소 코호몰로지의 보조정리
\ref{local-cohomology-lemma-depth-2-connected-punctured-spectrum}에 의해 이 함자는 충만충실하다
(깊이 조건을 확인하려면 대수학의 보조정리 \ref{algebra-lemma-apply-grothendieck}도 적용한다).

\medskip\noindent
$\pi : V \to U$를 유한 에탈 사상이라 하고, $Y \to X$를 보조정리 \ref{pione-lemma-extend-S2}에서 구성한 유한 사상이라 하자.
$Y \to X$가 유한 에탈임을 보여야 한다. 모든 점 $x \in X$에 대하여 그 점이 주어진 점 $s \in S$로 갈 때
이를 보일 때에는 뇌터 스킴의 평탄 사상 $S' \to S$로서 $s$를 상에 포함하는 것에 의해 밑변환해도 된다.
이는 보조정리 \ref{pione-lemma-extend-S2}의 구성이 평탄 밑변환과 가환하기 때문이다.

\medskip\noindent
$U$를 확대하여 $U \subset X$가 $Y \to X$가 유한 에탈인 최대 열린집합이라고 가정해도 된다.
$Z \subset X$를 $U$의 여집합이라 하자. 모순을 얻기 위해 $Z \not = \emptyset$이라고 가정한다.
$s \in S$를 $Z \to S$의 상에 속하는 점으로서 $s$의 어떤 진정한 일반화점도 그 상에 속하지 않는 것으로 택한다.
그러면 $\Spec(\mathcal{O}_{S, s})$로 밑변환하여 $S = \Spec(A)$이고,
$(A, \mathfrak m, \kappa)$가 깊이 $\geq 2$인 뇌터 국소환이라고 가정해도 된다. 또한 $Z$는 닫힌 올
$X_s$에 포함되며 $X_s$ 안에서 어디에서도 조밀하지 않다. 닫힌점 $z \in Z$를 택한다. 그러면
$\kappa(z)/\kappa$는 유한이다(Hilbert 영점정리에 의한다. 대수학의 정리
\ref{algebra-theorem-nullstellensatz}를 보라). 국소 스킴의 유한 평탄 사상
$(S', s') \to (S, s)$로서 잉여체 확대 $\kappa(z)/\kappa$를 실현하는 것을 택한다. 대수학의 보조정리
\ref{algebra-lemma-finite-free-given-residue-field-extension}을 보라. $S' \to S$로 밑변환하면
$\kappa(z) = \kappa$인 경우로 환원된다.

\medskip\noindent
스킴의 사상 심화의 보조정리 \ref{more-morphisms-lemma-slice-smooth}에 의해, 국소 닫힌 부분스킴
$S' \subset X$로서 $z$를 지나고 $S' \to S$가 $z$에서 에탈인 것이 존재한다. $S' \to S$로
밑변환하여 $\sigma : S \to X$이고 $\sigma(s) = z$인 단면이 존재한다고 가정해도 된다.
$\Spec(B) \subset X$를 $s$의 아핀 근방으로 택한다. 그러면 $A \to B$는 매끄러운 환 준동형이고 단면
$\sigma : B \to A$를 가진다. $I = \Ker(\sigma)$라 두고, $B^\wedge$를 $I$-진 완비화, 즉 환 $B$의
완비화라 하자. 그러면 $B^\wedge \cong A[[x_1, \ldots, x_d]]$가 어떤 $d \geq 0$에 대하여 성립한다. 대수학의 보조정리
\ref{algebra-lemma-section-smooth}를 보라. $d > 0$임에 유의하라. 실제로 그렇지 않다면
$X \to S$는 $z$에서 에탈이고, $z$는 $X_s$의 일반점이므로 가정 (4)에 의해
$z \in U$가 된다. 마찬가지로 $d > 0$이면 $\mathfrak m B^\wedge$는 $U$ 안으로, 사상
$\Spec(B^\wedge) \to X$에 의해 간다. $Y \to X$가 $\Spec(B^\wedge)$로 밑변환한 뒤
유한 에탈임을 보이면 충분하다. $B \to B^\wedge$가 평탄이므로
(대수학의 보조정리 \ref{algebra-lemma-completion-flat}), 이는 보조정리
\ref{pione-lemma-purity-power-series-over-depth2}와 $Y \to X$의 구성의 유일성에서 따른다.
\end{proof}

\begin{proposition}
\label{pione-proposition-purity-smooth-over-depth2}
$A \to B$를 뇌터 국소환의 국소 준동형이라 하자. $A$의 깊이가 $\geq 2$이고,
$A \to B$가 $\mathfrak m_B$-진 위상에 관하여 형식적으로 매끄러우며,
$\dim(B) > \dim(A)$라고 가정하자. 열린 부분집합 $V \subset Y = \Spec(B)$가 다음을 포함한다고 하자.
\begin{enumerate}
\item 모든 소 아이디얼 $\mathfrak q \subset B$로서 $\mathfrak q \cap A \not = \mathfrak m_A$인 것.
\item 소 아이디얼 $\mathfrak m_A B$.
\end{enumerate}
그러면 함자 $\textit{F\'Et}_Y \to \textit{F\'Et}_V$는 동치이다. 특히 $B$에서 순수성이 성립한다.
\end{proposition}

\begin{proof}
소 아이디얼 $\mathfrak q \subset B$가 $V$에 포함되지 않으면 $\mathfrak m_A$ 위에 놓인다. 이 경우
$A \to B_\mathfrak q$는 평탄 국소 준동형이므로 $\text{depth}(B_\mathfrak q) \geq 2$이다
(대수학의 보조정리 \ref{algebra-lemma-apply-grothendieck}). 따라서 보조정리
\ref{pione-lemma-quasi-compact-dense-open-connected-at-infinity-Noetherian}와 국소 코호몰로지의 보조정리
\ref{local-cohomology-lemma-depth-2-connected-punctured-spectrum}에 의해 이 함자는 충만충실하다.

\medskip\noindent
$A^\wedge$ 및 $B^\wedge$를 각각 $A$ 및 $B$의 극대 아이디얼에 관한 완비화라 하자.
이 명제의 가정들은 $A^\wedge \to B^\wedge$에 대하여 성립한다. 대수학 심화의 보조정리
\ref{more-algebra-lemma-completion-dimension},
\ref{more-algebra-lemma-completion-depth} 및
\ref{more-algebra-lemma-formally-smooth-completion}을 보라. 보조정리 \ref{pione-lemma-extend-S2}의 구성의 유일성과
평탄 밑변환과의 가환성에 의해, $A^\wedge \to B^\wedge$ 및 $V$의 역상에 대하여 본질적 전사성을
증명하면 충분하다(세부사항은 생략한다. $V$가 천공 스펙트럼인 경우는 보조정리
\ref{pione-lemma-purity-and-completion}과 비교하라). 대수학 심화의 명제
\ref{more-algebra-proposition-fs-regular}에 의해, 이는 $A \to B$가 정칙이라고 가정해도 됨을 뜻한다.

\medskip\noindent
$W \to V$를 유한 에탈 사상이라 하자. Popescu 정리(환 사상의 평활화의 정리
\ref{smoothing-theorem-popescu})에 의해 $B = \colim B_i$를 매끄러운 $A$-대수들의 여과 여극한으로
쓸 수 있다. 어떤 $i$와 열린 부분집합 $V_i \subset \Spec(B_i)$로서 그 역상이 $V$인 것을 택할 수 있다
(극한의 보조정리 \ref{limits-lemma-descend-opens}). $i$를 크게 한 뒤에는 유한 에탈 사상
$W_i \to V_i$로서 그 $V$로의 밑변환이 $W \to V$인 것이 존재한다고 가정해도 된다.
극한의 보조정리 \ref{limits-lemma-descend-finite-presentation},
\ref{limits-lemma-descend-finite-finite-presentation} 및
\ref{limits-lemma-descend-etale}를 보라. $V_i$의 여집합이 $\Spec(B_i) \to \Spec(A)$의 닫힌 올에
포함된다고 가정해도 된다. 이는 $V$에 대하여 성립하기 때문이다($V_i$를 이와 같이 택하거나, 위의 보조정리를
사용하여 충분히 큰 $i$에 대하여 성립함을 보이면 된다). $\eta$를 $\Spec(B) \to \Spec(A)$의
닫힌 올의 일반점이라 하자. $\eta \in V$이므로 $\eta$의 상은 $V_i$에 속한다. 따라서 $V_i$를
$\Spec(B)$의 닫힌점의 상의 아핀 열린 근방으로 바꾼 뒤에는, $\Spec(B_i) \to \Spec(A)$의 닫힌 올이 기약이고
그 일반점이 $V_i$에 포함된다고 가정해도 된다(세부사항은 생략한다. 체 위의 매끄러운 스킴이 기약 스킴들의
서로소 합임을 사용한다). 이제 보조정리 \ref{pione-lemma-purity-smooth-over-depth2}를 적용하면
$W_i \to V_i$는 유한 에탈 사상 $\Spec(C_i) \to \Spec(B_i)$로 연장된다. 이를 $\Spec(B)$로
당겨오면 $W$는 원하는 대로 함자 $\textit{F\'Et}_Y \to \textit{F\'Et}_V$의 본질적 상에 속한다.
\end{proof}





\section{기본군에 대한 Lefschetz}
\label{pione-section-global-lefschetz}

\noindent
물론 국소적인 경우에는 이러한 종류의 결과를 이미 여러 개 증명하였다. 이 절에서는 사영적인 경우를 논한다.

\begin{proposition}
\label{pione-proposition-lefschetz-fully-faithful}
$k$를 체라 하고, $X$를 $k$ 위의 고유 스킴이라 하자. $\mathcal{L}$을 풍부한 가역
$\mathcal{O}_X$-가군이라 하자. $s \in \Gamma(X, \mathcal{L})$라 하고, $Y = Z(s)$를 $s$의 영점 스킴이라 하자.
모든 $x \in X \setminus Y$에 대하여
$$
\text{depth}(\mathcal{O}_{X, x}) + \dim(\overline{\{x\}}) > 1
$$
이라고 가정하자. 그러면 제한 함자 $\textit{F\'Et}_X \to \textit{F\'Et}_Y$는 충만충실하다.
실제로 $V \subset X$를 $Y$를 포함하는 임의의 열린 부분스킴이라 하면, 제한 함자
$\textit{F\'Et}_V \to \textit{F\'Et}_Y$는 충만충실하다.
\end{proposition}

\begin{proof}
첫째 주장은 보조정리 \ref{pione-lemma-restriction-fully-faithful-special}과
대수적 및 형식적 기하의 명제 \ref{algebraization-proposition-lefschetz}의 형식적 귀결이다.
둘째 주장은 보조정리 \ref{pione-lemma-restriction-fully-faithful-special}과
대수적 및 형식적 기하의 보조정리 \ref{algebraization-lemma-lefschetz-addendum}에서 따른다.
\end{proof}

\begin{proposition}
\label{pione-proposition-lefschetz-equivalence-general}
$k$를 체라 하고, $X$를 $k$ 위의 고유 스킴이라 하자. $\mathcal{L}$을 풍부한 가역
$\mathcal{O}_X$-가군이라 하자. $s \in \Gamma(X, \mathcal{L})$라 하고, $Y = Z(s)$를 $s$의 영점 스킴이라 하자.
$\mathcal{V}$를 $X$의 열린 부분스킴으로서 $Y$를 포함하는 것 전체의 집합에 포함관계의 역순으로 순서를 준 것이라 하자.
모든 $x \in X \setminus Y$에 대하여
$$
\text{depth}(\mathcal{O}_{X, x}) + \dim(\overline{\{x\}}) > 2
$$
이라고 가정하자. 그러면 제한 함자
$$
\colim_\mathcal{V} \textit{F\'Et}_V \to \textit{F\'Et}_Y
$$
는 동치이다.
\end{proposition}

\begin{proof}
이는 보조정리 \ref{pione-lemma-restriction-equivalence-general}과
대수적 및 형식적 기하의 명제 \ref{algebraization-proposition-lefschetz-equivalence}의 형식적 귀결이다.
\end{proof}

\begin{proposition}
\label{pione-proposition-lefschetz-equivalence}
$k$를 체라 하고, $X$를 $k$ 위의 고유 스킴이라 하자. $\mathcal{L}$을 풍부한 가역
$\mathcal{O}_X$-가군이라 하자. $s \in \Gamma(X, \mathcal{L})$라 하고, $Y = Z(s)$를 $s$의 영점 스킴이라 하자.
모든 $x \in X \setminus Y$에 대하여
$$
\text{depth}(\mathcal{O}_{X, x}) + \dim(\overline{\{x\}}) > 2
$$
이고, 또한 닫힌점 $x \in X \setminus Y$에 대하여 $\mathcal{O}_{X, x}$에서 순수성이 성립한다고 가정하자.
그러면 제한 함자 $\textit{F\'Et}_X \to \textit{F\'Et}_Y$는 동치이다. $X$, 또는 이와 동치로 $Y$가 연결이면
$$
\pi_1(Y, \overline{y}) \to \pi_1(X, \overline{y})
$$
는 임의의 기하점 $\overline{y}$가 $Y$ 위에 놓일 때 동형이다.
\end{proposition}

\begin{proof}
명제 \ref{pione-proposition-lefschetz-fully-faithful}에 의해 충만충실성이 성립한다. 명제
\ref{pione-proposition-lefschetz-equivalence-general}에 의해 $\textit{F\'Et}_Y$의 모든 대상은
섬유곱 $U \times_V Y$와 동형이다. 여기서 어떤 유한 에탈 사상 $U \to V$가 있고,
$V \subset X$는 $Y$를 포함하는 열린 부분스킴이다. 여집합 $T = X \setminus V$는\footnote{실제로 $T$는 $k$ 위에서
고유하고($X$에서 닫혀 있으므로), 아핀이며(아핀 스킴 $X \setminus Y$에서 닫혀 있기 때문이다.
스킴의 사상의 보조정리 \ref{morphisms-lemma-proper-ample-delete-affine}를 보라). 따라서 $k$ 위에서 유한이다
(스킴의 사상의 보조정리 \ref{morphisms-lemma-finite-proper}). 그러므로 $T$는 닫힌점들의 유한집합이다.}
$X \setminus Y$의 닫힌점들의 유한집합이다. $T = \{x_1, \ldots, x_n\}$이라 하자. 가정에 의해 유한 에탈 사상
$V'_i \to \Spec(\mathcal{O}_{X, x_i})$로서 $U \to V$와
$V \times_X \Spec(\mathcal{O}_{X, x_i})$ 위에서 일치하는 것을 택할 수 있다. 극한의 보조정리
\ref{limits-lemma-glueing-near-closed-point}을 $n$번 적용하면 $U \to V$가 원하는 유한 에탈 사상
$U' \to X$로 연장됨을 알 수 있다. 마지막 주장에 대해서는 보조정리 \ref{pione-lemma-what-equivalence-gives}를 보라.
\end{proof}










\section{분기 자취의 순수성}
\label{pione-section-purity-ramification}

\noindent
이 절에서는 일반적으로 유한한 사상에 대한 분기 자취의 순수성의 유사 결과를 논한다. 이 결과는 Gabber의 것으로 보인다.
특별한 경우는 van der Waerden의 순수성 정리로, 정규 다양체에서 매끄러운 다양체로 가는 쌍유리 사상이
동형이 아닌 자취에 관한 것이다.

\begin{lemma}
\label{pione-lemma-characterize-rational-singularity}
$A$를 차원 $2$인 뇌터 정규 국소 정역이라 하자. $A$가 Nagata이고, 쌍대화 가군 $\omega_A$를 가지며,
특이점 해소 $f : X \to \Spec(A)$를 가진다고 가정하자. $\omega_X$를 곡면의 해소의 주의
\ref{resolve-remark-dualizing-setup}에 나오는 것으로 하자. 어떤 유효 Cartier 인자 E에 대하여
$\omega_X \cong \mathcal{O}_X(E)$이고, 이 $E \subset X$의 받침이 예외 올에 포함되면,
$A$는 유리 특이점을 정의한다.
$f$가 극소 해소이면 $E = 0$이다.
\end{lemma}

\begin{proof}
대각합 사상 $Rf_*\omega_X \to \omega_A$가 존재한다. 스킴의 쌍대성의 절
\ref{duality-section-trace}을 보라. Grauert--Riemenschneider에 의해
(곡면의 해소의 명제 \ref{resolve-proposition-Grauert-Riemenschneider}),
$R^1f_*\omega_X = 0$이다. 따라서 대각합 사상은 $f_*\omega_X \to \omega_A$인 사상이다. 이제
$$
\mathcal{O}_{\Spec(A)} = f_*\mathcal{O}_X \to f_*\omega_X \to \omega_A
$$
를 생각할 수 있다. 첫째 사상은 보조정리의 명제에서 존재한다고 가정한 사상
$\mathcal{O}_X \to \mathcal{O}_X(E) = \omega_X$에서 온다. 합성은 인자의 보조정리
\ref{divisors-lemma-check-isomorphism-via-depth-and-ass}에 의해 동형이다. 실제로 이는 $A$의 천공 스펙트럼 위에서
동형이고(보조정리의 가정 및 $f$가 천공 스펙트럼 위에서 동형이라는 사실에 의해), $A$와 $\omega_A$는
$A$-가군으로서 깊이 $2$를 가진다(대수학의 보조정리 \ref{algebra-lemma-criterion-normal} 및
쌍대화 복합체의 보조정리 \ref{dualizing-lemma-depth-dualizing-module}에 의해). 따라서
$f_*\omega_X \to \omega_A$는 전사이고, 그러므로 동형이다. 이에 따라 $Rf_*\omega_X = \omega_A$이고,
쌍대성에 의해 $Rf_*\mathcal{O}_X = \mathcal{O}_{\Spec(A)}$이다. 따라서
$H^1(X, \mathcal{O}_X) = 0$이고, 이는 $A$가 유리 특이점을 정의함을 뜻한다
(곡면의 해소의 절 \ref{resolve-section-bounded}에서 한 논의, 특히 보조정리
\ref{resolve-lemma-regular-rational} 및 \ref{resolve-lemma-exact-sequence}를 보라). $f$가 극소이면
$E = 0$이다. 이는 사상 $f^*\omega_A \to \omega_X$가 곡면의 해소의 보조정리
\ref{resolve-lemma-dualizing-blow-up-rational}을 반복하여 적용함으로써 전사이고, 또한 위에서 보았듯이
$\omega_A \cong A$이기 때문이다.
\end{proof}

\begin{lemma}
\label{pione-lemma-key-purity-ramification}
$f : X \to \Spec(A)$를 유한형 사상이라 하고, 점 $x \in X$를 택하자. 다음을 가정한다.
\begin{enumerate}
\item $A$는 우수 정칙 국소환이다.
\item $\mathcal{O}_{X, x}$는 차원 $2$인 정규환이다.
\item $f$는 $\overline{\{x\}}$ 밖에서 에탈이다.
\end{enumerate}
그러면 $f$는 $x$에서 에탈이다.
\end{lemma}

\begin{proof}
먼저 $X$를 $x$의 아핀 열린 근방으로 바꾼다. $\mathcal{O}_{X, x}$는 우수 국소환이다
(대수학 심화의 보조정리 \ref{more-algebra-lemma-finite-type-over-excellent}).
따라서 특이점의 극소 해소 $W \to \Spec(\mathcal{O}_{X, x})$를 택할 수 있다.
곡면의 해소의 정리 \ref{resolve-theorem-resolve}를 보라. 필요하면 $X$를 $x$의
아핀 열린 근방으로 바꾸어, 고유 사상 $b : X' \to X$로서
$X' \times_X \Spec(\mathcal{O}_{X, x}) = W$를 만족하는 것을 찾을 수 있다.
극한의 보조정리 \ref{limits-lemma-glueing-near-closed-point}을 보라.
또한 $X$를 축소하여 $X'$가 정칙이라고 가정해도 된다. 실제로 $W$는 정칙이고 $X'$는 우수하며,
우수환의 스펙트럼의 정칙 자취는 열려 있다. 또한 $W \to \Spec(\mathcal{O}_{X, x})$는 사영적이므로
(정규화된 부풀리기들의 열이기 때문이다), $X$를 축소한 뒤 $b$가 사영적이라고 가정해도 된다
(세부사항은 생략한다). $U = X \setminus \overline{\{x\}}$라 두자.
$W \to \Spec(\mathcal{O}_{X, x})$가 천공 스펙트럼 위에서 동형이므로,
$b : X' \to X$가 $U$ 위에서 동형이라고 가정해도 된다. 따라서 $U$를 $X'$의 열린 부분스킴으로도
간주한다. $f' = f \circ b : X' \to \Spec(A)$라 두자.

\medskip\noindent
$A$가 정칙이므로 $\mathcal{O}_Y$는 $Y$의 쌍대화 복합체이다. 따라서
$f^!\mathcal{O}_Y$는 $X$ 위의 쌍대화 복합체이다
(스킴의 쌍대성의 보조정리 \ref{duality-lemma-shriek-dualizing}).
스킴의 쌍대성의 보조정리 \ref{duality-lemma-dualizing-module-CM-scheme}에 의해
$X$의 Cohen--Macaulay 자취는 열려 있다(이는 우수성을 사용해서도 증명할 수 있다).
$\mathcal{O}_{X, x}$가 Cohen--Macaulay이므로 $X$를 축소한 뒤에는 $X$가 Cohen--Macaulay이라고 가정해도 된다.
에탈 사상은 국소 완전교차임에 유의하라. 따라서 스킴의 쌍대성의 보조정리
\ref{duality-lemma-fundamental-class-almost-lci}를 $r = 0$으로 적용하여 사상
$$
\mathcal{O}_X \longrightarrow \omega_{X/Y} = H^0(f^!\mathcal{O}_Y)
$$
을 얻는다. 이는 $X \setminus \overline{\{x\}}$ 위에서 동형이다.
스킴의 쌍대성의 보조정리 \ref{duality-lemma-shriek-over-CM}에 의해 $\omega_{X/Y}$는 $(S_2)$이므로,
인자의 보조정리 \ref{divisors-lemma-check-isomorphism-via-depth-and-ass}에 의해 이 사상은 동형이다.
이로써 이미 $X$, 특히 $\mathcal{O}_{X, x}$가 Gorenstein임을 알 수 있다.

\medskip\noindent
$\omega_{X'/Y} = H^0((f')^!\mathcal{O}_Y)$라 두자. 위와 완전히 같은 방식으로 논하면,
$(f')^!\mathcal{O}_Y = \omega_{X'/Y}[0]$는 $X'$의 쌍대화 복합체이다. $X'$가 정칙이므로
사상 $X' \to Y$는 국소 완전교차 사상이다. 스킴의 사상 심화의 보조정리
\ref{more-morphisms-lemma-morphism-regular-schemes-is-lci}를 보라.
스킴의 쌍대성의 보조정리 \ref{duality-lemma-fundamental-class-lci}에 의해 사상
$$
\mathcal{O}_{X'} \longrightarrow \omega_{X'/Y}
$$
이 존재하고, 이는 $U$ 위에서 동형이다. 따라서 어떤 유효 Cartier 인자에 대하여
$\omega_{X'/Y} = \mathcal{O}_{X'}(E)$이고, 이 $E \subset X'$는 $U$와 서로 만나지 않는다.

\medskip\noindent
$\omega_{X/Y} = \mathcal{O}_Y$이므로
$\omega_{X'/Y} = b^! f^!\mathcal{O}_Y = b^!\mathcal{O}_X$이다.
$W \to \Spec(\mathcal{O}_{X, x})$로 돌아가면 $\omega_W = \mathcal{O}_W(E|_W)$임을 알 수 있다.
보조정리 \ref{pione-lemma-characterize-rational-singularity}에 의해 $E|_W = 0$이다. 따라서 이미 사용한
스킴의 쌍대성의 보조정리 \ref{duality-lemma-fundamental-class-lci}에 의해
$f' : X' \to Y$는 에탈이다. $f'$가 에탈이면 $b$는 동형이어야 하므로, 이로써 곧 증명이 끝난다.
\end{proof}

\begin{lemma}[분기 자취의 순수성]
\label{pione-lemma-purity-ramification}
\begin{reference}
복소공간에 대한 이 결과는 \cite{Fischer}의 170쪽에서 찾을 수 있다. 일반적인 경우에는
Gabber의 것으로 알려진 \cite[Theorem 2.4]{Zong}이다.
\end{reference}
$f : X \to Y$를 국소 뇌터 스킴의 사상이라 하자. $x \in X$라 하고, $y = f(x)$라 두자. 다음을 가정한다.
\begin{enumerate}
\item $\mathcal{O}_{X, x}$는 차원 $\geq 1$인 정규환이다.
\item $\mathcal{O}_{Y, y}$는 정칙이다.
\item $f$는 국소 유한형이다.
\item 특수화 $x' \leadsto x$가 $\dim(\mathcal{O}_{X, x'}) = 1$을 만족하면,
$f$는 $x'$에서 에탈이다.
\end{enumerate}
그러면 $f$는 $x$에서 에탈이다.
\end{lemma}

\begin{proof}
이 보조정리를 $d = \dim(\mathcal{O}_{X, x})$에 관한 귀납법으로 증명한다.

\medskip\noindent
$d = 1$인 경우에는 가정에 의해 사상 $f$가 $x$에서 에탈이므로 자명하다. $d \geq 2$라고 가정하자.

\medskip\noindent
보조정리의 결론에 영향을 주지 않고 $\Spec(\mathcal{O}_{Y, y}) \to Y$로 밑변환할 수 있다.
스킴의 사상의 보조정리 \ref{morphisms-lemma-set-points-where-fibres-etale}를 보라. 따라서
$Y = \Spec(A)$라고 가정해도 된다. 여기서 $A$는 정칙 국소환이고, $y$는 극대 아이디얼
$\mathfrak m$에 대응하며, 이는 $A$의 극대 아이디얼이다.

\medskip\noindent
특수화 $x' \leadsto x$로서 $x' \not = x$인 것을 택하자. 그러면 $\mathcal{O}_{X, x'}$는
$\mathcal{O}_{X, x}$의 국소화이므로 정규이다. $x'$가 $X$의 일반점이 아니면
$1 \leq \dim(\mathcal{O}_{X, x'}) < d$이고, 귀납 가정에 의해 $f$는 $x'$에서 에탈이다.
따라서 $f$가 $x$로 특수화되는 모든 점에서 에탈이라고 가정해도 된다. $f$가 에탈인
점들의 집합은 정의에 의해 $X$에서 열려 있으므로, $X$를 $x$의 열린 근방으로 바꾼 뒤에는
$f$가 $\overline{\{x\}}$ 밖에서 에탈이라고 가정해도 된다. 특히 $f$는 닫힌점
$y \in Y = \Spec(A)$ 위에 놓이는 점들을 제외하면 에탈이다.

\medskip\noindent
$X' = X \times_{\Spec(A)} \Spec(A^\wedge)$라 두고, $x' \in X'$를 $x$ 위에 놓이는 유일한 점이라 하자.
위의 논의에 의해 $X'$는 닫힌 올 밖에서 $\Spec(A^\wedge)$ 위에 에탈이고, 따라서 $X'$는
닫힌 올 밖에서 정규이다. $X$가 정규이므로 곡면의 해소의 보조정리
\ref{resolve-lemma-normalization-completion}에 의해 $X'$는 정규이다. 이제
$X' \to \Spec(A^\wedge)$가 $x'$에서 에탈임을 보일 수 있다면, 앞서 말한 스킴의 사상의 보조정리
\ref{morphisms-lemma-set-points-where-fibres-etale}에 의해 $f$는 $x$에서 에탈이다.
따라서 $A$가 정칙 완비 국소환이라고 가정해도 되며, 실제로 그렇게 가정한다.

\medskip\noindent
$d = 2$인 경우는 보조정리 \ref{pione-lemma-key-purity-ramification}에서 따른다.

\medskip\noindent
$d > 2$라고 가정하자. $t \in \mathfrak m$, $t \not \in \mathfrak m^2$를 잡는다.
$Y_0 = \Spec(A/tA)$ 및 $X_0 = X \times_Y Y_0$로 놓자. 그러면 $X_0 \to Y_0$는
닫힌점 위의 올을 제외한 곳에서 에탈이다. $d > 2$이므로
$\dim(\mathcal{O}_{X_0, x}) = d - 1$은 $\geq 2$이다. 정규화 $X_0' \to X_0$는 전사이고 유한이다
(완비 국소환 위에서 작업하고 있고, 그러한 환은 Nagata이기 때문이다). $x$로 가는 점
$x' \in X_0'$를 잡는다. 귀납 가정에 의해 사상 $X'_0 \to Y$는 $x'$에서 에탈이다. 포함관계
$\kappa(y) \subset \kappa(x) \subset \kappa(x')$로부터 $\kappa(x)$가 $\kappa(y)$의 유한 확대임을 안다.
따라서 $x$는 $X \to Y$의 $y$ 위 올의 닫힌점이다. 그러나 $x$는 이 올의
일반점이기도 하므로, $f$가 $x$에서 준유한임을 알 수 있고, 분기 자취의 순수성의 경우로 환원된다.
보조정리 \ref{pione-lemma-purity}를 보라.
\end{proof}




\section{분기 자취의 여집합의 아핀성}
\label{pione-section-stronger-purity}

\noindent
$f : X \to Y$를 Noether 스킴들의 유한형 사상이라 하고, $X$는 정규이며 $Y$는 정칙이라고 하자.
$V \subset X$를 $f$가 에탈인 최대 열린부분스킴이라 하자.
\cite[Chapter IV, Section 21.12]{EGA}의 논의는 $V \to X$가 아핀 사상일 가능성을 시사한다.
$V \to X$가 아핀이면, 인자의 보조정리
\ref{divisors-lemma-complement-affine-open-immersion}을 사용하여 분기 자취의 순수성
(보조정리 \ref{pione-lemma-purity-ramification})을 얻는다는 점에 주의하자.
따라서 $V \to X$의 아핀성은 분기 자취에 대한 순수성의 “강한” 형태이다.
이 절에서는 $X$와 $Y$가 등표수이고 우수할 때 $V \to X$가 아핀임을 증명한다.
정리 \ref{pione-theorem-global}을 보라. $X$와 $Y$가 혼합표수인 경우에도(여전히 우수하다면)
이 결과가 참일 것이라고 추측하는 것이 타당해 보인다.

\begin{lemma}
\label{pione-lemma-structure-cohomology}
$(A, \mathfrak m)$를 체를 포함하는 정칙 국소환이라 하자.
$f : V \to \Spec(A)$를 에탈이고 준콤팩트인 사상이라 하자.
$\mathfrak m \not \in f(V)$이고 $g : V \to \Spec(A) \setminus \{\mathfrak m\}$가
아핀이라고 가정하자. 그러면 $i > 0$에 대하여 $H^i(V, \mathcal{O}_V)$는
$A$의 잉여체의 단사 포락의 복사본들의 직합과 동형이다.
\end{lemma}

\begin{proof}
천공 스펙트럼을 $U = \Spec(A) \setminus \{\mathfrak m\}$라 쓰자. 그러면 $g : V \to U$는
아핀이다. 스킴의 코호몰로지의 보조정리
\ref{coherent-lemma-relative-affine-cohomology}에 의해
$H^i(V, \mathcal{O}_V) = H^i(U, g_*\mathcal{O}_V)$이다.
스킴의 보조정리 \ref{schemes-lemma-push-forward-quasi-coherent}에 의해
$\mathcal{O}_U$-가군 $g_*\mathcal{O}_V$는 준연접이다. 임의의 준연접
$\mathcal{O}_U$-가군 $\mathcal{F}$에 대하여, $i > 0$이면 코호몰로지
$H^i(U, \mathcal{F})$는 $\mathfrak m$-거듭제곱 꼬임이다. 예를 들어 국소 코호몰로지의 보조정리
\ref{local-cohomology-lemma-local-cohomology}를 보라. 특히 $i > 0$이면 $A$-가군
$H^i(V, \mathcal{O}_V)$는 $\mathfrak m$-거듭제곱 꼬임이다.
임의의 평탄 환 준동형 $A \to A'$에 대하여, 평탄 기저변환에 관한 스킴의 코호몰로지의 보조정리
\ref{coherent-lemma-flat-base-change-cohomology}에 의해
$H^i(V, \mathcal{O}_V) \otimes_A A' = H^i(V', \mathcal{O}_{V'})$이다. 여기서
$V' = V \times_{\Spec(A)} \Spec(A')$이다. $A'$을 $A$의 완비화로 잡으면
(대수학 심화의 절 \ref{more-algebra-section-permanence-completion}에 의해 평탄하다) 다음을 얻는다.
$$
H^i(V, \mathcal{O}_V) = H^i(V, \mathcal{O}_V) \otimes_A A' =
H^i(V', \mathcal{O}_{V'}),\quad\text{여기서 } i > 0
$$
첫 번째 등호는 방금 증명한 꼬임 성질과 대수학 심화의 보조정리
\ref{more-algebra-lemma-neighbourhood-equivalence}에서 따른다. 또한 잉여체 $k$의 단사 포락은
$A$와 $A'$에 대해 같다. 쌍대화 복합체의 보조정리
\ref{dualizing-lemma-compare}를 보라. 이와 같이
$A = k[[x_1, \ldots, x_d]]$인 경우로 환원된다. 대수학의 절
\ref{algebra-section-cohen-structure-theorem}을 보라.

\medskip\noindent
$k$의 표수가 $p > 0$이라고 가정하자. $F : A \to A$, $a \mapsto a^p$는 평탄하고
(국소 코호몰로지의 보조정리 \ref{local-cohomology-lemma-frobenius-flat-regular}),
또 에탈 사상의 보조정리 \ref{etale-lemma-relative-frobenius-etale}에 의해
$V \times_{\Spec(A), \Spec(F)} \Spec(A) \cong V$가 $\Spec(A)$ 위의 스킴으로서 동형이므로,
위의 논의에서
$H^i(V, \mathcal{O}_V) \otimes_{A, F} A \cong H^i(V, \mathcal{O}_V)$를 얻는다. 따라서
국소 코호몰로지의 보조정리
\ref{local-cohomology-lemma-structure-torsion-Frobenius-regular}에 의해 결과가 따른다.

\medskip\noindent
$k$의 표수가 $0$이라고 가정하자. 국소 코호몰로지의 보조정리
\ref{local-cohomology-lemma-etale-derivation}에 의해 $H^i(V, \mathcal{O}_V)$ 위에 가법 작용소
$D_j$, $j = 1, \ldots, d$가 존재하고,
$\partial_j = \partial/\partial x_j$에 관한 라이프니츠 법칙을 만족한다. 따라서 국소 코호몰로지의 보조정리
\ref{local-cohomology-lemma-structure-torsion-D-module-regular}에 의해 결과를 얻는다.
\end{proof}

\begin{lemma}
\label{pione-lemma-conclude}
보조정리 \ref{pione-lemma-structure-cohomology}의 상황에서
$i \geq \dim(A) - 1$에 대하여 $H^i(V, \mathcal{O}_V) = 0$이라고 가정하자.
그러면 $V$는 아핀이다.
\end{lemma}

\begin{proof}
$k = A/\mathfrak m$로 놓자. $V \times_{\Spec(A)} \Spec(k) = \emptyset$이므로,
코호몰로지와 기저변환에 의해
$$
R\Gamma(V, \mathcal{O}_V) \otimes_A^\mathbf{L} k = 0
$$
이다. 스킴의 유도 범주의 보조정리 \ref{perfect-lemma-compare-base-change}를 보라.
따라서 스펙트럼 열(대수학 심화의 예 \ref{more-algebra-example-tor})
$$
E_2^{p, q} = \text{Tor}_{-p}(k, H^q(V, \mathcal{O}_V)),\quad
d_2^{p, q} : E_2^{p, q} \to E_2^{p + 2, q - 1}
$$
이 있고, $d_r^{p, q} : E_r^{p, q} \to E_r^{p + r, q - r + 1}$는 영으로 수렴한다.
보조정리 \ref{pione-lemma-structure-cohomology}, 쌍대화 복합체의 보조정리
\ref{dualizing-lemma-tor-injective-hull}, 그리고
$i \geq \dim(A) - 1$에 대하여 $H^i(V, \mathcal{O}_V) = 0$이라는 가정으로부터,
위치 $(p, q) = (0, 0)$으로 들어오거나 그곳에서 나가는 영이 아닌 미분은 없다고 결론내릴 수 있다. 따라서
$H^0(V, \mathcal{O}_V) \otimes_A k = 0$이다. 이는
$\mathfrak m = (x_1, \ldots, x_d)$이면, 아핀 열린부분스킴
$V \times_{\Spec(A)} \Spec(A_{x_i})$들에 의한 열린 덮개
$V = \bigcup V \times_{\Spec(A)} \Spec(A_{x_i})$가 존재하고
($V$가 $A$의 천공 스펙트럼 위에서 아핀이기 때문이다),
$x_1, \ldots, x_d$가 $\Gamma(V, \mathcal{O}_V)$의 단위 아이디얼을 생성한다는 뜻이다.
성질의 보조정리 \ref{properties-lemma-characterize-affine}에 의해 이는 $V$가 아핀임을 뜻한다.
\end{proof}

\begin{theorem}
\label{pione-theorem-global}
$Y$를 한 체 위의 우수한 정칙 스킴이라 하자. $f : X \to Y$를 스킴들의 유한형 사상이라 하고,
$X$는 정규라고 하자. $V \subset X$를 $f$가 에탈인 최대 열린부분스킴이라 하자.
그러면 포함사상 $V \to X$는 아핀이다.
\end{theorem}

\begin{proof}
$x \in X$를 그 상이 $y \in Y$인 점이라 하자. $V \cap W$가 아핀이 되는
$x$의 아핀 열린 근방 $W$를 찾으면 충분하다. $\Spec(\mathcal{O}_{X, x})$가 스킴 $W$들의
극한이므로, 이는
$$
V_x = V \times_X \Spec(\mathcal{O}_{X, x})
$$
가 아핀인 것과 동치이다(극한의 보조정리 \ref{limits-lemma-limit-affine}).
따라서 사상
$X \times_Y \Spec(\mathcal{O}_{Y, y}) \to \Spec(\mathcal{O}_{Y, y})$에 대해 정리가 성립하면
원래 정리도 성립한다. 특히 $Y$가 유한차원 정칙 스킴이라고 가정할 수 있으므로 차원
$d = \dim(Y)$에 관해 귀납할 수 있다. 같은 논증을 다시 결합하면,
$Y$는 닫힌점 $y$를 갖는 국소 스킴이고,
$V \cap (X \setminus f^{-1}(\{y\}) \to X \setminus f^{-1}(\{y\})$는
아핀이라고 가정해도 된다.

\medskip\noindent
$x \in X$를 $y$ 위의 점이라 하자. $x \in V$이면 증명할 것이 없다.
$f^{-1}(\{y\}) \cap V$는 닫힌점들의 유한집합임에 주의하자
(에탈 사상의 올은 이산적이기 때문이다). 따라서 $X$를 $x$의 아핀 열린 근방으로 바꾼 뒤,
$y \not \in f(V)$라고 가정해도 된다. $V$가 아핀임을 증명해야 한다.

\medskip\noindent
$e(V)$를 $H^i(V, \mathcal{O}_V) \not = 0$인 $i$들의 최댓값이라 하자. $X$가 아핀이므로,
정수 $e(V)$는 수 $e(V_x)$들($x \in X \setminus V$)의 최댓값이다. 국소 코호몰로지의 보조정리
\ref{local-cohomology-lemma-cd-local} 및 국소 코호몰로지의 보조정리
\ref{local-cohomology-lemma-cd}에 나오는 코호몰로지 차원의 특성화를 보라.
국소 코호몰로지의 보조정리 \ref{local-cohomology-lemma-cd-dimension}에 의해
$e(V_x) \leq \dim(\mathcal{O}_{X, x}) - 1$이다.
$\dim(\mathcal{O}_{X, x}) \geq 2$이면, 분기 자취의 순수성
(보조정리 \ref{pione-lemma-purity-ramification})에 의해 $V_x$는 $\mathcal{O}_{X, x}$의
천공 스펙트럼보다 진정으로 작다. $\mathcal{O}_{X, x}$가 정규이고 우수하므로,
Hartshorne--Lichtenbaum 소멸
(국소 코호몰로지의 보조정리 \ref{local-cohomology-lemma-affine-complement})에 의해
$e(V_x) \leq \dim(\mathcal{O}_{X, x}) - 2$가 따른다. 한편 $X \to Y$는 유한형이고,
$V \subset X$는 조밀하므로(필요하면 $X$를 $V$의 폐포로 바꾼다),
차원 공식(스킴의 사상의 보조정리 \ref{morphisms-lemma-dimension-formula})에 의해
$\dim(\mathcal{O}_{X, x}) \leq d$이다. 따라서 $e(V) \leq \max(0, d -  2)$이다.
$d \geq 2$이면 보조정리 \ref{pione-lemma-conclude}에 의해 $V$는 아핀이다.
$d = 1$ 또는 $d = 0$이면 $\mathcal{O}_{Y, y}$의 천공 스펙트럼은 아핀이므로
$V$도 아핀이다.
\end{proof}




\section{매끄럽고 고유한 경우의 특수화 사상}
\label{pione-section-specialization-smooth-proper}

\noindent
이 절에서는 다음 결과를 논한다. $f : X \to S$를 스킴들의 고유하고 매끄러운 사상이라 하자.
$s \leadsto s'$를 $S$의 점들의 특수화라 하자. 그러면 절 \ref{pione-section-specialization-map}의 특수화 사상
$$
sp : \pi_1(X_{\overline{s}}) \longrightarrow \pi_1(X_{\overline{s}'})
$$
은 전사이고 다음이 성립한다.
\begin{enumerate}
\item $\kappa(s')$의 표수가 영이면 이는 동형사상이다.
\item $\kappa(s')$의 표수가 $p > 0$이면 이는 $p$와 서로소인 최대 몫들 사이의 동형사상을 유도한다.
\end{enumerate}

\begin{lemma}
\label{pione-lemma-specialization-map-surjective}
$f : X \to S$를 기하학적으로 연결인 올들을 갖는 평탄 고유 사상이라 하자.
$s' \leadsto s$를 특수화라 하자. $X_s$가 기하학적으로 기약이면, 특수화 사상
$sp : \pi_1(X_{\overline{s}'}) \to \pi_1(X_{\overline{s}})$는 전사이다.
\end{lemma}

\begin{proof}
$X_s$가 기하학적으로 기약이므로, 필요하면 $S$를 축소하여 모든 올이 기하학적으로 기약이라고
가정해도 된다. 스킴의 사상 심화의 보조정리
\ref{more-morphisms-lemma-geometrically-reduced-open}을 보라.
$\mathcal{O}_{S, s} \to A \to \kappa(\overline{s}')$를 특수화 사상의 구성에 나타나는 사상들이라 하자.
절 \ref{pione-section-specialization-map}을 보라. 따라서
$$
\pi_1(X_{\overline{s}'}) \to \pi_1(X_A)
$$
가 전사임을 보이면 충분하다. 이는 명제
\ref{pione-proposition-first-homotopy-sequence}와 $\pi_1(\Spec(A)) = \{1\}$에서 따른다.
\end{proof}

\begin{proposition}
\label{pione-proposition-specialization-map-isomorphism}
$f : X \to S$를 기하학적으로 연결인 올들을 갖는 매끄럽고 고유한 사상이라 하자.
$s' \leadsto s$를 특수화라 하자. $\kappa(s)$의 표수가 영이면 특수화 사상
$$
sp : \pi_1(X_{\overline{s}'}) \to \pi_1(X_{\overline{s}})
$$
은 동형사상이다.
\end{proposition}

\begin{proof}
이 사상은 보조정리 \ref{pione-lemma-specialization-map-surjective}에 의해 전사이다. 따라서 단사임을 보이면 충분하다.

\medskip\noindent
$S$는 아핀이라고 가정해도 된다. 그러면 $S$는 $\mathbf{Z}$ 위 유한형인 아핀 스킴들의
여과 역극한이다. 따라서 $X \to S$가 $X_0 \to S_0$의 기저변환이라고 가정할 수 있다.
여기서 $S_0$는 유한형 $\mathbf{Z}$-대수의 스펙트럼이고, $X_0 \to S_0$는 매끄럽고 고유하다.
극한의 보조정리 \ref{limits-lemma-descend-finite-presentation},
\ref{limits-lemma-descend-smooth}, 그리고 \ref{limits-lemma-eventually-proper}를 보라. 보조정리
\ref{pione-lemma-specialization-map-base-change}에 의해 기저가 Noether인 경우로 환원된다.

\medskip\noindent
보조정리 \ref{pione-lemma-specialization-map-discrete-valuation-ring}을 적용하면, 기저 $S$가 엄밀 헨젤 이산 부치환환
$A$의 스펙트럼이고 $A$ 위의 특수화 사상을 다루는 경우로 환원된다. $K$를 $A$의 분수체라 하자.
$\overline{K}$를 기하학적 일반점 $\overline{\eta}$($\Spec(A)$의 점)에 대응하는 대수적 폐포로 택한다.
$\overline{K}/L/K$가 유한 분리 확대일 때, $B \subset L$을 $L$ 안에서 취한 $A$의 정수적 폐포라 하자.
이는 대수학 심화의 주의 \ref{more-algebra-remark-finite-separable-extension}에 의해 이산 부치환환이다.

\medskip\noindent
$X \to \Spec(A)$를 앞 단락의 사상이라 하자. 특수화 사상의 단사성을 보이려면,
$X_{\overline{\eta}}$의 임의의 유한 에탈 덮개 $V$가 어떤 유한 에탈 덮개 $Y \to X$의 기저변환임을
증명하면 충분하다. 실제로 이 경우 보조정리 \ref{pione-lemma-functoriality-galois-injective}에 의해
$\pi_1(X_{\overline{\eta}}) \to \pi_1(X) = \pi_1(X_s)$는 단사이다.

\medskip\noindent
$V$가 주어졌을 때 먼저 보조정리 \ref{pione-lemma-perfection}에 의해 $V$를
$V' \to X_{K^{sep}}$로 강하시킬 수 있고, 이어서 보조정리 \ref{pione-lemma-limit}에 의해
$V'' \to X_L$로 강하시킬 수 있다. $Z \to X_B$를 $V''$ 안에서 취한 $X_B$의 정규화라 하자.
$Z$는 정규이고 $Z_L = V''$가 $X_L$ 위의 스킴으로서 성립함에 주의하자. 따라서 $Z \to X_B$는
일반올 위에서 유한 에탈이다. 문제는 $Z \to X_B$가 모든 곳에서 에탈인지 알 수 없다는 데 있다.
$X \to \Spec(A)$는 기하학적으로 연결인 매끄러운 올들을 가지므로 특수올 $X_s$는
기하학적으로 기약이다. 따라서 $X_B \to \Spec(B)$의 특수올은 기약이다. 그 일반점을
$\xi_B$라 하자. $\xi_B$로 가는 $Z$의 점들을 $\xi_1, \ldots, \xi_r$라 하자. 이제 첫 번째 문제
(그리고 이것이 유일한 문제임을 알게 될 것이다)는 이산 부치환환의 확대
$$
\mathcal{O}_{X_B, \xi_B} \subset \mathcal{O}_{Z, \xi_i}
$$
가 분기할 수 있다는 점이다. 이 확대의 분기지수를 $e_i$라 하자. $\kappa(s)$의 표수가 영이므로
이 분기는 순분기임에 주의하자.

\medskip\noindent
분기를 제거하기 위해, 유도된 확대 $B'/B$의 분기지수 $e$가 $e_i$로 나누어지는
더 큰 유한 분리 확대 $K^{sep}/L'/L/K$를 택한다. $Z'$를
$\Spec(B') \to \Spec(B)$에 관한 $Z$의 정규화 기저변환이라 하자. 스킴의 사상 심화의 절
\ref{more-morphisms-section-reduced-fibre-theorem}의 논의를 보라.
$\xi_{i, j}$를 $\xi_{B'}$ 및 $Z$의 점 $\xi_i$로 가는 $Z'$의 점이라 하자. 그러면 국소환
$$
\mathcal{O}_{Z', \xi_{i, j}}
$$
은 $\mathcal{O}_{Z, \xi_i}$의 분수체를 $F_i$라 할 때,
$L' \otimes_L F_i$ 안에서 취한 $\mathcal{O}_{Z, \xi_i}$의 정수적 폐포의 국소화이다. 세부사항은 생략한다.
따라서 Abhyankar 보조정리
(대수학 심화의 보조정리 \ref{more-algebra-lemma-abhyankar})에 의해
$$
\mathcal{O}_{X_{B'}, \xi_{B'}} \subset \mathcal{O}_{Z', \xi_{i, j}}
$$
은 비분기이다. 따라서 사상 $Z' \to X_{B'}$는 여차원 $1$의 바깥에서 에탈이다.
따라서 분기 자취의 순수성(보조정리 \ref{pione-lemma-purity})에 의해
$Z' \to X_{B'}$는 유한 에탈이다.

\medskip\noindent
그러나 $A \to B'$가 유도하는 잉여체 확대는 자명하다($A$의 잉여체는 표수 영인
분리적으로 닫힌 체이므로 대수적으로 닫혀 있기 때문이다). 따라서 보조정리
\ref{pione-lemma-finite-etale-on-proper-over-henselian}을 두 번 적용하면,
$Z'$가 어떤 유한 에탈 덮개 $Y \to X$의 기저변환임을 알 수 있다
(먼저 $A$ 위에서 $Y$를 얻고, 다음으로 $B$로의 당김이 $Z'$와 동형임을 보인다).
이로써 증명이 끝난다.
\end{proof}

\noindent
$G$를 프로유한군, $p$를 소수라 하자. 정의에 따라 {\it $p$와 서로소인 최대 몫}은
$$
G' = \lim_{U \subset G\text{는 열리고 정규이며 그 지수가 }p\text{와 서로소}} G/U
$$
이다. $X$가 연결 스킴이고 $p$가 주어졌을 때, $\pi_1(X)$의 $p$와 서로소인 최대 몫을
$\pi'_1(X)$로 쓴다.

\begin{theorem}
\label{pione-theorem-specialization-map-isomorphism-prime-to-p}
$f : X \to S$를 기하학적으로 연결인 올들을 갖는 매끄럽고 고유한 사상이라 하자.
$s' \leadsto s$를 특수화라 하자. $\kappa(s)$의 표수가 $p$이면 특수화 사상
$$
sp : \pi_1(X_{\overline{s}'}) \to \pi_1(X_{\overline{s}})
$$
은 전사이고, p와 서로소인 최대 몫들 사이의 동형사상
$$
\pi'_1(X_{\overline{s}'}) \cong \pi'_1(X_{\overline{s}})
$$
을 유도한다.
\end{theorem}

\begin{proof}
증명은 명제 \ref{pione-proposition-specialization-map-isomorphism}의 증명과 완전히 같지만,
다음 차이들이 있다.
\begin{enumerate}
\item $X/A$가 주어졌을 때, 함자
$\textit{F\'Et}_X \to \textit{F\'Et}_{X_{\overline{\eta}}}$가 본질상 전사임을 보이는 데까지 나아가지는 않는다.
Galois 군의 위수가 $p$와 서로소인 Galois 대상이 본질상에 속한다는 것만 보인다. 같은 논증으로
$\pi'_1(X_{\overline{s}'}) \to \pi'_1(X_{\overline{s}})$의 단사성을 결론내리기에는 이것으로 충분하다.
\item 확대
$\mathcal{O}_{X_B, \xi_B} \subset \mathcal{O}_{Z, \xi_i}$
는 순분기이다. 대응하는 분수체 확대가 Galois이고 그 군의 위수가 $p$와 서로소이기 때문이다.
대수학 심화의 보조정리 \ref{more-algebra-lemma-galois-conclusion}을 보라.
\item 확대 $\kappa_B/\kappa_A$는 더 이상 반드시 자명하지는 않지만 순수 비분리이다. 따라서 사상
$X_{\kappa_B} \to X_{\kappa_A}$는 보편 위상동형사상이고, 명제
\ref{pione-proposition-universal-homeomorphism}에 의해 기본군들의 동형사상을 유도한다.
\end{enumerate}
\end{proof}













\section{순분기}
\label{pione-section-tame}

\noindent
$X \to Y$를 $\mathbf{Z}$ 위 유한형인 스킴들의 유한 에탈 사상이라 하자. $f$가 $\infty$에서
순분기한다는 것을 정의하는 방법은 여러 가지가 있다. 논문 \cite{Kerz-Schmidt}에서는
이 개념들이 어느 정도 일치하는지 논한다.

\medskip\noindent
이 절에서는 무한대에서의 순분기성 개념보다 앞서는, 더 초등적인 별개의 문제를 논한다.
\cite{Grothendieck-Murre}의 다소 다른 논의와 비교하라. 다음이 주어졌다고 가정하자.
\begin{enumerate}
\item 국소 Noether 스킴 $X$,
\item 조밀 열린집합 $U \subset X$,
\item 유한 에탈 사상 $f : Y \to U$.
\end{enumerate}
또한 $Z \cap U = \emptyset$인 모든 소인자 $Z \subset X$에 대하여,
$Z$의 일반점 $\xi$에서의 $X$의 국소환 $\mathcal{O}_{X, \xi}$가
이산 부치환환이라고 하자.
$K_\xi$를 $\mathcal{O}_{X, \xi}$의 분수체라 놓으면 스킴들의 데카르트 도식
$$
\xymatrix{
\Spec(K_\xi) \ar[r] \ar[d] & U \ar[d] \\
\Spec(\mathcal{O}_{X, \xi}) \ar[r] & X
}
$$
을 얻는다. 특히 $Y \times_U \Spec(K_\xi)$는 유한 분리 대수 $L_\xi/K_\xi$의 스펙트럼이다.
이때 $L_\xi/K_\xi$가 위와 같은 모든 $(Z, \xi)$에 대해 $\mathcal{O}_{X, \xi}$에 관하여 각각
비분기이거나 순분기라는 뜻에서, 각각
{\it $Y$가 $X$ 위 여차원 $1$에서 비분기이다},
{\it $Y$가 $X$ 위 여차원 $1$에서 순분기이다}라고 한다. 대수학 심화의 정의
\ref{more-algebra-definition-types-of-extensions}를 보라. 더 정확히 말하면, $L_\xi$를
$K_\xi$의 유한 분리 체 확대들의 곱으로 분해하고, 그 각각이
$\mathcal{O}_{X, \xi}$에 관하여 각각 비분기이거나 순분기일 것을 요구한다.

\begin{lemma}
\label{pione-lemma-pullback-tame-codim1}
$X' \to X$를 국소 Noether 스킴들의 사상이라 하고, $U \subset X$를 조밀 열린집합이라 하자. 다음을 가정하자.
\begin{enumerate}
\item $U' = f^{-1}(U)$는 $X'$의 조밀 열린집합이다.
\item $Z \cap U = \emptyset$인 모든 소인자 $Z \subset X$에 대하여,
$Z$의 일반점 $\xi$에서의 $X$의 국소환 $\mathcal{O}_{X, \xi}$는 이산 부치환환이다.
\item $Z' \cap U' = \emptyset$인 모든 소인자 $Z' \subset X'$에 대하여,
$Z'$의 일반점 $\xi'$에서의 $X'$의 국소환 $\mathcal{O}_{X', \xi'}$는 이산 부치환환이다.
\item $\xi' \in X'$이 (3)과 같은 점이면 $\xi = f(\xi')$는 (2)와 같은 점이다.
\end{enumerate}
이때 $f : Y \to U$가 유한 에탈이고 $Y$가 $X$ 위 여차원 $1$에서 각각 비분기이거나 순분기이면,
$Y' = Y \times_X X' \to U'$는 유한 에탈이고, $Y'$는 $X'$ 위 여차원 $1$에서 각각
비분기이거나 순분기이다.
\end{lemma}

\begin{proof}
이 보조정리에서 유일하게 흥미로운 사실은 대수학 심화의 보조정리
\ref{more-algebra-lemma-tame-goes-up}에 주어진 가환대수학의 결과이다.
\end{proof}

\noindent
위에서 도입한 용어를 사용하면 앞에서 얻은 순수성 결과를 다음과 같이 간결하게 다시 말할 수 있다.

\begin{lemma}
\label{pione-lemma-purity-one-divisor-general}
$X$를 국소 Noether 스킴이라 하자. $U \subset X$를 조밀 열린부분스킴이라 하자.
$Y \to U$를 유한 에탈 사상이라 하자. 다음을 가정한다.
\begin{enumerate}
\item $Y$는 $X$ 위 여차원 $1$에서 비분기이다.
\item $\mathcal{O}_{X, x}$는 모든 $x \in X \setminus U$에 대하여 정칙이다.
\end{enumerate}
그러면 유한 에탈 사상 $Y' \to X$로서 그 $U$로의 제한이 $Y$인 것이 존재한다.
\end{lemma}

\begin{proof}
$\xi \in X \setminus U$를 $X \setminus U$의 여차원 $1$인 기약 성분의 일반점이라 하자. 그러면
$\mathcal{O}_{X, \xi}$는 이산 부치환환이다. 위의 논의와 마찬가지로
$Y \times_U \Spec(K_\xi) = \Spec(L_\xi)$라 쓰자. $B_\xi$를
$\mathcal{O}_{X, \xi}$의 $L_\xi$ 안에서의 정수적 폐포라 하자. $Y$가 $X$ 위 여차원 $1$에서 비분기라는
가정은 $\mathcal{O}_{X, \xi} \to B_\xi$가 유한 에탈이라는 뜻이다.
따라서 유한 에탈 사상 $Y_\xi \to \Spec(\mathcal{O}_{X, \xi})$과 동형사상
$$
Y \times_U \Spec(K_\xi) \cong
Y_\xi \times_{\Spec(\mathcal{O}_{X, \xi})} \Spec(K_\xi)
$$
을 얻는다. 이는 $\Spec(K_\xi)$ 위의 동형사상이다. 극한의 보조정리
\ref{limits-lemma-glueing-near-point}에 의해 $\xi$를 포함하는 열린부분스킴 $U \subset U' \subset X$와
유한 표시 사상 $Y' \to U'$로서, $U$로 제한하면 $Y$를 되찾고
$\Spec(\mathcal{O}_{X, \xi})$로 제한하면 $Y_\xi$를 되찾는 것을 얻는다. 마지막으로 극한의 보조정리
\ref{limits-lemma-glueing-near-point-properties}에 의해 사상 $Y' \to U'$는, 필요하면 $U'$를 더 작은
열린집합으로 축소한 뒤 유한 에탈이라고 가정해도 된다. $X \setminus U$의 여차원 $1$인
다른 일반점들에 대하여 이 논의를 반복하면, 유한 에탈 사상 $Y' \to U'$로서 $Y \to U$를 연장하는 것이 있다고
가정해도 된다. 여기서 열린부분스킴 $U' \subset X$는 $U$와 모든 여차원 $1$인 점
($X \setminus U$의 점들)을 포함한다.
마지막으로 보조정리 \ref{pione-lemma-extend-pure}를 $Y' \to U'$에 적용한다. 실제로 모든 국소환
$\mathcal{O}_{X, x}$, $x \in X \setminus U'$는 정칙이고
$\dim(\mathcal{O}_{X, x}) \geq 2$를 만족한다. 따라서 보조정리 \ref{pione-lemma-local-purity}에 의해
$\mathcal{O}_{X, x}$에 대하여 순수성이 성립한다.
\end{proof}

\begin{lemma}
\label{pione-lemma-purity-one-divisor}
$X$를 국소 Noether 스킴이라 하자. $D \subset X$를 유효 Cartier 인자라 하고, $D$가 정칙 스킴이라고 하자.
$Y \to X \setminus D$를 유한 에탈 사상이라 하자. $Y$가 $X$ 위 여차원 $1$에서 비분기이면,
유한 에탈 사상 $Y' \to X$로서 그 $X \setminus D$로의 제한이 $Y$인 것이 존재한다.
\end{lemma}

\begin{proof}
이는 보조정리 \ref{pione-lemma-purity-one-divisor-general}의 특수한 경우이다. 우선 $D$는 $X$에서
어디에서도 조밀하지 않은 부분집합이므로(인자의 절
\ref{divisors-section-effective-Cartier-divisors}의 논의를 보라),
$X \setminus D$는 $X$에서 조밀하다. 다음으로 환 $\mathcal{O}_{X, x}$는 모든 $x \in D$에 대하여
정칙 국소환이다. 이는 대수학의 보조정리 \ref{algebra-lemma-regular-mod-x}와
$\mathcal{O}_{D, x}$가 정칙이라는 가정에서 따른다.
\end{proof}

\begin{example}[표준 순분기 사상]
\label{pione-example-tamely-ramified}
$A$를 Noether 환이라 하자. $f \in A$를 비영인자라 하고 $A/fA$가 기약이라고 하자. 그러면
$A_\mathfrak p$는 소원 $f$를 갖는 이산 부치환환이다. 이는 $f$ 위에 놓인 임의의 극소 소아이디얼
$\mathfrak p$에 대하여 성립한다. $e \geq 1$을 $A$에서 가역인 정수라 하자.
$$
C = A[x]/(x^e - f)
$$
로 놓자. 그러면 $\Spec(C) \to \Spec(A)$는 유한 국소 자유 사상이고 $A_f$의 스펙트럼 위에서 에탈이다.
유한 에탈 사상
$$
\Spec(C_f) \longrightarrow \Spec(A_f)
$$
은 $\Spec(A)$ 위 여차원 $1$에서 순분기한다. 이 순분기성은 대수학 심화의 보조정리
\ref{more-algebra-lemma-characterize-tame}에 나오는 순분기 확대의 특성화에서 즉시 따른다.
\end{example}

\noindent
정칙 인자에 대한 Abhyankar 보조정리의 한 형태를 진술한다.

\begin{lemma}[정칙 인자에 대한 Abhyankar 보조정리]
\label{pione-lemma-abhyankar-one-divisor}
$X$를 국소 Noether 스킴이라 하자. $D \subset X$를 유효 Cartier 인자라 하고, $D$가 정칙 스킴이라고 하자.
$Y \to X \setminus D$를 유한 에탈 사상이라 하자. $Y$가 $X$ 위 여차원 $1$에서 순분기하면,
$X$ 위 에탈 국소적으로 사상 $Y \to X$는 예
\ref{pione-example-tamely-ramified}에 기술한 표준 순분기 사상들의 유한 서로소 합으로 주어진다.
\end{lemma}

\begin{proof}
모든 $x \in X$에 대하여 에탈 근방 $(U, u) \to (X, x)$로서
$U = \Spec(A)$이고 기저변환 $Y \times_X U \to U$가 예
\ref{pione-example-tamely-ramified}에서와 같은 표준 순분기 사상들의 유한 서로소 합이 되는 것을 찾아야 한다.
$x \in D$라고 가정하자. $x \not \in D$인 경우는 에탈 사상의 보조정리
\ref{etale-lemma-finite-etale-etale-local}과 예 \ref{pione-example-tamely-ramified}에서
$e = 1$ 및 $f = 1$로 택하는 것에서 따른다.

\medskip\noindent
이 단락에서는 $X$가 엄밀 헨젤 국소환의 스펙트럼이고 $x$가 닫힌점인 경우로 환원한다. 실제로 $X$를
축소하여 $X = \Spec(A)$이고 $D \subset X$가 비영인자 $f \in A$로 주어진다고 가정해도 된다.
그러면 $Y$는 $\Spec(A_f)$의 유한 에탈 덮개이므로 아핀이다. $Y = \Spec(B)$라 쓰자.
$A^{sh}$를 $\mathcal{O}_{X, x}$의 엄밀 헨젤화라 하자. $A \to A^{sh}$가 평탄이고 $x \in D$이므로,
$f$는 $A^{sh}$의 극대 아이디얼 안의 비영인자로 보내진다. 또한
$A^{sh}/fA^{sh} = (A/fA)^{sh}$는(대수학의 보조정리
\ref{algebra-lemma-quotient-strict-henselization}), $\mathcal{O}_{D, x}$의 엄밀 헨젤화로서
정칙 국소환이다. 대수학 심화의 보조정리
\ref{more-algebra-lemma-henselization-regular}을 보라. 보조정리
\ref{pione-lemma-pullback-tame-codim1}에 의해 $Y$의 $\Spec(A^{sh})$로의 기저변환은 여차원 $1$에서
순분기한다. 이 $Y$의 $\Spec(A^{sh})$로의 기저변환에 대하여 주장을 증명했다고 가정하자.
$A^{sh}$가 엄밀 헨젤이므로 모든 에탈 근방은 단면을 가지며, 동형사상
$$
A^{sh} \otimes_A B \cong
\prod\nolimits_{i \in I} A^{sh}_f[x]/(x^{e_i} - f)
$$
을 얻는다. 여기서 $I$는 유한집합이고 각 $e_i$는 $A^{sh}$에서 가역인 정수이다.
$A^{sh} = \colim A_\lambda$를 에탈 $A$-대수 $A_\lambda$들의 여과 여극한으로 쓰자.
위의 동형사상은 충분히 큰 $\lambda$에 대하여 동형사상
$$
A_\lambda \otimes_A B \cong \prod\nolimits_{i \in I}
(A_\lambda)_f[x]/(x^{e_i} - f)
$$
으로 내려간다. 예를 들어 대수학의 보조정리
\ref{algebra-lemma-colimit-category-fp-algebras}를 보라. $\lambda$를 조금 더 크게 잡으면
$e_i$가 $A_\lambda$에서 가역이라고 가정해도 된다. 그러면
$\Spec(A_\lambda) \to \Spec(A)$는 우리가 찾는 $x$의 에탈 근방이다.

\medskip\noindent
$X = \Spec(A)$라고 가정하자. 여기서 $A$는 엄밀 헨젤 국소환이고, $x \in X$는 $A$의 극대 아이디얼
$\mathfrak m$에 대응한다. $f \in \mathfrak m$을 $D$를 잘라내는 비영인자라 하자.
그러면 $A/fA$가 정칙이고 대수학의 보조정리 \ref{algebra-lemma-regular-mod-x}에 의해 $A$도 정칙이다.
정칙 국소환의 몇몇 성질, 예를 들어 정규 정역이라는 사실을 사용할 것이다. 대수학의 보조정리
\ref{algebra-lemma-regular-domain} 및 \ref{algebra-lemma-regular-normal}을 보라. 특히 $A$는 정역이다.
위와 마찬가지로 $Y = \Spec(B)$는 아핀이고 $A_f \to B$는 유한 에탈이다. 따라서 $B$도 정규이다
(대수학의 보조정리 \ref{algebra-lemma-normal-goes-up}). 그러므로 대수학의 보조정리
\ref{algebra-lemma-characterize-reduced-ring-normal}에 의해 $Y$는 정규 정역들의 스펙트럼의 유한 서로소 합이다.
따라서 $B$도 정역이라고 가정해도 된다. $A$와 $A/fA$가 정역이므로 $f$는 $A$의 소원이고,
높이 $1$인 소아이디얼 $\mathfrak p = fA$를 생성한다. $A_\mathfrak p$는 $f$를 소원으로 갖는
이산 부치환환임에 주의하자.

\medskip\noindent
$K$를 $A$의 분수체, $L$을 $B$의 분수체라 하자. 순분기라는 가정은 $L$이
$A_\mathfrak p$에 대하여 순분기한다는 뜻이다. 대수학 심화의 주의
\ref{more-algebra-remark-finite-separable-extension}에서와 같이 $e \geq 1$을 $L$의 $A_\mathfrak p$ 위의
분기지수들의 곱이라 하자. 그러면 $e$는 $\kappa(\mathfrak p)$에서 가역이지만, 이 시점에서는 $e$가
$A/fA$ 또는 $A$에서 가역인지 알 수 없다.

\medskip\noindent
유한 자유 $A$-대수
$$
A' = A[x]/(x^e - f)
$$
를 생각하자. 이는 엄밀 헨젤 국소환의 국소 유한 확대이므로 엄밀 헨젤이다.
$f' = x$는 $A'$의 비영인자이고 $A'/f'A' \cong A/fA$는 정칙환임에 주의하자.
따라서 앞에서와 마찬가지로 $A'$는 정칙이고, 특히 정역이다.
$B' = B \otimes_A A' = B \otimes_{A_f} A'_{f'}$.
로 놓자. Abhyankar 보조정리(대수학 심화의 보조정리 \ref{more-algebra-lemma-abhyankar})에 의해
$\Spec(B')$는 $\Spec(A')$ 위 여차원 $1$에서 비분기이다. 실제로 보조정리
\ref{pione-lemma-pullback-tame-codim1}에 의해 $\Spec(B')$는 적어도 여전히 $\Spec(A')$ 위 여차원 $1$에서
순분기한다. 그러나 Abhyankar 보조정리에 의해 모든 분기지수가 $1$이 되었다. 보조정리
\ref{pione-lemma-purity-one-divisor}에 의해 유한 에탈 사상 $\Spec(B') \to \Spec(A'_{f'})$는
유한 에탈 사상 $\Spec(C) \to \Spec(A')$로 연장된다. 그런데 $A'$가 엄밀 헨젤이므로 $\Spec(C)$는
$\Spec(A')$의 복사본들의 유한 서로소 합이다. 따라서 적어도 하나의 사상
$\Spec(A'_{f'}) \to \Spec(B')$가 존재한다. 그러므로 $A_f$-대수로서의 포함
$B \subset A'_{f'}$가 존재한다.

\medskip\noindent
따라서 포함관계 $K \subset L \subset K[f^{1/e}]$를 갖는다. 대수학 심화의 보조정리
\ref{more-algebra-lemma-pull-root-uniformizer}에 의해 $L$은 $K[f^{1/n}]$과 같다. 여기서 $n$은
$e$의 어떤 약수이다. 이제 $B$의 정규성에서 $B \cong A_f[y]/(y^n - f)$가 따른다. 세부사항은 생략한다.
그러나 환 준동형 $A_f \to A_f[y]/(y^n - f)$의 분기 자취는 $ny^{n - 1}$로 잘라지므로,
$n$은 $A_f$의 단위원이어야 한다. $f$가 소원이므로 이는 $n = u f^s$라는 뜻이다.
여기서 $u$는 $A$의 어떤 단위원이고 $s \in \mathbf{Z}$이다. $n$이
$\kappa(\mathfrak p)$의 단위원으로 보내지므로 $s = 0$이고, 증명이 끝난다.
\end{proof}

\begin{lemma}
\label{pione-lemma-extend-tame-covering-normal}
보조정리 \ref{pione-lemma-abhyankar-one-divisor}의 상황에서 $Y$ 안에서 취한 $X$의 정규화는 유한 국소 자유 사상
$\pi : Y' \to X$이고 다음을 만족한다.
\begin{enumerate}
\item $Y'$의 $X \setminus D$로의 제한은 $Y$와 동형이다.
\item $D' = \pi^{-1}(D)_{red}$는 $Y'$ 위의 유효 Cartier 인자이다.
\item $D'$는 정칙 스킴이다.
\end{enumerate}
더 나아가 $X$ 위 에탈 국소적으로 사상 $Y' \to X$는 사상
$$
\Spec(A[x]/(x^e - f)) \to \Spec(A)
$$
들의 유한 서로소 합이다. 여기서 $A$는 Noether 환이고, $f \in A$는 비영인자이며 $A/fA$는 정칙이고,
$e \geq 1$은 $A$에서 가역이다.
\end{lemma}

\begin{proof}
이는 보조정리 \ref{pione-lemma-abhyankar-one-divisor}에 대한 보충일 뿐이며, 실제로 그 보조정리의 증명을 읽으면
이 보조정리가 참이라는 것은 거의 즉시 따른다. 그러나 보조정리
\ref{pione-lemma-abhyankar-one-divisor}의 결과에서 직접 유도할 수도 있다. 실제로 $Y$ 안에서 취한 $X$의
정규화는 에탈 기저변환과 가환한다. 스킴의 사상 심화의 보조정리
\ref{more-morphisms-lemma-normalization-smooth-localization}을 보라. 따라서 $Y' \to X$의 국소 구조에 관한
주장은 $X$ 위 에탈 국소적으로 증명해도 된다. 그러므로 보조정리
\ref{pione-lemma-abhyankar-one-divisor}에 의해 $X = \Spec(A)$이고 $A$는 Noether 환이며, 비영인자
$f\in A$에 대하여 $A/fA$가 정칙이고, $Y$는 환 $A_f[x]/(x^e - f)$의 스펙트럼들의 유한 서로소 합이며,
$e$가 $A$에서 가역이라고 가정해도 된다. $A_f[x]/(x^e - f)$ 안에서 취한 $A$의 정수적 폐포가
$A' = A[x]/(x^e - f)$와 같다는 검증은 생략한다
(이를 보려면 $(f)$ 위에 놓인 소아이디얼들에서의 $A'$의 국소화가 정칙임을 보이면 된다).
세부사항은 생략한다.
\end{proof}

\begin{lemma}
\label{pione-lemma-tame-covering-split}
보조정리 \ref{pione-lemma-abhyankar-one-divisor}의 상황에서 $Y' \to X$를 보조정리
\ref{pione-lemma-extend-tame-covering-normal}의 사상이라 하자. $R$을 분수체 $K$를 갖는 이산 부치환환이라 하자.
$$
t : \Spec(R) \to X
$$
를 사상이라 하고, 스킴론적 역상 $t^{-1}D$가 $\Spec(R)$의 기약 닫힌점이라고 하자.
\begin{enumerate}
\item $t|_{\Spec(K)}$가 $Y$의 점으로 올려지면, 올림
$t' : \Spec(R) \to Y'$로서 $Y' \to X$가 $t'(\Spec(R))$을 따라 에탈인 것을 얻는다.
\item $\Spec(K) \times_X Y$가 $\Spec(K)$의 복사본들의 서로소 합과 동형이면, $Y' \to X$는
$t(\Spec(R))$의 어떤 열린 근방 위에서 유한 에탈이다.
\end{enumerate}
\end{lemma}

\begin{proof}
유한 사상 $Y' \to X$에 고유성에 대한 부치환 판정법을 적용하면, $Y$의 $\Spec(K)$-값 점 중에서
$X$로 가는 사상으로서 $t|_{\Spec(K)}$와 일치하는 것은 사상 $t' : \Spec(R) \to Y'$로 유일하게 올려진다.
따라서 주장 (1)은 의미가 있다.

\medskip\noindent
에탈 근방 $(U, u) \to (X, t(\mathfrak m_R))$로서 $U = \Spec(A)$이고
$Y' \times_X U \to U$가 어떤 $f \in A$에 대하여 보조정리 \ref{pione-lemma-extend-tame-covering-normal}과 같은
기술을 갖는 것을 택하자. 그러면 $\Spec(R) \times_X U \to \Spec(R)$는 에탈이고 전사이다.
$R'$을 $\Spec(R) \times_X U$의 국소환 중 $\Spec(R)$의 닫힌점 위에 놓인 것이라 하자. 그러면 $R'$은
이산 부치환환이고 $R \subset R'$은 이산 부치환환들의 비분기 확대이다
(대수학 심화의 보조정리 \ref{more-algebra-lemma-Dedekind-etale-extension}). $t$에 대한 가정은
사상 $A \to R'$로서
$$
\Spec(R') \to \Spec(R) \times_X U \to U
$$
에 대응하는 것이 $f$를 소원 $\pi \in R'$로 보낸다는 뜻이다. 이제 어떤 $e_i \geq 1$에 대하여
$$
Y' \times_X U =
\coprod\nolimits_{i \in I} \Spec(A[x]/(x^{e_i} - f))
$$
라고 가정하자. 그러면
$$
\Spec(R') \times_U (Y' \times_X U) =
\coprod\nolimits_{i \in I} \Spec(R'[x]/(x^{e_i} - \pi))
$$
를 얻는다. 환 $R'[x]/(x^{e_i} - f)$는 이산 부치환환이다
(대수학 심화의 보조정리 \ref{more-algebra-lemma-pull-root-uniformizer}). 따라서 $R'$의 분수체로 가는
사상을 가지려면 $e_i = 1$이어야 한다.

\medskip\noindent
(1)의 증명. 이 경우 사상 $t' : \Spec(R) \to Y'$는 기저변환에 의해 대응하는 사상
$t'' : \Spec(R') \to Y' \times_X U$를 정하고, 위의 논의에 의해 이는 $i \in I$이고 $e_i = 1$인
성분으로 가야 한다. 따라서 $Y' \times_X U \to U$가 $t''$의 상을 따라 에탈이라는 것은 명백하다.
에탈성은 에탈 기저변환 뒤에 판정할 수 있는 성질이므로, 이것으로 (1)이 증명되었다.

\medskip\noindent
(2)의 증명. 이 경우 가정은 모든 $i \in I$에 대하여 $e_i = 1$이라는 뜻이다.
따라서 $Y' \times_X U \to U$는 유한 에탈이고 앞에서와 마찬가지로 결론을 얻는다.
\end{proof}

\begin{lemma}
\label{pione-lemma-extend-covering}
$S$를 일반점 $\eta$를 갖는 정역 정규 Noether 스킴이라 하자.
$f : X \to S$를 기하학적으로 연결인 올들을 갖는 매끄러운 사상이라 하자.
$\sigma : S \to X$를 $f$의 단면이라 하자. $Z \to X_\eta$를 유한 에탈 Galois 덮개
(절 \ref{pione-section-finite-etale-under-galois})라 하고, 그 군 $G$의 위수가 $S$ 위에서 가역이며
$Z$가 $\sigma(\eta)$로 가는 $\kappa(\eta)$-유리점을 가진다고 하자. 그러면 군 $G$를 갖고
$X_\eta$로의 제한이 $Z$인 유한 에탈 Galois 덮개 $Y \to X$가 존재한다.
\end{lemma}

\begin{proof}
먼저 $S = \Spec(R)$가 이산 부치환환 $R$의 스펙트럼이고 닫힌점을 $s \in S$라 하자.
그러면 $X_s$는 $X$의 유효 Cartier 인자이고, 체 위의 매끄러운 스킴으로서 $X_s$는 정칙이다.
또 일반올 $X_\eta$는 열린부분스킴 $X \setminus X_s$이다.
대수학 심화의 보조정리 \ref{more-algebra-lemma-galois-conclusion}과 $G$에 대한 가정에 의해 $Z$는
$X$ 위 여차원 $1$에서 순분기한다. $Z' \to X$를 보조정리
\ref{pione-lemma-extend-tame-covering-normal}의 사상이라 하자. $G$의 $Z$ 위의 작용이 $G$의 $Z'$ 위의 작용으로
연장됨에 주의하자. 보조정리 \ref{pione-lemma-tame-covering-split}에 의해 $Z' \to X$는
$\sigma(y)$의 어떤 열린 근방 위에서 유한 에탈이다. $X_s$가 기약이므로, 이는 $Z \to X_\eta$가
$X$ 위 여차원 $1$에서 비분기라는 뜻이다. 그러면 보조정리
\ref{pione-lemma-purity-one-divisor}에 의해 $X_\eta$로의 제한이 $Z$인 유한 에탈 사상 $Y \to X$를 얻는다.
물론 $Y \cong Z'$이고(세부사항은 생략한다. 힌트: 에탈 국소적으로 계산하라),
따라서 $Y$는 군 $G$를 갖는 Galois 덮개이다.

\medskip\noindent
일반적인 경우. $U \subset S$를 다음 성질을 갖는 최대 열린부분스킴이라 하자. 즉, 군 $G$를 갖고
$X_\eta$로의 제한이 $Z$와 동형인 유한 에탈 Galois 덮개 $Y \to X \times_S U$가 존재한다고 하자.
모순을 얻기 위해 $U \not = S$라고 가정한다. $s \in S \setminus U$를
$S \setminus U$의 한 기약 성분의 일반점이라 하자. $U_s$를 $U$의
$\Spec(\mathcal{O}_{S, s})$ 안에서의 역상이라 하면, 이는 $\mathcal{O}_{S, s}$의 천공 스펙트럼이다. 사상
$Y \times_S U_s \to X \times_S U_s$는 군 $G$를 갖는 유한 에탈 Galois 덮개
$Y'_s \to X \times_S \Spec(\mathcal{O}_{S, s})$의 제한이라고 주장한다.

\medskip\noindent
먼저 이 주장이 원하는 모순을 준다는 것을 보이자. 극한의 보조정리
\ref{limits-lemma-glueing-near-point}에 의해 $s$를 포함하는 열린부분스킴
$U \subset U' \subset S$와 유한 표시 사상 $Y'' \to U'$로서, $U$로 제한하면 $Y' \to U$를 되찾고
$\Spec(\mathcal{O}_{S, s})$로 제한하면 $Y'_s$를 되찾는 것을 얻는다. 더 나아가 이 보조정리가 주는 범주들의 동치에 의해,
$U'$를 축소한 뒤 사상 $Y'' \to U' \times_S X$와 $U' \times_S X$ 위의 $G$의 $Y''$ 위 작용이 존재하고,
이들이 $U$ 및 $\Spec(\mathcal{O}_{S, s})$로 기저변환한 뒤 주어진 사상과 작용에 맞는다고 가정해도 된다.
필요하면 $U'$를 더 축소하여 $Y'' \to U \times_S X$가 유한 에탈이라고 가정해도 된다.
극한의 보조정리 \ref{limits-lemma-glueing-near-point-properties}를 보라.
이는 $Y$가 군 $G$를 갖는 유한 에탈 Galois 덮개로 연장되는 $S$의 진정으로 더 큰 열린집합을
찾았다는 뜻이며, 우리가 원한 모순을 준다.

\medskip\noindent
주장의 증명. $S$를 $\Spec(\mathcal{O}_{S, s})$로 바꾸어도 되고, 실제로 그렇게 하자. 그러면
$S = \Spec(A)$이고 $(A, \mathfrak m)$은 정규 국소 정역이다. 또한 $U \subset S$는 천공 스펙트럼이고,
군 $G$를 갖는 유한 에탈 Galois 덮개 $Y \to X \times_S U$가 있다. $\dim(A) = 1$이면
증명의 첫 단락에 의해 $Y$를 $X$의 Galois 덮개로 연장할 수 있다. 따라서
$\dim(A) \geq 2$라고 가정해도 되고, $S$가 정규이므로 $\text{depth}(A) \geq 2$이다.
대수학의 보조정리 \ref{algebra-lemma-criterion-normal}을 보라. $X \to S$가 평탄이므로,
$s$로 가는 모든 점 $x \in X$에 대하여 $\text{depth}(\mathcal{O}_{X, x}) \geq 2$라고 결론내릴 수 있다.
대수학의 보조정리 \ref{algebra-lemma-apply-grothendieck}을 보라.
$$
Y' \longrightarrow X
$$
를 $Y \to X \times_S U$를 사용하여 보조정리 \ref{pione-lemma-extend-S2}에서 구성한 유한 사상이라 하자.
$G$의 표준 작용을 $Y$ 위에 얻는다는 점에 주의하자. 따라서 남은 것은 $Y'$가 $X$ 위에서 에탈임을
보이는 것뿐이다. 실제로 예를 들어 보조정리 \ref{pione-lemma-purity-smooth-over-depth2}에 의해
$Y' \to X$가 올 $X_s$의 유일한 일반점 위에서 에탈임을 보이는 것만으로도 충분하다.
이는 $\sigma(s)$의 형식적 근방에서의 국소 계산으로 보인다.

\medskip\noindent
$\sigma(s)$를 포함하는 아핀 열린집합 $\Spec(B) \subset X$를 택하자. 그러면 $A \to B$는 매끄러운 환 준동형이고
단면 $\sigma : B \to A$를 갖는다. $I = \Ker(\sigma)$라 쓰고 $B^\wedge$를 $B$의 $I$-진 완비화라 하자.
그러면 어떤 $d \geq 0$에 대하여 $B^\wedge \cong A[[x_1, \ldots, x_d]]$이다. 대수학의 보조정리
\ref{algebra-lemma-section-smooth}를 보라. 물론 $B \to B^\wedge$는 평탄이고
(대수학의 보조정리 \ref{algebra-lemma-completion-flat}), $\Spec(B^\wedge) \to X$의 상은 $X_s$의 일반점을 포함한다.
$V \subset \Spec(B^\wedge)$를 $U$의 역상이라 하자. 유한 에탈 사상
$$
W = Y \times_{(X \times_S U)} V \longrightarrow V
$$
를 생각하자. 보조정리 \ref{pione-lemma-extend-S2}에서 $Y'$의 구성이 평탄 기저변환과 양립하므로, 기저변환
$Y' \times_X \Spec(B^\wedge) \to \Spec(B^\wedge)$는 $W \to V$로부터
$\Spec(B^\wedge)$ 위에서 보조정리 \ref{pione-lemma-extend-S2}의 절차로 구성된다.
$V_0 = V \cap V(x_1, \ldots, x_d) \subset V$ 및 $W_0 = W \times_V V_0$로 놓자.
이는 정규 정역 스킴이고 사상 $\Spec(B^\wedge) \to X$에 의해 $\sigma(S)$ 안으로 가며, 실제로
$\sigma(U)$와 동일시된다. 따라서 $W_0 \to V_0 = U$는 완전히 분해된다. 실제로 그 일반올에 대해서는
$Z \to X_\eta$가 $\sigma(\eta)$ 위에 놓인 $\kappa(\eta)$-유리점을 갖는다는 가정에서 이 사실이 따른다
(그리고 물론 $G$-작용으로부터 전체 올 $Z_{\sigma(\eta)}$는 스킴
$\eta = \Spec(\kappa(\eta))$의 복사본들의 서로소 합이다). 마지막으로 보조정리
\ref{pione-lemma-fully-faithful-power-series-over-depth2}에 의해
$$
W_0 \times_U V \cong W
$$
이다. 이는 $W$가 $V$의 복사본들의 서로소 합임을 보인다. 따라서
$Y' \times_X \Spec(B^\wedge)$는 $\Spec(B^\wedge)$의 복사본들의 서로소 합이고, 증명이 끝난다.
\end{proof}

\begin{lemma}
\label{pione-lemma-extend-covering-general}
$S$를 일반점 $\eta$를 갖는 준콤팩트 준분리 정역 정규 스킴이라 하자.
$f : X \to S$를 준콤팩트하고 준분리이며 기하학적으로 연결인 올들을 갖는 매끄러운 사상이라 하자.
$\sigma : S \to X$를 $f$의 단면이라 하자. $Z \to X_\eta$를 유한 에탈 Galois 덮개
(절 \ref{pione-section-finite-etale-under-galois})라 하고, 그 군 $G$의 위수가 $S$ 위에서 가역이며
$Z$가 $\sigma(\eta)$로 가는 $\kappa(\eta)$-유리점을 가진다고 하자. 그러면 군 $G$를 갖고
$X_\eta$로의 제한이 $Z$인 유한 에탈 Galois 덮개 $Y \to X$가 존재한다.
\end{lemma}

\begin{proof}
$S$가 Noether이면 이는 보조정리 \ref{pione-lemma-extend-covering}의 결과이다.
일반적인 경우는 표준적인 극한 논법으로부터 따른다.
독자에게 이 증명은 건너뛸 것을 강하게 권한다.

\medskip\noindent
$S = \lim S_i$를 아핀 전이사상을 갖는 스킴들의 계의 유향 역극한으로 쓰되,
$S_i$들은 $\mathbf{Z}$ 위 유한형이라고 하자. 극한의 명제
\ref{limits-proposition-approximate}를 보라.
각 $i$에 대하여 $S \to S'_i \to S_i$를 $S$ 안에서 취한 $S_i$의 정규화에서 오는 사상들이라 하자.
스킴의 사상의 절 \ref{morphisms-section-normalization-X-in-Y}을 보라.
대수학의 명제 \ref{algebra-proposition-ubiquity-nagata}, 스킴의 사상의 보조정리
\ref{morphisms-lemma-nagata-normalization-finite} 및
\ref{morphisms-lemma-normal-normalization}을 함께 쓰면, $S'_i$는 $\mathbf{Z}$ 위 유한형이고
$S_i$ 위 유한이며, 더 나아가 $S'_i$는 정역 정규 스킴이고 $S \to S'_i$는 우세하다고 결론내릴 수 있다.
스킴의 사상의 보조정리 \ref{morphisms-lemma-functoriality-normalization}에 의해 전이사상
$S'_{i'} \to S'_i$로서 전이사상 $S_{i'} \to S_i$ 및 정의역이 $S$인 사상들과 양립하는 것을 얻는다.
$S = \lim S'_i$라고 주장한다. 이 주장의 증명은 생략한다(힌트: 어떤 $i$에 대하여 택한 $S_i$의
아핀 열린집합 위의 아핀 열린집합들을 생각하면, 이는 소박한 대수학 문제로 바뀐다).
따라서 $S = \lim S_i$를 정규 정역 스킴 $S_i$들의 계의 유향 역극한으로 쓰되, 전이사상들은 아핀이고
$S_i$들은 $\mathbf{Z}$ 위 유한형이라고 가정해도 된다.

\medskip\noindent
어떤 $i$에 대하여 유한 표시인 매끄러운 사상 $X_i \to S_i$로서 $S$로의 기저변환이 $X \to S$인 것을
찾을 수 있다. 극한의 보조정리 \ref{limits-lemma-descend-finite-presentation} 및
\ref{limits-lemma-descend-smooth}를 보라.
$i$를 크게 한 뒤 단면 $\sigma$가 단면 $\sigma_i : S_i \to X_i$로 내려간다고 가정해도 된다
(극한의 보조정리 \ref{limits-lemma-descend-finite-presentation}에 나오는 범주들의 동치에 의한다).
$X_i$를 스킴의 사상 심화의 절 \ref{more-morphisms-section-connected-components}에서 연구한
열린부분스킴 $X_i^0$로 바꾸어도 된다. 실제로 $X \to X_i$의 상은 분명히 그 안에 놓인다
(열림성에 대해서는 스킴의 사상 심화의 보조정리
\ref{more-morphisms-lemma-connected-along-section-open}을 보라).
따라서 $X_i \to S_i$의 올들이 기하학적으로 연결이라고 가정해도 된다.
$i$를 크게 한 뒤 $|G|$가 $S_i$ 위에서 가역이라고 가정해도 된다.
$\eta_i \in S_i$를 일반점이라 하자. $X_\eta$가 스킴 $X_{i, \eta_i}$들의 극한이므로,
똑같은 논법에 의해 필요하면 $i$를 크게 하여 $Z \to X_\eta$를 어떤 유한 에탈 Galois 덮개
$Z_i \to X_{i, \eta_i}$로 내릴 수 있다. 보조정리 \ref{pione-lemma-limit}을 보라.
필요하면 $i$를 한 번 더 크게 한 뒤 $Z_i$가 $\sigma_i(\eta_i)$로 가는
$\kappa(\eta_i)$-유리점을 가진다고 가정해도 된다.
이제 Noether인 경우의 보조정리를 적용한 뒤 $X$로 당겨오면 결론을 얻는다.
\end{proof}











\section{양의 표수에서의 기법}
\label{pione-section-achinger}

\noindent
Piotr Achinger의 논문 \cite{Achinger}에서는 양의 표수인 아핀 스킴이 항상
$K(\pi, 1)$임을 증명한다. 이 절에서는 \cite{Achinger}의 더 초등적인 부분을 설명한다.
즉, 양의 표수인 체 $k$에 대하여 $\mathbf{A}^n_k$ 위에서 에탈인 아핀 스킴이 실제로는
다른 사상에 의하여 $\mathbf{A}^n_k$ 위에서 유한 에탈임을 보인다. 또한 양의 표수인
연결 아핀 스킴들의 닫힌 몰입이 에탈 기본군들 사이의 단사를 유도함도 보인다.

\medskip\noindent
$k$를 표수 $p > 0$인 체라 하자.
전사
$$
k[x_1, \ldots, x_n] \longrightarrow A
$$
를 유한형 $k$-대수들의 전사라 하되, 그 정의역은 $x_1, \ldots, x_n$에 관한 다항식환이라고 하자.
그 핵을 $I \subset k[x_1, \ldots, x_n]$라 쓰면
$A = k[x_1, \ldots, x_n]/I$이다. $A$가 비영이라고 가정하지 않는다
(달리 말하면 $A$가 영환이고 $I = k[x_1, \ldots, x_n]$인 경우도 허용한다). 마지막으로
유한 에탈 환 준동형 $\pi : A \to B$가 주어졌다고 하자.

\medskip\noindent
$k, n, k[x_1, \ldots, x_n] \to A, I, \pi : A \to B$가 주어졌다고 하자.
$C$를 $k$-대수라 하자. 가환 도식
$$
\xymatrix{
& B \\
C \ar[r] & C/\varphi(I)C \ar[u]^\tau \\
k[x_1, \ldots, x_n] \ar[u]^\varphi \ar[r] &
A \ar[u] \ar@/_3em/[uu]_\pi
}
$$
를 생각하자. 여기서 $\varphi$는 에탈 $k$-대수 준동형이고 $\tau$는 전사 $k$-대수 준동형이다.
$C, \varphi, \tau$가 주어졌다고 하자. 임의의 $r \geq 0$ 및 $y_1, \ldots, y_r \in C$로서
$C$를 $\Im(\varphi)$ 위의 대수로 생성하는 것에 대하여,
$s = s(r, y_1, \ldots, y_r) \in \{0, \ldots, r\}$를 $1 \leq i \leq s$에 대하여
$y_i$가 $\Im(\varphi)$ 위 정수가 되는 가장 큰 정수라 하자. $NF(C, \varphi, \tau)$를
위와 같은 $r$ 및 $y_1, \ldots, y_r$의 모든 선택에 걸쳐
$r - s = r - s(r, y_1, \ldots, y_r)$가 취하는 최솟값으로 정의한다.
$NF(C, \varphi, \tau)$가 $0$인 것과 $\varphi$가 유한인 것은 동치이다.

\begin{lemma}
\label{pione-lemma-lower-invariant}
위의 상황에서 $NF(C, \varphi, \tau) > 0$이면, 가환 도식
$$
\xymatrix{
& B \\
C \ar[r] & C/\varphi'(I)C \ar[u]_{\tau'} \\
k[x_1, \ldots, x_n] \ar[u]^{\varphi'} \ar[r] &
A \ar[u] \ar@/_3em/[uu]_\pi
}
$$
에 들어가고 $NF(C, \varphi', \tau') < NF(C, \varphi, \tau)$를 만족하는 에탈 $k$-대수 준동형 $\varphi'$과
전사 $k$-대수 준동형 $\tau'$가 존재한다.
\end{lemma}

\begin{proof}
$C$를 $\Im(\varphi)$ 위에서 생성하는 $r \geq 0$ 및 $y_1, \ldots, y_r \in C$를 택하고,
$y_1, \ldots, y_s$가 $\Im(\varphi)$ 위에서 정수적이며
$r - s = NF(C, \varphi, \tau) > 0$이 되도록 $0 \leq s \leq r$를 택한다.
$B$가 $A$ 위 유한이므로 $B$에서의 $y_{s + 1}$의 상은 $A$ 위의 일계수 다항식을 만족한다.
따라서 $d \geq 1$ 및 $f_1, \ldots, f_d \in k[x_1, \ldots, x_n]$로서
$$
z = y_{s + 1}^d + \varphi(f_1) y_{s + 1}^{d - 1} + \ldots + \varphi(f_d) \in
J = \Ker(C \to C/\varphi(I)C \xrightarrow{\tau} B)
$$
가 되는 것을 찾을 수 있다. $\varphi : k[x_1, \ldots, x_n] \to C$가 에탈이므로,
영이 아니고 상수가 아닌 다항식 $g \in k[T_1, \ldots, T_{n + 1}]$로서
$$
g(\varphi(x_1), \ldots, \varphi(x_n), z) = 0
\quad\text{(다음 환에서)}\quad
C
$$
가 되는 것을 찾을 수 있다. 이를 보려면 예를 들어
$C \otimes_{\varphi, k[x_1, \ldots, x_n]} k(x_1, \ldots, x_n)$가
$k(x_1, \ldots, x_n)$의 유한차 분리 체 확대들의 유한 직접곱이라는 사실
(대수학의 보조정리 \ref{algebra-lemma-etale-over-field}을 보라)을 사용하면 된다.
따라서 $z$는 $k(x_1, \ldots, x_n)$ 위의 일계수 다항식을 만족한다. 분모를 없애면 $g$를 얻는다.

\medskip\noindent
$g$의 존재와 대수학의 보조정리 \ref{algebra-lemma-helper-polynomial}에 의해, $z$가 $C$의 부분환 $C'$, 즉
$t_1 = \varphi(x_1) + z^{pe_1}, \ldots, t_n = \varphi(x_n) + z^{pe_n}$로 생성되는 부분환 위에서
정수적이 되게 하는 정수 $e_1, e_2, \ldots, e_n \geq 1$을 얻는다. 물론 원소
$\varphi(x_1), \ldots, \varphi(x_n)$도 $C'$ 위에서 정수적이고, 원소 $y_1, \ldots, y_s$도 마찬가지이다.
마지막으로 $z$를 택한 방식에 의해 원소 $y_{s + 1}$도 $C'$ 위에서 정수적이다.

\medskip\noindent
환 준동형
$$
\varphi' : k[x_1, \ldots, x_n] \longrightarrow C, \quad
x_i \longmapsto t_i
$$
를 생각하자. 그 상은 $C'$이다.
$\text{d}(\varphi(x_i)) = \text{d}(t_i) = \text{d}(\varphi'(x_i))$가
$\Omega_{C/k}$에서 성립하므로(여기서 $k$의 표수가 $p > 0$이라는 사실을 사용한다),
$\varphi$가 에탈인 것에서 $\varphi'$도 에탈이라고 결론내릴 수 있다. 대수학의 보조정리
\ref{algebra-lemma-characterize-etale-over-polynomial-ring}을 보라.
$\varphi'(x_i) - \varphi(x_i) = t_i - \varphi(x_i) = z^{pe_i}$는 $z$를 $J$의 원소로 택한 것에 의해
핵 $J$, 즉 사상 $C \to C/\varphi(I)C \to B$의 핵에 속함에 주의하자.
따라서 $f \in I$에 대하여 원소
$$
\varphi'(f) =
f(t_1, \ldots, t_n) =
f(\varphi(x_1) + z^{pe_1}, \ldots, \varphi(x_n) + z^{pe_n}) =
\varphi(f) + \text{다음 아이디얼의 원소 }(z)
$$
도 $J$에 속한다. 달리 말해 $\varphi'(I)C \subset J$이고, $A$ 위에서 에탈인 대수들 사이의 전사
$$
\tau' : C/\varphi'(I)C \longrightarrow C/J \cong B
$$
를 얻는다. 마지막으로 대수 $C$는 원소
$\varphi(x_1), \ldots, \varphi(x_n), y_1, \ldots, y_r$에 의해
$C' = \Im(\varphi')$ 위에서 생성되고,
$\varphi(x_1), \ldots, \varphi(x_n), y_1, \ldots, y_{s + 1}$는
$C' = \Im(\varphi')$ 위에서 정수적이다. 따라서
$NF(C, \varphi', \tau') < r - s = NF(C, \varphi, \tau)$이다. 이것으로 증명이 끝난다.
\end{proof}

\begin{lemma}
\label{pione-lemma-affine-etale-over-affine-space}
$k$를 표수 $p > 0$인 체라 하자. $X \to \mathbf{A}^n_k$를 에탈 사상이라 하고 $X$가 아핀이라고 하자.
그러면 유한 에탈 사상 $X \to \mathbf{A}^n_k$가 존재한다.
\end{lemma}

\begin{proof}
$X = \Spec(C)$라 쓰자. $A = 0$으로 놓고 $I = k[x_1, \ldots, x_n]$라 쓰자.
가정에 의해 어떤 에탈 $k$-대수 준동형
$\varphi : k[x_1, \ldots, x_n] \to C$가 존재한다. $\tau : C/\varphi(I)C \to 0$를 유일한 전사라 쓰자.
$N(C, \varphi, \tau)$가 최소가 되도록 $\varphi$와 $\tau$를 택할 수 있다. 보조정리
\ref{pione-lemma-lower-invariant}에 의해 $N(C, \varphi, \tau) = 0$을 얻는다. 따라서 $\varphi$는 유한 에탈이다.
\end{proof}

\begin{lemma}
\label{pione-lemma-dominate-affine-space}
$k$를 표수 $p > 0$인 체라 하자. $Z \subset \mathbf{A}^n_k$를 닫힌부분스킴이라 하자.
$Y \to Z$가 유한 에탈이라고 하자. 그러면 유한 에탈 사상 $f : U \to \mathbf{A}^n_k$로서,
$Z$ 위의 열린 동시에 닫힌 몰입 $Y \to f^{-1}(Z)$가 존재하게 하는 것이 존재한다.
\end{lemma}

\begin{proof}
이 문제를 대수학으로 바꾸어 보자.
$\mathbf{A}^n_k = \Spec(k[x_1, \ldots, x_n])$.
라고 쓰자. 그러면 $Z = \Spec(A)$이고 $A = k[x_1, \ldots, x_n]/I$이다. 여기서
$I \subset k[x_1, \ldots, x_n]$는 어떤 아이디얼이다.
$Y = \Spec(B)$라 쓰면 $Y \to Z$는 유한 에탈 $k$-대수 준동형 $A \to B$에 대응한다.

\medskip\noindent
대수학의 보조정리 \ref{algebra-lemma-lift-etale}에 의해 에탈 환 준동형
$$
\varphi : k[x_1, \ldots, x_n] \to C
$$
와 전사 $A$-대수 준동형 $\tau : C/\varphi(I)C \to B$가 존재한다.
($C, \varphi, \tau$는 $\tau$가 동형사상이 되도록 택할 수도 있지만, 이를 사용하지 않는다.)
$N(C, \varphi, \tau)$가 최소가 되도록 $\varphi$와 $\tau$를 택할 수 있다. 보조정리
\ref{pione-lemma-lower-invariant}에 의해 $N(C, \varphi, \tau) = 0$을 얻는다. 따라서 $\varphi$는 유한 에탈이다.

\medskip\noindent
$f : U = \Spec(C) \to \mathbf{A}^n_k$를 $\varphi$에 대응하는 유한 에탈 사상이라 하자.
사상 $Y \to f^{-1}(Z) = \Spec(C/\varphi(I)C)$는 $\tau$에 의해 유도되고, $\tau$가 전사이므로 닫힌 몰입이며
스킴의 사상의 보조정리 \ref{morphisms-lemma-etale-permanence}에 의해 에탈 사상이므로 열리기도 하다.
이것으로 증명이 끝난다.
\end{proof}

\noindent
다음이 주요 결과이다.

\begin{proposition}
\label{pione-proposition-injective}
$p$를 소수라 하자. $i : Z \to X$를 $\mathbf{F}_p$ 위의 연결 아핀 스킴들의 닫힌 몰입이라 하자.
$Z$의 임의의 기하점 $\overline{z}$에 대하여 사상
$$
\pi_1(Z, \overline{z}) \to \pi_1(X, \overline{z})
$$
은 단사이다.
\end{proposition}

\begin{proof}
$Y \to Z$를 유한 에탈 사상이라 하자. $Y$가 $f^{-1}(Z)$의 열린 동시에 닫힌 부분스킴과 동형이 되게 하는
유한 에탈 사상 $f : U \to X$를 구성하면 충분하다.
보조정리 \ref{pione-lemma-functoriality-galois-injective}를 보라.
$Y = \Spec(A)$ 및 $X = \Spec(R)$라 쓰면 닫힌 몰입 $Y \to X$는 전사 $R \to A$로 주어진다.
$A = \colim A_i$를 유한형 $\mathbf{F}_p$-부분대수들의 여과 여극한으로 쓸 수 있다.
보조정리 \ref{pione-lemma-limit}에 의해 어떤 $i$와 유한 에탈 사상 $Y_i \to Z_i = \Spec(A_i)$로서
$Y = Z \times_{Z_i} Y_i$가 되는 것을 찾을 수 있다.

\medskip\noindent
전사 $\mathbf{F}_p[x_1, \ldots, x_n] \to A_i$를 택하자. 이는 닫힌 몰입
$$
Z_i = \Spec(A_i)
\longrightarrow
X_i = \mathbf{A}^n_{\mathbf{F}_p} = \Spec(\mathbf{F}_p[x_1, \ldots, x_n])
$$
을 정한다. 다항식 대수의 보편 성질과 $R \to A$가 전사라는 사실에 의해 가환 도식
$$
\xymatrix{
\mathbf{F}_p[x_1, \ldots, x_n] \ar[r] \ar[d] &
A_i \ar[d] \\
R \ar[r] &
A
}
$$
을 $\mathbf{F}_p$-대수들 사이에서 찾을 수 있다. 따라서 가환 도식
$$
\xymatrix{
Y_i \ar[r] &
Z_i \ar[r] &
X_i \\
Y \ar[u] \ar[r] &
Z \ar[u] \ar[r] &
X \ar[u]
}
$$
을 얻고, 그 오른쪽 정사각형은 데카르트이다. 유한 에탈 사상 $f_i : U_i \to X_i$로서 $Y_i$가
$f_i^{-1}(Z_i)$의 열린 동시에 닫힌 부분스킴과 동형이 되는 것을 찾을 수 있다면,
$f_i$를 $X \to X_i$에 의해 기저변환한 $f : U \to X$가 이 문제의 해라는 것은 명백하다.
따라서 보조정리 \ref{pione-lemma-dominate-affine-space}를
$Y_i \to Z_i \to X_i = \mathbf{A}^n_{\mathbf{F}_p}$에 적용하면 결론을 얻는다.
\end{proof}
